% Time-stamp: <[choco-doc.tex] last changed 28-10-2010 13:18 by sofdem>
\documentclass[a4paper,10pt]{book}
\usepackage{hyperref}
\usepackage{fancyhdr,lastpage,fullpage,fancybox}
\usepackage{amsmath,amssymb}

%%%% Glossary  %%%%%%%%%%%%%%%%%%%%%%%%%%%%%%%%%%%%%%%%%%%%%%%%%%%%%%%%%
\usepackage[acronym]{glossaries} % make a separate list of acronyms
% to use call
% $> latex choco-doc 
% $> makeglossaries choco-doc 
% $> latex choco-doc 

\makeglossaries

%%%% Web Tralics  %%%%%%%%%%%%%%%%%%%%%%%%%%%%%%%%%%%%%%%%%%%%%%%%%%%%%%%%%
\newif\ifweb
%\webtrue
\webfalse

\ifweb
\else
%%%% Refs, formats, styles  %%%%%%%%%%%%%%%%%%%%%%%%%%%%%%%%%%%%%%%%%%%%%%%%%%%%%%%%%
\hypersetup{plainpages=false, colorlinks=true, urlcolor=blue, linkcolor=cyan, citecolor=orange, filecolor=red}
\fancypagestyle{plain}{\fancyhf{}
    \fancyhead[LE]{\slshape \leftmark} 
    \fancyhead[RO]{\slshape \rightmark} 
    \lfoot{\footnotesize\emph{CHOCO solver documentation \\BSD licence \number\year}}
    \cfoot{\footnotesize\textsf{-\thepage/\pageref{LastPage}-}}
    \rfoot{\footnotesize\emph{\number\day/\number\month/\number\year}}}
\fancypagestyle{fancyfront}{\fancyhf{}
    \fancyhead[LE]{\slshape \leftmark} 
    \fancyhead[RO]{\slshape \rightmark} 
    \lfoot{\footnotesize\emph{CHOCO solver documentation \\BSD licence \number\year}}
    \cfoot{\footnotesize\textsf{-\thepage-}}
    \rfoot{\footnotesize\emph{\number\day/\number\month/\number\year}}}
%    \renewcommand{\headrulewidth}{0pt}}
\addtolength{\headsep}{\baselineskip} 

\makeatletter 
\def\cleardoublepage{\clearpage\if@twoside \ifodd\c@page\else 
\hbox{} 
\thispagestyle{empty} 
\newpage 
\if@twocolumn\hbox{}\newpage\fi\fi\fi} 
\makeatother 

%%%% Fonts, Figures, Colors %%%%%%%%%%%%%%%%%%%%%%%%%%%%%%%%%%%%%%%%%%%%%%%%%%%%%%%%%
\newcommand{\algofont}{\fontencoding{U}\fontfamily{eur}\fontseries{b}\fontshape{n}\selectfont}
\usepackage[usenames]{color}
\usepackage[pdftex]{graphicx,xcolor} %\pdfcompresslevel=9  %\pdfimageresolution=600
\def\Graphique#1#2{\includegraphics[#2]{#1.pdf}}
\def\Graph#1#2{\includegraphics[#2]{#1}}
%\def\insertLogo#1{\includegraphics{#1}}
\newcommand{\insertGraphique}[3]{\begin{figure}[htp]\centerline{\Graph{#2}{width=#1}}\caption[]{#3}\label{fig:#2}\end{figure}}
\newcommand{\newrgbcolor}[2]{\definecolor{#1}{rgb}{#2}}
\definecolor{darkred}{rgb}{0.5372549, 0.27843137, 0.27058824}
\definecolor{darkblue}{rgb}{0.27843137, 0.27058824, 0.5372549}
\newcommand{\gr}[1]{{\color{lightgray} #1}}
\newcommand{\dr}[1]{{\color{darkblue} #1}}
%%%% listings, env  %%%%%%%%%%%%%%%%%%%%%%%%%%%%%%%%%%%%%%%%%%%%%%%%%%%%%%%%%%%%%%%%%%%%%%%
\usepackage{listings}
\definecolor{Brown}{cmyk}{0,0.81,1,0.60}
\definecolor{OliveGreen}{cmyk}{0.64,0,0.95,0.40}
\definecolor{CadetBlue}{cmyk}{0.62,0.57,0.23,0}
\lstset{language=java,frame=ltrb,framesep=5pt,
    numbers=none,
    tabsize=3,
    breaklines=true,
    aboveskip=1ex,
    belowskip=1ex,
    basicstyle=\small\ttfamily,
    backgroundcolor=\color{gray!25},
    columns=fullflexible,
    keywordstyle=\ttfamily\color{OliveGreen},
    identifierstyle=\ttfamily\color{CadetBlue}\bfseries, 
    commentstyle=\color{Brown},
    stringstyle=\ttfamily,
    showstringspaces=true}

\newcommand{\mylst}[1]{\lstinline|#1|}

\newenvironment{note}{%
\setlength{\fboxsep}{15pt}%
\Sbox% 
\minipage{.7\linewidth}% 
}% 
{\endminipage\endSbox% 
\begin{center}\doublebox{\TheSbox}\end{center}}

\newenvironment{notedef}{%
\setlength{\fboxsep}{15pt}%
\Sbox% 
\minipage{.95\linewidth}% 
}% 
{\endminipage\endSbox% 
\par\hspace{-1em}\shadowbox{\TheSbox}\par\noindent}

\newenvironment{myquote}{\begin{notedef}}{\end{notedef}}
%\newenvironment{notedef}{\begin{note}}{\end{note}}
\fi
%%%%%%%%%%%%%%%%%%%%%%%%%%%%%%%%%%%%%%%%%%%%%%%%%%%%%%%%%%%%%%%%%%%%%%%%%%%%%%%%%%%%%
\newcommand{\collec}[2]{\ensuremath{\langle #1,..,#2\rangle}}
\newcommand{\coll}[1]{\ensuremath{\langle #1\rangle}}
\newcommand{\N}{\ensuremath{\mathbb{N}}}
\newcommand{\Z}{\ensuremath{\mathbb{Z}}}
\newcommand{\R}{\ensuremath{\mathbb{R}}}
\newcommand{\defarrow}{\ensuremath{\stackrel{\text{\tiny def}}{\longleftrightarrow}}}
\usepackage{theorem}
\newtheorem{definition}{Definition}
\newenvironment{proof}{\par\noindent\textbf{Proof.} }{\ \rule{0.5em}{0.5em}\par\bigskip}
\newtheorem{theorem}{Theorem}
\newtheorem{lemma}[theorem]{Lemma}
\newtheorem{proposition}[theorem]{Proposition}
\newtheorem{corollary}[theorem]{Corollary}
\newtheorem{example}{Example}
\newtheorem{question}{Question}

\newcommand{\gui}[1]{\flqq\,#1\,\frqq}
\newcommand{\todo}[1]{\textcolor{orange}{#1}}

\newcommand{\HRule}{\rule{\linewidth}{0.5mm}}

\newcommand{\eg}{\textit{e.g.},~}
\newcommand{\ie}{\textit{i.e.},~}
%%%%%%%%%%%%%%%%%%%%%%%%%%%%%%%%%%%%%%%%%%%%%%%%%%%%%%%%%%%%%%%%%%%%%%%%%%%%%%%%%%%%%

\begin{document}

\pagestyle{fancyfront}
\frontmatter

\begin{titlepage}
\begin{center}
\vspace*{\stretch{1}}
% Upper part of the page
\textsc{\LARGE Choco solver}\\
\url{http://choco.mines-nantes.fr/}\\[1.5cm]

% Title
\HRule \\[0.4cm]
{ \Huge \bfseries Documentation}\\
\HRule \\[1cm]
\vspace*{\stretch{1}} 
\includegraphics[width=0.20\textwidth]{../shared/media/logo.pdf} 

\vfill
% Bottom of the page
{\large \today}
\end{center}
\end{titlepage}

\clearpage
\vspace*{\stretch{1}}
\begin{myquote}
Copyright (C)  1999, 2010  F. Laburthe, N. Jussien.\\
Permission is granted to copy, distribute and/or modify this document\\
under the terms of the GNU Free Documentation License, Version 1.3\\
or any later version published by the Free Software Foundation;\\
with no Invariant Sections, no Front-Cover Texts, and no Back-Cover Texts.\\
A copy of the license is included in the section entitled \hyperlink{licence}{"GNU \\Free Documentation License"}.
\end{myquote}
\vspace*{\stretch{1}}

\tableofcontents

\mainmatter
\pagestyle{plain}

%!TEX root = choco-doc.tex
\addcontentsline{toc}{chapter}{Preface}
\chapter*{Preface}
Choco is a java library for constraint satisfaction problems (CSP) and constraint programming (CP). It is built on a event-based propagation mechanism with backtrackable structures.
Choco is an open-source software, distributed under a \textbf{BSD licence} and hosted by \href{http://sourceforge.net/projects/choco/}{sourceforge.net}.
For any informations visit \url{http://choco.emn.fr}.
\bigskip

\noindent This document is organized as follows:
\begin{itemize}
\item \hyperlink{ch:doc}{Documentation} is the user-guide of Choco. After a short \hyperlink{doc:introduction}{introduction} to constraint programming and to the Choco solver, it presents the basics of \hyperlink{doc:model}{modeling} and \hyperlink{doc:solver}{solving} with Choco, and some \hyperlink{doc:advanced}{advanced usages} (customizing propagation and search).
\item \hyperlink{part:elements}{Elements of Choco} gives a detailed description of the \hyperlink{ch:vars}{variables}, \hyperlink{ch:operators}{operators}, \hyperlink{ch:constraints}{constraints} currently available in Choco.
\item \hyperlink{ch:extra}{Extras} presents future works, only available on the beta version or extension of the current jar, such as the \hyperlink{chocoandvisu:chocoandvisu}{visualization module of Choco}. The section dedicated to \hyperlink{sudokuandcp:sudokuandconstraintprogramming}{Sudoku} aims at explaining the basic principles of Constraint Programming (propagation and search) on this famous game.


\end{itemize}

\part{Documentation}\label{ch:doc}\hypertarget{ch:doc}{}

The documentation of Choco is organized as follows:
\begin{itemize}
\item 
The concise \hyperlink{doc:introduction}{introduction} provides some informations \hyperlink{introduction:aboutconstraintprogramming}{about constraint programming} concepts and a ``Hello world''-like \hyperlink{introduction:myfirstchocoprogram}{first Choco program}.
\item 
The \hyperlink{doc:model}{model} section gives informations on \hyperlink{doc:model}{how to create a model} and introduces \hyperlink{model:variables}{variables} and \hyperlink{model:constraints}{constraints}.
\item 
The \hyperlink{doc:solver}{solver} section gives informations on \hyperlink{doc:solver}{how to create a solver}, to \hyperlink{doc:solver}{read a model}, to define a \hyperlink{solver:searchstrategy}{search strategy}, and finally to \hyperlink{solver:solveaproblem}{solve a problem}.
\item 
The \hyperlink{doc:advanced}{advanced use} section explains how to define your own \hyperlink{advanced:defineyourownlimitsearchspace}{limit search space}, \hyperlink{advanced:defineyourownsearchstrategy}{search strategy}, \hyperlink{advanced:defineyourownconstraint}{constraint}, \hyperlink{advanced:defineyourownoperator}{operator}, \hyperlink{advanced:defineyourownvariable}{variable}, or \hyperlink{advanced:backtrackablestructures}{backtrackable structure}.
\item 
The \hyperlink{doc:applications}{applications} section shows the use of Choco defined global constraints on \hyperlink{schedulinganduseofthecumulative:schedulinganduseofthecumulativeconstraint}{scheduling} or \hyperlink{geostdescription:placementanduseofthegeostconstraint}{placement} problems.
%\item 
%Lastly, the catalog of Choco defined \hyperlink{ch:constraints}{constraints} is presented.
\end{itemize}

%\section*{Beta}\label{documentation:beta}\hypertarget{documentation:beta}{}
%Here you can find short documentation concerning futur works, only available on the beta version or extension of the current jar:
%\begin{itemize}
%%	\item \hyperlink{chocoandgraphviz}{Choco and Graphviz} \emph{not yet available. Has been included in \hyperlink{chocoandvisu}{Choco and Visu}.}
%	\item \hyperlink{chocoandvisu}{Choco and Visu}
%\end{itemize}

%\section*{Material}\label{documentation:material}\hypertarget{documentation:material}{}
%Here you can find some materials at your disposal. If you create any, you can also send it to us, so we can add it to the list.

%\subsection{Presentations}\label{documentation:presentations}\hypertarget{documentation:presentations}{}
%\begin{itemize}
%	\item \textbf{August 2008} : The \emph{CHOCO : an Open Source Java Constraint Constraint Programming} white paper presentation send to the \href{http://www.cril.univ-artois.fr/cpai08/}{CSP'08 competition}. \href{media/pdf/choco-presentation.pdf}{PDF} of white paper presentation of Choco
%	\item \textbf{june 2008} : \emph{The CHOCO constraint programming solver} presentation held at the \href{https://projects.coin-or.org/events/wiki/cpaior2008}{Workshop on Open-Source Software for Integer and Constraint Programming} during the last \href{http://contraintes.inria.fr/cpaior08/}{CPAIOR} conference. \href{media/slides/cpaior-choco.pdf}{PDF slides} of the presentation given by Guillaume Rochart
%\end{itemize}
%\Graph{media/banner.png}{width=\linewidth}

\chapter{Introduction to constraint programming and Choco}\label{doc:introduction}\hypertarget{doc:introduction}{}

\section{About constraint programming}\label{introduction:aboutconstraintprogramming}\hypertarget{introduction:aboutconstraintprogramming}{}

\begin{myquote}
Constraint programming represents one of the closest approaches computer science has yet made to the Holy Grail of programming: the user states the problem, the computer solves it.
\begin{flushright}\bf E. C. Freuder, Constraints, 1997.\end{flushright}
\end{myquote}


Fast increasing computing power in the 1960s led to a wealth of works around problem solving, at the root of Operational Research, Numerical Analysis, Symbolic Computing, Scientific Computing, and a large part of Artificial Intelligence and programming languages. Constraint Programming is a discipline that gathers, interbreeds, and unifies ideas shared by all these domains to tackle decision support problems.

Constraint programming has been successfully applied in numerous domains. Recent applications include computer graphics (to express geometric coherence in the case of scene analysis), natural language processing (construction of efficient parsers), database systems (to ensure and/or restore consistency of the data), operations research problems (scheduling, routing), molecular biology (DNA sequencing), business applications (option trading), electrical engineering (to locate faults), circuit design (to compute layouts), etc.

Current research in this area deals with various fundamental issues, with implementation aspects and with new applications of constraint programming.

\subsection{Constraints}\label{introduction:constraints}\hypertarget{introduction:constraints}{}
A constraint is simply a logical relation among several unknowns (or variables), each taking a value in a given domain. A constraint thus restricts the possible values that variables can take, it represents some partial information about the variables of interest. For instance, the circle is inside the square relates two objects without precisely specifying their positions, i.e., their coordinates. Now, one may move the square or the circle and he or she is still able to maintain the relation between these two objects. Also, one may want to add another object, say a triangle, and to introduce another constraint, say the square is to the left of the triangle. From the user (human) point of view, everything remains absolutely transparent.

Constraints naturally meet several interesting properties:
\begin{itemize}
	\item constraints may specify partial information, i.e. constraint need not uniquely specify the values of its variables,
	\item constraints are non-directional, typically a constraint on (say) two variables $X, Y$ can be used to infer a constraint on $X$ given a constraint on $Y$ and vice versa,
	\item constraints are declarative, i.e. they specify what relationship must hold without specifying a computational procedure to enforce that relationship,
	\item constraints are additive, i.e. the order of imposition of constraints does not matter, all that matters at the end is that the conjunction of constraints is in effect,
	\item constraints are rarely independent, typically constraints in the constraint store share variables.
\end{itemize}

Constraints arise naturally in most areas of human endeavor. The three angles of a triangle sum to 180 degrees, the sum of the currents floating into a node must equal zero, the position of the scroller in the window scrollbar must reflect the visible part of the underlying document, these are some examples of constraints which appear in the real world. Thus, constraints are a natural medium for people to express problems in many fields. 

\subsection{Constraint Programming}\label{introduction:constraintprogramming}\hypertarget{introduction:constraintprogramming}{}
Constraint programming is the study of computational systems based on constraints. The idea of constraint programming is to solve problems by stating constraints (conditions, properties) which must be satisfied by the solution.

Work in this area can be tracked back to research in Artificial Intelligence and Computer Graphics in the sixties and seventies. Only in the last decade, however, has there emerged a growing realization that these ideas provide the basis for a powerful approach to programming, modeling and problem solving and that different efforts to exploit these ideas can be unified under a common conceptual and practical framework, constraint programming. 

\begin{note}
If you know \textbf{sudoku}, then you know \textbf{constraint programming}. See why \hyperlink{sudokuandcp}{here}.
\end{note}


\section{Modeling with Constraint programming}\label{introduction:modelingwithconstraintprogramming}\hypertarget{introduction:modelingwithconstraintprogramming}{}
The formulation and the resolution of combinatorial problems are the two main goals of the constraint programming domain. This is an essential way to solve many interesting industrial problems such as scheduling, planning or design of timetables. The main interest of constraint programming is to propose to the user to model a problem without being interested in the way the problem is solved.

\subsection{The Constraint Satisfaction Problem}\label{introduction:csp}\hypertarget{introduction:csp}{}

Constraint programming allows to solve combinatorial problems modeled by a Constraint Satisfaction Problem (CSP). Formally, a CSP is defined by a triplet $(X,D,C)$:
\begin{itemize}
	\item \textbf{Variables}: $X = \{X_1,X_2,\ldots,X_n\}$ is the set of variables of the problem.
	\item \textbf{Domains}: $D$ is a function which associates to each variable $X_i$ its domain $D(X_i)$, i.e. the set of possible values that can be assigned to $X_i$. The domain of a variable is usually a finite set of integers: $D(X_i)\subset\Z$ (\emph{integer variable}). But a domain can also be continuous ($D(X_i)\subseteq\R$ for a \emph{real variable}) or made of discrete set values ($D(X_i)\subseteq\mathcal{P}(\Z)$ for a \emph{set variable}).
	\item \textbf{Constraints}: $C = \{C_1,C_2,\ldots,C_m\}$ is the set of constraints. A constraint $C_j$ is a relation defined on a subset $X^j = \{X^j_1,X^j_2,\ldots,X^j_{n^j}\}\subseteq X$ of variables which restricts the possible tuples of values $(v_1,\ldots,v_{n^j})$ for these variables:
$$(v_1,\ldots,v_{n^j})\in C_j\cap (D(X^j_1)\times D(X^j_2)\times\cdots\times D(X^j_{n^j})).$$
Such a relation can be defined explicitely (ex: $(X_1,X_2)\in\{(0,1),(1,0)\}$) or implicitely (ex: $X_1+X_2\le 1$).
\end{itemize}

Solving a CSP is to find a tuple $v=(v_1,\ldots,v_{n})\in D(X)$ on the set of variables which satisfies all the constraints:
$$(v_1,\ldots,v_{n^j})\in C_j,\quad\forall j\in\{1,\ldots,m\}.$$

For optimization problems, one need to define an \textbf{objective function} $f:D(X)\rightarrow\R$. An optimal solution is then a solution tuple of the CSP that minimizes (or maximizes) function $f$.

\subsection{Examples of CSP models}\label{introduction:examples}\hypertarget{introduction:examples}{}
This part provides three examples using different types of variables in different problems. These examples are used throughout this tutorial to illustrate their modeling with Choco.

\subsubsection{Example 1: the n-queens problem.}\label{introduction:example1:nqueens}\hypertarget{introduction:example1:nqueens}{}
Let us consider a chess board with $n$ rows and $n$ columns. A queen can move as far as she pleases, horizontally, vertically, or diagonally. The standard $n$-queens problem asks how to place $n$ queens on an $n$-ary chess board so that none of them can hit any other in one move.

The $n$-queens problem can be modeled by a CSP in the following way:
\begin{itemize}
	\item \textbf{Variables}: $X = \{X_{i}\ |\ i\in [1,n]\}$.
	\item \textbf{Domain}: for all variable $X_{i}\in X$, $D(X_{i}) = \{1,2,\ldots, n\}$.
	\item \textbf{Constraints}: the set of constraints is defined by the union of the three following constraints,
	\begin{itemize}
		\item queens have to be on distinct lines:
		\begin{itemize}
			\item $C_{lines} = \{X_{i}\neq X_{j}\ |\ i,j\in [1,n], i\neq j\}$.
		\end{itemize}
		\item queens have to be on distinct diagonals:
		\begin{itemize}
			\item $C_{diag1} = \{X_{i}\neq X_{j+j-i}\ |\ i,j\in [1,n], i\neq j\}$.
			\item $C_{diag2} = \{X_{i}\neq X_{j+i-j}\ |\ i,j\in [1,n], i\neq j\}$.
		\end{itemize}
	\end{itemize}
\end{itemize}

\subsubsection{Example 2: the ternary Steiner problem.}\label{introduction:example2:theternarysteinerproblem}\hypertarget{introduction:example2:theternarysteinerproblem}{}
A ternary Steiner system of order $n$ is a set of $n*(n-1)/6$ triplets of distinct elements taking their values in $[1,n]$, such that all the pairs included in two distinct triplets are different.
See \url{http://mathworld.wolfram.com/SteinerTripleSystem.html} for details. 

The ternary Steiner problem can be modeled by a CSP in the following way:
\begin{itemize}
	\item let $t = n*(n-1)/6$.
	\item \textbf{Variables}: $X = \{X_{i}\ |\ i\in [1,t]\}$.
	\item \textbf{Domain}: for all $i\in [1,t]$, $D(X_{i}) = \{1,...,n\}$.
	\item \textbf{Constraints}:
	\begin{itemize}
		\item every set variable $X_i$ has a cardinality of 3:
		\begin{itemize}
			\item for all $i\in [1,t]$, $|X_{i}| = 3$.
		\end{itemize}
		\item the cardinality of the intersection of every two distinct sets must not exceed 1:
		\begin{itemize}
			\item for all $i,j\in [1,t]$, $i\neq j$, $|X_{i}\cap X_{j}|\le 1$.
		\end{itemize}
	\end{itemize}
\end{itemize}

\subsubsection{Example 3: the CycloHexane problem.}\label{introduction:example3:thecyclohexaneproblem}\hypertarget{introduction:example3:thecyclohexaneproblem}{}
The problem consists in finding the 3D configuration of a cyclohexane molecule. It is described with a system of three non linear equations:
\begin{itemize}
	\item \textbf{Variables}: $x,y,z$.
	\item \textbf{Domain}: $]-\infty;+\infty[$.
	\item \textbf{Constraints}:
	\begin{align*}
		y^{2} * (1 + z^{2}) + z * (z - 24 * y) &= -13\\
		x^{2} * (1 + y^{2}) + y * (y - 24 * x) &= -13\\
		z^{2} * (1 + x^{2}) + x * (x - 24 * z) &= -13
	\end{align*}
\end{itemize}

\section{My first Choco program: the magic square}\label{introduction:myfirstchocoprogram}\hypertarget{introduction:myfirstchocoprogram}{}

\subsection{The magic square problem}\label{introduction:amagicsquareproblem}\hypertarget{introduction:amagicsquareproblem}{}
In the following, we will address the magic square problem of order 3 to illustrate step-by-step how to model and solve this problem using choco. 

\subsubsection{Definition:}
A magic square of order $n$ is an arrangement of $n^{2}$ numbers, usually distinct integers, in a square, such that the $n$ numbers in all rows, all columns, and both diagonals sum to the same constant. A standard magic square contains the integers from 1 to $n^{2}$.

The constant sum in every row, column and diagonal is called the magic constant or magic sum $M$. The magic constant of a classic magic square depends only on $n$ and has the value:
$M(n)=n(n^2 +1)/2$.

\href{http://en.wikipedia.org/wiki/magicsquare}{More details on the magic square problem.}


\subsection{A mathematical model}\label{introduction:mathematicalmodeling}\hypertarget{introduction:mathematicalmodeling}{}

Let $x_{ij}$ be the variable indicating the value of the $j^{th}$ cell of row $i$. 
Let $C$ be the set of constraints modeling the magic square as:
\begin{align*}
&x_{ij} \in [1,n^2],\ &&\forall i,j \in [1, n]\\
&x_{ij}\ne x_{kl},\ &&\forall i,j,k,l \in [1,n], i\ne k, j\ne l\\
&\sum_{j=1}^{n} x_{ij} = n^2,\ &&\forall i \in [1,n]\\
&\sum_{i=1}^{n} x_{ij} = n^2,\ &&\forall j \in [1,n]\\
&\sum_{i=1}^{n} x_{ii} = n^2&&\\
&\sum_{i=n}^{1} x_{i(n-i)} = n^2&&\\
\end{align*}

We have all the required information to model the problem with Choco.
\begin{note}
	For the moment, we just talk about \emph{model translation} from a mathematical representation to Choco.
	Choco can be used as a \emph{black box}, that means we just need to define the problem without knowing the way it will be solved. We can therefore focus on the modeling not on the solving.
\end{note}

\subsection{To Choco...}\label{introduction:inchoco}\hypertarget{introduction:inchoco}{}

First, we need to know some of the basic Choco objects:
\begin{itemize}
\item 
The \textbf{model} (object \texttt{Model} in Choco) is one of the central elements of a Choco program. Variables and constraints are associated to it.
\item
The \textbf{variables} (objects \texttt{IntegerVariable}, \texttt{SetVariable}, and \texttt{RealVariable} in Choco) are the \emph{unknown} of the problem. Values of variables are taken from a \textbf{domain} which is defined by a set of values or quite often simply by a lower bound and an upper bound of the allowed values. The domain is given when creating the variable.
\begin{note}
Do not forget that we manipulate \textbf{variables} in the mathematical sense (as opposed to classical computer science). Their effective value will be known only once the problem has been solved.
\end{note}
\item
The \textbf{constraints} define relations to be satisfied between variables and constants.
In our first model, we only use the following constraints provided by Choco:
\begin{itemize}
	\item \texttt{eq(var1, var2)} which ensures that \texttt{var1} equals \texttt{var2}.
	\item \texttt{neq(var1, var2)} which ensures that \texttt{var1} is not equal to \texttt{var2}.
	\item \texttt{sum(var[])} which returns expression \texttt{var[0]+var[1]+...+var[n]}.
\end{itemize}
\end{itemize}

\subsection{The program}\label{introduction:theprogram}\hypertarget{introduction:theprogram}{}
After having created your java class file, import the Choco class to use the API:
\begin{lstlisting}
  import choco.Choco;
\end{lstlisting}
First of all, let's create a Model:
\lstinputlisting{java/imagicsquare1.j2t}
We create an instance of \texttt{CPModel()} for \textbf{C}onstraint \textbf{P}rogramming Model.
Do not forget to add the following imports:
\begin{lstlisting}
  import choco.cp.model.CPModel;
\end{lstlisting}
Then we declare the variables of the problem:
\lstinputlisting{java/imagicsquare2.j2t}
Add the import:
\begin{lstlisting}
  import choco.kernel.model.variables.integer.IntegerVariable;
\end{lstlisting}
We have defined the variable using the \texttt{makeIntVar} method which creates an enumerated domain: all the values are stored in the java object (beware, it is usually not necessary to store all the values and it is less efficient than to create a bounded domain).

\noindent Now, we are going to state a constraint ensuring that all variables must have a different value:
\lstinputlisting{java/imagicsquare3.j2t}
Add the import:
\begin{lstlisting}
  import choco.kernel.model.constraints.Constraint;
\end{lstlisting}
Then, we add the constraint ensuring that the magic sum is respected:
\lstinputlisting{java/imagicsquare4.j2t}
Then we define the constraint ensuring that each column is equal to the magic sum.
Actually, \texttt{var} just denotes the rows of the square. So we have to declare a temporary array of variables that defines the columns.
\lstinputlisting{java/imagicsquare5.j2t}
It is sometimes useful to define some temporary variables to keep the model simple or to reorder array of variables. That is why we also define two other temporary arrays for diagonals.
\lstinputlisting{java/imagicsquare6.j2t}
Now, we have defined the model. The next step is to solve it.
For that, we build a Solver object:
\lstinputlisting{java/imagicsquare7.j2t}

with the imports:
\begin{lstlisting}
  import choco.cp.solver.CPSolver;
\end{lstlisting}
We create an instance of \texttt{CPSolver()} for Constraint Programming Solver.
Then, the solver reads (translates) the model and solves it:
\lstinputlisting{java/imagicsquare8.j2t}
The only variables that need to be printed are the ones in \texttt{var} (all the others are only references to these ones). 
\begin{note}
We have to use the Solver to get the value of each variable of the model. The Model only declares the objects, the Solver finds their value.
\end{note}
We are done, we have created our first Choco program. 
The complete source code can be found here: \href{media/zip/exmagicsquare.zip}{ExMagicSquare.zip}


\subsection{In summary}\label{introduction:whatisimportant}\hypertarget{introduction:whatisimportant}{}
\begin{itemize}
	\item A Choco Model is defined by a set of Variables with a given domain and a set of Constraints that link Variables:
it is necessary to add both Variables and Constraints to the Model.
	\item temporary Variables are useful to keep the Model readable, or necessary when reordering arrays.
	\item The value of a Variable can be known only once the Solver has found a solution.
	\item To keep the code readable, you can avoid the calls to the static methods of the Choco classes, by importing the static classes, i.e. instead of:
\begin{lstlisting}
  import choco.Choco;
  ...
  IntegerVariable v = Choco.makeIntVar("v", 1, 10);
  ...
  Constraint c = Choco.eq(v, 5);
\end{lstlisting}
you can use:
\begin{lstlisting}
  import static choco.Choco.*;
  ...
  IntegerVariable v = makeIntVar("v", 1, 10);
  ...
  Constraint c = eq(v, 5);
\end{lstlisting}
\end{itemize}

\section{Complete examples}\label{model:completeexamples}\hypertarget{model:completeexamples}{}
We provide now the complete Choco model for the three examples \hyperlink{introduction:examples}{previously described}.

\subsection{Example 1: the n-queens problem with Choco}\label{model:example1:nqueenschoco}\hypertarget{model:example1:nqueenschoco}{}
This first model for the \hyperlink{introduction:example1:nqueens}{n-queens problem} only involves binary constraints of differences between integer variables. One can immediately recognize the 4 main elements of any Choco code. First of all, create the model object. Then create the variables by using the Choco API (One variable per queen giving the row (or the column) where the queen will be placed). Finally, add the constraints and solve the problem. 

\lstinputlisting{java/inqueen.j2t}

\subsection{Example 2: the ternary Steiner problem with Choco}\label{model:example2:ternarysteinerchoco}\hypertarget{model:example2:ternarysteinerchoco}{}
The \hyperlink{introduction:example2:theternarysteinerproblem}{ternary Steiner problem} is entirely modeled using set variables and set constraints. 
\lstinputlisting{java/iternarysteiner.j2t}

\subsection{Example 3: the CycloHexane problem with Choco}\label{model:example3:thecyclohexaneproblemwithchoco}\hypertarget{model:example3:thecyclohexaneproblemwithchoco}{}
Real variables are illustrated on the problem of finding the 3D configuration of a cyclohexane molecule. 
\lstinputlisting{java/icyclohexane.j2t}



\chapter{The model}\label{doc:model}\hypertarget{doc:model}{}

The {\tt Model}, along with the {\tt Solver}, is one of the two key elements of any Choco program. The Choco {\tt Model} allows to describe a problem in an easy and declarative way: it simply records the variables and the constraints defining the problem.

This section describes the large API provided by Choco to create different types of \hyperlink{model:variables}{variables} and \hyperlink{model:constraints}{constraints}.

%\begin{note}
\textbf{Note that a static import is required to use the Choco API:}
\begin{lstlisting}
  import static choco.Choco.*;
\end{lstlisting}
%It is mandatory in order to compile !
%\end{note}

%\section{How to create a model}\label{model:howtocreateamodel}\hypertarget{model:howtocreateamodel}{}
First of all, a {\tt Model} object is created as follows:
\begin{lstlisting}
Model model = new CPModel();
\end{lstlisting}
In that specific case, a Constraint Programming (CP) {\tt Model} object has been created. 


%%%%%%%%%%%%%%%%%%%%%%%%%%%%%%%%%%%%%%%%%%%%%%%%%%%%%%%%%%%%%%%%%%%%%%%%%%%%%%%%%%%%%%%%%%%%%%%%%%%%%%%%%%%%%%%%%%%%%%%%%%%%%%%%%%%%%%%%%%%%%%%%%%%
%%%%%%%%%%%%%%%%%%%%%%%%%%%%%%%%%%%%%%%%%%%%%%%%%%%%% VARIABLE %%%%%%%%%%%%%%%%%%%%%%%%%%%%%%%%%%%%%%%%%%%%%%%%%%%%%%%%%%%%%%%%%%%%%%%%%%%%%%%%%%%%
%%%%%%%%%%%%%%%%%%%%%%%%%%%%%%%%%%%%%%%%%%%%%%%%%%%%%%%%%%%%%%%%%%%%%%%%%%%%%%%%%%%%%%%%%%%%%%%%%%%%%%%%%%%%%%%%%%%%%%%%%%%%%%%%%%%%%%%%%%%%%%%%%%%


\section{Variables}\label{model:variables}\hypertarget{model:variables}{}

%Choco provides a large API to create different types of variables : \textbf{integer}, \textbf{real} and \textbf{set}. 

A Variable is defined by a type (\hyperlink{integervariable}{integer}, \hyperlink{realvariable}{real}, or \hyperlink{setvariable}{set} variable), a name, and the values of its domain. When creating a simple variable, some options can be set to specify its domain representation (ex: enumerated or bounded) within the {\tt Solver}.
%Some kinds of variables have options for their domain, it may have an effect on what kind of specific object is created when the model is read by the solver.
\begin{note}
The choice of the domain should be considered. The efficiency of the solver often depends on judicious choice of the domain type.
\end{note}
Variables can be combined as \hyperlink{model:expressionvariables}{expression variables} using operators.

One or more variables can be added to the model using the following methods of the \texttt{Model} class:
\lstinputlisting{java/mvariabledeclaration1.j2t}

\begin{note}
Explictly addition of variables is not mandatory. See \hyperlink{model:constraints}{\tt Constraint} for more details.
\end{note}

Specific role of variables \emph{var} can be defined with \emph{options}:  \hyperlink{model:decisionvariables}{non-decision} variables or  \hyperlink{model:objectivevariable}{objective} variable;
\lstinputlisting{java/mvariabledeclaration2.j2t}

%%%%%%%%%%%%%%%%%%%%%%%%%%%%%%%%%%%%%%%%%%%%%%%%%%%%%%%%%%%%%%%%%%%%%%%%%%%%%%%%%%%%%%%%%%%%%%%%%%%%%%%%%%%%%%%%%%%%%%%%%%%%%%%%%%%%%%%%%%%%%%%%%%%
%%%%%%%%%%%%%%%%%%%%%%%%%%%%%%%%%%%%%%%%%%%%%%%%%%%%% SIMPLE VARIABLE  %%%%%%%%%%%%%%%%%%%%%%%%%%%%%%%%%%%%%%%%%%%%%%%%%%%%%%%%%%%%%%%%%%%%%%%%%%%%

\subsection{Simple Variables}\label{model:simplevariables}\hypertarget{model:simplevariables}{}
See Section \hyperlink{ch:vars}{Variables} for details:

\begin{notedef}\tt
\hyperlink{integervariable}{IntegerVariable}, \hyperlink{setvariable}{SetVariable}, \hyperlink{realvariable}{RealVariable}
\end{notedef}

%%%%%%%%%%%%%%%%%%%%%%%%%%%%%%%%%%%%%%%%%%%%%%%%%%%%%%%%%%%%%%%%%%%%%%%%%%%%%%%%%%%%%%%%%%%%%%%%%%%%%%%%%%%%%%%%%%%%%%%%%%%%%%%%%%%%%%%%%%%%%%%%%%%
%%%%%%%%%%%%%%%%%%%%%%%%%%%%%%%%%%%%%%%%%%%%%%%%%%%%% CONSTANT VARIABLE  %%%%%%%%%%%%%%%%%%%%%%%%%%%%%%%%%%%%%%%%%%%%%%%%%%%%%%%%%%%%%%%%%%%%%%%%%%

\subsection{Constants}\label{model:constants}\hypertarget{model:constants}{}
A constant is a variable with a fixed domain. An \hyperlink{integervariable}{\tt IntegerVariable} declared with a unique value is automatically set as constant. A constant declared twice or more is only stored once in a model.

\lstinputlisting{java/mconstant.j2t}

%%%%%%%%%%%%%%%%%%%%%%%%%%%%%%%%%%%%%%%%%%%%%%%%%%%%%%%%%%%%%%%%%%%%%%%%%%%%%%%%%%%%%%%%%%%%%%%%%%%%%%%%%%%%%%%%%%%%%%%%%%%%%%%%%%%%%%%%%%%%%%%%%%%
%%%%%%%%%%%%%%%%%%%%%%%%%%%%%%%%%%%%%%%%%%%%%%%%%%%%% EXPRESSION VARIABLE  %%%%%%%%%%%%%%%%%%%%%%%%%%%%%%%%%%%%%%%%%%%%%%%%%%%%%%%%%%%%%%%%%%%%%%%%
\subsection{Expression variables and operators}\label{model:expressionvariables}\hypertarget{model:expressionvariables}{}
Expression variables represent the result of combinations between variables of the same type made by operators. Two types of expression variables exist : 
\begin{notedef}
\textbf{\tt IntegerExpressionVariable} and \textbf{\tt RealExpressionVariable}.
\end{notedef}
One can define a buffered expression variable to make a constraint easy to read, for example:
\lstinputlisting{java/mexpressionvariable.j2t}

%\section{Operators}\label{model:operators}\hypertarget{model:operators}{}

To construct expressions of variables, simple operators can be used. Each returns a \texttt{ExpressionVariable} object:
\begin{notedef}\tt
\hyperlink{abs:absoperator}{abs}, \hyperlink{cos:cosoperator}{cos}, \hyperlink{disteq:disteqoperator}{distEq}, \hyperlink{distgt:distgtoperator}{distGt}, \hyperlink{distlt:distltoperator}{distLt}, \hyperlink{distneq:distneqoperator}{distNeq}, \hyperlink{div:divoperator}{div}, \hyperlink{ifthenelse:ifthenelseoperator}{ifThenElse}, \hyperlink{max:maxoperator}{max}, \hyperlink{min:minoperator}{min}, \hyperlink{minus:minusoperator}{minus}, \hyperlink{mod:modoperator}{mod}, \hyperlink{mult:multoperator}{mult}, \hyperlink{neg:negoperator}{neg}, \hyperlink{plus:plusoperator}{plus}, \hyperlink{power:poweroperator}{power}, \hyperlink{scalar:scalaroperator}{scalar}, \hyperlink{sin:sinoperator}{sin}, \hyperlink{sum:sumoperator}{sum}.
\end{notedef}
Note that these operators are not considered as constraints: they do not return a \texttt{Constraint} objet but a \texttt{Variable} object.

%%%%%%%%%%%%%%%%%%%%%%%%%%%%%%%%%%%%%%%%%%%%%%%%%%%%%%%%%%%%%%%%%%%%%%%%%%%%%%%%%%%%%%%%%%%%%%%%%%%%%%%%%%%%%%%%%%%%%%%%%%%%%%%%%%%%%%%%%%%%%%%%%%%
%%%%%%%%%%%%%%%%%%%%%%%%%%%%%%%%%%%%%%%%%%%%%%%%%%%%% MULTIPLE VARIABLE  %%%%%%%%%%%%%%%%%%%%%%%%%%%%%%%%%%%%%%%%%%%%%%%%%%%%%%%%%%%%%%%%%%%%%%%%%%

\subsection{MultipleVariables}\label{model:multiplevariables}\hypertarget{model:multiplevariables}{}
These are syntaxic sugar. To make their declaration easier, \hyperlink{tree:treeconstraint}{\tt tree}, \hyperlink{geost:geostconstraint}{\tt geost}, and scheduling constraints allow or require to use multiple variables, like \texttt{TreeParametersObject}, \texttt{GeostObject} or \hyperlink{taskvariable}{\tt TaskVariable}.
See also the code examples for these constraints.

%%%%%%%%%%%%%%%%%%%%%%%%%%%%%%%%%%%%%%%%%%%%%%%%%%%%%%%%%%%%%%%%%%%%%%%%%%%%%%%%%%%%%%%%%%%%%%%%%%%%%%%%%%%%%%%%%%%%%%%%%%%%%%%%%%%%%%%%%%%%%%%%%%%
%%%%%%%%%%%%%%%%%%%%%%%%%%%%%%%%%%%%%%%%%%%%%%%%%%%%% OPTIONS %%%%%%%%%%%%%%%%%%%%%%%%%%%%%%%%%%%%%%%%%%%%%%%%%%%%%%%%%%%%%%%%%%%%%%%%%%%%%%%%%%%%%

\subsection{Decision/non-decision variables}\label{model:decisionvariables}\hypertarget{model:decisionvariables}{}

By default, each variable added to a model is a decision variable, \textit{i.e.} is included in the default search strategy. A variable can be stated as a non decision one if its value can be computed by side-effect. To specify non decision variables, one can 
\begin{itemize}
\item exclude them from its search strategies (see \hyperlink{solver:searchstrategy}{search strategy} for more details);
\item specify non-decision variables (adding \hyperlink{vnodecision:vnodecisionoptions}{\tt Options.V\_NO\_DECISION} to their options) and keep the default search strategy.
\end{itemize}
\lstinputlisting{java/mnodecision1.j2t}
Each of these options can also be set within a single instruction for a group of variables, as follows: 
\lstinputlisting{java/mnodecision2.j2t}

\begin{note}
 The declaration of a \hyperlink{solver:searchstrategy}{search strategy} will erase setting \hyperlink{vnodecision:vnodecisionoptions}{\tt Options.V\_NO\_DECISION}.
\end{note}
  \todo{more precise: user-defined/pre-defined, variable and/or value heuristics ?}

\subsection{Objective variable}\label{model:objectivevariable}\hypertarget{model:objectivevariable}{}
You can define an objective variable directly within the model, by using option \hyperlink{vobjective:vobjectiveoptions}{\tt Options.V\_OBJECTIVE}:
\lstinputlisting{java/mobjective.j2t}

Only one variable can be defined as an objective. If more than one objective variable is declared, then only the last one will be taken into account.

Note that optimization problems can be declared without defining an objective variable within the model (see the \hyperlink{solver:optimization}{optimization example}.)

%%%%%%%%%%%%%%%%%%%%%%%%%%%%%%%%%%%%%%%%%%%%%%%%%%%%%%%%%%%%%%%%%%%%%%%%%%%%%%%%%%%%%%%%%%%%%%%%%%%%%%%%%%%%%%%%%%%%%%%%%%%%%%%%%%%%%%%%%%%%%%%%%%%
%%%%%%%%%%%%%%%%%%%%%%%%%%%%%%%%%%%%%%%%%%%%%%%%%%%%% CONSTRAINT  %%%%%%%%%%%%%%%%%%%%%%%%%%%%%%%%%%%%%%%%%%%%%%%%%%%%%%%%%%%%%%%%%%%%%%%%%%%%%%%%%
%%%%%%%%%%%%%%%%%%%%%%%%%%%%%%%%%%%%%%%%%%%%%%%%%%%%%%%%%%%%%%%%%%%%%%%%%%%%%%%%%%%%%%%%%%%%%%%%%%%%%%%%%%%%%%%%%%%%%%%%%%%%%%%%%%%%%%%%%%%%%%%%%%%

\section{Constraints}\label{model:constraints}\hypertarget{model:constraints}{}
Choco provides a large number of simple and global constraints and allows the user to easily define its own new constraint.
% Either basic, global (a \hyperlink{constraints}{large set of global constraints} are available) or \hyperlink{advanced:defineyourownconstraint}{user-defined} constraints, they are used to specify conditions to be held on variables to the model. 
A constraint deals with one or more variables of the model and specify conditions to be held on these variables. 
A constraint is stated into the model by using the following methods available from the \texttt{Model} API: 

\lstinputlisting{java/mconstraintdeclaration1.j2t}

\begin{note}\
Adding a constraint automatically adds its variables to the model (explicit declaration of variables addition is not mandatory).
\end{note}


\subsubsection{Example:} adding a difference (disequality) constraint between two variables of the model

\lstinputlisting{java/mconstraintdeclaration2.j2t}

Available \emph{options} depend on the kind of constraint \emph{c} to add: they allow, for example, to choose the filtering algorithm to run during propagation. See \hyperlink{optionssettings}{Section options ans settings} for more details, specific APIs exist for declaring options constraints.

This section presents the constraints available in the Choco API sorted by type or by domain. Related sections:
\begin{itemize}
\item a detailed description (with options, examples, references) of each constraint is given in Section \hyperlink{ch:constraints}{constraints}
\item Section \hyperlink{doc:applications}{applications} shows how to apply some specific global constraints
\item Section \hyperlink{advanced:defineyourownconstraint}{user-defined constraint} explains how to create its own constraint.
\end{itemize}

%%%%%%%%%%%%%%%%%%%%%%%%%%%%%%%%%%%%%%%%%%%%%%%%%%%%%%%%%%%%%%%%%%%%%%%%%%%%%%%%%%%%%%%%%%%%%%%%%%%%%%%%%%%%%%%%%%%%%%%%%%%%%%%%%%%%%%%%%%%%%%%%%%%
%%%%%%%%%%%%%%%%%%%%%%%%%%%%%%%%%%%%%%%%%%%%%%%%%%%%% BINARY CONSTRAINT  %%%%%%%%%%%%%%%%%%%%%%%%%%%%%%%%%%%%%%%%%%%%%%%%%%%%%%%%%%%%%%%%%%%%%%%%%%

\subsection{Binary constraints}\label{model:comparisonconstraints}\hypertarget{model:comparisonconstraints}{}
%The simplest constraints are comparisons which are defined over expressions of variables such as linear combinations. The following comparison constraints can be accessed through the \texttt{Model} API:
Constraints involving two integer variables
\begin{notedef}\tt
  \begin{itemize}
  \item \hyperlink{eq:eqconstraint}{eq}, \hyperlink{geq:geqconstraint}{geq}, \hyperlink{gt:gtconstraint}{gt}, \hyperlink{leq:leqconstraint}{leq}, \hyperlink{lt:ltconstraint}{lt}, \hyperlink{neq:neqconstraint}{neq}
  \item \hyperlink{abs:absconstraint}{abs}, \hyperlink{oppositesign:oppositesignconstraint}{oppositeSign}, \hyperlink{samesign:samesignconstraint}{sameSign}
  \end{itemize}
\end{notedef}

%%%%%%%%%%%%%%%%%%%%%%%%%%%%%%%%%%%%%%%%%%%%%%%%%%%%%%%%%%%%%%%%%%%%%%%%%%%%%%%%%%%%%%%%%%%%%%%%%%%%%%%%%%%%%%%%%%%%%%%%%%%%%%%%%%%%%%%%%%%%%%%%%%%
%%%%%%%%%%%%%%%%%%%%%%%%%%%%%%%%%%%%%%%%%%%%%%%%%%%%% TERNARY CONSTRAINT  %%%%%%%%%%%%%%%%%%%%%%%%%%%%%%%%%%%%%%%%%%%%%%%%%%%%%%%%%%%%%%%%%%%%%%%%%

\subsection{Ternary constraints}\label{model:ternaryconstraints}\hypertarget{model:ternaryconstraints}{}
Constraints involving three integer variables
\begin{notedef}\tt
  \begin{itemize}
  \item \hyperlink{distanceeq:distanceeqconstraint}{distanceEQ}, \hyperlink{distanceneq:distanceneqconstraint}{distanceNEQ}, \hyperlink{distancegt:distancegtconstraint}{distanceGT}, \hyperlink{distancelt:distanceltconstraint}{distanceLT}
  \item \hyperlink{intdiv:intdivconstraint}{intDiv}, \hyperlink{mod:modconstraint}{mod}, \hyperlink{times:timesconstraint}{times}
  \end{itemize}
\end{notedef}

%%%%%%%%%%%%%%%%%%%%%%%%%%%%%%%%%%%%%%%%%%%%%%%%%%%%%%%%%%%%%%%%%%%%%%%%%%%%%%%%%%%%%%%%%%%%%%%%%%%%%%%%%%%%%%%%%%%%%%%%%%%%%%%%%%%%%%%%%%%%%%%%%%%
%%%%%%%%%%%%%%%%%%%%%%%%%%%%%%%%%%%%%%%%%%%%%%%%%%%%% REAL CONSTRAINT  %%%%%%%%%%%%%%%%%%%%%%%%%%%%%%%%%%%%%%%%%%%%%%%%%%%%%%%%%%%%%%%%%%%%%%%%%%%%

\subsection{Constraints involving real variables}\label{model:realconstraints}\hypertarget{model:realconstraints}{}
%The simplest constraints are comparisons which are defined over expressions of variables such as linear combinations. The following comparison constraints can be accessed through the \texttt{Model} API:
Constraints involving two real variables
\begin{notedef}\tt
  \begin{itemize}
  \item \hyperlink{eq:eqconstraint}{eq}, \hyperlink{geq:geqconstraint}{geq}, \hyperlink{leq:leqconstraint}{leq}
  \end{itemize}
\end{notedef}

%%%%%%%%%%%%%%%%%%%%%%%%%%%%%%%%%%%%%%%%%%%%%%%%%%%%%%%%%%%%%%%%%%%%%%%%%%%%%%%%%%%%%%%%%%%%%%%%%%%%%%%%%%%%%%%%%%%%%%%%%%%%%%%%%%%%%%%%%%%%%%%%%%%
%%%%%%%%%%%%%%%%%%%%%%%%%%%%%%%%%%%%%%%%%%%%%%%%%%%%% SET CONSTRAINT  %%%%%%%%%%%%%%%%%%%%%%%%%%%%%%%%%%%%%%%%%%%%%%%%%%%%%%%%%%%%%%%%%%%%%%%%%%%%%

\subsection{Constraints involving set variables}\label{model:setconstraints}\hypertarget{model:setconstraints}{}
%The simplest constraints are comparisons which are defined over expressions of variables such as linear combinations. The following comparison constraints can be accessed through the \texttt{Model} API:
Set constraints are illustrated on the \hyperlink{model:example2:ternarysteinerchoco}{ternary Steiner problem}. 
\begin{notedef}\tt
  \begin{itemize}
  \item \hyperlink{eqcard:eqcardconstraint}{eqCard}, \hyperlink{geqcard:geqcardconstraint}{geqCard}, \hyperlink{leqcard:leqcardconstraint}{leqCard}
  \item \hyperlink{member:memberconstraint}{member}, \hyperlink{notmember:notmemberconstraint}{notMember}
  \item \hyperlink{isincluded:isincludedconstraint}{isIncluded}, \hyperlink{isnotincluded:isnotincludedconstraint}{isNotIncluded}, \hyperlink{setdisjoint:setdisjointconstraint}{setDisjoint}
  \item \hyperlink{setinter:setinterconstraint}{setInter}, \hyperlink{setunion:setunionconstraint}{setUnion}
  \item \hyperlink{max:maxofaset}{max}, \hyperlink{min:minofaset}{min}
  \item \hyperlink{pack:packconstraint}{pack}
  \end{itemize}
\end{notedef}

%\hyperlink{max:maxconstraint}{max}, \hyperlink{min:minconstraint}{min},

%%%%%%%%%%%%%%%%%%%%%%%%%%%%%%%%%%%%%%%%%%%%%%%%%%%%%%%%%%%%%%%%%%%%%%%%%%%%%%%%%%%%%%%%%%%%%%%%%%%%%%%%%%%%%%%%%%%%%%%%%%%%%%%%%%%%%%%%%%%%%%%%%%%
%%%%%%%%%%%%%%%%%%%%%%%%%%%%%%%%%%%%%%%%%%%%%%%%%%%%% CHANNELING CONSTRAINT  %%%%%%%%%%%%%%%%%%%%%%%%%%%%%%%%%%%%%%%%%%%%%%%%%%%%%%%%%%%%%%%%%%%%%%

\subsection{Channeling constraints}\label{model:channelingconstraints}\hypertarget{model:channelingconstraints}{}
The use of a redundant model is a frequent technique to strengthen propagation or to get more freedom to design dedicated search heuristics. The following constraints allow to ensure integrity of different models:
\begin{notedef}\tt
  \begin{itemize}
  \item \hyperlink{inversechanneling:inversechannelingconstraint}{inverseChanneling}, \hyperlink{boolchanneling:boolchannelingconstraint}{boolChanneling}, \hyperlink{domainconstraint:domainconstraintconstraint}{domainConstraint}
  \end{itemize}
\end{notedef}
More complex channeling can be done using reified constraints (see Section \hyperlink{model:reifiedconstraints}{reification}) although they are less efficient. For example, to ensure that two variables are equal or not, one can reify the equality into a boolean variables :
\lstinputlisting{java/cchannelingreified.j2t}

%%%%%%%%%%%%%%%%%%%%%%%%%%%%%%%%%%%%%%%%%%%%%%%%%%%%%%%%%%%%%%%%%%%%%%%%%%%%%%%%%%%%%%%%%%%%%%%%%%%%%%%%%%%%%%%%%%%%%%%%%%%%%%%%%%%%%%%%%%%%%%%%%%%
%%%%%%%%%%%%%%%%%%%%%%%%%%%%%%%%%%%%%%%%%%%%%%%%%%%%% EXTENSIONS CONSTRAINT  %%%%%%%%%%%%%%%%%%%%%%%%%%%%%%%%%%%%%%%%%%%%%%%%%%%%%%%%%%%%%%%%%%%%%%

\subsection{Constraints in extension and relations}\label{model:arbitraryconstraintsinextension}\hypertarget{model:arbitraryconstraintsinextension}{}
Choco supports the statement of constraints defining arbitrary relations over two or more variables.
Such a relation may be defined by three means:
\begin{itemize}
	\item \textbf{feasible table:} the list of allowed tuples of values (that belong to the relation),
	\item \textbf{infeasible table:} the list of forbidden tuples of values (that do not belong to the relation),
	\item \textbf{predicate:} a method to be called in order to check whether a tuple of values belongs or not to the relation.
\end{itemize}
On the one hand, constraints based on tables may be rather memory consuming in case of large domains, although one relation table may be shared by several constraints. On the other hand, predicate constraints require little memory as they do not cache truth values, but imply some run-time overhead for calling the feasibility test. Table constraints are thus well suited for constraints over small domains; while predicate constraints are well suited for situations with large domains. 

Different levels of consistency can be enforce on constraints in extension: 
\begin{itemize}
\item several arc-consistency (AC) algorithms for binary relations
\item two AC algorithms for n-ary relations dedicated either to positive or to negative tables (relation defined by the allowed or forbidden tuples)
\item a weaker forward-checking (FC) algorithm for n-ary relations.
\end{itemize}

The Choco API for creating constraints in extension are as follows:
\begin{notedef}\tt
  \begin{itemize}
  \item \hyperlink{feaspairac:feaspairacconstraint}{feasPairAC}, \hyperlink{infeaspairac:infeaspairacconstraint}{infeasPairAC}, \hyperlink{relationpairac:relationpairacconstraint}{relationPairAC}
  \item \hyperlink{feastupleac:feastupleacconstraint}{feasTupleAC}, \hyperlink{infeastupleac:infeastupleacconstraint}{infeasTupleAC}, \hyperlink{relationtupleac:relationtupleacconstraint}{relationTupleAC}
  \item \hyperlink{feastuplefc:feastuplefcconstraint}{feasTupleFC}, \hyperlink{infeastuplefc:infeastuplefcconstraint}{infeasTupleFC}, \hyperlink{relationtuplefc:relationtuplefcconstraint}{relationTupleFC}
  \end{itemize}
\end{notedef}

\subsubsection{Relations.}
A same relation might be shared among several constraints, in this case it is highly recommended to create it first and then use the \hyperlink{relationpairac:relationpairacconstraint}{relationPairAC}, \hyperlink{relationtupleac:relationtupleacconstraint}{relationTupleAC}, or \hyperlink{relationtuplefc:relationtuplefcconstraint}{relationTupleFC} API  on the same relation for each constraint.

For binary relations, the following Choco API is provided:
\mylst{makeBinRelation(int[] min, int[] max, List<int[]>pairs, boolean feas)}

It returns a \texttt{BinRelation} giving a list of compatible (\texttt{feas=true}) or incompatible (\texttt{feas=false}) pairs of values. This relation can be applied to any pair of variables $(x_1,x_2)$ whose domains are included in the \texttt{min/max} intervals, i.e. such that:
$$\mathtt{min}[i] \le x_i.\mathtt{getInf}() \le x_i.\mathtt{getSup}() \le  \mathtt{max}[i],\quad \forall i.$$
Bounds \texttt{min/max} are mandatory in order to allow to compute the opposite of the relation if needed.

For n-ary relations, the corresponding Choco API is:
\mylst{makeLargeRelation(int[] min, int[] max, List<int[]> tuples, boolean feas);}
It returns a \texttt{LargeRelation}. If \texttt{feas=true}, the returned relation matches also the \texttt{IterLargeRelation} interface which provides constant time iteration abilities over tuples (for compatibility with the GAC algorithm used over feasible tuples).
\lstinputlisting{java/mlargerelation.j2t}

Lastly, some specific relations can be defined without storing the tuples, as in the following example (\texttt{TuplesTest} extends \texttt{LargeRelation}):
\lstinputlisting{java/mnotallequal.j2t}
Then, a \emph{NotAllEqual} constraint can be stated within the problem by:
\lstinputlisting{java/mrelationtuplefc.j2t}
%Again, for compatibility with the GAC algorithm invoked by relationTupleAC, such a relation has to match the \texttt{IterLargeRelation} interface for feasible tuples.


%%%%%%%%%%%%%%%%%%%%%%%%%%%%%%%%%%%%%%%%%%%%%%%%%%%%%%%%%%%%%%%%%%%%%%%%%%%%%%%%%%%%%%%%%%%%%%%%%%%%%%%%%%%%%%%%%%%%%%%%%%%%%%%%%%%%%%%%%%%%%%%%%%%
%%%%%%%%%%%%%%%%%%%%%%%%%%%%%%%%%%%%%%%%%%%%%%%%%%%%% REIFIED CONSTRAINT  %%%%%%%%%%%%%%%%%%%%%%%%%%%%%%%%%%%%%%%%%%%%%%%%%%%%%%%%%%%%%%%%%%%%%%%%%

\subsection{Reified constraints}\label{model:reifiedconstraints}\hypertarget{model:reifiedconstraints}{}
Constraints involved in another constraint are usually called reified constraints. Typical examples of reified constraints are
 constraints combined with logical operators, such as $(x \neq y) \lor (z \le 9)$.

%\subsubsection{To reify a constraint into a boolean variable.}\label{model:toreifyaconstraintintoabooleanvariable}\hypertarget{model:toreifyaconstraintintoabooleanvariable}{}
Choco provides a generic constraint to reify any constraints on integer variables or set variables into a boolean variable expressing its truth value:
\begin{notedef}\tt
  \begin{itemize}
  \item \hyperlink{reifiedconstraint:reifiedconstraintconstraint}{reifiedConstraint}, \hyperlink{reifiedand:reifiedandconstraint}{reifiedAnd}, \hyperlink{reifiedleftimp:reifiedleftimpconstraint}{reifiedLeftImp}, \hyperlink{reifiednot:reifiednotconstraint}{reifiedNot}, \hyperlink{reifiedor:reifiedorconstraint}{reifiedOr}, \hyperlink{reifiedrightimp:reifiedrightimpconstraint}{reifiedRightImp}, \hyperlink{reifiedxnor:reifiedxnorconstraint}{reifiedXnor}, \hyperlink{reifiedxor:reifiedxorconstraint}{reifiedXor}
  \end{itemize}
\end{notedef}
This mechanism can be used for example to model MaxCSP problems where the number of satisfied constraints has to be maximized.
It is also intended to give the freedom to the user to build complex reified constraints. However, Choco provides a more simple and direct API to build complex expressions using boolean operators:
\begin{notedef}\tt
  \begin{itemize}
  \item \hyperlink{and:andconstraint}{and}, \hyperlink{or:orconstraint}{or}, \hyperlink{implies:impliesconstraint}{implies}, \hyperlink{ifonlyif:ifonlyifconstraint}{ifOnlyIf}, \hyperlink{ifthenelse:ifthenelseconstraint}{ifThenElse}, \hyperlink{not:notconstraint}{not}
  \end{itemize}
\end{notedef}
Such an expression is represented as a tree of operators. The leaves of this tree are made of variables, constants or even traditional constraints. Variables and constants can be combined as \texttt{ExpressionVariable} using \hyperlink{model:expressionvariables}{operators} (e.g, \texttt{mult(10,abs(w))}), or using simple constraints (e.g., \texttt{leq(z,9)}), or even using global constraints (e.g, \texttt{alldifferent(vars)}).
The language available on expressions is therefore slightly richer and matches the language used in the \href{http://cpai.ucc.ie/08/}{Constraint Solver Competition 2008} of the CPAI workshop.

For example, the following expression
$$((x = 10 * |y|) \lor (z \le 9))\quad \iff\quad \texttt{alldifferent}(a,b,c)$$
could be represented by :
\begin{lstlisting}
	Constraint exp = ifOnlyIf( or( eq(x, mult(10, abs(y))), leq(z, 9) ), 
                               alldifferent(new IntegerVariable[]{a,b,c}) );
\end{lstlisting}


\subsubsection{Handling complex expressions.}\label{model:handlingcomplexexpressions}\hypertarget{model:handlingcomplexexpressions}{}
Expressions offer a more powerful modeling language than the one available via standard constraints. However, they 
can not be handled as efficiently as the standard constraints that embed a dedicated propagation algorithm. We therefore
recommend you to carefully check that you can not model the expression using the intensional constraints of Choco before using
expressions.
Inside the solver, expressions can be represented in two different ways that can be decided at the modeling level, using the following {\tt Model} API:
\begin{lstlisting}
  setDefaultExpressionDecomposition(boolean decomp);
\end{lstlisting}
or the option \hyperlink{edecomp:edecompoptions}{\tt Options.E\_DECOMP}.
\begin{itemize}
\item The first way (\texttt{decomp=false}) is to handle them as \hyperlink{model:arbitraryconstraintsinextension}{constraints in extension}. The expression is then used to check a tuple in a dynamic way just like a n-ary relation that is defined without listing all the possible tuples. The expression is then propagated using the GAC3rm algorithm. This is very powerful as arc-consistency is obtained on the corresponding constraints.
\item The second way (\texttt{decomp=true}) is to decompose the expression automatically by introducing intermediate variables and eventually the generic \hyperlink{reifiedintconstraint:reifiedintconstraintconstraint}{\tt reifiedIntConstraint}. By doing so, the level of pruning decreases but expressions of larger arity involving large domains can be represented.
\end{itemize}

%\subsubsection{Tell the solver how to consider an expression.}
%The default representation of expressions can be enforced using the following API  on the model object: 
%Parameter \emph{decomp} tells the solver whether expressions shoud be considered as extensional constraints (\texttt{decomp=false}) or decomposed 
Once the default representation is chosen, one can also make exception for a particular expression using options on \texttt{addConstraint}. 
For example, the following code tells the solver to decompose e1 and not e2 :
\begin{lstlisting}
	model.setDefaultExpressionDecomposition(false);
	IntegerVariable x = makeIntVar("x", 1, 3, Options.V_BOUND);
	IntegerVariable y = makeIntVar("y", 1, 3, Options.V_BOUND);
	IntegerVariable z = makeIntVar("z", 1, 3, Options.V_BOUND);

	Constraint e1 = or(lt(x, y), lt(y, x));
	model.addConstraint(Options.E_DECOMP, e1);
	
	Constraint e2 = or(lt(y, z), lt(z, y));
	model.addConstraint(e2);
\end{lstlisting}

\subsubsection{When and how should I use expressions ?}\label{model:whenshouldiuseexpressions}\hypertarget{model:whenshouldiuseexpressions}{}
An expression (represented in extension) should be used in the case of a complex logical relationship that involves \textbf{few different variables}, each of \textbf{small domain}, and if \textbf{arc consistency} is desired on those variables.
In such a case, an expression can even be more powerful than a model using intermediate variables and intensional constraints.
Imagine the following ``crazy'' example :
\begin{lstlisting}
 or( and( eq( abs(sub(div(x,50),div(y,50))),1), eq( abs(sub(mod(x,50),mod(y,50))),2)),
     and( eq( abs(sub(div(x,50),div(y,50))),2), eq( abs(sub(mod(x,50),mod(y,50))),1)))
\end{lstlisting}
This expression has a small arity: it involves only two variables $x$ and $y$.
Let assume that their domains has no more than 300 values, then such an expression should typically not be decomposed. Indeed, arc consistency will create many holes in the domains and filter much more than if the relation was decomposed.

Conversely, an expression should be decomposed as soon as it involves a large number of variables, or at least one variable with a large domain.

%%%%%%%%%%%%%%%%%%%%%%%%%%%%%%%%%%%%%%%%%%%%%%%%%%%%%%%%%%%%%%%%%%%%%%%%%%%%%%%%%%%%%%%%%%%%%%%%%%%%%%%%%%%%%%%%%%%%%%%%%%%%%%%%%%%%%%%%%%%%%%%%%%%
%%%%%%%%%%%%%%%%%%%%%%%%%%%%%%%%%%%%%%%%%%%%%%%%%%%%% GLOBAL CONSTRAINT  %%%%%%%%%%%%%%%%%%%%%%%%%%%%%%%%%%%%%%%%%%%%%%%%%%%%%%%%%%%%%%%%%%%%%%%%%%

\subsection{Global constraints}\label{model:advancedconstraints}\hypertarget{model:advancedconstraints}{}
Choco includes several \href{http://www.emn.fr/x-info/sdemasse/gccat/}{global constraints}. Those constraints accept any number of variables and offer dedicated filtering algorithms which are able to make deductions where a decomposed model would not.
For instance, constraint \texttt{alldifferent}$(a,b,c,d)$ with $a,b\in[1,4]$ and $c,d\in[3,4]$ allows to deduce that $a$ and $b$ cannot be instantiated to $3$ or $4$; such rule cannot be inferred by simple binary constraints. 

The up-to-date list of global constraints available in Choco can be found within the Javadoc API.
Most of these global constraints are listed below according to their application fields.
Details and examples can be found in Section \hyperlink{ch:constraints}{Elements of Choco/Constraints}.
\subsubsection{Value constraints}\label{model:valueconstraints}\hypertarget{model:valueconstraints}{}
Constraints that put a restriction on how values can be assigned to usually one or several collections of variables.
See also in Global Constraint Catalog: \href{http://www.emn.fr/x-info/sdemasse/gccat/Kvalue_constraint.html}{value constraint}.

\vspace{1em}\noindent\begin{notedef}\tt
  \begin{itemize}
  \item counting distinct values: 
\hyperlink{alldifferent:alldifferentconstraint}{allDifferent}, 
\hyperlink{atmostnvalue:atmostnvalueconstraint}{atMostNValue},
\hyperlink{increasingnvalue:increasingnvalueconstraint}{increasingnvalue},
  \item counting values: 
\hyperlink{occurrence:occurrenceconstraint}{occurrence},
\hyperlink{occurrencemax:occurrencemaxconstraint}{occurrenceMax},
\hyperlink{occurrencemin:occurrenceminconstraint}{occurrenceMin},
\hyperlink{globalcardinality:globalcardinalityconstraint}{globalCardinality},
  \item indexing values: 
\hyperlink{nth:nthconstraint}{nth} (element),
\hyperlink{max:maxconstraint}{max},
\hyperlink{min:minconstraint}{min},
  \item ordering: 
\hyperlink{sorting:sortingconstraint}{sorting},
\hyperlink{increasingnvalue:increasingnvalueconstraint}{increasingnvalue},
\hyperlink{lex:lexconstraint}{lex}, 
\hyperlink{lexeq:lexeqconstraint}{lexeq},
\hyperlink{leximin:leximinconstraint}{leximin},
\hyperlink{lexchain:lexchainconstraint}{lexChain},
\hyperlink{lexchaineq:lexchaineqconstraint}{lexChainEq},
  \item tuple matching: 
\hyperlink{feastupleac:feastupleacconstraint}{feasTupleAC},
\hyperlink{feastuplefc:feastuplefcconstraint}{feasTupleFC},
\hyperlink{infeastupleac:infeastupleacconstraint}{infeasTupleAC},
\hyperlink{infeastuplefc:infeastuplefcconstraint}{infeasTupleFC},
\hyperlink{relationtupleac:relationtupleacconstraint}{relationTupleAC},
\hyperlink{relationtuplefc:relationtuplefcconstraint}{relationTupleFC},
  \item pattern matching: 
\hyperlink{regular:regularconstraint}{regular},
\hyperlink{costregular:costregularconstraint}{costRegular},
\hyperlink{multicostregular:multicostregularconstraint}{multiCostRegular}, 
\hyperlink{stretchcyclic:stretchcyclicconstraint}{stretchCyclic}, 
\hyperlink{stretchpath:stretchpathconstraint}{stretchPath}, 
\hyperlink{tree:treeconstraint}{tree},
  \end{itemize}
\end{notedef}

\subsubsection{Boolean constraints}\label{model:logicconstraints}\hypertarget{model:logicconstraints}{}
Logical operations on boolean expressions.
See also in Global Constraint Catalog: \href{http://www.emn.fr/x-info/sdemasse/gccat/KBoolean_constraint.html}{boolean constraint}.

\vspace{1em}\noindent\begin{notedef}\tt
\hyperlink{and:andconstraint}{and},
\hyperlink{or:orconstraint}{or},
\hyperlink{clause:clauseconstraint}{clause},
\end{notedef}

\subsubsection{Channelling constraints}\label{model:channellingconstraints}\hypertarget{model:channellingconstraints}{}
Constraints linking two collections of variables (many-to-many) or indexing one among many variables (one-to-many).
See also in Global Constraint Catalog: \href{http://www.emn.fr/x-info/sdemasse/gccat/Kchannelling_constraint.html}{channelling constraint}.

 \vspace{1em}\noindent\begin{notedef}\tt
   \begin{itemize}
   \item one-to-many: 
 \hyperlink{domainconstraint:domainconstraintconstraint}{domainConstraint},
 \hyperlink{nth:nthconstraint}{nth} (element),
 \hyperlink{max:maxconstraint}{max},
 \hyperlink{min:minconstraint}{min},
   \item many-to-many: 
 \hyperlink{inversechanneling:inversechannelingconstraint}{inverseChanneling},
 \hyperlink{inverseset:inversesetconstraint}{inverseset},
 \hyperlink{sorting:sortingconstraint}{sorting},
 \end{itemize}
 \end{notedef}

\subsubsection{Optimization constraints}\label{model:optimizationconstraints}\hypertarget{model:optimizationconstraints}{}
Constraints channelling a variable to the sum of the weights of a collection of variable-value assignments.
See also in Global Constraint Catalog: \href{http://www.emn.fr/x-info/sdemasse/gccat/Kcost_filtering_constraint.html}{cost-filtering constraint}.
\vspace{1em}\noindent\begin{notedef}\tt
 \begin{itemize}
  \item one cost: 
\hyperlink{occurrence:occurrenceconstraint}{occurrence},
\hyperlink{occurrencemax:occurrencemaxconstraint}{occurrenceMax},
\hyperlink{occurrencemin:occurrenceminconstraint}{occurrenceMin},
\hyperlink{knapsackproblem:knapsackproblemconstraint}{knapsackProblem},
\hyperlink{equation:equationconstraint}{equation},
\hyperlink{costregular:costregularconstraint}{costRegular},
\hyperlink{tree:treeconstraint}{tree},
 \item several costs:
\hyperlink{globalcardinality:globalcardinalityconstraint}{globalCardinality},
\hyperlink{multicostregular:multicostregularconstraint}{multiCostRegular}, 
 \end{itemize}
\end{notedef}

\subsubsection{Packing constraints (capacitated resources)}\label{model:packingconstraints}\hypertarget{model:packingconstraints}{}
Constraints involving items to be packed in bins without overlapping. More generaly, any constraints modelling the concurrent assignment of objects to one or several capacitated resources.
See also in Global Constraint Catalog: \href{http://www.emn.fr/x-info/sdemasse/gccat/Kresource_constraint.html}{resource constraint}.

\vspace{1em}\noindent\begin{notedef}\tt
   \begin{itemize}
   \item packing problems: 
\hyperlink{equation:equationconstraint}{equation},
\hyperlink{knapsackproblem:knapsackproblemconstraint}{knapsackProblem},
\hyperlink{pack:packconstraint}{pack} (bin-packing),
   \item geometric placement problems: 
\hyperlink{geost:geostconstraint}{geost}, 
   \item scheduling problems: 
\hyperlink{disjunctive:disjunctiveconstraint}{disjunctive}, 
\hyperlink{cumulative:cumulativeconstraint}{cumulative}, 
   \item timetabling problems: 
\hyperlink{costregular:costregularconstraint}{costRegular},
\hyperlink{multicostregular:multicostregularconstraint}{multiCostRegular}, 
 \end{itemize}
\end{notedef}

\subsubsection{Scheduling constraints (time assignment)}\label{model:schedulingconstraints}\hypertarget{model:schedulingconstraints}{}
Constraints involving tasks to be scheduled over a time horizon.
See also \hyperlink{schedulinganduseofthecumulative:schedulinganduseofthecumulativeconstraint}{scheduling application} and in Global Constraint Catalog: \href{http://www.emn.fr/x-info/sdemasse/gccat/Kscheduling_constraint.html}{scheduling constraint}.

\vspace{1em}\noindent\begin{notedef}\tt
   \begin{itemize}
   \item temporal constraints: 
\hyperlink{preceding:precedingconstraint}{preceding}, 
\hyperlink{precedencedisjoint:precedencedisjointconstraint}{precedenceDisjoint}, 
\hyperlink{precedenceimplied:precedenceimpliedconstraint}{precedenceImplied}, 
\hyperlink{precedencereified:precedencereifiedconstraint}{precedenceReified},
\hyperlink{forbiddeninterval:forbiddenintervalconstraint}{forbiddenInterval},
\hyperlink{tree:treeconstraint}{tree},
   \item resource constraints: 
\hyperlink{cumulative:cumulativeconstraint}{cumulative}, 
\hyperlink{disjunctive:disjunctiveconstraint}{disjunctive}, 
\hyperlink{geost:geostconstraint}{geost}, 
 \end{itemize}
\end{notedef}

%\part{solver}
\label{solver}
\hypertarget{solver}{}


\chapter{The solver}\label{solver:thesolver}\hypertarget{solver:thesolver}{}

%\section{How to create a solver}\label{solver:howtocreateasolver}\hypertarget{solver:howtocreateasolver}{}

To create a {\tt Solver}, one just needs to create a new object as follow:
\begin{lstlisting}
Solver solver = new CPSolver();
\end{lstlisting}
By this, a Constraint Programming (CP) {\tt Solver} object is created. 

%\section{Read a model}\label{solver:readamodel}\hypertarget{solver:readamodel}{}
The solver gives an API to read a model. The reading of a model is compulsory and must be done after the entire definition of the model. 
\begin{lstlisting}
solver.read(model);
\end{lstlisting}
The reading is divided in 2 parts: \hyperlink{solver:variablesreading}{variables reading} and \hyperlink{solver:constraintsreading}{constraints reading}.

\section{Variables reading}\label{solver:variablesreading}\hypertarget{solver:variablesreading}{}
The solver iterates over the variables of the Model to create solver-specific variables and domains (as defined in the model). 
Thus, three types of variables can be created: integer variables, real variables and set variables. 
Depending on the constructor, the correct domain is created (like bounded domain or enumerated domain for integer variables). 

\begin{note}
\textbf{Bound variables} are related to large domains which are only represented by their lower and upper bounds. The domain is encoded in a space efficient way and propagation events only concern bound updates. Value removals between the bounds are therefore ignored (\emph{holes} are not considered). The level of consistency achieved by most constraints on these variables is called \emph{bound-consistency}.

On the contrary, the domain of an \textbf{enumerated variable} is explicitly represented and every value is considered while pruning. Basic constraints are therefore often able to achieve \emph{arc-consistency} on enumerated variables (except for NP global constraint such as the cumulative constraint). Remember that switching from an enumerated variable to a bounded variables decrease the level of propagation achieved by the system.
\end{note}

%\begin{note}
Model variables and Solver variables are distinct. Solver variables are solver representation of the model variables. One can't access to variable value directly from the model variable. To access to a model variable thanks to the solver, use the following \texttt{Solver} API: \mylst{getVar(Variable v);}
%\end{note}

\subsection{Solver and IntegerVariables}\label{solver:solverandintegervariables}\hypertarget{solver:solverandintegervariables}{}

A model integer variable can be accessed by the method \textbf{\tt getVar(IntegerVariable v)} which returns a \textbf{\tt IntDomainVar} object:
\begin{lstlisting}
  IntegerVariable x = makeEnumIntVar("x", 1, 100);  // model variable
  IntDomainVar xOnSolver = solver.getVar(x);  // solver variable
\end{lstlisting}

The state of an \texttt{IntDomainVar} can be accessed through the main following public methods :

\noindent\begin{tabular}{p{.3\linewidth}p{.7\linewidth}}
  \hline
  \texttt{IntDomainVar} API &  description \\
  \hline
	\mylst{hasEnumeratedDomain()} &checks if the variable is an enumerated or a bound one\\
	\mylst{getInf()} &returns the lower bound of the variable\\
	\mylst{getSup()} &returns the upper bound of the variable\\
	\mylst{getVal()} &returns the value if it is instantiated\\
	\mylst{isInstantiated()} &checks if the domain is reduced to a singleton\\
	\mylst{canBeInstantiatedTo(int v)} &checks if the value \emph{v} is contained in the domain of the variable\\
	\mylst{getDomainSize()} &returns the current size of the domain\\
  \hline\\
\end{tabular}

For more informations on advanced uses of such \texttt{IntDomainVar}, see \hyperlink{advanced}{advanced uses}.

\subsection{Solver and SetVariables}\label{solver:solverandsetvariables}\hypertarget{solver:solverandsetvariables}{}

A model set variable can be access by the method \textbf{\tt getVar(SetVariable v)} which returns a \textbf{\tt SetVar} object:
\begin{lstlisting}
	SetVariable x = makeBoundSetVar("x", 1, 40); // model variable
	SetVar xOnSolver = solver.getVar(x); // solver variable
\end{lstlisting}
A set variable on integer values between $[1,n]$ has $2^{n}$ values (every possible subsets of $\{1..n\}$). This makes an exponential number of values and the domain is represented with two bounds corresponding to the intersection of all possible sets (called the kernel) and the union of all possible sets (called the envelope) which are the possible candidate values for the variable.

The state of a \texttt{SetVar} can be accessed through the main following public methods on the SetVar class:

\noindent\begin{tabular}{p{.3\linewidth}p{.7\linewidth}}
  \hline
  \texttt{SetVar} API &  description \\
  \hline
	\mylst{getCard()} &returns the \texttt{IntDomainVar} representing the cardinality of the set variable\\
	\mylst{isInDomainKernel(int v)} &checks if value \emph{v} is contained in the current kernel\\
	\mylst{isInDomainEnveloppe(int v)} &checks if value \emph{v} is contained in the current envelope\\
	\mylst{getDomain()} &returns the domain of the variable as a \texttt{SetDomain}. Iterators on envelope or kernel can than be called\\
	\mylst{getKernelDomainSize()} &returns the size of the kernel\\
	\mylst{getEnveloppeDomainSize()} &returns the size of the envelope\\
	\mylst{getEnveloppeInf()} &returns the first available value of the envelope\\
	\mylst{getEnveloppeSup()} &returns the last available value of the envelope\\
	\mylst{getKernelInf()} &returns the first available value of the kernel\\
	\mylst{getKernelSup()} &returns the last available value of the kernel\\
	\mylst{getValue()} &returns a table of integers \texttt{int[]} containing the current domain\\
  \hline\\
\end{tabular}


For more informations on advanced uses of such \texttt{SetVar}, see \hyperlink{advanced}{advanced uses}.

\subsection{Solver and RealVariables}\label{solver:solverandrealvariables}\hypertarget{solver:solverandrealvariables}{}

\begin{note}
\emph{Real variables are still under development but can be used to solve toy problems such as small systems of equations.}
\end{note}
 
A model real variable can be access by the method \textbf{\tt getVar(RealVariable v)} which returns a \texttt{RealVar} object:
\begin{lstlisting}
	RealVariable x = makeRealVar("x", 1.0, 3.0); // model variable
	RealVar xOnSolver = s.getVar(x); // solver variable
\end{lstlisting}

Continuous variables are useful for non linear equation systems which are encountered in physics for example.

\noindent\begin{tabular}{p{.3\linewidth}p{.7\linewidth}}
  \hline
  \texttt{RealVar} API &  description \\
  \hline
	\mylst{getInf()} &returns the lower bound of the variable (\texttt{double})\\
	\mylst{getSup()} &returns the upper bound of the variable (\texttt{double})\\
	\mylst{isInstantiated()} &checks if the domain of a variable is reduced to a canonical interval. A canonical interval indicates that the domain has reached the precision given by the user or the solver\\
  \hline\\
\end{tabular}


For more informations on advanced uses of such \texttt{RealVar}, see \hyperlink{advanced}{advanced uses}.

\section{Constraints reading}\label{solver:constraintsreading}\hypertarget{solver:constraintsreading}{}
After variables, the Solver iterates over the constraints added to the Model. It creates Solver constraints that encapsulates a filtering algorithm which are called when a propagation step occur or when external events happen on the variables belonging to the constraint, such as value removals or bounds modifications. And it add it to the constraint network. 

\section{Search Strategy}\label{solver:searchstrategy}\hypertarget{solver:searchstrategy}{}

A key ingredient of any constraint approach is a clever branching strategy. The construction of the search tree is done according to a series of \textit{branching objects} (that plays the role of achieving intermediate goals in logic programming). The user may specify the sequence of branching objects to be used to build the search tree. A common way to branch in CP is by assigning variables to values. We will present in this section how to define your branching strategies with existing variables and values selectors/iterators. But first...

\subsection{Override the default search stragtegy}\label{solver:overridethedefaultsearchstrategy}\hypertarget{solver:overridethedefaultsearchstrategy}{}


Basically, a search strategy is the composition of three objects: a branching strategy, a variable selector and a value selector. Some branching simply assign a selected value to a selected variable, like \hyperlink{assignvar:assignvarbranchstrat}{AssignVar}, others branching strategies embed the variable selector, like \hyperlink{domoverwdeg:domoverwdegbranchstrat}{DomOverWDegBranchingNew}, or more, like  \hyperlink{impact:impactbranchstrat}{ImpactBasedBranching}.

The default branchings are: 

\noindent\begin{tabular}{p{.4\linewidth}p{.6\linewidth}}
\hline
Variable &  Default strategy \\
\hline
Integer & \hyperlink{domoverwdeg:domoverwdegbranchstrat}{DomOverWDegBranchingNew} +\hyperlink{increasingdomain:increasingdomainvaliterator}{IncreasingDomain}\\
Set &   \hyperlink{assignsetvar:assignsetvarbranchstrat}{AssignSetVar} + \hyperlink{mindomset:mindomsetvarselector}{MinDomainSet} + \hyperlink{minenv:minenvvalselector}{MinEnv} \\
 Real &  \hyperlink{assigninterval:assignintervalbranchstrat}{AssignInterval} + \hyperlink{cyclicrealvarselector:cyclicrealvarselectorvarselector}{CyclicRealVarSelector}+ \hyperlink{realincreasingdomain:realincreasingdomainvaliterator}{RealIncreasingDomain} \\
\hline\\
\end{tabular}

There are two ways to custom a search strategy: use the ones define in the factory  \texttt{BranchingFactory} or compose it. 
A branching strategy can be set or add to the previous one using the following API (must be done before calling the method \mylst{solve()}):

  \mylst{solver.addGoal(AbstractIntBranchingStrategy branching)} 

\noindent to clear the list of goals, use:

  \mylst{solver.clearGoals()} 

You might want to apply different heuristics to different set of variables of the problem. In that case, the search is viewed as a sequence of branching objects (or goals). Up to now, we only had one branching or one goal including all the variables of the problem but several goals can be used.

Adding a new goal is made through the solver with the \mylst{solver.addGoal(AbstractIntBranchingStrategy branching)} method. 

The following example add three branching objects on integer variables \emph{vars1}, \emph{vars2} and set variables \emph{svars} to solver \emph{s}. The first two branchings are both \texttt{AssignVar} but use two different variable/values selection strategies:
\begin{lstlisting}
  s.addGoal(new AssignVar(new MinDomain(s,s.getVar(vars1)), new IncreasingDomain()));
  s.addGoal(new AssignVar(new DomOverDeg(s,s.getVar(vars2)),new DecreasingDomain());
  s.addGoal(new AssignSetVar(new MinDomSet(s,s.getVar(svars)), new MinEnv(s)));
  s.solve();
\end{lstlisting}

\begin{note}
Strategies are made of \texttt{Solver} variables (not \texttt{Model} variables).
\end{note}

\subsubsection{Branching strategy.}\label{solver:branchstrat}\hypertarget{solver:branchstrat}{}
It defines the way to take a decision in a tree search node.
  
\noindent The \textbf{branching strategies} currently available in Choco are the following: 
\begin{notedef}\tt
\hyperlink{assigninterval:assignintervalbranchstrat}{AssignInterval}, \hyperlink{assignorforbidintvarval:assignorforbidintvarvalbranchstrat}{AssignOrForbidIntVarVal}, \hyperlink{assignorforbidintvarvalpair:assignorforbidintvarvalpairbranchstrat}{AssignOrForbidIntVarValPair}, \hyperlink{assignsetvar:assignsetvarbranchstrat}{AssignSetVar}, \hyperlink{assignvar:assignvarbranchstrat}{AssignVar}, \hyperlink{domoverwdeg:domoverwdegbranchstrat}{DomOverWDegBranchingNew}, \hyperlink{domoverwdegbin:domoverwdegbinbranchstrat}{DomOverWDegBinBranchingNew}, \hyperlink{impact:impactbranchstrat}{ImpactBasedBranching}, \hyperlink{packdynremovals:packdynremovalsbranchstrat}{PackDynRemovals}, \hyperlink{settimes:settimesbranchstrat}{SetTimes}, \hyperlink{taskdomoverwdeg:taskdomoverwdegbranchstrat}{TaskOverWDegBinBranching}.
\end{notedef}    


\subsubsection{Variable selector.}\label{solver:variableselector}\hypertarget{solver:variableselector}{}
It defines the way to choose the non instantiated variable on which the next decision will be made.

\noindent The \textbf{integer variable selectors} currently available in Choco are the following: 
\begin{notedef}\tt
\hyperlink{compositeintvarselector:compositeintvarselectorvarselector}{CompositeIntVarSelector}, \hyperlink{lexintvarselector:lexintvarselectorvarselector}{LexIntVarSelector}, \hyperlink{maxdomain:maxdomainvarselector}{MaxDomain}, \hyperlink{maxregret:maxregretvarselector}{MaxRegret}, \hyperlink{maxvaldomain:maxvaldomainvarselector}{MaxValueDomain}, \hyperlink{mindomain:mindomainvarselector}{MinDomain}, \hyperlink{minvaldomain:minvaldomainvarselector}{MinValueDomain}, \hyperlink{mostconstrained:mostconstrainedvarselector}{MostConstrained},  \hyperlink{randomvarint:randomvarintvarselector}{RandomIntVarSelector},  \hyperlink{staticvarorder:staticvarordervarselector}{StaticVarOrder}
\end{notedef}

\noindent The \textbf{set variable selectors} currently available in Choco are the following: 
\begin{notedef}\tt
\hyperlink{maxdomset:maxdomsetvarselector}{MaxDomainSet}, \hyperlink{maxregretset:maxregretsetvarselector}{MaxRegretSet}, \hyperlink{maxvaldomainset:maxvaldomainsetvarselector}{MaxValueDomainSet}, \hyperlink{mindomset:mindomsetvarselector}{MinDomainSet}, \hyperlink{minvaldomainset:minvaldomainsetvarselector}{MinValueDomainSet}, \hyperlink{mostconstrainedset:mostconstrainedsetvarselector}{MostConstrainedSet},  \hyperlink{randomvarset:randomvarsetvarselector}{RandomSetVarSelector},  \hyperlink{staticsetvarorder:staticsetvarordervarselector}{StaticSetVarOrder}
\end{notedef}

\noindent The \textbf{real variable selector} currently available in Choco is the following: 
\begin{notedef}\tt
\hyperlink{cyclicrealvarselector:cyclicrealvarselectorvarselector}{CyclicRealVarSelector}
\end{notedef}

\subsubsection{Value iterator}\label{solver:valueiterator}\hypertarget{solver:valueiterator}{}
Once the variable has been choosen, the Solver has to compute its value. The first way to do it is to schedule the value once and give an iterator to the solver.

\noindent The \textbf{integer value iterator} currently available in Choco are the following: 
\begin{notedef}\tt
\hyperlink{decreasingdomain:decreasingdomainvaliterator}{DecreasingDomain}, \hyperlink{increasingdomain:increasingdomainvaliterator}{IncreasingDomain}
\end{notedef}

\noindent The \textbf{real value iterator} currently available in Choco is the following: 
\begin{notedef}\tt
\hyperlink{realincreasingdomain:realincreasingdomainvaliterator}{RealIncreasingDomain}
\end{notedef}




\subsubsection{Value selector}\label{solver:valueselector}\hypertarget{solver:valueselector}{}
The second way to do it is to compute the following value at each call.

\noindent The \textbf{integer value selector} currently available in Choco are the following: 
\begin{notedef}\tt
\hyperlink{bestfit:bestfitvalselector}{BestFit}, \hyperlink{costregularvalselector:costregularvalselectorvalselector}{CostRegularValSelector}, \hyperlink{fcostregularvalselector:fcostregularvalselectorvalselector}{FCostRegularValSelector}, \hyperlink{maxval:maxvalvalselector}{MaxVal}, \hyperlink{mcrvalselector:mcrvalselectorvalselector}{MCRValSelector}, \hyperlink{midval:midvalvalselector}{MidVal}, \hyperlink{minval:minvalvalselector}{MinVal}
\end{notedef}

\noindent The \textbf{set value selector} currently available in Choco is the following: 
\begin{notedef}\tt
\hyperlink{minenv:minenvvalselector}{MinEnv}, \hyperlink{randomsetvalselector:randomsetvalselectorvalselector}{RandomSetValSelector}
\end{notedef}

\subsection{Why is it important to define a search strategy ?}\label{solver:whyisitimportanttodefineasearchstrategy}\hypertarget{solver:whyisitimportanttodefineasearchstrategy}{}

In a partial instantiation, when a fix point has been reached, the Solver needs to take a decision to resume the search. The way decisions are chosen has a \textbf{real impact on the resolution step efficient}. 
\begin{note}
\emph{The search strategy should not be overlooked!!}
An adapted search strategy can reduce: the execution time, the number of node expanded, the number of backtrack done.
\end{note}
Let see that small example:
\begin{lstlisting}
	Model m = new CPModel();
        int n = 1000;
        IntegerVariable var = Choco.makeIntVar("var", 0, 2);
        IntegerVariable[] bi = Choco.makeBooleanVarArray("b", n);
        m.addConstraint(Choco.eq(var, Choco.sum(bi)));

        Solver badStrat = new CPSolver();
        badStrat.read(m);
        badStrat.addGoal(
                new AssignVar(
                        new MinDomain(badStrat), 
                        new IncreasingDomain()
                ));
        badStrat.solve();
        badStrat.printRuntimeStatistics();

        Solver goodStrat = new CPSolver();
        goodStrat.read(m);
        goodStrat.addGoal(
                new AssignVar(
                        new MinDomain(goodStrat, goodStrat.getVar(new IntegerVariable[]{var})), 
                        new DecreasingDomain()
                ));
        goodStrat.solve();
        goodStrat.printRuntimeStatistics();
\end{lstlisting}

This model ensures that $var = b_{0} + b_{1} + \ldots + b_{1000}$ where \emph{var} has a small domain and $b_{i}$ is a binary variable. The propagation has no effect on any domain and a fix point is reached at the beginning of the search. So, a decision has to be done choosing a variable and its value. If the variable selector is set to \texttt{MinDomain} (see below), the solver will iterate over the variables, starting by the 1000 binary variables and ending with \emph{var}, and 1001 nodes will be created.

\subsection{Restarts}\label{solver:restarts}\hypertarget{solver:restarts}{}

You can set geometric restarts by using the following API available on the solver:
\begin{lstlisting}
setGeometricRestart(int base, double grow);
setGeometricRestart(int base, double grow, int restartLimit);
\end{lstlisting}
It performs a search with restarts regarding the number of backtrack. An initial allowed number of backtrack is given (parameter base) and once this limit is reached a restart is performed and the new limit imposed to the search is increased by multiplying the previous limit with the parameter grow. restartLimit parameter states the maximum number of restarts. Restart strategies makes really sense with strategies that make choices based on the past experience of the search : \texttt{DomOverWdeg} or Impact based search. It could also be used with a random heuristic
\begin{lstlisting}
	CPSolver s = new CPSolver();
	s.read(model);
	
	s.setGeometricRestart(14, 1.5d);
	s.setFirstSolution(true);
	s.generateSearchStrategy();
	s.attachGoal(new DomOverWDegBranching(s, new IncreasingDomain()));
	s.launch();
\end{lstlisting}

You can also set Luby restarts by using the following API available on the solver:
\begin{lstlisting}
setLubyRestart(int base);
setLubyRestart(int base, int grow);
setLubyRestart(int base, int grow, int restartLimit);
\end{lstlisting}
it performs a search with restarts regarding the number of backtracks. One way to describe this strategy is to say that all run lengths are power of two, and that each time a pair of runs of a given length has been completed, a run of twice that length is immediatly executed. The limit is equals to \emph{length*base}.
\begin{itemize}
	\item \textbf{example with growing factor of 2 : [1, 1, 2, 1, 1, 2, 4, 1, 1, 2, 1, 1, 2, 4, 8, 1,...]}
	\item \textbf{example with growing factor of 3 : [1, 1, 1, 3, 1, 1, 1, 3, 9,...]}
\end{itemize}

\begin{lstlisting}
	CPSolver s = new CPSolver();
	s.read(model);
	
	s.setLubyRestart(50, 2, 100);
	s.setFirstSolution(true);
	s.generateSearchStrategy();
	s.attachGoal(new DomOverWDegBranching(s, new IncreasingDomain()));
	s.launch();
\end{lstlisting}

\section{Limiting Search Space}\label{solver:limitingsearchspace}\hypertarget{solver:limitingsearchspace}{}
The Solver class provides some limits on the search strategy that you can fix or just monitor.
Limits may be imposed on the search algorithm to avoid spending too much time in the exploration. The limits are updated and checked each time a new node is created. It has to be specified before the resolution. 
After having created the solver, you can specify whether or not you want to fix a limit:

\begin{description}
\item[time limit] State a time limit on tree search. When the execution time is equal to the time limit, the search stops whatever a solution is found or not. You can define a time limit with the following API : \mylst{setTimeLimit(int timeLimit)} where unit is millisecond. Or just monitor (or not) the search time with the API : \mylst{monitorTimeLimit(boolean b)}. The default value is set to \texttt{true}. Finally, you can get the time limit, once the solve method has been called, with the API: \mylst{getTimeCount()} 
\item[node limit] State a node limit on tree search. When the number of nodes explored is equal to the node limit, the search stops whatever a solution is found or not. You can define a node limit with the following API: \mylst{setNodeLimit(int nodeLimit)} where unit is the number of nodes. Or just monitor (or not) the number of nodes explored with the API: \mylst{monitorNodeLimit(boolean b)}. The default value is set to \texttt{true}. Finally, you can get the node limit, once the solve method has been called, with the API: \mylst{getNodeCount()} 
\item[backtrack limit] State a backtrack limit on tree search. When the number of backtracks done is equal to the backtrack limit, the search stops whatever a solution is found or not. You can define a backtrack limit with the following API: \mylst{setBackTrackLimit(int backtrackLimit)} where unit is the number of backtracks. Or just monitor (or not) the number of backtrack done with the API: \mylst{monitorBackTrackLimit(boolean b)}. The default value is set to \texttt{false}. Finally, you can get the backtrack limit, once the solve method has been called, with the API: \mylst{getBackTrackCount()} 
\item[fail limit] State a fail limit on tree search. When the number of failure is equal to the fail limit, the search stops whatever a solution is found or not. You can define a fail limit with the following API : \mylst{setFailLimit(int failLimit)} where unit is the number of failure. Or just monitor (or not) the number of failure encountered with the API : \mylst{monitorFailLimit(boolean b)}. The default value is set to \texttt{false}. Finally, you can get the fail limit, once the solve method has been called, with the API : \mylst{getFailCount()} 
%\item[CPU time limit] State a CPU limit on tree search. When the CPU time (user + system) is equal to the CPU time limit, the search stops whatever a solution is found or not. You can define a CPU time limit with the following API: \mylst{setCpuTimeLimit(int timeLimit)} where unit is millisecond. Or just monitor (or not) the search time with the API: \mylst{monitorCpuTimeLimit(boolean b)}. The default value is set to \texttt{false}. Finally, you can get the CPU time limit, once the solve method has been called, with the API: \mylst{getCpuTimeCount()} 
\end{description}

\todo{add example}

\section{Solve a problem}\label{solver:solveaproblem}\hypertarget{solver:solveaproblem}{}
As Solver is the second element of a Choco program, the control of the search process without using predefined tools is made on the Solver.

\noindent\begin{tabular}{p{.4\linewidth}p{.6\linewidth}}
  \hline
  \texttt{Solver} API & description \\
  \hline
      \mylst{solve()} &  Compute the first solution of the Model, if the Model is feasible. \\
      \mylst{solve(boolean all)} &  If \emph{all} is set to true, computes all solutions of the Model, if the Model is feasible. \\
      \mylst{solveAll()} &  Computes all the solution of the Model, if the Model is feasible. \\
      \mylst{propagate()} &  Computes initial propagation of the Model, and reachs the first Fix Point. It reduces variables Domain through constraints linked and other variables domain. Can throw a \texttt{ContradictionException} if the Solver detects a contradiction in the Model. \\
      \mylst{maximize(Var obj, boolean restart)} &  Allows user to find a solution that maximizing the objective varible \emph{obj}. The optimization finds a first solution then finds a new solution that improves \emph{obj} and so on till no other solution can be found that improves \emph{obj}. Parameter \emph{restart} is a boolean indicating whether the Solver will restart the search after each solution found (if set to \texttt{true}) or if it will keep backtracking from the leaf of the last solution found. See \hyperlink{solver:optimization}{example}. \textbf{Beware}: the variable \emph{obj} expected must be a Solver variable and not a Model variable. \\
      \mylst{minimize(Var obj, boolean restart)} &  Allows user to find a solution that minimizing the objective varible \emph{obj}. The optimization finds a first solution then finds a new solution that improves \emph{obj} and so on till no other solution can be found that improves \emph{obj}. Parameter \emph{restart} is a boolean indicating whether the Solver will restart the search after each solution found (if set to \texttt{true}) or if it will keep backtracking from the leaf of the last solution found. See \hyperlink{solver:optimization}{example}. \textbf{Beware}: the variable \emph{obj} expected must be a Solver variable and not a Model variable. \\
      \mylst{nextSolution()} &  Allows the Solver to find the next solution, if one or more solution have already been find with \texttt{solve()} or \texttt{nextSolution()}. \\
      \mylst{isFeasible()} &  Indicates whether or not the Model has at least one solution. \\
      \hline\\
	\end{tabular}

\subsection{Solver settings}\label{solver:solversettings}\hypertarget{solver:solversettings}{}

\subsubsection{Logs}\label{solver:logs}\hypertarget{solver:logs}{}
A logging class is instrumented in order to produce trace statements throughout search: ChocoLogging. The verbosity level of the solver can be set, by the following static method
\begin{lstlisting}
	ChocoLogging.toVerbose();
	// And after solver.solve()
	ChocoLogging.flushLogs();
\end{lstlisting}

The code above ensure that messages are printed in order to describe the construction of the search tree.

Six verbosity levels are available:

\noindent\begin{tabular}{p{.4\linewidth}p{.6\linewidth}}
  \hline
  Level & prints... \\
  \hline
 \texttt{ChocoLogging.toSilent()} & display only severe messages from core loggers and warning messages otherwise\\
 \ \texttt{ChocoLogging.toQuiet()} & display only severe messages from core loggers and info messages otherwise\\
 \texttt{ChocoLogging.toDefault()} & display information about initial and final state of the search\\
 \texttt{ChocoLogging.toVerbose()} & display search information at regular node intervals\\
 \texttt{ChocoLogging.toSolution()} & display all solutions\\
 \texttt{ChocoLogging.toSearch()} & display the search tree\\
\hline\\
\end{tabular}

Note that in the case of a verbosity greater or equals to \texttt{toVerbose()}, the regular search information step is set to 1000, by default. You can change this value, using:
\begin{lstlisting}
  ChocoLogging.setEveryXNodes(20000);
\end{lstlisting}
 

Note that in the case of verbosity \texttt{toSearch()}, trace statements are printed up to a maximal depth in the search tree. The default value is set to 25, but you can change the value of this threshold, say to 10, with the following setter method:
\begin{lstlisting}
  ChocoLogging.setLoggingMaxDepth(10);
\end{lstlisting}

\subsection{Optimization}\label{solver:optimization}\hypertarget{solver:optimization}{}
\todo{to introduce}
\begin{lstlisting}
  Model m = new CPModel();
  IntegerVariable obj1 = makeEnumIntVar("obj1", 0, 7);
  IntegerVariable obj2 = makeEnumIntVar("obj1", 0, 5);
  IntegerVariable obj3 = makeEnumIntVar("obj1", 0, 3);
  IntegerVariable cost = makeBoundIntVar("cout", 0, 1000000);
  int capacity = 34;
  int[] volumes = new int[]{7, 5, 3};
  int[] energy = new int[]{6, 4, 2};
  // capacity constraint
  m.addConstraint(leq(scalar(volumes, new IntegerVariable[]{obj1, obj2, obj3}), capacity));
	
  // objective function
  m.addConstraint(eq(scalar(energy, new IntegerVariable[]{obj1, obj2, obj3}), cost));
  
  Solver s = new CPSolver();
  s.read(m);
  
  s.maximize(s.getVar(cost), false);
\end{lstlisting}
\label{doc:solver}\hypertarget{doc:solver}{}
%%\part{constraints}
\label{constraints}
\hypertarget{constraints}{}



\chapter{Constraints in alphabetical order}\label{constraints:constraintsinalphabeticalorder}\hypertarget{constraints:constraintsinalphabeticalorder}{}

Choco offers a large number of available constraints.
In this part, you will find a description of each defined constraint, with API, options and examples.
\begin{note}
	 \hyperlink{abs:absconstraint}{abs},
	 \hyperlink{alldifferent:alldifferentconstraint}{allDifferent},
	 \hyperlink{and:andconstraint}{and},
	 \hyperlink{atmostnvalue:atmostnvalueconstraint}{atMostNValue},
	 \hyperlink{boolchanneling:boolchannelingconstraint}{boolChanneling},
	 \hyperlink{cumulative:cumulativeconstraint}{cumulative},
	 \hyperlink{disjunctive:disjunctiveconstraint}{disjunctive},
	 \hyperlink{distanceeq:distanceeqconstraint}{distanceEQ},
	 \hyperlink{distancegt:distancegtconstraint}{distanceGT},
	 \hyperlink{distancelt:distanceltconstraint}{distanceLT},
	 \hyperlink{distanceneq:distanceneqconstraint}{distanceNEQ},
	 \hyperlink{eq:eqconstraint}{eq},
	 \hyperlink{eqcard:eqcardconstraint}{eqCard},
	 \hyperlink{equation:equationconstraint}{equation},
	 \hyperlink{false:falseconstraint}{FALSE},
	 \hyperlink{feaspairac:feaspairacconstraint}{feasPairAC},
	 \hyperlink{feastupleac:feastupleacconstraint}{feasTupleAC},
	 \hyperlink{feastuplefc:feastuplefcconstraint}{feasTupleFC},
	 \hyperlink{geost:geostconstraint}{geost},
	 \hyperlink{geq:geqconstraint}{geq},
	 \hyperlink{geqcard:geqcardconstraint}{geqCard},
	 \hyperlink{globalcardinality:globalcardinalityconstraint}{globalCardinality},
	 \hyperlink{gt:gtconstraint}{gt},
	 \hyperlink{ifonlyif:ifonlyifconstraint}{ifOnlyIf},
	 \hyperlink{ifthenelse:ifthenelseconstraint}{ifThenElse},
	 \hyperlink{implies:impliesconstraint}{implies},
	 \hyperlink{infeaspairac:infeaspairacconstraint}{infeasPairAC},
	 \hyperlink{infeastupleac:infeastupleacconstraint}{infeasTupleAC},
	 \hyperlink{infeastuplefc:infeastuplefcconstraint}{infeasTupleFC},
	 \hyperlink{intdiv:intdivconstraint}{intDiv},
	 \hyperlink{inversechanneling:inversechannelingconstraint}{inverseChanneling},
	 \hyperlink{isincluded:isincludedconstraint}{isIncluded},
	 \hyperlink{isnotincluded:isnotincludedconstraint}{isNotIncluded},
	 \hyperlink{leq:leqconstraint}{leq},
	 \hyperlink{leqcard:leqcardconstraint}{leqCard},
	 \hyperlink{lex:lexconstraint}{lex},
	 \hyperlink{lexchain:lexchainconstraint}{lexChain},
	 \hyperlink{lexchaineq:lexchaineqconstraint}{lexChainEq},
	 \hyperlink{lexeq:lexeqconstraint}{lexeq},
	 \hyperlink{leximin:leximinconstraint}{leximin},
	 \hyperlink{lt:ltconstraint}{lt},
	 \hyperlink{max:maxconstraint}{max},
	 \hyperlink{member:memberconstraint}{member},
	 \hyperlink{min:minconstraint}{min},
	 \hyperlink{mod:modconstraint}{mod},
	 \hyperlink{multicostregular:multicostregularconstraint}{multiCostRegular},
	 \hyperlink{neq:neqconstraint}{neq},
	 \hyperlink{neqcard:neqcardconstraint}{neqCard},
	 \hyperlink{not:notconstraint}{not},
	 \hyperlink{notmember:notmemberconstraint}{notMember},
	 \hyperlink{nth:nthconstraint}{nth},
	 \hyperlink{occurrencemax:occurrencemaxconstraint}{occurrenceMax},
	 \hyperlink{occurrencemin:occurrenceminconstraint}{occurrenceMin},
	 \hyperlink{occurrence:occurrenceconstraint}{occurrence},
	 \hyperlink{oppositesign:oppositesignconstraint}{oppositeSign},
	 \hyperlink{or:orconstraint}{or},
	 \hyperlink{pack:packconstraint}{pack},
	 \hyperlink{precedencereified:precedencereifiedconstraint}{precedenceReified},
	 \hyperlink{preceding:precedingconstraint}{preceding},
	 \hyperlink{regular:regularconstraint}{regular},
	 \hyperlink{reifiedintconstraint:reifiedintconstraintconstraint}{reifiedIntConstraint},
	 \hyperlink{relationpairac:relationpairacconstraint}{relationPairAC},
	 \hyperlink{relationtupleac:relationtupleacconstraint}{relationTupleAC},
	 \hyperlink{relationtuplefc:relationtuplefcconstraint}{relationTupleFC},
	 \hyperlink{samesign:samesignconstraint}{sameSign},
	 \hyperlink{setdisjoint:setdisjointconstraint}{setDisjoint},
	 \hyperlink{setinter:setinterconstraint}{setInter},
	 \hyperlink{setunion:setunionconstraint}{setUnion},
	 \hyperlink{sorting:sortingconstraint}{sorting},
	 \hyperlink{stretchpath:stretchpathconstraint}{stretchPath},
	 \hyperlink{times:timesconstraint}{times},
	 \hyperlink{tree:treeconstraint}{tree},
	 \hyperlink{true:trueconstraint}{TRUE}.
   \end{note}


\chapter{Operators in alphabetical order}\label{constraints:operatorsinalphabeticalorder}\hypertarget{constraints:operatorsinalphabeticalorder}{}

  
\begin{note}
	 \hyperlink{abs:absoperator}{abs},
	 \hyperlink{cos:cosoperator}{cos},
	 \hyperlink{div:divoperator}{div},
	 \hyperlink{false:falseoperator}{FALSE},
	 \hyperlink{ifthenelse:ifthenelseoperator}{ifThenElse},
	 \hyperlink{max:maxoperator}{max},
	 \hyperlink{min:minoperator}{min},
	 \hyperlink{minus:minusoperator}{minus},
	 \hyperlink{mod:modoperator}{mod},
	 \hyperlink{mult:multoperator}{mult},
	 \hyperlink{neg:negoperator}{neg},
	 \hyperlink{plus:plusoperator}{plus},
	 \hyperlink{power:poweroperator}{power},
	 \hyperlink{scalar:scalaroperator}{scalar},
	 \hyperlink{sin:sinoperator}{sin},
	 \hyperlink{sum:sumoperator}{sum},
	 \hyperlink{true:trueoperator}{TRUE}. 
\end{note}
\label{doc:constraints}\hypertarget{doc:constraints}{}
%\part{advanced}
\label{advanced}
\hypertarget{advanced}{}


\chapter{Advanced uses of Choco}\label{advanced:advancedusesofchoco}\hypertarget{advanced:advancedusesofchoco}{}

\section{Environment}\label{advanced:environment}\hypertarget{advanced:environment}{}

Environment is a central object of the backtracking system. It defines the notion of \textit{world}. A world contains values of storable objects or operations that permit to \textit{backtrack} to its state. The environment \textit{pushes} and \textit{pops} worlds.

There are \textit{primitive} data types (\texttt{IstateBitSet, IStateBool, IStateDouble, IStateInt, IStateLong}) and \textit{objects} data types (\texttt{IStateBinarytree, IStateIntInterval, IStateIntProcedure, IStateIntVector, IStateObject, IStateVector}).

There are two different environments: \textit{EnvironmentTrailing} and \textit{EnvironmentCopying}.

\subsection{Copying}\label{advanced:copying}\hypertarget{advanced:copying}{}
In that environment, each data type is defined by a value (primitive or object) and a timestamp. Every time a world is pushed, each value is copied in an array (one array per data type), with finite indice. When a world is popped, every value is restored. 

\subsection{Trailing}\label{advanced:trailing}\hypertarget{advanced:trailing}{}
In that environment, data types are defined by its value. Every operation applied to a data type is pushed in a \textit{trailer}. When a world is pushed, the indice of the last operation is stored. When a world is popped, these operations are popped and \textit{unapplied} until reaching the last operation of the previous world.\\\textit{Default one in CPSolver}

\section{Define your own search strategy}\label{advanced:defineyourownsearchstrategy}\hypertarget{advanced:defineyourownsearchstrategy}{}
%A key ingredient of any constraint approach is a clever branching strategy. The construction of the search tree is done according to a series of Branching objects (that plays the role of achieving intermediate goals in logic programming). The user may specify the sequence of branching objects to be used to build the search tree. 
Section~\hyperlink{solver:searchstrategy}{Search strategy} presented the default branching strategies available in Choco and showed how to post them or to compose them as goals.
In this section, we will start with a very simple and common way to branch by choosing values for variables and specially how to define its own variable/value selection strategy. We will then focus on more complex branching such as dichotomic or n-ary choices. Finally we will show how to control the search space in more details with well known strategy such as LDS (Limited discrepancy search).

For integer variables, the variable and value selection strategy objects are based on the following interfaces:
\begin{itemize}
	\item \texttt{AbstractIntBranchingStrategy}: abstract class for the branching strategy,
	\item \texttt{VarSelector<V>} : Interface for the variable selection (V extends Var),
	\item \texttt{ValIterator<V>} : Interface to describes an iteration scheme on the domain of a variable,
	\item \texttt{ValSelector<V>} : Interface for a value selection.
\end{itemize}

Concrete examples of these interfaces are respectively,  \hyperlink{assignvar:assignvarbranchstrat}{AssignVar}, \hyperlink{mindomain:mindomainvarselector}{MinDomain}, \hyperlink{increasingdomain:increasingdomainvaliterator}{IncreasingDomain}, \hyperlink{maxval:maxvalvalselector}{MaxVal}.

\subsection{How to define your own Branching object}\label{advanced:beyondvariable/valueselection,howtodefineyourownbranchingobject}\hypertarget{advanced:beyondvariable/valueselection,howtodefineyourownbranchingobject}{}

Beyond Variable/value selection...

\subsection{Define your own variable selection}\label{advanced:defineyourownvariableselection}\hypertarget{advanced:defineyourownvariableselection}{}
You may extend this small library of branching schemes and heuristics by defining your own concrete classes of \texttt{AbstractIntVarSelector}. We give here an example of an \texttt{IntVarSelector} with the implementation of a static variable ordering :
\begin{lstlisting}
  public class StaticVarOrder extends AbstractIntVarSelector {

      // the sequence of variables that need be instantiated
	  protected IntDomainVar[] vars;
	
	  public StaticVarOrder(IntDomainVar[] vars) {
          this.vars = vars;
	  }
	
	  public IntDomainVar selectIntVar() {
          for (int i = 0; i < vars.length; i++)
              if (!vars[i].isInstantiated())
                  return vars[i];
          return null;
	  }
  }
\end{lstlisting}

Notice on this example that you only need to implement method \texttt{selectIntVar()} which belongs to the contract of \texttt{IntVarSelector}. This method should return a non instantiated variable or \texttt{null}. Once the branching is finished, the next branching (if one exists) is taken by the search algorithm to continue the search, otherwise, the search stops as all variable are instantiated. To avoid the loop over the variables of the branching, a backtrackable integer (\texttt{StoredInt}) could be used to remember the last instantiated variable and to directly select the next one in the table. Notice that backtrackable structures could be used in any of the code presented in this chapter to speedup the computation of dynamic choices.

You can add your variable selector as a part of a search strategy, using \mylst{solver.addGoal()}.

\subsection{Define your own value selection}\label{advanced:defineyourownvalueselection}\hypertarget{advanced:defineyourownvalueselection}{}
You may also define your own concrete classes of \texttt{ValIterator} or \texttt{ValSelector}. 

\subsubsection{Value selector}\label{advanced:valueselector}\hypertarget{advanced:valueselector}{}

\insertGraphique{.5\linewidth}{media/valselector.pdf}{ValSelector interface and implementation}

We give here an example of an \texttt{IntValSelector} with the implementation of a minimum value selecting:
\begin{lstlisting}
  public class MinVal extends AbstractSearchHeuristic implements ValSelector {
	 /**
      * selecting the lowest value in the domain
      * @param x the variable under consideration
      * @return what seems the most interesting value for branching
      */
	  public int getBestVal(IntDomainVar x) {
          return x.getInf();
	  }
  }
\end{lstlisting}
Only \texttt{getBestVal()} method must be implemented, returning the best value \emph{in the domain} according to the heuristic.

You can add your value selector as a part of a search strategy, using \mylst{solver.addGoal()}.

\begin{note}
Using a value selector with bounded domain variable is strongly inadvised, except if it pick up bounds value. If the value selector pick up a value that is not a bound, when it goes up in the tree search, that value could be not removed and picked twice (or more)!
\end{note} 

\subsubsection{Values iterator}\label{advanced:valuesiterator}\hypertarget{advanced:valuesiterator}{}
We give here an example of an \texttt{ValIterator} with the implementation of an increasing domain iterator:
\begin{lstlisting}
  public final class IncreasingDomain implements ValIterator {
	  /**
	   * testing whether more branches can be considered after branch i, 
       * on the alternative associated to variable x
	   * @param x the variable under scrutiny
	   * @param i the index of the last branch explored
	   * @return true if more branches can be expanded after branch i
	   */
       public boolean hasNextVal(Var x, int i) {
           return (i < ((IntDomainVar) x).getSup());
       }
	
	  /**
	   * accessing the index of the first branch for variable x
	   * @param x the variable under scrutiny
	   * @return the index of the first branch: first value to be assigned to x
	   */
       public int getFirstVal(Var x) {
           return ((IntDomainVar) x).getInf();
       }
	
	  /**
	   * generates the index of the next branch after branch i, 
       * on the alternative associated to variable x
	   * @param x the variable under scrutiny
	   * @param i the index of the last branch explored
	   * @return the index of the next branch to be expanded after branch i
	   */
       public int getNextVal(Var x, int i) {
           return ((IntDomainVar) x).getNextDomainValue(i);
       }
   }
\end{lstlisting}
%Works as an basic \texttt{Iterator} object, implementing the three main methods \texttt{hasNextVal()}, \texttt{getFirstVal()} and \texttt{getNextVal()}.

You can add your value iterator as a part of a search strategy, using \mylst{solver.addGoal()}.

\todo{under development} See \href{http://choco-solver.net/index.phptitle=userguide:beyondvariable.2fvalueselection.2chowtodefineyourownbranchingobject}{old version}

\section{Define your own limit search space}\label{advanced:defineyourownlimitsearchspace}\hypertarget{advanced:defineyourownlimitsearchspace}{}

To define your own limits/statistics (notice that a limit object can be used only to get statistics about the search), you can create a limit object by extending the \texttt{AbstractGlobalSearchLimit} class or implementing directly the interface \texttt{IGlobalSearchLimit}. Limits are managed at each node of the tree search and are updated each time a node is open or closed. Notice that limits are therefore time consuming. Implementing its own limit need only to specify to the following interface :

\begin{lstlisting}
	/**
	 * The interface of objects limiting the global search exploration
	 */
	public interface GlobalSearchLimit {

	  /**
	   * resets the limit (the counter run from now on)
	   * @param first true for the very first initialization, false for subsequent ones
	   */
	  public void reset(boolean first);
	
	  /**
	   * notify the limit object whenever a new node is created in the search tree
	   * @param solver the controller of the search exploration, managing the limit
	   * @return true if the limit accepts the creation of the new node, false otherwise
	   */
	  public boolean newNode(AbstractGlobalSearchSolver solver);
	
	  /**
	   * notify the limit object whenever the search closes a node in the search tree
	   * @param solver the controller of the search exploration, managing the limit
	   * @return true if the limit accepts the death of the new node, false otherwise
	   */
	  public boolean endNode(AbstractGlobalSearchSolver solver);
	}
\end{lstlisting}

Look at the following example to see a concrete implementation of the previous interface. We define here a limit on the depth of the search (which is not found by default in choco). The \texttt{getWorldIndex()} is used to get the current world, i.e the current depth of the search or the number of choices which have been done from baseWorld. 

\begin{lstlisting}
	public class DepthLimit extends AbstractGlobalSearchLimit {
	
	  public DepthLimit(AbstractGlobalSearchSolver theSolver, int theLimit) {
	    super(theSolver,theLimit);
	    unit = "deep";
	  }
	
	  public boolean newNode(AbstractGlobalSearchSolver solver) {
	    nb = Math.max(nb, this.getProblem().getWorldIndex() –
	    this.getProblem().getSolver().getSearchSolver().baseWorld);
	    return (nb < nbMax);
	  }
	
	  public boolean endNode(AbstractGlobalSearchSolver solver) {
	    return true;
	  }
	
	  public void reset(boolean first) {
	   if (first) {
	    nbTot = 0;
	   } else {
	    nbTot = Math.max(nbTot, nb);
	   }
	   nb = 0;
	  }
\end{lstlisting}

Once you have implemented your own limit, you need to tell the search solver to take it into account. Instead of using a call to the \texttt{solve()} method, you have to create the search solver by yourself and add the limit to its limits list such as in the following code :
\begin{lstlisting}
	Solver s = new CPSolver();
	s.read(model);
	s.setFirstSolution(true);
	s.generateSearchStrategy();
	s.getSearchStrategy().limits.add(new DepthLimit(s.getSearchStrategy(),10));
	s.launch();
\end{lstlisting}


\subsection{How does a search loop work ?}\label{advanced:howdoesasearchloopwork}\hypertarget{advanced:howdoesasearchloopwork}{}
The seach loop is created when a \texttt{solve()} method is called. It goes down and up in the branches in order to cover the tree search. 

%\subsubsection{Search loop}\label{advanced:searchloop}\hypertarget{advanced:searchloop}{}
\begin{lstlisting}[title={Algorithm of the search loop in Choco}]
  next_move = new node
  WHILE no solution AND in search limit
        IF next_move is new node
	    THEN
	        create a new node : variable/value selection ;
	        IF node exists 
	        THEN
	            next_move <-- go down branch ;            
	        ELSE 
	            next_move <-- go up branch ;
	            solution is found ;
	
	    ELSE IF next_move is go down branch
	        propagate ;
	        IF no contradiction 
	        THEN
	            next_move <-- new node ;                        
	        ELSE
	            next_move <-- go up branch ;
	        
	    ELSE IF next_move is go up branch 
	        find next branch ;
	        propagate ;
	        IF has next branch AND no contradiction
	        THEN
	            next_move <-- go down branch ;                            
	        ELSE
	            next_move <-- go up branch ;
	
	    END IF
	
  END WHILE
\end{lstlisting}
%\subsubsection{Search loop with recomputation}\label{advanced:searchloopwithrecomputation}\hypertarget{advanced:searchloopwithrecomputation}{}

\section{Define your own constraint}\label{advanced:defineyourownconstraint}\hypertarget{advanced:defineyourownconstraint}{}

This section describes how to add you own constraint, with specific propagation algorithms. Note that this section is only useful in case you want to express a constraint for which the basic propagation algorithms (using tables of tuples, or boolean predicates) are not efficient enough to propagate the constraint.

The general process consists in defining a new constraint class and implementing the various propagation methods. We recommend the user to follow the examples of existing constraint classes (for instance, such as \texttt{GreaterOrEqualXYC} for a binary inequality) 

\subsection{The constraint hierarchy}\label{advanced:theconstrainthierarchy}\hypertarget{advanced:theconstrainthierarchy}{}

Each new constraint must be represented by an object implementing the \texttt{\bf SConstraint} interface (\texttt{S} for solver constraint). To help the user defining new constraint classes, several abstract classes defining \texttt{SConstraint} have been implemented. These abstract classes provide the user with a management of the constraint network and the propagation engineering. They should be used as much as possible.

For constraints on integer variables, the easiest way to implement your own constraint is to inherit from one of the following classes, depending of the number of solver integer variables (\texttt{IntDomainVar}) involved:

\centerline{\begin{tabular}{ll}
      \hline
  Default class to implement &  number of solver integer variables \\
  \hline
  \mylst{AbstractUnIntSConstraint} &  \textbf{one} variable \\
  \mylst{AbstractBinIntSConstraint} &  \textbf{two} variables \\
  \mylst{AbstractTernIntSConstraint} &  \textbf{three} variables \\
  \mylst{AbstractLargeIntSConstraint} &  any number of variables. \\
  \hline\\
\end{tabular}}

\noindent Constraints over integers must implement the following methods (grouped in the \texttt{IntSConstraint} interface):

\noindent\begin{tabular}{lp{.6\linewidth}}
  \hline
  Method to implement &  description \\
  \hline
  \mylst{pretty()} &Returns a pretty print of the constraint \\
  \mylst{propagate()} &The main propagation method (propagation from scratch). Propagating the constraint until local consistency is reached. \\
  \mylst{awake()} &Propagating the constraint for the very first time until local consistency is reached. The awake is meant to initialize the data structures contrary to the propagate. Specially, it is important to avoid initializing the data structures in the constructor. \\
  \mylst{awakeOnInst(int x)} &Default propagation on instantiation: full constraint re-propagation. \\
  \mylst{awakeOnBounds(int x)} &Default propagation on improved bounds: propagation on domain revision. \\
  \mylst{awakeOnRemovals(int x, IntIterator v)} &Default propagation on mutliple values removal: propagation on domain revision. The iterator allow to iterate over the values that have been removed. \\
&\\
\hline
\multicolumn{2}{l}{Methods \texttt{awakeOnBounds} and \texttt{awakeOnRemovals} can be replaced by more fine grained methods:}\\
\hline
%Alternative Method &  description \\
%  \hline
  \mylst{awakeOnInf(int x)} &Default propagation on improved lower bound: propagation on domain revision. \\
  \mylst{awakeOnSup(int x)} &Default propagation on improved upper bound: propagation on domain revision. \\
  \mylst{awakeOnRem(int x, int v)} &Default propagation on one value removal: propagation on domain revision.  \\
&\\
  \hline
\multicolumn{2}{l}{To use the constraint in expressions or reification, the following minimum API is mandatory:}\\
  \hline
  \mylst{isSatisfied(int[] x)} &Tests if the constraint is satisfied when the variables are instantiated. \\
	\mylst{isEntailed()} &Checks if the constraint must be checked or must fail. It returns true if the constraint is known to be satisfied whatever happend on the variable from now on, false if it is violated. \\
	\mylst{opposite()} &It returns an AbstractSConstraint that is the opposite of the current constraint. \\
    \hline\\
	\end{tabular}

In the same way, a \textbf{set constraint} can inherit from \texttt{AbstractUnSetSConstraint}, \texttt{AbstractBinSetSConstraint}, \texttt{AbstractTernSetSConstraint} or \texttt{AbstractLargeSetSConstraint}.

A \textbf{real constraint} can inherit from \texttt{AbstractUnRealSConstraint}, \texttt{AbstractBinRealSConstraint} or \texttt{AbstractLargeRealSConstraint}.

A mixed constraint between \textbf{set and integer variables} can inherit from \texttt{AbstractBinSetIntSConstraint} or \texttt{AbstractLargeSetIntSConstraint}.

\begin{note}
A simple way to implement its own constraint is to:
\begin{itemize}
	\item create an empty constraint with only \texttt{propagate()} method implemented and every \texttt{awakeOnXxx()} ones set to \texttt{this.constAwake(false);}
	\item when the propagation filter is sure, separate it into the \texttt{awakeOnXxx()} methods in order to have finer granularity
	\item finally, if necessary, use backtrackables objects to improve the efficient of your constraint
\end{itemize}

\end{note}

\subsubsection{How do I add my constraint to the Model ?}\label{advanced:howdoiaddmyconstrainttothemodel}\hypertarget{advanced:howdoiaddmyconstrainttothemodel}{}

Adding your constraint to the model requires you to definite a specific constraint manager (that can be a inner class of your Constraint).
This manager need to implement:
\begin{lstlisting}
makeConstraint(Solver s, Variable[] vars, Object params, HashSet<String> options)
\end{lstlisting}
This method allows the Solver to create an instance of your constraint, with your parameters and Solver objects.

\begin{note}
If you create your constraint manager as an inner class, you must declare this class as \textbf{public and static}.
If you don't, the solver can't instantiate your manager.
\end{note}

Once this manager has been implemented, you simply add your constraint to the model using the \texttt{addConstraint()} API with a \texttt{ComponentConstraint} object:
\begin{lstlisting}
  model.addConstraint( new ComponentConstraint(MyConstraintManager.class, params, vars) );
  // OR
  model.addConstraint( new ComponentConstraint("package.of.MyConstraint", params, vars) );
\end{lstlisting}
Where \emph{params} is whatever you want (\texttt{Object[], int, String},...) and \emph{vars} is an array of Model Variables (or more specific) objects.

\subsection{Example: implement and add the \texttt{IsOdd} constraint}
One creates the constraint by implementing the \texttt{AbstractUnIntSConstraint} (one integer variable) class:
\lstinputlisting{java/isodd.j2t}

To add the constraint to the model, one creates the following class (or inner class):
\lstinputlisting{java/isoddmanager.j2t}
It calls the constructor of the constraint, with every \emph{vars}, \emph{params} and \emph{options} needed.

Then, the constraint can be added to a model as follows:
\begin{lstlisting}
	// Creation of the model
	Model m = new CPModel();
	
	// Declaration of the variable
	IntegerVariable aVar = Choco.makeIntVar("a_variable", 0, 10);
	
	// Adding the constraint to the model, 1st solution:
	m.addConstraint(new ComponentConstraint(IsOddManager.class, null, new IntegerVariable[]{aVar}));
	// OR 2nd solution:
	m.addConstraint(new ComponentConstraint("myPackage.Constraint.IsOddManager", null, new IntegerVariable[]{aVar}));
	
	Solver s = new CPSolver();
	s.read(m);
	s.solve();
\end{lstlisting}
And that's it!!

\subsection{Example of an empty constraint}\label{advanced:anexempleofemptyconstraint}\hypertarget{advanced:anexempleofemptyconstraint}{}

See the complete code: \href{media/zip/constraintpattern.zip}{ConstraintPattern.zip}

\begin{lstlisting}
  public class ConstraintPattern extends AbstractLargeIntSConstraint {
      
      public ConstraintPattern(IntDomainVar[] vars) {
          super(vars);
      }
	
      /**
      * pretty print. The String is not constant and may depend on the context.
      * @return a readable string representation of the object
      */
      public String pretty() {
          return null;
      }
	
      /**
      * check whether the tuple satisfies the constraint
      * @param tuple values
      * @return true if satisfied
      */
      public boolean isSatisfied(int[] tuple) {
          return false;
      }

      /**
      * propagate until local consistency is reached
      */
      public void propagate() throws ContradictionException {
          // elementary method to implement
      }
	    
      /**
      * propagate for the very first time until local consistency is reached.
      */
      public void awake() throws ContradictionException {
          constAwake(false);        // change if necessary
      }
	
	
      /**
      * default propagation on instantiation: full constraint re-propagation
      * @param var index of the variable to reduce
      */
      public void awakeOnInst(int var) throws ContradictionException {
          constAwake(false);        // change if necessary
      }
	
      /**
      * default propagation on improved lower bound: propagation on domain revision
      * @param var index of the variable to reduce
      */
      public void awakeOnInf(int var) throws ContradictionException {
          constAwake(false);        // change if necessary
      }
	
	
      /**
      * default propagation on improved upper bound: propagation on domain revision
      * @param var index of the variable to reduce
      */
      public void awakeOnSup(int var) throws ContradictionException {
          constAwake(false);        // change if necessary
      }
	
      /**
      * default propagation on improve bounds: propagation on domain revision
      * @param var index of the variable to reduce
      */
      public void awakeOnBounds(int var) throws ContradictionException {
          constAwake(false);        // change if necessary
      }
	
      /**
      * default propagation on one value removal: propagation on domain revision
      * @param var index of the variable to reduce
      * @param val the removed value
      */
      public void awakeOnRem(int var, int val) throws ContradictionException {
          constAwake(false);        // change if necessary
      }
	
      /**
      * default propagation on one value removal: propagation on domain revision
      * @param var index of the variable to reduce
      * @param delta iterator over remove values
      */
      public void awakeOnRemovals(int var, IntIterator delta) throws ContradictionException {
          constAwake(false);        // change if necessary
      }
  }
\end{lstlisting}

The first step to create a constraint in Choco is to implement all \texttt{awakeOn...} methods with \texttt{constAwake(false)} and to put your propagation algorithm in the \texttt{propagate()} method. 

A constraint can choose not to react to fine grained events such as the removal of a value of a given variable but instead delay its propagation at the end of the fix point reached by ``fine grained events'' and fast constraints that deal with them incrementally (that's the purpose of the constraints events queue). 

To do that, you can use \texttt{constAwake(false)} that tells the solver that you want this constraint to be called only once the variables events queue is empty. This is done so that heavy propagators can delay their action after the fast one to avoid doing a heavy processing at each single little modification of domains.

\section{Define your own operator}\label{advanced:defineyourownoperator}\hypertarget{advanced:defineyourownoperator}{}
\todo{to complete}

\section{Define your own variable}\label{advanced:defineyourownvariable}\hypertarget{advanced:defineyourownvariable}{}
\todo{to complete}

\section{Backtrackable structures}\label{advanced:backtrackablestructures}\hypertarget{advanced:backtrackablestructures}{}
\todo{to complete}

\section{Logging System}\label{advanced:loggingsystem}\hypertarget{advanced:loggingsystem}{}

Choco logging system is based on the \texttt{java.util.logging} package and located in the package \texttt{common.logging}.
Most Choco abstract classes or interfaces propose a static field \texttt{LOGGER}.
The following figures present the architecture of the logging system with the default verbosity.

\insertGraphique{.9\linewidth}{media/logger-default.png}{Logger Tree with the default verbosity}

The shape of the node depicts the type of logger:
\begin{itemize}
	\item The \emph{house} loggers represent private loggers. Do not use directly these loggers because their level are low and all messages would always be displayed.
	\item The \emph{octagon} loggers represent critical loggers. These loggers are provided in the variables, constraints and search classes and could have a huge impact on the global performances.
	\item The \emph{box} loggers are provided for dev and users.
\end{itemize}
The color of the node gives its logging level with DEFAULT verbosity:
\texttt{Level.FINEST} (\textcolor{yellow}{gold}),
\texttt{Level.INFO} (\textcolor{orange}{orange}),
\texttt{Level.WARNING} (\textcolor{red}{red}).

\subsubsection{Verbosity and messages.}\label{advanced:verbosityandmessages}\hypertarget{advanced:verbosityandmessages}{}
The following table summarizes the verbosities available in choco: 

\begin{itemize}
	\item \textbf{OFF -- level 0:} Disable logging.
	\item \textbf{SILENT -- level 1:} Display only severe messages.
	\item \textbf{DEFAULT -- level 2:} Display informations on final search state.
		\begin{itemize}
			\item ON START
				\lstset{language={sh},columns=fixed}
\begin{lstlisting}
 ** CHOCO : Constraint Programming Solver
 ** CHOCO v2.1.1 (April, 2009), Copyleft (c) 1999-2010
 \end{lstlisting}
			\item ON COMPLETE SEARCH:
				\begin{lstlisting}
- Search completed -
 [Maximize		: {0},]
 [Minimize		: {1},]
  Solutions		: {2},
  Times (ms)	: {3},
  Nodes			: {4},
  Backtracks	: {5},
  Restarts		: {6}.
  \end{lstlisting}
	brackets [\textit{line}] indicate \textit{line} is optionnal,\\
 	\texttt{Maximize} --resp. \texttt{Minimize}-- indicates the best known value before exiting of the objective value in \textit{maximize()} -- --resp. \textit{minimize()}-- strategy.

			\item ON COMPLETE SEARCH WITHOUT SOLUTIONS :
				\begin{lstlisting}
- Search completed - No solutions
 [Maximize		: {0},]
 [Minimize		: {1},]
  Solutions		: {2},
  Times (ms)	: {3},
  Nodes			: {4},
  Backtracks	: {5},
  Restarts		: {6}.
\end{lstlisting}
	brackets [\textit{line}] indicate \textit{line} is optionnal,\\
 	\texttt{Maximize} --resp. \texttt{Minimize}-- indicates the best known value before exiting of the objective value in \textit{maximize()} -- --resp. \textit{minimize()}-- strategy.

			\item ON INCOMPLETE SEARCH:
				\begin{lstlisting}
- Search incompleted - Exiting on limit reached
  Limit			: {0},
 [Maximize		: {1},]
 [Minimize		: {2},]
  Solutions		: {3},
  Times (ms)	: {4},
  Nodes			: {5},
  Backtracks	: {6},
  Restarts		: {7}.
  
  \end{lstlisting}
	brackets [\textit{line}] indicate \textit{line} is optionnal,\\
 	\texttt{Maximize} --resp. \texttt{Minimize}-- indicates the best known value before exiting of the objective value in \textit{maximize()} -- --resp. \textit{minimize()}-- strategy.
		\end{itemize}			

	\item \textbf{VERBOSE -- level 3:} Display informations on search state.
		\begin{itemize}
			\item EVERY X (default=1000) NODES:
			\begin{lstlisting}
- Partial search - [Objective : {0}, ]{1} solutions, {2} Time (ms), {3} Nodes, {4} Backtracks, {5} Restarts.
			\end{lstlisting}
			\texttt{Objective} indicates the best known value.

			\item ON RESTART : 
			\begin{lstlisting}
- Restarting search - {0} Restarts.
			\end{lstlisting}
		\end{itemize}

	\item \textbf{SOLUTION -- level 4:} display all solutions.
		\begin{itemize}
			\item AT EACH SOLUTION:
			\begin{lstlisting}
- Solution #{0} found. [Objective: {0}, ]{1} Solutions, {2} Time (ms), {3} Nodes, {4} Backtracks, {5} Restarts.
  X_1:v1, x_2:v2...
			\end{lstlisting}
		\end{itemize}

	\item \textbf{SEARCH -- level 5:} Display the search tree.
		\begin{itemize}
			\item AT EACH NODE, ON DOWN BRANCH:
			\begin{lstlisting}
...[w] down branch X==v branch b
			\end{lstlisting}
where \texttt{w} is the current world index, \texttt{X} the branching variable, \texttt{v} the branching value and \texttt{b} the branch index. This message can be adapted on variable type and search strategy.

			\item AT EACH NODE, ON UP BRANCH:
			\begin{lstlisting}
...[w] up branch X==v branch b
			\end{lstlisting}
where \texttt{w} is the current world index, \texttt{X} the branching variable, \texttt{v} the branching value and \texttt{b} the branch index. This message can be adapted on variable type and search strategy.
		\end{itemize}

	\item \textbf{FINEST -- level 6:} display all logs.

\end{itemize}

More precisely, if the verbosity level is greater than DEFAULT, then the verbosity levels of the model and of the solver are increased to INFO, and the verbosity levels of the search and of the branching are slightly modified to display the solution(s) and search messages.

The verbosity level can be changed as follows:
\begin{lstlisting}
	ChocoLogging.setVerbosity(Verbosity.VERBOSE);
\end{lstlisting}


\subsubsection{How to write logging statements ?}\label{advanced:howtowriteloggingstatements}\hypertarget{advanced:howtowriteloggingstatements}{}

\begin{itemize}
	\item Critical Loggers are provided to display error or warning. Displaying too much message really \textbf{impacts the performances}.
	\item Check the logging level before creating arrays or strings.
	\item Avoid multiple calls to \texttt{Logger} functions. Prefer to build a \texttt{StringBuilder} then call the \texttt{Logger} function.
	\item Use the \texttt{Logger.log} function instead of building string in \texttt{Logger.info()}.
\end{itemize}

\subsubsection{Handlers.}\label{advanced:handlers}\hypertarget{advanced:handlers}{}
Logs are displayed on \texttt{System.out} but warnings and severe messages are also displayed on \texttt{System.err}.
\texttt{ChocoLogging.java} also provides utility functions to easily change handlers:
\begin{itemize}
	\item Functions \texttt{set...Handler} remove current handlers and replace them by a new handler.
	\item Functions \texttt{add...Handler} add new handlers but do not touch existing handlers.
\end{itemize}

\subsubsection{Define your own logger.}\label{advanced:defineyourownlogger}\hypertarget{advanced:defineyourownlogger}{}
\begin{lstlisting}
ChocoLogging.makeUserLogger(String suffix);
\end{lstlisting}

% \subsubsection{Figure source}\label{advanced:figuresource}\hypertarget{advanced:figuresource}{}
% \begin{lstlisting}
%   digraph G {
%       node [style=filled, shape=box];
%       choco [shape=house,fillcolor=gold];
	
%       kernel [shape=house,fillcolor=gold];
%       engine [shape=octagon,fillcolor=indianred];
%       search [shape=octagon,fillcolor=darkorange];
%       branching [shape=octagon,fillcolor=indianred];
	
%       api [shape=house,fillcolor=gold];
%       model [fillcolor=indianred];
%       solver [fillcolor=indianred];
%       parser [fillcolor=darkorange];
      
%       user [fillcolor=darkorange];
%       samples [fillcolor=darkorange];
      
%       test [fillcolor=indianred];
      
%       choco -> kernel;
%       choco -> API;
%       choco -> user;
%       choco -> test;
      
%       kernel -> engine;
%       kernel -> search;
      
%       api -> model;
%       api -> solver
%       api -> parser;
      
%       user  -> samples;
      
%       search -> branching;
% 	}
% \end{lstlisting}
\label{doc:advanced}\hypertarget{doc:advanced}{}
\chapter{Applications with global constraints}\label{doc:applications}\hypertarget{doc:applications}{}
%\part{scheduling and use of the cumulative}
%\label{schedulinganduseofthecumulative}
%\hypertarget{schedulinganduseofthecumulative}{}

\section{Scheduling and use of the cumulative constraint}\label{schedulinganduseofthecumulative:schedulinganduseofthecumulativeconstraint}\hypertarget{schedulinganduseofthecumulative:schedulinganduseofthecumulativeconstraint}{}

\emph{This tutorial is the Choco2 version of \href{http://choco-solver.net/index.phptitle=schedulinganduseofthecumulative}{this one}}

We present a simple example of a scheduling problem solved using the cumulative global constraint.

%The problem is to schedule a given set of tasks on a single resource and in a given time horizon. 
The problem is to maximize the number of tasks that can be scheduled on  a single resource within a given time horizon.
% and to provide a corresponding valid schedule.

The following picture summarizes the instance that we will use as example. 
It shows the resource profile on the left and on the right, the set of tasks to be scheduled. Each task is represented here as a rectangle whose height is the resource consumption of the task and whose length is the duration of the task. Notice that the profile is not a straight line but varies in time. 

\insertGraphique{\linewidth}{media/schedulinstance.jpg}{A cumulative scheduling problem instance}

This tutorial might help you to deal with :

\begin{itemize}
	\item The profile of the resource that varies in time whereas the API of the cumulative only accepts a constant capacity
	\item The objective function that implies optional tasks which is a priori not allowed by cumulative
	\item A search heuristic that will first assign the tasks to the resource and then try to schedule them while maximizing the number of tasks
\end{itemize}

The first point is easy to solve by adding fake tasks at the right position to simulate the consumption of the resource. The second point is possible thanks to the ability of the cumulative to handle variable heights. We shall explain it in more details soon.

Let's have a look at the source code and first start with the representation of the instance. We need three fake tasks to decrease the profile accordindly to the instance capacity. There is otherwise 11 tasks. Their heights and duration are fixed and given in the two following int[] tables. The three first tasks correspond to the fake one are all of duration 1 and heights 2, 1, 4. 
\begin{lstlisting}
  CPModel m = new CPModel();
  // data
  int n = 11 + 3; //number of tasks (include the three fake tasks)
  int[] heights_data = new int[]{2, 1, 4, 2, 3, 1, 5, 6, 2, 1,3, 1, 1, 2};
  int[] durations_data = new int[]{1, 1, 1, 2, 1, 3, 1, 1, 3, 4,2, 3, 1, 1};
\end{lstlisting}

The variables of the problem consist of four variables for each task (start, end, duration, height). We recall here that the scheduling convention is that a task is active on the interval [start, end-1] so that the upper bound of the start and end variables need to be 5 and 6 respectively. Notice that start and end variables are BoudIntVar variables. Indeed, the cumulative is only performing bound reasonning so it would be a waste of efficiency to declare here EnumVariables. Duration and heights are constant in this problem. However, our plan is to simulate the allocation of a task to the resource by using variable height. In other word, we will define the height of the task i as a variable of domain $\{0, \mathtt{height[i]}\}$. The height of the task takes its normal value if the task is assigned to the resource and 0 otherwise. The duration is really constant and is therefore created as a ConstantIntVar.

Moreover, we add a boolean variable per task to specify if the task is assigned to the resource or not. The objective variable is created as a BoundIntVar. 

\begin{lstlisting}
  IntegerVariable capa = constant(7);
  IntegerVariable[] starts = makeIntVarArray("start", n, 0, 5, "cp:bound");
  IntegerVariable[] ends = makeIntVarArray("end", n, 0, 6, "cp:bound");
  
  IntegerVariable[] duration = new IntegerVariable[n];
  IntegerVariable[] height = new IntegerVariable[n];
  for (int i = 0; i < height.length; i++) {
      duration[i] = constant(durations_data[i]);
      height[i] = makeIntVar("height " + i, new int[]{0, heights_data[i]});
  }
  
  IntegerVariable[] bool = makeIntVarArray("taskIn?", n, 0, 1);
  IntegerVariable obj = makeIntVar("obj", 0, n, "cp:bound", "cp:objective");
\end{lstlisting}

We then add the constraints to the model. Three constraints are needed. First, the cumulative ensures that the resource consumption is respected at any time. Then we need to make sure that if a task is assigned to the cumulative, its height can not be null which is done by the use of boolean channeling constraints. Those constraints ensure that :

$$\mathtt{bool[i]} = 1 \quad\iff\quad \mathtt{height[i]} = \mathtt{heights\_data[i]}$$

We state the objective function to be the sum of all boolean variables. 
\begin{lstlisting}
  //post the cumulative
  m.addConstraint(cumulative(starts, ends, duration, height, capa, ""));                                                                            
  //post the channeling to know if the task is scheduled or not
  for (int i = 0; i < n; i++) {
      m.addConstraint(boolChanneling(bool[i], height[i], heights_data[i]));
  }
  
  //state the objective function
  m.addConstraint(eq(sum(bool), obj));
\end{lstlisting}

Finally we fix the fake task at their position to simulate the profil: 
\begin{lstlisting}
  CPSolver s = new CPSolver();
  s.read(m);
  
  //set the fake tasks to establish the profile capacity of the resource
  try {
      s.getVar(starts[0]).setVal(1); s.getVar(ends[0]).setVal(2); s.getVar(height[0]).setVal(2);
      s.getVar(starts[1]).setVal(2); s.getVar(ends[1]).setVal(3); s.getVar(height[1]).setVal(1);
      s.getVar(starts[2]).setVal(3); s.getVar(ends[2]).setVal(4); s.getVar(height[2]).setVal(4);
  } catch (ContradictionException e) {
      System.out.println("error, no contradiction expected at this stage");
  }
\end{lstlisting}
We are now ready to solve the problem. We could call a maximize(false) but we want to add a specific heuristic that first assigned the tasks to the cumulative and then tries to schedule them.
\begin{lstlisting}
  s.maximize(s.getVar(obj),false);
\end{lstlisting}
We now want to print the solution and will use the following code : 
\begin{lstlisting}
	System.out.println("Objective : " + (s.getVar(obj).getVal() - 3));
	for (int i = 3; i < starts.length; i++) {
	   if (s.getVar(height[i]).getVal() != 0)
	      System.out.println("[" + s.getVar(starts[i]).getVal() + " - " 
                                 + (s.getVar(ends[i]).getVal() - 1) + "]:" 
                                 + s.getVar(height[i]).getVal());
	}
\end{lstlisting}
Choco gives the following solution : 
\begin{lstlisting}
	Objective : 9
	[1 - 2]:2
	[0 - 0]:3
	[2 - 4]:1
	[4 - 4]:5
	[5 - 5]:6
	[0 - 2]:2
	[0 - 3]:1
	[3 - 5]:1
	[0 - 0]:1
\end{lstlisting}
This solution could be represented by the following picture :

\insertGraphique{\linewidth}{media/schedulsolution.jpg}{Cumulative profile of a solution.}

Notice that the cumulative gives a necesserary condition for packing (if no schedule exists then no packing exists) but this condition is not sufficient as shown on the picture because it only ensures that the capacity is respected at each time point. Specially, the tasks might be splitted to fit in the profile as in the previous solution.
The complete code can be found \hyperlink{cumulative:cumulativeconstraint}{here}.

%\part{geost description}
\label{geostdescription}
\hypertarget{geostdescription}{}


\section{Placement and use of the Geost constraint}\label{geostdescription:placementanduseofthegeostconstraint}\hypertarget{geostdescription:placementanduseofthegeostconstraint}{}

The global constraint \texttt{\bf geost}($k,O,S,C$) handles in a generic way a variety of geometrical constraints $C$ in space and time between polymorphic $k\in\N$ dimensional objects $O$, each of which taking a shape among a set of shapes $S$ during a given time interval and at a given position in space. Each shape from $S$ is defined as a finite set of shifted boxes, where each shifted box is described by a box in a $k$-dimensional space at the given offset with the given sizes.

More precisely a \emph{shifted box} $s$= \emph{shape(sid,t[],l[])} is an entity defined by a shape id \emph{sid}, an shift offset \emph{s.t[d]}, $0\le d < k$, and a size \emph{s.l[d]}$>0$, $0\le d<k$. All attributes of a shifted box are integer values. Then, a \emph{shape} from $S$ is a collection of shifted boxes sharing all the same shape id. Note that the shifted boxes associated with a given shape may or may not overlap. This sometimes allows a drastic reduction in the number of shifted boxes needed to describe a shape.
Each \emph{object} \emph{o= object( id, sid,x[], start, duration,end)} from $O$ is an entity defined by a unique object id \emph{o.id} (an integer), a shape id \emph{o.sid}, an origin \emph{o.x[d]}, $0\le d<k$, a starting time \emph{o.start}, a duration \emph{o.duration}$>0$, and a finishingxs time \emph{o.end}.

All attributes \emph{sid, x[0],x[1],...,x[k-1], start, duration, end} correspond to domain variables. Typical constraints from the
list of constraints $C$ are for instance the fact that a given subset of objects from $O$ do not pairwise overlap.
Constraints of the list of constraints $C$ have always two first arguments $A_i$ and $O_i$ (followed by possibly some additional arguments) which respectively specify:
\begin{itemize}
	\item A list of dimensions (integers between 0 and k-1), or attributes of the objects of $O$ the constraint considers.
	\item A list of identifiers of the objects to which the constraint applies.
\end{itemize}

\subsection{Example and way to implement it}\label{geostdescription:exampleandwaytoimplementit}\hypertarget{geostdescription:exampleandwaytoimplementit}{}

We will explain how to use \emph{geost} although a 2D example. Consider we have 3 objects $o_0, o_1, o_2$ to place them inside a box $B$ (3x4) such that they don't overlap (see Figure bellow). The first object $o_0$ has two potential shapes while $o_1$ and $o_2$ have one shape. Given that the placement of the objects should be totally inside $B$ this means that the domain of the origins of objects are as follows (we start from 0 this means that the placement space is from 0 to 2 on $x$ and from 0 to 3 on $y$:
\begin{itemize}
	\item $o_0$: $x$ in 0..1, $y$ in 0..1,
	\item $o_1$: $x$ in 0..1, $y$ in 0..1,
	\item $o_2$: $x$ in 0..1, $y$ in 0..3.
\end{itemize}

\insertGraphique{.9\linewidth}{media/exp_geost.png}{Geost objects and shapes}

We describe now how to solve this problem by using Choco.

\subsubsection{Build a CP model.}\label{geostdescription:buildacpmodel}\hypertarget{geostdescription:buildacpmodel}{}
To begin the implementation we build a CP model:
\begin{lstlisting}
  Model m = new CPModel();
\end{lstlisting}
\subsubsection{Set the Dimension.}\label{geostdescription:setthedimension}\hypertarget{geostdescription:setthedimension}{}
Then we first need to specify the dimension k we are working in. This is done by assigning the dimension to a local variable that we will use later:
\begin{lstlisting}
  int dim = 2;
\end{lstlisting}
\subsubsection{Create the Objects.}\label{geostdescription:createtheobjects}\hypertarget{geostdescription:createtheobjects}{}
Then we start by creating the objects and store them in a vector as such:
\begin{lstlisting}
  Vector<GeostObject> objects = new Vector<GeostObject>();
\end{lstlisting}
Now we create the first object $o_0$ by creating all its attributes.
\begin{lstlisting}
  int objectId = 0; // object id
  IntegerVariable shapeId = Choco.makeIntVar("sid", 0, 1); // shape id (2 possible values)
  IntegerVariable coords[] = new IntegerVariable[dim]; // coordinates of the origin 
  coords[0] = Choco.makeIntVar("x", 0, 1);  
  coords[1] =  Choco.makeIntVar("y", 0, 1);
\end{lstlisting}
We need to specify 3 more Integer Domain Variables representing the temporal attributes (start, duration and end), which for the current implementation of \texttt{geost} are not working, however we need to give them dummy values. 
\begin{lstlisting}
  IntegerVariable start = Choco.makeIntVar("start", 0, 0);
  IntegerVariable duration = Choco.makeIntVar("duration", 1, 1);
  IntegerVariable end = Choco.makeIntVar("end", 1, 1);
\end{lstlisting}
Finally we are ready to add the object 0 to our \emph{objects} Vector:
\begin{lstlisting}
  objects.add(new GeostObject(dim, objectId, shapeId, coords, start, duration, end));
\end{lstlisting}
Now we do the same for the other  object $o_1$  and  $o_2$ and add them to our \emph{objects} vector.


\subsubsection{Create the Shifted Boxes.}\label{geostdescription:createtheshiftedboxes}\hypertarget{geostdescription:createtheshiftedboxes}{}
To create the shapes and their shifted boxes we create the shifted boxes and associate them with the corresponding shapeId. This is done as follows, first we create a Vector called \emph{sb} for example 
\begin{lstlisting}
Vector<ShiftedBox> sb = new Vector<ShiftedBox> ();
\end{lstlisting}
To create the shifted boxes for the shape 0 (that corresponds to $o_0$), we start by the first shifted box by creating the sid and 2 arrays one to specify the offset of the box in each dimension and one for the size of the box in each dimension:
\begin{lstlisting}
  int sid = 0; 
  int[] offset = {0,0};
  int[] sizes = {1,3};
\end{lstlisting} 
Now we add our shiftedbox to the \emph{sb} Vector:
\begin{lstlisting}
  sb.add(new ShiftedBox(sid, offset, sizes));
\end{lstlisting} 
We do the same with second shifted box:
\begin{lstlisting}
  sb.add(new ShiftedBox(0, new int[]{0,0}, new int[]{2,1}));
\end{lstlisting} 
By the same way we create the shifted boxes corresponding to second shape $S_1$:
\begin{lstlisting}
  sb.add(new ShiftedBox(1, new int[]{0,0}, new int[]{2,1}));
  sb.add(new ShiftedBox(1, new int[]{1,0}, new int[]{1,3}));
\end{lstlisting}
and the third shape $S_2$ consisting of three shifted boxes:
\begin{lstlisting}
  sb.add(new ShiftedBox(2, new int[]{0,0}, new int[]{2,1})) ;
  sb.add(new ShiftedBox(2, new int[]{1,0}, new int[]{1,3})); 
  sb.add(new ShiftedBox(2, new int[]{0,2}, new int[]{2,1}));
\end{lstlisting}
and finally the last shape $S_3$
\begin{lstlisting}
  sb.add(new ShiftedBox(3, new int[]{0,0}, new int[]{2,1}));
\end{lstlisting}

\subsubsection{Create the constraints.}\label{geostdescription:createtheconstraints}\hypertarget{geostdescription:createtheconstraints}{}
First we create a Vector called ectr that will contain the external constraints. 
\begin{lstlisting}
  Vector <ExternalConstraint> ectr = new Vector <ExternalConstraint>();
\end{lstlisting}
In order to create the non-overlapping constraint we first create an array containing all the dimensions the constraint will be active in (in our example it is all dimensions) and lets name this array \emph{ectrDim} and a list of objects \emph{objOfEctr} that this constraint will apply to (in our example it is all objects).

\begin{note}
Note that in the current implementation of \texttt{geost} only the non-overlapping constraint is available. Moreover, 
\emph{ectrDim} should contain all dimensions and \emph{objOfEctr} should contain all the objects, i.e. the non-overlapping constraint applies to all the objects in all dimensions.
\end{note}
 
After that we add the constraint to a vector \emph{ectr} that contains all the constraints we want to add.
The code for these steps is as follows: 
\begin{lstlisting}
  int[] ectrDim = new int[dim]; 
  for(i = 0; i < dim; i ++)  
      ectrDim[i] = i; 
  int[] objOfEctr = new int[3]; 
  for(i = 0; i < 3; i ++) 
      objOfEctr[i] = objects.elementAt(i).getObjectId();
\end{lstlisting}
All we need to do now is create the non-overlapping constraint and add it to the \emph{ectr} vector
that holds all the constraints. this is done as follows:
\begin{lstlisting}
  //Constants.NON_OVERLAPPING indicates the id of the non-overlapping constraint
  NonOverlapping n = new NonOverlapping(Constants.NON_OVERLAPPING, ectrDim, objOfEctr);
  ectr.add(n);
\end{lstlisting}

\subsubsection{Create the \texttt{geost} constraint and add it to the model.}\label{geostdescription:createthegeostconstraintandaddittothemodel}\hypertarget{geostdescription:createthegeostconstraintandaddittothemodel}{}
\begin{lstlisting}
  Constraint geost = Choco.geost(dim, objects, sb, ectr);
  m.addConstraint(geost);
\end{lstlisting}
	
\subsubsection{Solve the problem.}\label{geostdescription:solvetheproblem}\hypertarget{geostdescription:solvetheproblem}{}
\begin{lstlisting}
  Solver s = new CPSolver();
  s.read(m);
  s.solve();
\end{lstlisting}

The full java code can be found here: \href{media/zip/geostexp.java}{geostexp.java}

\subsection{Support for Greedy Assignment within the geost Kernel}\label{geostdescription:supportforgreedyassignmentwithinthegeostkernel}\hypertarget{geostdescription:supportforgreedyassignmentwithinthegeostkernel}{}

\subsubsection{Motivation and functionality description.}\label{geostdescription:motivationandfunctionalitydescription}\hypertarget{geostdescription:motivationandfunctionalitydescription}{}
Since, for performance reasons, the \texttt{geost} kernel offers a mode where he tries to fix all objects during one single propagation step, we provide a way to specify a preferred order on how to fix all the objects in one single propagation step. This is achieved by:
\begin{itemize}
	\item Fixing the objects according to the order they were passed to the \texttt{geost} kernel.
	\item When considering one object, fixing its shape variable as well as its coordinates:
	\begin{itemize}
		\item According to an order on these variables that can be explicitly specified.
		\item A value to assign that can either be the smallest or the largest value, also specified by the user.
	\end{itemize}
\end{itemize}

\begin{note}
Note that the use of the greedy mode assumes that no other constraint is present in the problem.
\end{note}

	This is encoded by a term that has exactly the same structure as the term associated to an object of  \texttt{geost}. The only difference consists of the fact that a variable is replaced by an expression \_ (\emph{The character \_ denotes the fact that the corresponding attribute is irrelevant, since for instance, we know that it is always fixed}), $\min(I)$ (respectively, $\max(I)$), where $I$ is a strictly positive integer. The meaning is that the corresponding variable should be fixed to its minimum (respectively maximum value) in the order $I$.   We can in fact give a list of vectors  $v_1,v_2,\ldots,v_p$ in order to specify how to fix objects $o_{(1+pa)},o_{(2+pa)},...,o_{(p+pa)}$.

This is illustrated by Figure bellow: for instance, Part(\emph{I}) specifies that we alternatively:
\begin{itemize}
	\item fix the shape variable of an object to its maximum value (i.e., by using max(1) ), fix the $x$-coordinate of an object to its its minimum value (i.e., by using min(2)), fix the $y$-coordinate of an object to its its minimum value (i.e., by using min(3)) and
	\item fix the shape variable of an object to its maximum value (i.e., by using max(1)), fix the $x$-coordinate of an object to its its maximum value (i.e., by using max(2)), fix the $y$-coordinate of an object to its its maximum value (i.e., by using max(3)).
\end{itemize}

In the example associated with Part (I) we successively fix objects $o_1, o_2, o_3, o_4, o_5, o_6$ by alternatively using strategies (1) \mylst{object(_,max(1),x[min(2),min(3)])} and (2) \mylst{object(_,max(1),x[max(2),max(3))}. 

\insertGraphique{\linewidth}{media/greedy.png}{Greedy placement}

\subsubsection{Implementation.}\label{geostdescription:implementation}\hypertarget{geostdescription:implementation}{}
The greedy algorithm for fixing an object o is controlled by a vector \emph{v} of length \emph{k+1} such that:
\begin{itemize}
	\item The shape variable \emph{o.sid} should be set to its minimum possible value if \emph{v}[0]$<0$, and to its maximum possible value otherwise.
	\item abs(\emph{v}[1])-2 is the most significant dimension (the one that varies the slowest) during the sweep. The values are tried in ascending order if \emph{v}[1]$<0$, and in descending order otherwise.
	\item abs(\emph{v}[2])-2 is the next most significant dimension, and its sign indicates the value order, and so on.
\end{itemize}

For example, a term \mylst{object(_,min(1),[max(3),min(4),max(2)])} is encoded as the vector $[-1,4,2,-3]$.
	                                      
\subsubsection{Second example.}\label{geostdescription:secondexample}\hypertarget{geostdescription:secondexample}{}
We will explain although a 2D example how to take into account of greedy mode. Consider we have 12 identical objects $o_0,o_1,\ldots,o_{11}$ having 4 potential shapes and we want to place them in a box $B$ (7x6) (see Figure bellow). Given that the placement of the objects should be totally inside B this means that the domain of the origins of objects are as follows $x\in [0,5]$, $y\in [0,4]$. Moreover, suppose that we want use two strategies when greedy algorithm is called: the term \mylst{object(_,min(1),[min(2),min(3)])} for objects $o_0, o_2, o_4, o_6, o_8, o_{10}$, and the term \mylst{object(_,max(1),[max(2),max(3)])} for objects $o_1, o_3, o_5, o_7, o_9, o_{11}$. These strategies are encoded respectively as $[-1,-2,-3]$ and $[1, 2, 3]$.

\insertGraphique{\linewidth}{media/exp2_geost.png}{A second Geost instance.}

We comment only the additional step w.r.t. the preceding example. In fact we just need to create the list of controlling vectors before creating the geost constraint. Each controlling vector is an array:
\begin{lstlisting}
  Vector <int[]> ctrlVs = new Vector<int[]>();
  int[] v0 = {-1, -2, -3};
  int[] v1 = {1, 2, 3};
  ctrlVs.add(v0);
  ctrlVs.add(v1);
\end{lstlisting}
and then create the \texttt{geost} constraint, by adding the list of controlling vectots as an another argument, as follows:
\begin{lstlisting}
  Constraint geost = Choco.geost(dim, objects, sb, ectr,ctrlVs);
  m.addConstraint(geost);
\end{lstlisting}

The full java code can be reached \href{media/zip/greddyexp.java}{here}.
\lstinputlisting{media/zip/greedyexp.java}

The placement obtained using the preceding strategies is displayed in the following figure (right side).

\insertGraphique{\linewidth}{media/solution_exp2.png}{A solution placement}





\part{Elements of Choco}\label{part:elements}\hypertarget{part:elements}{}
\chapter{Variables (Model)}\label{ch:vars}\hypertarget{ch:vars}{}
This section describes the three kinds of \hyperlink{model:variables}{variables} that can be used within a Choco Model, and an object-variable.
\section{Integer variables}\label{integervariable}\hypertarget{integervariable}{}
\texttt{IntegerVariable} is a variable whose associated domain is made of integer values. 

\subsubsection{constructors:}
      \noindent\begin{tabular}{p{.8\linewidth}p{.15\linewidth}}
        Choco method & return type \\
        \hline
        \mylst{makeIntVar(String name, int lowB, int uppB, String... options)} &\texttt{IntegerVariable}\\
		\mylst{makeIntVar(String name, List<Integer> values, String... options)} &\texttt{IntegerVariable}\\
		\mylst{makeIntVar(String name, int[] values, String... options)} &\texttt{IntegerVariable}\\
        \mylst{makeBooleanVar(String name, String... options)}  &\texttt{IntegerVariable}\\
        \mylst{makeIntVarArray(String name, int dim, int lowB, int uppB, String... options)} &\texttt{IntegerVariable[]}\\
        \mylst{makeIntVarArray(String name, int dim, int[] values, String... options)} &\texttt{IntegerVariable[]}\\
        \mylst{makeBooleanVarArray(String name, int dim, String... options)}  &\texttt{IntegerVariable[]}\\
        \mylst{makeIntVarArray(String name, int dim1, int dim2, int lowB, int uppB, String... options)}  &\texttt{IntegerVariable[][]}\\
        \mylst{makeIntVarArray(String name, int dim1, int dim2, int[] values, String... options)}  &\texttt{IntegerVariable[][]}\\
      \end{tabular}
% 	\begin{itemize}
% 		\item to create an \textbf{IntegerVariable} object:
% 		\begin{itemize}
% 			\item \mylst{makeIntVar(String name, int lowB, int uppB, String... options)}
% 			\item \mylst{makeIntVar(String name, List<Integer> values, String... options)}
% 			\item \mylst{makeIntVar(String name, int[] values, String... options)}
% 		\end{itemize}
% 		\item to create an \textbf{array of IntegerVariable} object:
% 		\begin{itemize}
% 			\item \mylst{makeIntVarArray(String name, int dim, int lowB, int uppB, String... options)}
% 			\item \mylst{makeIntVarArray(String name, int dim, int[] values, String... options)}
% 		\end{itemize}
% 		\item to create a \textbf{matrix of IntegerVariable} object:
% 		\begin{itemize}
% 			\item \mylst{makeIntVarArray(String name, int dim1, int dim2, int lowB, int uppB, String... options)}
% 			\item \mylst{makeIntVarArray(String name, int dim1, int dim2, int[] values, String... options)}
% 		\end{itemize}
% 		\item to create an \textbf{IntegerVariable} object with pre defined domain [0,1]:
% 		\begin{itemize}
% 			\item \mylst{makeBooleanVar(String name, String... options)}
% 		\end{itemize}
% 		\item to create an \textbf{array of IntegerVariable} object with pre defined domain [0,1]:
% 		\begin{itemize}
% 			\item \mylst{makeBooleanVarArray(String name, int dim, String... options)}
% 		\end{itemize}
% 	\end{itemize}
% 	\item \textbf{return type} : \texttt{IntegerVariable} \emph{or} \texttt{IntegerVariable[]} \emph{or} \texttt{IntegerVariable[][]}
\subsubsection{options:}
	\begin{itemize}
		\item \emph{no option} : equivalent to option \hyperlink{venum:venumoptions}{\tt Options.V\_ENUM}
		\item \hyperlink{venum:venumoptions}{\tt Options.V\_ENUM} : to force Solver to create enumerated domain for the variable.
		\item \hyperlink{vbound:vboundoptions}{\tt Options.V\_BOUND} : to force Solver to create bounded domain for the variable.
		\item \hyperlink{vlink:vlinkoptions}{\tt Options.V\_LINK} : to force Solver to create linked list domain for the variable.
		\item \hyperlink{vbtree:vbtreeoptions}{\tt Options.V\_BTREE} : to force Solver to create binary tree domain for the variable.
		\item \hyperlink{vblist:vblistoptions}{\tt Options.V\_BLIST} : to force Solver to create bipartite list domain for the variable.
		\item \hyperlink{vmakespan:vmakespanoptions}{\tt Options.V\_MAKEPSAN} : declare the current variable as makespan.
		\item \hyperlink{vnodecision:vnodecisionoptions}{\tt Options.V\_NO\_DECISION} : to force variable to be removed from the pool of decisional variables.
		\item \hyperlink{vobjective:vobjectiveoptions}{\tt Options.V\_OBJECTIVE} : to define the variable to be the one to optimize.
	\end{itemize}
\subsubsection{methods:}
      \begin{itemize}
      \item \mylst{removeVal(int val)}: remove value \emph{val} from the domain of the current variable
      \end{itemize}

A variable with $\{0,1\}$ domain is automatically considered as boolean domain.

\subsubsection{Example:}
\lstinputlisting{java/vintegervariable.j2t}

Integer variables are illustrated on the \hyperlink{model:example1:nqueenschoco}{n-Queens problem}. 

\section{Real variables}\label{realvariable}\hypertarget{realvariable}{}
\texttt{RealVariable} is a variable whose associated domain is made of real values. Only enumerated domain is available for real variables. 

Such domain are memory consuming. In order to minimize the memory use and to have the precision you need, the model offers a way to set a precision (default value is 1.0e-6):
\lstinputlisting{java/vprecision.j2t}

\subsubsection{constructor:}
      \noindent\begin{tabular}{p{.8\linewidth}p{.15\linewidth}}
        Choco method & return type \\
        \hline
        \mylst{makeRealVar(String name, double lowB, double uppB, String... options)} &\texttt{RealVariable}\\
      \end{tabular}
%	\begin{itemize}
%		\item to create a \textbf{RealVariable} object:
%		\begin{itemize}
%			\item \mylst{makeRealVar(String name, double lowB, double uppB, String... options)}
%		\end{itemize}
%	\end{itemize}
%	\item \textbf{return type} : \texttt{RealVariable}
\subsubsection{options:}
	\begin{itemize}
		\item \emph{no option} : no particular choice on decision or objective.
		\item \hyperlink{vnodecision:vnodecisionoptions}{\tt CPOptions.V\_NO\_DECISION} : to force variable to be removed from the pool of decisional variables.
		\item \hyperlink{vobjective:vobjectiveoptions}{\tt CPOptions.V\_OBJECTIVE} : to define the variable to be the one to optimize.
	\end{itemize}

\subsubsection{Example:}
\lstinputlisting{java/vrealvariable.j2t}

Real variables are illustrated on the \hyperlink{model:example3:thecyclohexaneproblemwithchoco}{CycloHexan problem}. 

\section{Set variables}\label{setvariable}\hypertarget{setvariable}{}
\texttt{SetVariable} is high level modeling tool. It allows to represent variable whose values are sets. A SetVariable on integer values between $[1,n]$ has $2*n$ values (every possible subsets of $\{1..n\}$). This makes an exponential number of values and the domain is represented with two bounds corresponding to the intersection of all possible sets (called the kernel) and the union of all possible sets (called the envelope) which are the possible candidate values for the variable. The consistency achieved on SetVariables is therefore a kind of bound consistency.

\subsubsection{constructors:}
      \noindent\begin{tabular}{p{.8\linewidth}p{.15\linewidth}}
        Choco method & return type \\
        \hline
        \mylst{makeSetVar(String name, int lowB, int uppB, String... options)} &\texttt{SetVariable}\\
        \mylst{makeSetVarArray(String name, int dim, int lowB, int uppB, String... options)} &\texttt{SetVariable[]}
      \end{tabular}
%	\begin{itemize}
%		\item to create an \textbf{SetVariable} object:
%		\begin{itemize}
%			\item \mylst{makeSetVar(String name, int lowB, int uppB, String... options)}
%		\end{itemize}
%		\item to create an \textbf{array of SetVariable} object:
%		\begin{itemize}
%			\item \mylst{makeSetVarArray(String name, int dim, int lowB, int uppB, String... options)}
%		\end{itemize}
%	\end{itemize}
%	\item \textbf{return type} : \texttt{SetVariable} \emph{or} \texttt{SetVariable[]}
\subsubsection{options:}
	\begin{itemize}
		\item \emph{no option} : equivalent to option \hyperlink{venum:venumoptions}{\tt Options.V\_ENUM}
		\item \hyperlink{venum:venumoptions}{\tt Options.V\_ENUM} : to force Solver to create \texttt{SetVariable} with enumerated domain for the caridinality variable.
		\item \hyperlink{vbound:vboundoptions}{\tt Options.V\_BOUND} : to force Solver to create \texttt{SetVariable} with bounded cardinality.
		\item \hyperlink{vnodecision:vnodecisionoptions}{\tt Options.V\_NO\_DECISION} : to force variable to be removed from the pool of decisional variables.
		\item \hyperlink{vobjective:vobjectiveoptions}{\tt Options.V\_OBJECTIVE} : to define the variable to be the one to optimize.
	\end{itemize}

The variable representing the cardinality can be accessed and constrained using method \texttt{getCard()} that returns an \hyperlink{integervariable}{\tt IntegerVariable} object.

\subsubsection{Example:}
\lstinputlisting{java/vsetvariable.j2t}

Set variables are illustrated on the \hyperlink{model:example2:ternarysteinerchoco}{ternary Steiner problem}. 



\section{Task variables}\label{taskvariable}\hypertarget{taskvariable}{}
\texttt{TaskVariable} is an object-variable composed of three \hyperlink{integervariable}{\tt IntegerVariable}: a starting time integer variable $start$, an ending time integer variable $end$ and a duration integer variable $duration$. To create a \texttt{TaskVariable}, one can creates the $start$, $end$ and $duration$ before, or indicates the earliest starting time ($int$), the latest completion time ($int$) and the duration({\tt int} or {\tt IntegerVariable}).

\subsubsection{constructors:}
      \noindent\begin{tabular}{p{.8\linewidth}p{.15\linewidth}}
        Choco method & return type \\
        \hline
        \mylst{makeTaskVar(String name, IntegerVariable start, IntegerVariable end, IntegerVariable duration, String... options)} &\texttt{TaskVariable}\\
		\mylst{makeTaskVar(String name, IntegerVariable start, IntegerVariable duration, String... options)} &\texttt{TaskVariable}\\
		\mylst{makeTaskVar(String name, int binf, int bsup, IntegerVariable duration, String... options)} &\texttt{TaskVariable}\\
        \mylst{makeTaskVar(String name, int binf, int bsup, int duration, String... options)}  &\texttt{TaskVariable}\\
        \mylst{makeTaskVar(String name, int bsup, IntegerVariable duration, String... options)} &\texttt{TaskVariable}\\
        \mylst{makeTaskVar(String name, int bsup, int duration, String... options)} &\texttt{TaskVariable}\\
        \mylst{makeTaskVarArray(String prefix, IntegerVariable[] starts, IntegerVariable[] ends, IntegerVariable[] durations, String... options)}  &\texttt{TaskVariable[]}\\
        \mylst{makeTaskVarArray(String name,  int binf, int bsup, IntegerVariable[] durations, String... options)}  &\texttt{TaskVariable[]}\\
        \mylst{makeTaskVarArray(String name,  int binf, int bsup, int[] durations, String... options)}  &\texttt{TaskVariable[][]}\\
        \mylst{makeTaskVarArray(String name,  int binf, int bsup, IntegerVariable[][] durations, String... options) }  &\texttt{TaskVariable[][]}\\
        \mylst{makeTaskVarArray(String name,  int binf, int bsup, int[][] durations, String... options)}  &\texttt{TaskVariable[][]}\\
      \end{tabular}

\subsubsection{options:}
Options are for the three \texttt{IntegerVariable}. See \hyperlink{integervariable}{\tt IntegerVariable} for more details about options.

\subsubsection{Example:}
\lstinputlisting{java/vtaskvariable.j2t} 

\chapter{Operators (Model)}\label{ch:operators}\hypertarget{ch:operators}{}
This section lists and details the \hyperlink{model:expressionvariables}{operators} that can be used within a Choco Model to combine variables in expressions.
\section{abs (operator)}\label{abs:absoperator}\hypertarget{abs:absoperator}{}
Returns an expression variable that represents the absolute value of the argument (\(|n|\)).

\begin{itemize}
	\item \textbf{API} : abs(IntegerExpressionVariable n)
	\item \textbf{return type} : IntegerExpressionVariable
	\item \textbf{options} : \emph{n/a}
	\item \textbf{favorite domain} : unknown
\end{itemize}

\textbf{Example}:
\begin{lstlisting}
	CPModel m = new CPModel();
	IntegerVariable x = makeIntVar("x", 1, 5, "cp:enum");
	IntegerVariable y = makeIntVar("y", -5, 5, "cp:enum");
	m.addConstraint(eq(abs(x), y));
	CPSolver s = new CPSolver();
	s.read(m);
	s.solve();
	Assert.assertEquals(s.getVar(x).getVal(),Math.abs(s.getVar(y).getVal()));
\end{lstlisting}

%%% Local Variables: 
%%% mode: latex
%%% TeX-master: t
%%% End: 

%\part{cos}
\label{cos}
\hypertarget{cos}{}

\section{cos (operator)}\label{cos:cosoperator}\hypertarget{cos:cosoperator}{}
Returns an expression variable corresponding to the cosinus value of the argument (\(cos(x)\)).

\begin{itemize}
	\item \textbf{API} : \mylst{cos(RealExpressionVariable exp)}
	\item \textbf{return type} : RealExpressionVariable
	\item \textbf{options} : \emph{n/a}
	\item \textbf{favorite domain} : real
\end{itemize}

\textbf{Example}:
\lstinputlisting{java/ocos.j2t}

%\part{div}
\label{div}
\hypertarget{div}{}

\section{div (operator)}\label{div:divoperator}\hypertarget{div:divoperator}{}
Returns an expression variable that represents the \(integer\) \(quotient\) of the division of the first argument variable by the second one (\(n_1/n_2\)).

\begin{itemize}
	\item \textbf{API} :
	\begin{itemize}
		\item \mylst{div(IntegerExpressionVariable n1, IntegerExpressionVariable n2)}
		\item \mylst{div(IntegerExpressionVariable n1, int n2)}
		\item \mylst{div(int n1, IntegerExpressionVariable n2)}
	\end{itemize}
	\item \textbf{return type} : IntegerExpressionVariable
	\item \textbf{options} : \emph{n/a}
	\item \textbf{favorite domain} : \emph{n/a}
\end{itemize}

\textbf{Example}:
\lstinputlisting{java/odiv.j2t}


\section{ifThenElse (operator)}\label{ifthenelse:ifthenelseoperator}\hypertarget{ifthenelse:ifthenelseoperator}{}
\texttt{ifThenElse}$(c,v_1,v_2)$ states that if the constraint $c$ is satisfied, it returns the second parameter $v_1$, otherwise it returns the third one $v_2$.

\begin{itemize}
	\item \textbf{API} : \mylst{ifThenElse(Constraint c, IntegerExpressionVariable v1, IntegerExpressionVariable v2)}
	\item \textbf{return type} : IntegerExpressionVariable
	\item \textbf{options} : \emph{n/a}
	\item \textbf{favorite domain} : unknown
\end{itemize}

\textbf{Example}:
\lstinputlisting{java/oifthenelse.j2t}

\section{max (operator)}\label{max:maxoperator}\hypertarget{max:maxoperator}{}
Returns an expression variable equals to the greater value of the argument (\(max(x_1, x_2, ..., x_n)\)).

\begin{itemize}
	\item \textbf{API} :
	\begin{itemize}
		\item max(IntegerExpressionVariable x1, IntegerExpressionVariable x2)
		\item max(int x1, IntegerExpressionVariable x2)
		\item max(IntegerExpressionVariable x1, int x2)
		\item max(IntegerExpressionVariable[] x)
	\end{itemize}
	\item \textbf{return type}: IntegerExpressionVariable
	\item \textbf{options} : \emph{n/a}
	\item \textbf{favorite domain} : \emph{to complete}
\end{itemize}

\textbf{Example}:
\begin{lstlisting}
	Model m = new CPModel();
	m.setDefaultExpressionDecomposition(true);
	IntegerVariable[] v = makeIntVarArray("v", 3, -3, 3);
	IntegerVariable minv = makeIntVar("min", -3, 3);
	Constraint c = eq(minv, max(v));
	m.addConstraint(c);
	Solver s = new CPSolver();
	s.read(m);
	s.solveAll();
\end{lstlisting}

%%% Local Variables: 
%%% mode: latex
%%% TeX-master: t
%%% End: 


\section{min (operator)}\label{min:minoperator}\hypertarget{min:minoperator}{}
Returns an expression variable equals to the smaller value of the argument (\(min(x_1, x_2, ..., x_n)\)).

\begin{itemize}
	\item \textbf{API} :
	\begin{itemize}
		\item \mylst{min(IntegerExpressionVariable x1, IntegerExpressionVariable x2)}
		\item \mylst{min(int x1, IntegerExpressionVariable x2)}
		\item \mylst{min(IntegerExpressionVariable x1, int x2)}
		\item \mylst{min(IntegerExpressionVariable[] x)}
	\end{itemize}
	\item \textbf{return type}: IntegerExpressionVariable
	\item \textbf{options} : \emph{n/a}
	\item \textbf{favorite domain} : \emph{to complete}
\end{itemize}

\textbf{Example}:
\lstinputlisting{java/omin.j2t}

%%% Local Variables: 
%%% mode: latex
%%% TeX-master: t
%%% End: 

%\part{minus}
\label{minus}
\hypertarget{minus}{}

\section{minus (operator)}\label{minus:minusoperator}\hypertarget{minus:minusoperator}{}
Returns an expression variable that corresponding to the difference between the two arguments (\(x-y\)).

\begin{itemize}
	\item \textbf{API} :
	\begin{itemize}
		\item \mylst{minus(IntegerExpressionVariable x, IntegerExpressionVariable y)}
		\item \mylst{minus(IntegerExpressionVariable x, int y)}
		\item \mylst{minus(int x, IntegerExpressionVariable y)}
		\item \mylst{minus(RealExpressionVariable x, RealExpressionVariable y)}
		\item \mylst{minus(RealExpressionVariable x, double y)}
		\item \mylst{minus(double x, RealExpressionVariable y)}
	\end{itemize}
	\item \textbf{return type} :
	\begin{itemize}
		\item \texttt{IntegerExpressionVariable}, if parameters are \texttt{IntegerExpressionVariable}
		\item \texttt{RealExpressionVariable}, if parameters are \texttt{RealExpressionVariable}
	\end{itemize}
	\item \textbf{options} : \emph{n/a}
	\item \textbf{favorite domain} : \emph{to complete}
\end{itemize}

\textbf{Example}
\begin{lstlisting}
	Model m = new CPModel();
	Solver s = new CPSolver();
	IntegerVariable a = makeIntVar("a", 0, 4);
	m.addConstraint(eq(minus(a, 1), 2));
	s.read(m);
	s.solve();
\end{lstlisting}


\section{mod (operator)}\label{mod:modoperator}\hypertarget{mod:modoperator}{}
Returns an expression variable that represents the integer remainder of the division of the first argument variable by the second one (\(x_1\%x_2\)).

\begin{itemize}
	\item \textbf{API}:
	\begin{itemize}
		\item mod(IntegerExpressionVariable x1, IntegerExpressionVariable x2)
		\item mod(int x1, IntegerExpressionVariable x2)
		\item mod(IntegerExpressionVariable x1, int x2)
	\end{itemize}
	\item \textbf{return type} : \texttt{IntegerExpressionVariable}
	\item \textbf{options} : \emph{n/a}
	\item \textbf{favorite domain} : \emph{n/a}
\end{itemize}

\textbf{Example}:
\lstinputlisting{java/omod.j2t}

%%% Local Variables: 
%%% mode: latex
%%% TeX-master: t
%%% End: 

%\part{mult}
\label{mult}
\hypertarget{mult}{}

\section{mult (operator)}\label{mult:multoperator}\hypertarget{mult:multoperator}{}
Returns an expression variable that corresponding to the product of variables in argument (\(x*y\)).

\begin{itemize}
	\item \textbf{API} :
	\begin{itemize}
		\item mult(IntegerExpressionVariable x, IntegerExpressionVariable y)
		\item mult(IntegerExpressionVariable x, int y)
		\item mult(int x, IntegerExpressionVariable y)
		\item mult(RealExpressionVariable x, RealExpressionVariable y)
		\item mult(RealExpressionVariable x, double y)
		\item mult(double x, RealExpressionVariable y)
	\end{itemize}
	\item \textbf{return type} :
	\begin{itemize}
		\item \texttt{IntegerExpressionVariable}, if parameters are \texttt{IntegerExpressionVariable}
		\item \texttt{RealExpressionVariable}, if parameters are \texttt{RealExpressionVariable}
	\end{itemize}
	\item \textbf{options} : \emph{n/a}
	\item \textbf{favorite domain} : \emph{to complete}
\end{itemize}

\textbf{Example}
\begin{lstlisting}
	CPModel m = new CPModel();
	IntegerVariable x = makeIntVar("x", -10, 10);
	IntegerVariable z = makeIntVar("z", -10, 10);
	IntegerVariable w = makeIntVar("w", -10, 10);
	m.addVariable(x, z, w);
	
	CPSolver s = new CPSolver();
	// x >= z * w
	Constraint exp = geq(x, mult(z,w));
	
	m.setDefaultExpressionDecomposition(true);
	m.addConstraint(exp);
	
	s.read(m);
	s.solveAll();
\end{lstlisting}

%\part{neg}
\label{neg}
\hypertarget{neg}{}

\section{neg (operator)}\label{neg:negoperator}\hypertarget{neg:negoperator}{}

Returns an expression variable that is the opposite of the expression integer variable in argument (\(-x\)).

\begin{itemize}
	\item \textbf{API} : neg(IntegerExpressionVariable x)
	\item \textbf{return type} : \texttt{IntegerExpressionVariable}
	\item \textbf{options} : \emph{n/a}
	\item \textbf{favorite domain} : \emph{n/a}
\end{itemize}

\textbf{Example}:
\begin{lstlisting}
	Model m = new CPModel();
	Solver s = new CPSolver();
	
	IntegerVariable x = makeIntVar("x", -10, 10);
	IntegerVariable w = makeIntVar("w", -10, 10);
	// -x = w - 20
	m.addConstraint(eq(neg(x), minus(w, 20)));
	
	s.read(m);
	s.solve();
\end{lstlisting}

%\part{plus}
\label{plus}
\hypertarget{plus}{}



\section{plus (operator)}\label{plus:plusoperator}\hypertarget{plus:plusoperator}{}
Returns an expression variable that corresponding to the sum of the two arguments (\(x+y\)).

\begin{itemize}
	\item \textbf{API} :
	\begin{itemize}
		\item \mylst{plus(IntegerExpressionVariable x, IntegerExpressionVariable y)}
		\item \mylst{plus(IntegerExpressionVariable x, int y)}
		\item \mylst{plus(int x, IntegerExpressionVariable y)}
		\item \mylst{plus(RealExpressionVariable x, RealExpressionVariable y)}
		\item \mylst{plus(RealExpressionVariable x, double y)}
		\item \mylst{plus(double x, RealExpressionVariable y)}
	\end{itemize}
	\item \textbf{return type} :
	\begin{itemize}
		\item \texttt{IntegerExpressionVariable}, if parameters are \texttt{IntegerExpressionVariable}
		\item \texttt{RealExpressionVariable}, if parameters are \texttt{RealExpressionVariable}
	\end{itemize}
	\item \textbf{options} : \emph{n/a}
	\item \textbf{favorite domain} : \emph{to complete}
\end{itemize}

\textbf{Example}
% \begin{itemize}
% 	\item example1:
% \end{itemize}

\lstinputlisting{java/oplus.j2t}


%\part{power}
\label{power}
\hypertarget{power}{}

\section{power (operator)}\label{power:poweroperator}\hypertarget{power:poweroperator}{}
Returns an expression variable that represents the first argument raised to the power of the second argument (\(x^y\)).

\begin{itemize}
	\item \textbf{API} :
	\begin{itemize}
		\item \mylst{power(IntegerExpressionVariable x, IntegerExpressionVariable y)}
		\item \mylst{power(int x, IntegerExpressionVariable y)}
		\item \mylst{power(IntegerExpressionVariable x, int y)}
		\item \mylst{power(RealExpressionVariable x, int y)}
	\end{itemize}
	\item \textbf{return type}:
	\begin{itemize}
		\item \texttt{IntegerExpressionVariable}, if parameters are \texttt{IntegerExpressionVariable}
		\item \texttt{RealExpressionVariable}, if parameters are \texttt{RealExpressionVariable}
	\end{itemize}
	\item \textbf{option} : \emph{n/a}
	\item \textbf{favorite domain} : \emph{to complete}
\end{itemize}

\textbf{Example} : 
\lstinputlisting{java/opower.j2t}

%\part{scalar}
\label{scalar}
\hypertarget{scalar}{}

\section{scalar (operator)}\label{scalar:scalaroperator}\hypertarget{scalar:scalaroperator}{}
Return an integer expression that corresponds to the scalar product of coefficients array and variables array (\(c_1*x_1+c_2*x_2+...+c_n*x_n\)).

\begin{itemize}
	\item \textbf{API} :
	\begin{itemize}
		\item scalar(int[] c, IntegerVariable[] x)
		\item scalar(IntegerVariable[] x, int[] c)
	\end{itemize}
	\item \textbf{return type} : IntegerExpressionVariable
	\item \textbf{options} : \emph{n/a}
	\item \textbf{favorite domain} : \emph{to complete}
\end{itemize}

\textbf{Example}:

\lstinputlisting{java/oscalar.j2t}

%\part{sin}
\label{sin}
\hypertarget{sin}{}

\section{sin (operator)}\label{sin:sinoperator}\hypertarget{sin:sinoperator}{}
Returns a real variable that corresponding to the sinus value of the argument (\(sin(x)\)).

\begin{itemize}
	\item \textbf{API} : sin(RealExpressionVariable exp)
	\item \textbf{return type} : RealExpressionVariable
	\item \textbf{options} : \emph{n/a}
	\item \textbf{favorite domain} : real
\end{itemize}

\textbf{Example}:

\emph{No valid example for the moment}

%\part{sum}
\label{sum}
\hypertarget{sum}{}

\section{sum (operator)}\label{sum:sumoperator}\hypertarget{sum:sumoperator}{}
Return an integer expression that corresponds to the sum of the variables given in argument (\(x_1+x_2+...+x_n\)).

\begin{itemize}
	\item \textbf{API}: sum(IntegerVariable... lv)
	\item \textbf{return type} : IntegerExpressionVariable
	\item \textbf{options} : \emph{n/a}
	\item \textbf{favorite domain} : \emph{to complete}
\end{itemize}

\textbf{Example} :
\lstinputlisting{java/osum.j2t}

%
\section{TRUE (operator)}\label{true:trueoperator}\hypertarget{true:trueoperator}{}
Returns an expression always equals to \emph{true}.

%%% Local Variables: 
%%% mode: latex
%%% TeX-master: t
%%% End: 

\chapter{Constraints (Model)}\label{ch:constraints}\hypertarget{ch:constraints}{}
This section lists and details the \hyperlink{model:constraints}{constraints} currently available in Choco.
%\part{abs}
%\label{abs}\hypertarget{abs}{}
\section{abs (constraint)}\label{abs:absconstraint}\hypertarget{abs:absconstraint}{}
\begin{notedef}
  \texttt{abs}$(x,y)$ states that $x$ is the absolute value of $y$:
$$x = |y|$$
\end{notedef}

\begin{itemize}
	\item \textbf{API} : \mylst{abs(IntegerVariable x, IntegerVariable y)}
	\item \textbf{return type} : \texttt{Constraint}
	\item \textbf{options} : \emph{n/a}
	\item \textbf{favorite domain} : enumerated
\end{itemize}

\textbf{Example}:
\begin{lstlisting}
	CPModel m = new CPModel();
	IntegerVariable x = makeIntVar("x", 1, 5, "cp:enum");
	IntegerVariable y = makeIntVar("y", -5, 5, "cp:enum");
	m.addConstraint(abs(x,y));
	CPSolver s = new CPSolver();
	s.read(m);
	s.solve();
	Assert.assertEquals(s.getVar(x).getVal(),Math.abs(s.getVar(y).getVal()));
\end{lstlisting}


%\part{alldifferent}
\label{alldifferent}
\hypertarget{alldifferent}{}

\section{allDifferent (constraint)}\label{alldifferent:alldifferentconstraint}\hypertarget{alldifferent:alldifferentconstraint}{}
\begin{notedef}
  \texttt{allDifferent}$(x_1,\ldots,x_n)$ states that the arguments have pairwise distinct values:
 $$x_i \neq x_j,\quad \forall\ i\neq j$$  
\end{notedef}
This constraint is useful for some matching problems.
Notice that the filtering algorithm used will depend on the nature (enumerated or bounded) of the variables: 
when \emph{enumerated}, the constraint refers to the alldifferent of \cite{ReginAAAI94};
when \emph{bounded}, a dedicated algorithm for bound propagation is used \cite{LopezIJCAI03}.

\begin{itemize}
	\item \textbf{API} :
	\begin{itemize}
		\item \mylst{allDifferent(IntegerVariable... x)}
		\item \mylst{allDifferent(String options, IntegerVariable... x)}
	\end{itemize}
	\item \textbf{return type} : \texttt{Constraint}
	\item \textbf{options} :
	\begin{itemize}
		\item \emph{no option} clever choice made on domains of given variables
		\item \texttt{CPOptions.C_ALLDIFFERENT_AC} for \cite{ReginAAAI94} implementation of arc consistency
		\item \texttt{CPOptions.C_ALLDIFFERENT_BC} for \cite{LopezIJCAI03} implementation of bound consistency
		\item \texttt{CPOptions.C_ALLDIFFERENT_CLIQUE} for propagating the clique of differences
	\end{itemize}
	\item \textbf{favorite domain} : depending of options.
	\item \textbf{references} :
      \begin{itemize}
      \item  \cite{ReginAAAI94}: \emph{A filtering algorithm for constraints of difference in CSPs}
      \item  \cite{LopezIJCAI03}: \emph{A fast and simple algorithm for bounds consistency of the alldifferent constraint}
      \item global constraint catalog: \href{http://www.emn.fr/x-info/sdemasse/gccat/Calldifferent.html}{\tt alldifferent}
      \end{itemize}
\end{itemize}



\textbf{Example}:
\lstinputlisting{java/calldifferent.j2t}

\label{among}
\hypertarget{among}{}

\section{among (constraint)}\label{among:amongconstraint}\hypertarget{among:amongconstraint}{}

% \subsection{among values}\label{among:amongvalues}\hypertarget{among:amongvalues}{}

% \begin{notedef}
%   \texttt{among}$(x, V)$ states that the variable $x$ takes its value in $V$:
%  $$x \subseteq V$$  
% \end{notedef}

% \begin{itemize}
% 	\item \textbf{API}: \mylst{among(IntegerVariable x, int[] v)}
% 	\item \textbf{return type}: \texttt{Constraint}
% \end{itemize}

% \textbf{Example}:
% \lstinputlisting{java/camong1.j2t}

%\subsection{among values with counter}\label{among:amongvaluescounter}\hypertarget{among:amongvaluescounter}{}

\begin{notedef}
\texttt{among}$(z, \collec{x_1}{x_n}, s)$ states that $z$ is the number of $x_i$ belonging to set $s$:
 $$ z = \vert\lbrace i\ |\ x_i \in s \rbrace\vert $$  
\end{notedef}

\begin{itemize}
\item \textbf{API}: 
  \begin{itemize}
  \item \mylst{among(IntegerVariable z, IntegerVariable[] x, int[] v)}
  \item \mylst{among(IntegerVariable z, IntegerVariable[] x, SetVariable s)}
  \end{itemize}
	\item \textbf{return type}: \texttt{Constraint}
	\item  \cite{Bessiere05ERCIM}: \emph{\textit{Among}, \textit{common} and \textit{disjoint} Constraints}
	\item  \cite{Bessiere06ERCIM}: \emph{\textit{Among}, \textit{common} and \textit{disjoint} Constraints}
	\item global constraint catalog: \href{http://www.emn.fr/x-info/sdemasse/gccat/Camong.html}{\tt among}
\end{itemize}

\textbf{Example}:

among with a collection of values
\lstinputlisting{java/camong2.j2t}

among with a set variable
\lstinputlisting{java/camong3.j2t}

%\part{and}
\label{and}
\hypertarget{and}{}

\section{and (constraint)}\label{and:andconstraint}\hypertarget{and:andconstraint}{}
\begin{notedef}
  \texttt{and}$(c_1,\ldots,c_n)$ states that every constraints in arguments are satisfied:
$$ c_1 \land c_2 \land\ldots\land c_n$$
\end{notedef}

\begin{itemize}
\item \textbf{API} : \mylst{and(Constraint... c)}
\item \textbf{return type} : \texttt{Constraint}
\item \textbf{options} : \emph{n/a}
\item \textbf{favorite domain} : \emph{n/a}
\item \textbf{references} :\\
  global constraint catalog: \href{http://www.emn.fr/x-info/sdemasse/gccat/Cand.html}{\tt and}
\end{itemize}

\textbf{Example}:
\lstinputlisting{java/cand.j2t}

%\part{atmostnvalue}
\label{atmostnvalue}
\hypertarget{atmostnvalue}{}

\section{atMostNValue (constraint)}\label{atmostnvalue:atmostnvalueconstraint}\hypertarget{atmostnvalue:atmostnvalueconstraint}{}
\begin{notedef}
\texttt{atMostNValue}$(x,z)$ states that the number of different values occurring in the array of variables $x$ is at most \emph{z}:
$$z\ge|\{x_1,\ldots,x_n\}|$$  
\end{notedef}

\begin{itemize}
	\item \textbf{API} : \mylst{atMostNValue(IntegerVariable[] x, IntegerVariable z)}
	\item \textbf{return type} : \texttt{Constraint}
	\item \textbf{options} : \emph{n/a}
	\item \textbf{favorite domain} : \emph{n/a}
	\item \textbf{references} :
      \begin{itemize}
      \item  \cite{BessiereCPAIOR05} \emph{Filtering algorithms for the NValue constraint}
      \item global constraint catalog: \href{http://www.emn.fr/x-info/sdemasse/gccat/Catmost_nvalue.html}{\tt atmost\_nvalue}
      \end{itemize}
    \end{itemize}

\textbf{Example}:
\lstinputlisting{java/catmostnvalue.j2t}

%\part{boolchanneling}
\label{boolchanneling}
\hypertarget{boolchanneling}{}

\section{boolChanneling (constraint)}\label{boolchanneling:boolchannelingconstraint}\hypertarget{boolchanneling:boolchannelingconstraint}{}
\begin{notedef}  
\texttt{boolChanneling}$(b,x,v)$ states that boolean $b$ is true if and only if $x$ has value $v$:
$$(b=1)\quad\iff\quad (x=v)$$ 
\end{notedef}

$b$ is an indicator variable acting as an observer of value $v$. 
See also \hyperlink{domainchanneling}{\texttt{domainChanneling}} for observing all the values of $x$.
%Imagine a bin packing problem where variable $x$ tells you on which a given bin object is placed. By stating the boolean channeling, $b$ is true if and only if the object is placed on bin $v$, the knapsack constraint for bin $v$ can then be easily stated as a scalar of the boolean variables.
\begin{itemize}
	\item \textbf{API} : \mylst{boolChanneling(IntegerVariable b, IntegerVariable x, int v)}
	\item \textbf{return type} : \texttt{Constraint}
	\item \textbf{options} : \emph{n/a}
	\item \textbf{favorite domain} : enumerated for $x$
\end{itemize}

\textbf{Example}:
\lstinputlisting{java/cboolchanneling.j2t}

\label{clause}
\hypertarget{clause}{}

\section{clause (constraint)}\label{clause:clauseconstraint}\hypertarget{clause:clauseconstraint}{}
\begin{notedef}
  \texttt{clause}$(\collec{b^+_1}{b^+_n}, \collec{b^-_1}{b^-_m})$ states that at least one boolean $b^+_i$ is true or one boolean $b^-_j$ is false.
% $$(x_1=1) \vee (x_2=1) \vee \dots \vee (x_n=1) \vee (y_1=0) \vee (y_2=0) \vee \dots \vee (y_m=0)$$  
 $$\bigvee_{i=1}^n (b^+_i=1) \vee \bigvee_{j=1}^m (b^-_j=0)$$
\end{notedef}

\begin{itemize}
	\item \textbf{API} :
	\begin{itemize}
		\item \mylst{clause(IntegerVariable[] bpos, IntegerVariable[] bneg)}
		\item \mylst{clause(String options, IntegerVariable[] bpos, IntegerVariable[] bneg)}
	\end{itemize}
	\item \textbf{return type} : \texttt{Constraint}
	\item \textbf{options} :
	\begin{itemize}
		\item \emph{no option} default filtering
		\item \hyperlink{cclause:cclauseoptions}{\tt Options.C\_CLAUSES\_ENTAIL} ensures quick entailment tests
	\end{itemize}
	\item \textbf{favorite domain} : \emph{n/a}.
	\item \textbf{references} : global constraint catalog: \href{http://www.emn.fr/x-info/sdemasse/gccat/Cclause_or.html}{\tt clause\_or}
\end{itemize}



\textbf{Example}:
\lstinputlisting{java/cclause.j2t}

\hypertarget{costregular}{}

\section{costRegular (constraint)}\label{costregular:costregularconstraint}\hypertarget{costregular:costregularconstraint}{}
\begin{notedef}
  \texttt{costRegular}$(z, \collec{x_1}{x_n},\mathcal{L}(\Pi), \coll{c_{i,j}})$ states that sequence \collec{x_1}{x_n} is a word belonging to the regular language $\mathcal{L}(\Pi)$ and that $z$ is its cost computed as the sum of the individual symbol weights $c_{i,x_i}$:
$$ \collec{x_1}{x_n} \in \mathcal{L}(\Pi)\quad\land\quad \sum_{i=1}^n c_{i,x_i} = z.$$
\end{notedef}

Like \hyperlink{regular}{\texttt{regular}}, this constraint is useful for modelling sequencing rules in personnel scheduling and rostering problems. Furthermore it allows to handle a linear counter (or cost) on the sequence. See \hyperlink{multicostregular}{\texttt{multiCostRegular}} to simultaneously handle several linear counters.

\texttt{costRegular} is the optimization variant of the \hyperlink{regular:regularconstraint}{\texttt{regular}} constraint. Enforcing GAC is NP-Hard, then the implemented algorithm~\cite{DemasseyC06} achieves an intermediate AC-BC level of consistency. Let $\mathcal{L}_x=\mathcal{L}(\Pi)\cap (D_1\times\ldots\times D_n)$ be the set of words of the language $\mathcal{L}(\Pi)$ that can be matched by \collec{x_1}{x_n} according to their current domains \collec{D_1}{D_n}, then:  
      \begin{itemize}
      \item Arc Consistency is enforced over $x$ regarding the language and the lower and upper bounds of $z$: for each value $v\in D_i$, there exists a word in $\mathcal{L}_x$, with $v$ as its $i$-th symbol and whose cost is between the bounds of $z$.
      \item Bound Consistency is enforced over $z$ regarding $x$ and the language: the lower and upper bounds of $z$ are set as the minimum and maximum costs of any words in $\mathcal{L}_x$.
      \end{itemize}
In summary, \texttt{costRegular}$(z, x,\mathcal{L}(\Pi), c)$ dominates its decomposition \hyperlink{regular}{\texttt{regular}}$(x,\mathcal{L}(\Pi))\land$\hyperlink{equation}{\texttt{equation}}$(z, x, c)$. Another decomposition proposed in~\cite{BeldiceanuC05} can easily be generated by introducing intermediary cost variables \collec{z_1}{z_n} and state variables \collec{q_0}{q_n}, then posting constraints in extension on each tuple $(q_{i-1},x_i,q_i,z_i)$ with the $\Pi$ transition table, and one linear sum $z=z_1+\cdots+z_n$. In terms of consistency, the two approaches are incomparable (words with costs out of the bounds of $z$ may not be filtered by the decomposition, see examples in~\cite{MenanaCPAIOR09}).  

Several API exists for defining the regular language:
\begin{itemize}
\item With a deterministic finite automaton (DFA) $\Pi$ weighted by a cost table $c$ with two dimensions, then $c[i][j]$ is the cost of any transition in $\Pi$ labeled by $j$ when processing the $i$-th symbol of a word: it models the cost of assigning variable $x_i$ to value $j$. The constraint ensures that $z=\sum_i c[i][x_i]$.
\item With a DFA $\Pi$ weighted by a cost table $c$ with three dimensions, then $c[i][j][s]$ is the cost of the transition in $\Pi$ outgoing from state $s$ and labeled by $j$ when processing the $i$-th symbol of a word: it models the cost of assigning variable $x_i$ to value $j$ if assignment sequence $\collec{x_1}{x_{i-1}}$ reaches state $s$ when processed by $\Pi$.  The constraint ensures that $z=\sum_i c[i][x_i][s_i]$ where $\collec{s_0, s_1}{s_n}$ is the sequence of states encountered when recognizing $\collec{x_1}{x_{n}}$ in $\Pi$.
\item With a weighted valued multi-graph $G(\Pi)$ and a node $s$, then $G(\Pi)$ must be a layered graph with $n+1$ layers and $s$ be the unique node in layer 0. Such a graph defines a valued DFA, by setting the arcs as the transitions, the arc values as the transition labels, the arc weights as the transition costs, the nodes as the states, and the nodes in the last layer as the accepting states. Note that this DFA recognizes only words of length $n$. The constraint ensures that $z$ is the total weight of the path in $G(\Pi)$ produced when recognizing \collec{x_1}{x_n}.  
\end{itemize}

Automaton $\Pi$ is encoded as an object of class \texttt{FiniteAutomaton}, whose API contains:
\begin{lstlisting}
  FiniteAutomaton();
  FiniteAutomaton(String regularExpression);
  int addState();
  void setInitialState(int state); 
  void setFinal(int state); 
  void addTransition(int state1, int state2, int.. labels);
  FiniteAutomaton union(FiniteAutomaton a);
  FiniteAutomaton intersection(FiniteAutomaton a);
  FiniteAutomaton complement();
  void minimize();
  int getNbStates();
  void toDotty(String dotFileName);
\end{lstlisting}

\begin{itemize}
	\item \textbf{API} :
	\begin{itemize}
		\item \mylst{costRegular(IntegerVariable z, IntegerVariable[] x, FiniteAutomaton pi, int[][] c)}
		\item \mylst{costRegular(IntegerVariable z, IntegerVariable[] x, FiniteAutomaton pi, int[][][] c)}
		\item \mylst{costRegular(IntegerVariable z, IntegerVariable[] x, DirectedMultigraph<Node,Arc> g, Node s)}
	\end{itemize}
	\item \textbf{return type} : \texttt{Constraint}
	\item \textbf{options} : \emph{n/a}
	\item \textbf{favorite domain} : \emph{n/a}
	\item \textbf{references} : \cite{DemasseyC06}: \emph{A \texttt{Cost-Regular} based hybrid column generation approach}
\end{itemize}



\textbf{Example}:

Build the \texttt{FiniteAutomaton} manually by adding states and transitions:
\lstinputlisting{java/ccostregular.j2t}

Build the \texttt{FiniteAutomaton} from a combination of several regular expressions: 
\lstinputlisting{java/ccostregular2.j2t}

%\part{cumulative}
\label{cumulative}
\hypertarget{cumulative}{}

\section{cumulative (constraint)}\label{cumulative:cumulativeconstraint}\hypertarget{cumulative:cumulativeconstraint}{}
\todo{to be cleaned.}

\begin{notedef}
  \texttt{cumulative(start,duration,height,capacity)} states that a set of tasks (defined by their starting times, finishing dates, durations and heights (or consumptions)) are executed on a cumulative resource of limited capacity. That is, the total height of the tasks which are executed at any time $t$ does not exceed the capacity of the resource:
$$\sum_{\{i\ |\ \mathtt{start}[i]\le t < \mathtt{start}[i]+\mathtt{duration}[i]\}} \mathtt{height}[i] \le \mathtt{capacity},\quad (\forall \text{ time } t)$$
\end{notedef}

The notion of task does not exist yet in Choco. The \texttt{cumulative} takes therefore as input, several arrays of integer variables (of same size $n$) denoting the starting, duration, and height of each task. When the array of finishing times is also specified, the constraint ensures that \texttt{start[i] + duration[i] = end[i]} for all task $i$.
As usual, a task is executed in the interval \texttt{[start,end-1]}.

For further informations, see the section devoted to this constraint in the Choco Tutorial document. 
%A tutorial on the use of this constraint is available \hyperlink{schedulinganduseofthecumulative:schedulinganduseofthecumulativeconstraint}{here}

\begin{itemize}
	\item \textbf{API} :
	\begin{itemize}
		\item \mylst{cumulative(IntegerVariable[] start, IntegerVariable[] end, IntegerVariable[] duration, IntegerVariable[] height, IntegerVariable capa, String... options)}
		\item \mylst{cumulative(IntegerVariable[] start, IntegerVariable[] end, IntegerVariable[] duration, int[] height, int capa, String... options)}
		\item \mylst{cumulative(IntegerVariable[] start, IntegerVariable[] duration, IntegerVariable[] height, IntegerVariable capa, String... options)}
	\end{itemize}
	\item \textbf{return type} : \texttt{Constraint}
	\item \textbf{options} :
	\begin{itemize}
		\item \emph{no option}
		\item \hyperlink{ccumulativeti:ccumulativetioptions}{SettingType.TASK\_INTERVAL.getOptionName()} for fast task intervals
		\item \hyperlink{ccumulativesti:ccumulativestioptions}{SettingType.SLOW\_TASK\_INTERVAL.getOptionName()} for slow task intervals
		\item \hyperlink{ccumulativecef:ccumulativecefoptions}{SettingType.VILIM\_CEF\_ALGO.getOptionName()} for Vilim theta lambda tree + lazy computation of the inner maximization of the edge finding rule of Van hentenrick and Mercier
		\item \hyperlink{ccumulativescef:ccumulativescefoptions}{SettingType.VHM\_CEF\_ALGO\_N2K.getOptionName()} for Simple $n^2 * k$ algorithm (lazy for R) (CalcEF -- Van Hentenrick)
	\end{itemize}
	\item \textbf{favorite domain} : \emph{n/a}
	\item \textbf{references} :
      \begin{itemize}
      \item  \cite{BeldiceanuCP02} \emph{A new multi-resource cumulatives constraint with negative heights}
      \item global constraint catalog: \href{http://www.emn.fr/x-info/sdemasse/gccat/Ccumulative.html}{\tt cumulative}
      \end{itemize}
\end{itemize}

\textbf{Example}:
\lstinputlisting{java/ccumulative.j2t} 

\section{cumulativeMax (constraint)}\label{cumulativemax:cumulativemaxconstraint}\hypertarget{cumulativemax:cumulativemaxconstraint}{}
Specific case of \hyperlink{cumulative:cumulativeconstraint}{Cumulative}, where the \textbf{capacity} is infinite.

\begin{itemize}
	\item \textbf{API} :
	\begin{itemize}
		\item \mylst{cumulativeMax(String name, TaskVariable[] tasks, IntegerVariable[] heights, IntegerVariable[] usages, IntegerVariable capacity, String... options)}
		\item \mylst{cumulativeMax(String name, TaskVariable[] tasks, IntegerVariable[] heights, IntegerVariable capacity, String... options)}
		\item \mylst{cumulativeMax(TaskVariable[] tasks, int[] heights, int capacity, String... options)}
	\end{itemize}
	\item \textbf{return type} : \texttt{Constraint}
	\item \textbf{options} :
	\begin{itemize}
		\item \emph{no option}
		\item \hyperlink{ccumulativeti:ccumulativetioptions}{SettingType.TASK\_INTERVAL.getOptionName()} for fast task intervals
		\item \hyperlink{ccumulativesti:ccumulativestioptions}{SettingType.SLOW\_TASK\_INTERVAL.getOptionName()} for slow task intervals
		\item \hyperlink{ccumulativecef:ccumulativecefoptions}{SettingType.VILIM\_CEF\_ALGO.getOptionName()} for Vilim theta lambda tree + lazy computation of the inner maximization of the edge finding rule of Van hentenrick and Mercier
		\item \hyperlink{ccumulativescef:ccumulativescefoptions}{SettingType.VHM\_CEF\_ALGO\_N2K.getOptionName()} for Simple $n^2 * k$ algorithm (lazy for R) (CalcEF -- Van Hentenrick)
	\end{itemize}
	\item \textbf{favorite domain} : \emph{n/a}
\end{itemize}
\section{cumulativeMin (constraint)}\label{cumulativemin:cumulativeminconstraint}\hypertarget{cumulativemin:cumulativeminconstraint}{}
Specific case of \hyperlink{cumulative:cumulativeconstraint}{Cumulative}, where the \textbf{consumption} is equal to 0.

\begin{itemize}
	\item \textbf{API} :
	\begin{itemize}
		\item \mylst{cumulativeMin(String name, TaskVariable[] tasks, IntegerVariable[] heights, IntegerVariable[] usages, IntegerVariable consumption, String... options)}
		\item \mylst{cumulativeMin(String name, TaskVariable[] tasks, IntegerVariable[] heights, IntegerVariable consumption, String... options)}
		\item \mylst{cumulativeMin(TaskVariable[] tasks, int[] heights, int consumption, String... options)}
	\end{itemize}
	\item \textbf{return type} : \texttt{Constraint}
	\item \textbf{options} :
	\begin{itemize}
		\item \emph{no option}
		\item \hyperlink{ccumulativeti:ccumulativetioptions}{SettingType.TASK\_INTERVAL.getOptionName()} for fast task intervals
		\item \hyperlink{ccumulativesti:ccumulativestioptions}{SettingType.SLOW\_TASK\_INTERVAL.getOptionName()} for slow task intervals
		\item \hyperlink{ccumulativecef:ccumulativecefoptions}{SettingType.VILIM\_CEF\_ALGO.getOptionName()} for Vilim theta lambda tree + lazy computation of the inner maximization of the edge finding rule of Van hentenrick and Mercier
		\item \hyperlink{ccumulativescef:ccumulativescefoptions}{SettingType.VHM\_CEF\_ALGO\_N2K.getOptionName()} for Simple $n^2 * k$ algorithm (lazy for R) (CalcEF -- Van Hentenrick)
	\end{itemize}
	\item \textbf{favorite domain} : \emph{n/a}
\end{itemize}
\label{disjoint}
\hypertarget{disjoint}{}

\section{disjoint (constraint)}\label{disjoint:disjointconstraint}\hypertarget{disjoint:disjointconstraint}{}

\begin{notedef}
  \texttt{disjoint}$(\collec{T^1_1}{T^1_n}, \collec{T^2_1}{T^2_m})$ states that each pair of tasks $(T^1_i, T^2_j)$ is in disjunction i.e., the processings of the two tasks do not overlap in time:
 $$T^1_i.end \le T^2_j.start\quad\lor\quad  T^2_j.end \le T^1_i.start,\quad \forall i=1..n, j=1..m $$  
\end{notedef}

CHOCO only provides a decomposition with reified precedences because the coloured cumulative is not available.

\begin{itemize}
	\item \textbf{API}: \mylst{disjoint(TaskVariable[] t1, TaskVariable[] t2)}
	\item \textbf{return type}: \texttt{Constraint[]}
	\item \textbf{favorite domain} : \emph{n/a}.
	\item \textbf{references} :
      \begin{itemize}
      \item global constraint catalog: \href{http://www.emn.fr/x-info/sdemasse/gccat/Cdisjoint_tasks.html}{\tt disjoint\_tasks}
      \end{itemize}
\end{itemize}
\textbf{Example}:
\lstinputlisting{java/cdisjoint2.j2t}

%\part{disjunctive}
\label{disjunctive}
\hypertarget{disjunctive}{}

\section{disjunctive (constraint)}\label{disjunctive:disjunctiveconstraint}\hypertarget{disjunctive:disjunctiveconstraint}{}

\begin{notedef}
  \texttt{disjunctive(start,duration)} states that a set of tasks (defined by their starting times and durations) are executed on a ddisjunctive resource, i.e. they do not overlap in time:
$$|\{i\ |\ \mathtt{start}[i]\le t < \mathtt{start}[i]+\mathtt{duration}[i]\}| \le 1,\quad (\forall \text{ time } t)$$
\end{notedef}

The notion of task does not exist yet in Choco. The \texttt{disjunctive} takes therefore as input arrays of integer variables (of same size $n$) denoting the starting and duration of each task. When the array of finishing times is also specified, the constraint ensures that \texttt{start[i] + duration[i] = end[i]} for all task $i$.
As usual, a task is executed in the interval \texttt{[start,end-1]}.

\begin{itemize}
	\item \textbf{API} :
	\begin{itemize}
		\item \mylst{disjunctive(IntegerVariable[] start, int[] duration, String...options)}
		\item \mylst{disjunctive(IntegerVariable[] start, IntegerVariable[] duration, String... options)}
		\item \mylst{disjunctive(IntegerVariable[] start, IntegerVariable[] end, IntegerVariable[] duration, String... options)}
		\item \mylst{disjunctive(IntegerVariable[] start, IntegerVariable[] end, IntegerVariable[] duration, IntegerVariable uppBound, String... options)}
	\end{itemize}
	\item \textbf{return type} : \texttt{Constraint}
	\item \textbf{options} :\emph{n/a}
	\item \textbf{favorite domain} : \emph{to complete}
	\item \textbf{references} :\\
      global constraint catalog: \href{http://www.emn.fr/x-info/sdemasse/gccat/Cdisjunctive.html}{\tt disjunctive}
\end{itemize}

\textbf{Example}:
%\lstinputlisting{java/cdisjunctive.j2t}
\mylst{//TODO: complete} 

%\part{distanceeq}
\label{distanceeq}
\hypertarget{distanceeq}{}

\section{distanceEQ (constraint)}\label{distanceeq:distanceeqconstraint}\hypertarget{distanceeq:distanceeqconstraint}{}
\begin{notedef}
  \texttt{distanceEQ}$(x_1,x_2,x_3,c)$ states that $x_3$ plus an offset $c$ (by default $c=0$) is equal to the distance between $x_1$ and $x_2$:
$$ x_3 + c = | x_1 - x_2 |$$
\end{notedef}

\begin{itemize}
	\item \textbf{API} :
	\begin{itemize}
		\item \mylst{distanceEQ(IntegerVariable x1, IntegerVariable x2, int x3)}
		\item \mylst{distanceEQ(IntegerVariable x1, IntegerVariable x2, IntegerVariable x3)}
		\item \mylst{distanceEQ(IntegerVariable x1, IntegerVariable x2, IntegerVariable x3, int c)}
	\end{itemize}
	\item \textbf{return type}: \texttt{Constraint}
	\item \textbf{options} : \emph{n/a}
	\item \textbf{favorite domain} : \emph{to complete}
	\item \textbf{references} :\\
      global constraint catalog: \href{http://www.emn.fr/x-info/sdemasse/gccat/Call_min_dist.html}{\tt all\_min\_dist} (variant)
\end{itemize}

\textbf{Example}:
\lstinputlisting{java/cdistanceeq.j2t}

%\part{distancegt}
\label{distancegt}
\hypertarget{distancegt}{}

\section{distanceGT (constraint)}\label{distancegt:distancegtconstraint}\hypertarget{distancegt:distancegtconstraint}{}
\begin{notedef}
  \texttt{distanceGT}$(x_1,x_2,x_3,c)$ states that $x_3$ plus an offset $c$ (by default $c=0$) is strictly greater than the distance between $x_1$ and $x_2$:
$$ x_3 + c > | x_1 - x_2 |$$
\end{notedef}

\begin{itemize}
	\item \textbf{API} :
	\begin{itemize}
		\item \mylst{distanceGT(IntegerVariable x1, IntegerVariable x2, int x3)}
		\item \mylst{distanceGT(IntegerVariable x1, IntegerVariable x2, IntegerVariable x3)}
		\item \mylst{distanceGT(IntegerVariable x1, IntegerVariable x2, IntegerVariable x3, int c)}
	\end{itemize}
	\item \textbf{return type}: \texttt{Constraint}
	\item \textbf{options} : \emph{n/a}
	\item \textbf{favorite domain} : \emph{to complete}
	\item \textbf{references} :\\
      global constraint catalog: \href{http://www.emn.fr/x-info/sdemasse/gccat/Call_min_dist.html}{\tt all\_min\_dist} (variant)
\end{itemize}

\textbf{Example}:
\begin{lstlisting}
	Model m = new CPModel();
	Solver s = new CPSolver();
	IntegerVariable v0 = makeIntVar("v0", 0, 5);
	IntegerVariable v1 = makeIntVar("v1", 0, 5);
	IntegerVariable v2 = makeIntVar("v2", 0, 5);
	m.addConstraint(distanceGT(v0, v1, v2, 0));
	s.read(m);
	s.solveAll();
\end{lstlisting}

%\part{distancelt}
\label{distancelt}
\hypertarget{distancelt}{}

\section{distanceLT (constraint)}\label{distancelt:distanceltconstraint}\hypertarget{distancelt:distanceltconstraint}{}
\begin{notedef}
  \texttt{distanceLT}$(x_1,x_2,x_3,c)$ states that $x_3$ plus an offset $c$ (by default $c=0$) is strictly smaller than the distance between $x_1$ and $x_2$:
$$ x_3 + c < | x_1 - x_2 |$$
\end{notedef}

\begin{itemize}
	\item \textbf{API} :
	\begin{itemize}
		\item \mylst{distanceLT(IntegerVariable x1, IntegerVariable x2, int x3)}
		\item \mylst{distanceLT(IntegerVariable x1, IntegerVariable x2, IntegerVariable x3)}
		\item \mylst{distanceLT(IntegerVariable x1, IntegerVariable x2, IntegerVariable x3, int c)}
	\end{itemize}
	\item \textbf{return type}: \texttt{Constraint}
	\item \textbf{options} : \emph{n/a}
	\item \textbf{favorite domain} : \emph{to complete}
\end{itemize}

\textbf{Example}:
\begin{lstlisting}
	Model m = new CPModel();
	Solver s = new CPSolver();
	IntegerVariable v0 = makeIntVar("v0", 0, 5);
	IntegerVariable v1 = makeIntVar("v1", 0, 5);
	IntegerVariable v2 = makeIntVar("v2", 0, 5);
	m.addConstraint(distanceLT(v0, v1, v2, 0));
	s.read(m);
	s.solveAll();
\end{lstlisting}

%\part{distanceneq}
\label{distanceneq}
\hypertarget{distanceneq}{}

\section{distanceNEQ (constraint)}\label{distanceneq:distanceneqconstraint}\hypertarget{distanceneq:distanceneqconstraint}{}
\begin{notedef}
  \texttt{distanceNEQ}$(x_1,x_2,x_3,c)$ states that $x_3$ plus an offset $c$ (by default $c=0$) is not equal to the distance between $x_1$ and $x_2$:
$$ x_3 + c \neq | x_1 - x_2 |$$
\end{notedef}

\begin{itemize}
	\item \textbf{API} :
	\begin{itemize}
		\item \mylst{distanceNEQ(IntegerVariable x1, IntegerVariable x2, int x3)}
		\item \mylst{distanceNEQ(IntegerVariable x1, IntegerVariable x2, IntegerVariable x3)}
		\item \mylst{distanceNEQ(IntegerVariable x1, IntegerVariable x2, IntegerVariable x3, int c)}
	\end{itemize}
	\item \textbf{return type}: \texttt{Constraint}
	\item \textbf{options} : \emph{n/a}
	\item \textbf{favorite domain} : \emph{to complete}
	\item \textbf{references} :\\
      global constraint catalog: \href{http://www.emn.fr/x-info/sdemasse/gccat/Call_min_dist.html}{\tt all\_min\_dist} (variant)
\end{itemize}

\textbf{Example}:
\begin{lstlisting}
	Model m = new CPModel();
	Solver s = new CPSolver();
	IntegerVariable v0 = makeIntVar("v0", 0, 5);
	IntegerVariable v1 = makeIntVar("v1", 0, 5);
	IntegerVariable v2 = makeIntVar("v2", 0, 5);
	m.addConstraint(distanceNEQ(v0, v1, v2, 0));
	s.read(m);
	s.solveAll();
\end{lstlisting}

%\part{domainchanneling}
\label{domainchanneling}
\hypertarget{domainchanneling}{}

\section{domainChanneling (constraint)}\label{domainchanneling:domainchannelingconstraint}\hypertarget{domainchanneling:domainchannelingconstraint}{}
\begin{notedef}  
\texttt{domainChanneling}$(x, \collec{b_1}{b_n})$ states that boolean $b_j$ is true if and only if $x$ has value $j$:
$$b_j=1\quad\iff\quad x=j,\qquad\forall j=1..n$$ 
\end{notedef}

It makes the link between a domain variable $x$ and those 0-1 variables $b$ that are associated with each potential value of $x$: the 0-1 variable $b_j$ associated with the value $j$ taken by $x$ is equal to 1, while the remaining 0-1 variables $b_i$ ($i\neq j$) are all equal to 0.

\begin{itemize}
	\item \textbf{API} : \mylst{domainChanneling(IntegerVariable x, IntegerVariable[] b)}
	\item \textbf{return type} : \texttt{Constraint}
	\item \textbf{options} : \emph{n/a}
	\item \textbf{favorite domain} : enumerated for $x$
	\item \textbf{references} :\\
	  global constraint catalog: \href{http://www.emn.fr/x-info/sdemasse/gccat/Cdomain_constraint.html}{\tt domain\_constraint}
\end{itemize}

\textbf{Example}:
\lstinputlisting{java/cdomainchanneling.j2t}

\section{element (constraint)}\label{element:elementconstraint}\hypertarget{element:elementconstraint}{}
See \hyperlink{nth:nthconstraint}{nth}.

\section{endsAfter (constraint)}\label{endsafter:endsafterconstraint}\hypertarget{endsafter:endsafterconstraint}{}
\begin{notedef}
  \texttt{endsAfter}$(t,c)$ states that the task variable $t$ ends after $c$:
  $$e_t \ge c$$
\end{notedef}

\begin{itemize}
	\item \textbf{API} :\mylst{endsAfter(TaskVariable t, int c) }
	\item \textbf{return type} : \texttt{Constraint}
	\item \textbf{options} : \emph{n/a}
	\item \textbf{favorite domain} : \emph{n/a}.
\end{itemize}

\textbf{Examples:}
%\lstinputlisting{java/ceq1.j2t}
\emph{to complete}
\section{endsAfterBegin (constraint)}\label{endsafterbegin:endsafterbeginconstraint}\hypertarget{endsafterbegin:endsafterbeginconstraint}{}

\begin{notedef}
\texttt{endsAfterBegin}$(t_1,t_2,c)$ states that the task variable $t_1$ ends after the beginning of $t_2$, minus $c$:
  $$e_{t_{1}} \ge s_{t_{2}} - c$$
\end{notedef}

\begin{itemize}
	\item \textbf{API} :\mylst{endsAfterBegin(TaskVariable t1, TaskVariable t2, int c)}
	\item \textbf{return type} : \texttt{Constraint}
	\item \textbf{options} : \emph{n/a}
	\item \textbf{favorite domain} : \emph{n/a}.
\end{itemize}

\textbf{Examples:}
%\lstinputlisting{java/ceq1.j2t}
\emph{to complete}

\section{endsAfterEnd (constraint)}\label{endsafterend:endsafterendconstraint}\hypertarget{endsafterend:endsafterendconstraint}{}
\begin{notedef}
\texttt{endsAfterEnd}$(T_1,T_2,c)$  states that task $T_1$ ends after the end time of $T_2$ minus $c$:
  $$T_{1}.end \ge T_{2}.end - c$$
\end{notedef}

\begin{itemize}
	\item \textbf{API} :\mylst{endsAfterEnd(TaskVariable t1, TaskVariable t2, int c)}
	\item \textbf{return type} : \texttt{Constraint}
	\item \textbf{options} : \emph{n/a}
	\item \textbf{favorite domain} : \emph{n/a}.
\end{itemize}

\textbf{Examples:}
%\lstinputlisting{java/ceq1.j2t}
\emph{to complete}

\section{endsBefore (constraint)}\label{endsbefore:endsbeforeconstraint}\hypertarget{endsbefore:endsbeforeconstraint}{}
\begin{notedef}
  \texttt{endsBefore}$(t,c)$ states that the task variable $t$ ends before $c$:
  $$e_t \le c$$
\end{notedef}

\begin{itemize}
	\item \textbf{API} :\mylst{endsBefore(final TaskVariable t, final int c) }
	\item \textbf{return type} : \texttt{Constraint}
	\item \textbf{options} : \emph{n/a}
	\item \textbf{favorite domain} : \emph{n/a}.
\end{itemize}

\textbf{Examples:}
%\lstinputlisting{java/ceq1.j2t}
\emph{to complete}

\section{endsBeforeBegin (constraint)}\label{endsbeforebegin:endsbeforebeginconstraint}\hypertarget{endsbeforebegin:endsbeforebeginconstraint}{}
\begin{notedef}
\texttt{endsBeforeBegin}$(t_1,t_2,c)$ states that the task variable $t_1$ ends before the beginning of $t_2$, minus $c$:
  $$e_{t_{1}} \le s_{t_{2}} - c$$
\end{notedef}

\begin{itemize}
	\item \textbf{API} :\mylst{endsBeforeBegin(TaskVariable t1, TaskVariable t2, int c)}
	\item \textbf{return type} : \texttt{Constraint}
	\item \textbf{options} : \emph{n/a}
	\item \textbf{favorite domain} : \emph{n/a}.
\end{itemize}

\textbf{Examples:}
%\lstinputlisting{java/ceq1.j2t}
\emph{to complete}

\section{endsBeforeEnd (constraint)}\label{endsbeforeend:endsbeforeendconstraint}\hypertarget{endsbeforeend:endsbeforeendconstraint}{}
\begin{notedef}
  \texttt{endsBeforeEnd}$(t_1,t_2,c)$ states that the task variable $t_1$ ends before the end of $t_2$, minus $c$:
  $$e_{t_{1}} \le e_{t_{2}} - c$$
\end{notedef}

\begin{itemize}
	\item \textbf{API} :\mylst{endsBeforeEnd(TaskVariable t1, TaskVariable t2, int c)}
	\item \textbf{return type} : \texttt{Constraint}
	\item \textbf{options} : \emph{n/a}
	\item \textbf{favorite domain} : \emph{n/a}.
\end{itemize}

\textbf{Examples:}
%\lstinputlisting{java/ceq1.j2t}
\emph{to complete}

\section{endsBetween (constraint)}\label{endsbetween:endsbetweenconstraint}\hypertarget{endsbetween:endsbetweenconstraint}{}
\begin{notedef}
  \texttt{endsBetween}$(t,c_1, c_2)$ states that the task variable $t$ ends between $c_1$ and $c_2$:
  $$c_1 \le e_t \le c_2$$
\end{notedef}

\begin{itemize}
	\item \textbf{API} :\mylst{endsBetween(TaskVariable t, int min, int max)}
	\item \textbf{return type} : \texttt{Constraint}
	\item \textbf{options} : \emph{n/a}
	\item \textbf{favorite domain} : \emph{n/a}.
\end{itemize}

\textbf{Examples:}
%\lstinputlisting{java/ceq1.j2t}
\emph{to complete}

%\part{eq}
\label{eq}
\hypertarget{eq}{}

\section{eq (constraint)}\label{eq:eqconstraint}\hypertarget{eq:eqconstraint}{}
\begin{notedef}
  \texttt{eq}$(x,y)$ states that the two arguments are equal:
$$x = y$$
\end{notedef}

\begin{itemize}
	\item \textbf{API} :
	\begin{itemize}
		\item \mylst{eq(IntegerExpressionVariable x, IntegerExpressionVariable y)}
		\item \mylst{eq(IntegerExpressionVariable x, int y)}
		\item \mylst{eq(int x, IntegerExpressionVariable y)}
		\item \mylst{eq(SetVariable x, SetVariable y)}
		\item \mylst{eq(RealExpressionVariable x, RealExpressionVariable y)}
		\item \mylst{eq(RealExpressionVariable x, double y)}
		\item \mylst{eq(double x, RealExpressionVariable y)}
		\item \mylst{eq(IntegerVariable x, RealVariable y)}
		\item \mylst{eq(RealVariable x, IntegerVariable y)}
	\end{itemize}
	\item \textbf{return type} : \texttt{Constraint}
	\item \textbf{options} : \emph{n/a}
	\item \textbf{favorite domain} : \emph{to complete}.
	\item \textbf{references} :\\
      global constraint catalog: \href{http://www.emn.fr/x-info/sdemasse/gccat/Ceq.html}{\tt eq} (on domain variables) and \href{http://www.emn.fr/x-info/sdemasse/gccat/Ceq_set.html}{\tt eq\_set} (on set variables). 
\end{itemize}

\textbf{Examples:}
\begin{itemize}
	\item example1:
\end{itemize}
\lstinputlisting{java/ceq1.j2t}
\begin{itemize}
	\item example2
\end{itemize}

\lstinputlisting{java/ceq2.j2t}

%\part{eqcard}
\label{eqcard}
\hypertarget{eqcard}{}

\section{eqCard (constraint)}\label{eqcard:eqcardconstraint}\hypertarget{eqcard:eqcardconstraint}{}
\begin{notedef}
  \texttt{eqCard}$(s,x)$ states that the cardinality of set $s$ is equal to $x$:
$$|s| = x$$
\end{notedef}

\begin{itemize}
	\item \textbf{API} :
	\begin{itemize}
		\item \mylst{eqCard(SetVariable s, IntegerVariable x)}
		\item \mylst{eqCard(SetVariable s, int x)}
	\end{itemize}
	\item \textbf{return type} : \texttt{Constraint}
	\item \textbf{options} : \emph{n/a}
	\item \textbf{favorite domain} : \emph{to complete}
\end{itemize}

\textbf{Example}:
\lstinputlisting{java/ceqcard.j2t}

%\part{equation}
\label{equation}
\hypertarget{equation}{}

\section{equation (constraint)}\label{equation:equationconstraint}\hypertarget{equation:equationconstraint}{}
\begin{notedef}
  \texttt{equation}$(z, \collec{x_1}{x_n},\collec{c_1}{c_n})$ states that $z$ is the weighted sum of $x$ by $c$:
$$c_1x_1+c_2x_2+...+c_nx_n = z$$
\end{notedef}
See also \hyperlink{knapsackproblem:knapsackproblemconstraint}{knapsackProblem}.

\begin{itemize}
	\item \textbf{API} :
	\begin{itemize}
		\item equation(int z, IntegerVariable[] x, int[] c)
		\item equation(String option, int z, IntegerVariable[] x, int[] c)
		\item equation(IntegerVariable z, IntegerVariable[] x, int[] c)
		\item equation(String option, IntegerVariable z, IntegerVariable[] x, int[] c)
	\end{itemize}
	\item \textbf{return type} : \texttt{Constraint}
	\item \textbf{options} :
	\begin{itemize}
		\item \emph{no option} or \texttt{"cp:ac"}: to enforce GAC  using \hyperlink{regular}{\texttt{regular}}
		\item \texttt{"cp:bc"}: to enforce bound consistency using \hyperlink{eq}{\texttt{eq}}$(z,$\hyperlink{scalar}{\texttt{scalar}}$(x,c))$
	\end{itemize}
	\item \textbf{favorite domain} : \emph{to complete}
	\item \textbf{global constraint catalog} : \href{http://www.emn.fr/z-info/sdemasse/gccat/Cscalar_product.html}{scalar\_product}
\end{itemize}

\textbf{Example}:
\lstinputlisting{java/cequation.j2t}

\section{exactly (constraint)}\label{exactly:exactlyconstraint}\hypertarget{exactly:exactlyconstraint}{}
\emph{To complete}

%\part{false}
%\label{false}\hypertarget{false}{}
\section{FALSE (constraint)}\label{false:falseconstraint}\hypertarget{false:falseconstraint}{}
\(FALSE\) always returns \emph{false}.

%\part{feaspairac}
\label{feaspairac}
\hypertarget{feaspairac}{}

\section{feasPairAC (constraint)}\label{feaspairac:feaspairacconstraint}\hypertarget{feaspairac:feaspairacconstraint}{}

\begin{notedef}
  \texttt{feasPairAC}$(x,y,feasTuples)$ states an extensional binary constraint on $(x,y)$ defined by the table $feasTuples$ of compatible pairs of values, and then enforces arc consistency. Two APIs are available to define the compatible pairs:
\begin{itemize}
	\item if $feasTuples$ is encoded as a list of pairs \texttt{List<int[2]>}, then:
      $$\exists \text{ tuple } i\ |\quad (x,y)=feasTuples[i]$$
	\item if $feasTuples$ is encoded as a boolean matrix \texttt{boolean[][]}, let $\underline{x}$ and  $\underline{y}$ be the initial minimum values of $x$ and $y$, then:
      $$\exists (u,v)\ |\quad (x,y)=(u+\underline{x},v+\underline{y})\ \land\ feasTuples[u][v]$$
\end{itemize}
\end{notedef}

The two APIs are duplicated to allow definition of options. 
\begin{itemize}
	\item \textbf{API} :
	\begin{itemize}
		\item \mylst{feasPairAC(IntegerVariable x, IntegerVariable y, List<int[]> feasTuples)}
		\item \mylst{feasPairAC(String options, IntegerVariable x, IntegerVariable y, List<int[]> feasTuples)}
		\item \mylst{feasPairAC(IntegerVariable x, IntegerVariable y, boolean[][] feasTuples)}
		\item \mylst{feasPairAC(String options, IntegerVariable x, IntegerVariable y, boolean[][] feasTuples)}
	\end{itemize}
	\item \textbf{return type} : \texttt{Constraint}
	\item \textbf{options} :
	\begin{itemize}
		\item \emph{no option}: use AC3 (default arc consistency)
		\item \texttt{cp:ac3}: to get AC3 algorithm (searching from scratch for supports on all values)
		\item \texttt{cp:ac2001}: to get AC2001 algorithm (maintaining the current support of each value)
		\item \texttt{cp:ac32}: to get AC3rm algorithm (maintaining the current support of each value in a non backtrackable way)
		\item \texttt{cp:ac322}: to get AC3 with the used of \texttt{BitSet} to know if a support still exists
	\end{itemize}
	\item \textbf{favorite domain} : \emph{to complete}
	\item \textbf{references} :\\
      global constraint catalog: \href{http://www.emn.fr/x-info/sdemasse/gccat/Celem.html}{elem}
\end{itemize}



\textbf{Example}:
\begin{lstlisting}
	Model m = new CPModel();
	Solver s = new CPSolver();
	
	ArrayList couples2 = new ArrayList();
	couples2.add(new int[]{1, 2});
	couples2.add(new int[]{1, 3});
	couples2.add(new int[]{2, 1});
	couples2.add(new int[]{3, 1});
	couples2.add(new int[]{4, 1});
	
	IntegerVariable v1 = makeIntVar("v1", 1, 4);
	IntegerVariable v2 = makeIntVar("v2", 1, 4);
	m.addConstraint(feasPairAC("cp:ac32",v1, v2, couples2));
	s.read(m);
	s.solveAll();
\end{lstlisting} 

%\part{feastupleac}
\label{feastupleac}
\hypertarget{feastupleac}{}

\section{feasTupleAC (constraint)}\label{feastupleac:feastupleacconstraint}\hypertarget{feastupleac:feastupleacconstraint}{}
\begin{notedef}
  \texttt{feasTupleAC}$(x,feasTuples)$ states an extensional constraint on $(x_1,\ldots,x_n)$ defined by the table $feasTuples$ of compatible tuples of values, and then enforces arc consistency:
      $$\exists \text{ tuple } i\ |\quad (x_1,\ldots,x_n)=feasTuples[i]$$
\end{notedef}

The API is duplicated to define options.
\begin{itemize}
	\item \textbf{API} :
	\begin{itemize}
		\item \mylst{feasTupleAC(List<int[]> feasTuples, IntegerVariable... x)}
		\item \mylst{feasTupleAC(String options, List<int[]> feasTuples, IntegerVariable... x)}
	\end{itemize}
	\item \textbf{return type}: \texttt{Constraint}
	\item \textbf{options} :
	\begin{itemize}
		\item \emph{no option}: use AC32 (default arc consistency)
		\item \texttt{cp:ac32}: to get AC3rm algorithm (maintaining the current support of each value in a non backtrackable way)
		\item \texttt{cp:ac2001}: to get AC2001 algorithm (maintaining the current support of each value)
		\item \texttt{cp:ac2008}: to get AC2008 algorithm (maintained by STR)
	\end{itemize}
	\item \textbf{favorite domain} : \emph{to complete}
	\item \textbf{references} :\\
      global constraint catalog: \href{http://www.emn.fr/x-info/sdemasse/gccat/Cin_relation.html}{in\_relation}
\end{itemize}

\textbf{Example}:
\begin{lstlisting}
	Model m = new CPModel();
	Solver s = new CPSolver();
	IntegerVariable v1 = makeIntVar("v1", 0, 2);
	IntegerVariable v2 = makeIntVar("v2", 0, 4);
	
	ArrayList feasTuple = new ArrayList();
	feasTuple.add(new int[]{1, 1}); // x*y = 1
	feasTuple.add(new int[]{2, 4}); // x*y = 1
	
	m.addConstraint(feasTupleAC("cp:ac2001",feasTuple, new IntegerVariable[]{v1, v2}));
	
	s.read(m);
	s.solve();
\end{lstlisting}

%\part{feastuplefc}
\label{feastuplefc}
\hypertarget{feastuplefc}{}

\section{feasTupleFC (constraint)}\label{feastuplefc:feastuplefcconstraint}\hypertarget{feastuplefc:feastuplefcconstraint}{}
\begin{notedef}
  \texttt{feasTupleFC}$(\collec{x_1}{x_n},feasTuples)$ states an extensional constraint on \collec{x_1}{x_n} defined by the table $feasTuples$ of compatible tuples of values, and then performs Forward Checking:
      $$\exists \text{ tuple } i\ |\quad \collec{x_1}{x_n}=feasTuples[i]$$
\end{notedef}


\begin{itemize}
	\item \textbf{API} : \mylst{feasTupleFC(List<int[]> tuples, IntegerVariable... x)}
	\item \textbf{return type}: \texttt{Constraint}
	\item \textbf{options} : \emph{n/a}
	\item \textbf{favorite domain}: \emph{to complete}
	\item \textbf{references} :\\
      global constraint catalog: \href{http://www.emn.fr/x-info/sdemasse/gccat/Cin_relation.html}{in\_relation}
\end{itemize}

\textbf{Example}:
\lstinputlisting{java/cfeastuplefc.j2t}

\section{forbiddenInterval (constraint)}\label{forbiddeninterval:forbiddenintervalconstraint}\hypertarget{forbiddeninterval:forbiddenintervalconstraint}{}
\begin{notedef}
  \texttt{forbiddenInterval}$(t)$ applies additionnal search tree reduction based on time intervals in which no operation can start or end in an optimal solution. \emph{The tasks must all belong to one disjunctive resource and have fixed durations}.
\end{notedef}

\begin{itemize}
	\item \textbf{API} :
	\begin{itemize}
		\item \mylst{forbiddenInterval(String name, TaskVariable[] tasks)}
		\item \mylst{forbiddenInterval(TaskVariable[] tasks)}
	\end{itemize}
	\item \textbf{return type} : \texttt{Constraint}
	\item \textbf{options} : \emph{n/a}
	\item \textbf{favorite domain} : \emph{n/a}
\end{itemize}

\textbf{Example}:
\emph{to complete}
%\lstinputlisting{java/ceqcard.j2t}


%\part{geost}
\label{geost}
\hypertarget{geost}{}

\section{geost (constraint)}\label{geost:geostconstraint}\hypertarget{geost:geostconstraint}{}
\todo{to be cleaned}

\begin{notedef}
\texttt{geost} is a global constraint that generically handles a variety of geometrical placement problems. 
It handles geometrical constraints (non-overlapping, distance, etc.) between polymorphic objects (ex: polymorphism can be used for representing rotation) in any dimension.
%The \texttt{geost}$(K, O, S, C)$ constraint is given set of parameters which will define the environment of \texttt{geost}. The parameters are as follows:
The parameters of \texttt{geost}$(dim, objects, shiftedBoxes, eCtrs)$ are respectively:
the space dimension, the list of geometrical objects, the set of boxes that compose the shapes of the objects, the set of geometrical constraints.
The greedy mode should be used without external constraints to have safe results, because it excludes external constraints from its exploration and look for instanciation of variables involved in geost which respect the geost constraint.
For further informations, see the section devoted to this constraint in the Choco Tutorial document. 
%visit the following \hyperlink{geostdescription:placementanduseofthegeostconstraint}{page}.
\end{notedef}

\begin{itemize}
	\item \textbf{API} :\\
\mylst{geost(int dim, Vector<GeostObject> objects, Vector<ShiftedBox> shiftedBoxes, Vector<ExternalConstraint> eCtrs)}\\
\mylst{geost(int dim, Vector<GeostObject> objects, Vector<ShiftedBox> shiftedBoxes, Vector<ExternalConstraint> eCtrs, Vector<int[]> ctrlVs)}
	\item \textbf{return type} : \texttt{Constraint}
	\item \textbf{options} :\emph{n/a}
	\item \textbf{favorite domain} : \emph{to complete}
	\item \textbf{references} :\\
      global constraint catalog: \href{http://www.emn.fr/x-info/sdemasse/gccat/Cgeost.html}{geost}
\end{itemize}

The geost constraint requires the creation of different objects:

\centerline{\begin{tabular}{p{3cm}p{5cm}p{6cm}}
parameter &type &description \\
\hline
\emph{objects} &\texttt{Vector<GeostObject>} &geometrical objects\\
\emph{shiftedBoxes} &\texttt{Vector<ShiftedBox>} &boxes that compose the object shapes\\
\emph{eCtrs} &\texttt{Vector<ExternalConstraint>} &geometrical constraints\\
\emph{ctrlVs} &\texttt{Vector<int[]>} &controlling vectors (for greedy mode)\\[1em]
\end{tabular}}

\noindent Where a \texttt{\bf GeostObject} is defined by:

\centerline{\begin{tabular}{p{3cm}p{4cm}p{7cm}}
attribute &type &description \\
\hline
\emph{dim} &\texttt{int} &dimension\\
\emph{objectId} &\texttt{int} &object id\\
\emph{shapeId} &\texttt{IntegerVariable} &shape id\\
\emph{coordinates} &\texttt{IntegerVariable[$dim$]} &coordinates of the origin\\
\emph{startTime} &\texttt{IntegerVariable} &starting time\\
\emph{durationTime} &\texttt{IntegerVariable} &duration\\
\emph{endTime} &\texttt{IntegerVariable} &finishing time\\[1em]
\end{tabular}}

\noindent Where a \texttt{\bf ShiftedBox} is a $dim$-box defined by the shape it belongs to, its origin (the coordinates of the lower left corner of the box) and its lengths in every dimensions:

\centerline{\begin{tabular}{p{3cm}p{4cm}p{7cm}}
attribute &type &description \\
\hline
\emph{sid} &\texttt{int} &shape id\\
\emph{offset} &\texttt{int[$dim$]} &coordinates of the offset (lower left corner)\\
\emph{size} &\texttt{int[$dim$]} &lengths in every dimensions\\[1em]
\end{tabular}}

\noindent Where an \texttt{\bf ExternalConstraint} contains informations and functionality common to all external constraints and is defined by:

\centerline{\begin{tabular}{p{3cm}p{4cm}p{7cm}}
attribute &type &description \\
\hline
 \emph{ectrID} &\texttt{int} &constraint id\\
 \emph{dimensions} &\texttt{int[]} &list of dimensions that the external constraint is active for\\
 \emph{objectIdentifiers} &\texttt{int[]} &list of object ids that this external constraint affects.\\[1em]
\end{tabular}}

\textbf{Example}:
\lstinputlisting{java/cgeost.j2t}

%\part{geq}
\label{geq}
\hypertarget{geq}{}

\section{geq (constraint)}\label{geq:geqconstraint}\hypertarget{geq:geqconstraint}{}
\begin{notedef}
  \texttt{geq}$(x,y)$ states that $x$ is greater than or equal to $y$:
$$x\ge y$$
\end{notedef}

\begin{itemize}
	\item \textbf{API} :
	\begin{itemize}
		\item \mylst{geq(IntegerExpressionVariable x, IntegerExpressionVariable y)}
		\item \mylst{geq(IntegerExpressionVariable x, int y)}
		\item \mylst{geq(int x, IntegerExpressionVariable y)}
		\item \mylst{geq(RealExpressionVariable x, RealExpressionVariable y)}
		\item \mylst{geq(RealExpressionVariable x, double y)}
		\item \mylst{geq(double x, RealExpressionVariable y)}
	\end{itemize}
	\item \textbf{return type} : \texttt{Constraint}
	\item \textbf{options} : \emph{n/a}
	\item \textbf{favorite domain} : \emph{to complete}.
	\item \textbf{references} :\\
      global constraint catalog: \href{http://www.emn.fr/x-info/sdemasse/gccat/Cgeq.html}{geq}
\end{itemize}

\textbf{Examples:}
\begin{itemize}
	\item example1:
\end{itemize}

\begin{lstlisting}
	Model m = new CPModel();
	Solver s = new CPSolver();
	int c = 1;
	IntegerVariable v = makeIntVar("v", 0, 2);
	m.addConstraint(geq(v, c));
	s.read(m);
	s.solve();
\end{lstlisting}
\begin{itemize}
	\item example2
\end{itemize}

\begin{lstlisting}
	Model m = new CPModel();
	Solver s = new CPSolver();
	IntegerVariable v1 = makeIntVar("v1", 0, 2);
	IntegerVariable v2 = makeIntVar("v2", 0, 2);
	IntegerExpressionVariable w1 = plus(v1, 1);
	IntegerExpressionVariable w2 = minus(v2, 1);
	m.addConstraint(geq(w1, w2));
	s.read(m);
	s.solve();
\end{lstlisting}

%\part{geqcard}
\label{geqcard}
\hypertarget{geqcard}{}

\section{geqCard (constraint)}\label{geqcard:geqcardconstraint}\hypertarget{geqcard:geqcardconstraint}{}
\begin{notedef}
  \texttt{geqCard}$(s,x)$ states that the cardinality of set $s$ is greater than or equal to $x$:
$$|s| \ge x$$
\end{notedef}

\begin{itemize}
	\item \textbf{API} :
	\begin{itemize}
		\item \mylst{geqCard(SetVariable s, IntegerVariable x)}
		\item \mylst{geqCard(SetVariable s, int x)}
	\end{itemize}
	\item \textbf{return type} : \texttt{Constraint}
	\item \textbf{options} : \emph{n/a}
	\item \textbf{favorite domain} : \emph{to complete}
\end{itemize}

\textbf{Example}:
\begin{lstlisting}
	Model m = new CPModel();
	Solver s = new CPSolver();
	SetVariable s = makeSetVar("s", 1, 5);
	IntegerVariable i = makeIntVar("card", 2, 3);
	m.addConstraint(member(x, 3));
	m.addConstraint(geqCard(x, i));
	s.read(m);
	s.solve();
\end{lstlisting}

%\part{globalcardinality}
\label{globalcardinality}
\hypertarget{globalcardinality}{}

\section{globalCardinality (constraint)}\label{globalcardinality:globalcardinalityconstraint}\hypertarget{globalcardinality:globalcardinalityconstraint}{}
\begin{notedef}
  \texttt{globalCardinality}$(x,low, up)$ states bounds on the occurrence numbers of any value $v$ in $x$ (here, offset $min$ is the minimum value over all variables in $x$) :
$$low[v-min]\le|\{i\ |\ x_i=v\}|\le up[v-min],\quad\forall \text{ value } v$$   
\end{notedef}

Mulitple APIs exist:
\begin{itemize}
%	\item \emph{bounds on cardinalities:} Given an array of variables $x$, $min$ the minimal value over all variables, and $max$ the maximal value over all variables, the constraint ensures that the number of occurrences of value $v$ among the variables is between $low[v-min]$ and $up[v-min]$. Note that the length of $low$ and $up$ should be $max - min + 1$. Use the propagator of \cite{QuimperCP03}.
	\item Given an array of variables $x$ and an $offset$, the constraint ensures that the number of occurrences of the value $v$ in $x$ is between $low[v-offset]$ and $up[v-offset]$. Use the propagator of \cite{QuimperCP03}.
	\item \emph{variable cardinalities:} Given an array of variables $x$, an array of variables $card$ to represent the cardinalities and an $offset$, the constraint ensures that the number of occurrences of the value $v$ among the variables is equal to $card[v-offset]$. This constraint:
      \begin{itemize}
      \item enforces Bound Consistency over $x$ regarding the lower and upper bounds of $card$, 
      \item maintains the upper bounds of $card$ by counting the number of variables in which each value can occur, 
      \item maintains the lower bounds of $card$ by counting the number of variables instantiated to each value, 
      \item enforces $card[0] + \cdots + card[m] = n$, where \emph{n} is the number of variables and \emph{m} the number of values.
      \end{itemize}
	Use the propagator of \cite{QuimperCP03}.      

\end{itemize}

The APIs are duplicated to define options. 

\begin{itemize}
	\item \textbf{API} :
      \begin{itemize}
	\item \mylst{globalCardinality(IntegerVariable[] x, int[] low, int[] up, int offset)}
	\item \mylst{globalCardinality(String options, IntegerVariable[] x, int[] low, int[] up, int offset)}
	\item \mylst{globalCardinality(IntegerVariable[] x, IntegerVariable[] card, int offset)}
      \end{itemize}
	\item \textbf{return type} : \texttt{Constraint}
	\item \textbf{options}:
	\begin{itemize}
		\item \emph{no option} :
		\item \hyperlink{cgccac:cgccacoptions}{\tt Options.C\_GCC\_AC} : for \cite{ReginAAAI96} implementation of arc consistency
		\item \hyperlink{cgccbc:cgccbcoptions}{\tt Options.C\_GCC\_BC} : for  \cite{QuimperCP03} implementation of bound consistency
	\end{itemize}
	\item \textbf{favorite domain} : \emph{to complete}
	\item \textbf{references} :
      \begin{itemize}
      \item \cite{ReginAAAI96}: \emph{Generalized arc consistency for global cardinality constraint},
      \item \cite{QuimperCP03}: \emph{An efficient bounds consistency algorithm for the global cardinality constraint}
      \item global constraint catalog: \href{http://www.emn.fr/x-info/sdemasse/gccat/Cglobal_cardinality.html}{global\_cardinality}
      \end{itemize}
\end{itemize}

\textbf{Examples:}
\begin{itemize}
	\item example1:
\end{itemize}

\lstinputlisting{java/cglobalcardinality1.j2t}

\begin{itemize}
	\item example2
\end{itemize}

\lstinputlisting{java/cglobalcardinality2.j2t}

%\part{gt}
\label{gt}
\hypertarget{gt}{}

\section{gt (constraint)}\label{gt:gtconstraint}\hypertarget{gt:gtconstraint}{}
\begin{notedef}
  \texttt{gt}$(x,y)$ states that $x$ is strictly greater than $y$:
$$x>y$$
\end{notedef}

\begin{itemize}
	\item \textbf{API} :
	\begin{itemize}
		\item \mylst{gt(IntegerExpressionVariable x, IntegerExpressionVariable y)}
		\item \mylst{gt(IntegerExpressionVariable x, int y)}
		\item \mylst{gt(int x, IntegerExpressionVariable y)}
	\end{itemize}
	\item \textbf{return type} : \texttt{Constraint}
	\item \textbf{options} : \emph{n/a}
	\item \textbf{favorite domain} : \emph{to complete}.
	\item \textbf{references} :\\
      global constraint catalog: \href{http://www.emn.fr/x-info/sdemasse/gccat/Cgt.html}{gt}
\end{itemize}

\textbf{Example:}
\begin{lstlisting}
	Model m = new CPModel();
	Solver s = new CPSolver();
	int c = 1;
	IntegerVariable v = makeIntVar("v", 0, 2);
	m.addConstraint(gt(v, c));
	s.read(m);
	s.solve();
\end{lstlisting}

%\part{ifonlyif}
\label{ifonlyif}
\hypertarget{ifonlyif}{}

\section{ifOnlyIf (constraint)}\label{ifonlyif:ifonlyifconstraint}\hypertarget{ifonlyif:ifonlyifconstraint}{}
\begin{notedef}
  \texttt{ifOnlyIf}$(c_1,c_2)$ states that $c_1$ holds if and only if $c_2$ holds:
$$c_1\iff c_2$$
\end{notedef}

\begin{itemize}
	\item \textbf{API} : \mylst{ifOnlyIf(Constraint c1, Constraint c2)}
	\item \textbf{return type} : \texttt{Constraint}
	\item \textbf{options} : \emph{n/a}
	\item \textbf{favorite domain} : \emph{n/a}
\end{itemize}

\textbf{Example}:
\lstinputlisting{java/cifonlyif.j2t}

%\part{ifthenelse}
\label{ifthenelse}
\hypertarget{ifthenelse}{}

\section{ifThenElse (constraint)}\label{ifthenelse:ifthenelseconstraint}\hypertarget{ifthenelse:ifthenelseconstraint}{}
\begin{notedef}
  \texttt{ifThenElse}$(c_1,c_2,c_3)$ states that if $c_1$ holds then $c_2$ holds, otherwise $c_3$ holds:
  $$(c_1\land c_2) \lor (\neg c_1 \land c_3)$$
\end{notedef}

\begin{itemize}
	\item \textbf{API} :\mylst{ifThenElse(Constraint c1, Constraint c2, Constraint c3)}
	\item \textbf{return type} : \texttt{Constraint}
	\item \textbf{options} : \emph{n/a}
	\item \textbf{favorite domain} : \emph{n/a}
\end{itemize}

\textbf{Example}:
\lstinputlisting{java/cifthenelse.j2t}

%\part{implies}
\label{implies}
\hypertarget{implies}{}

\section{implies (constraint)}\label{implies:impliesconstraint}\hypertarget{implies:impliesconstraint}{}
\begin{notedef}
  \texttt{implies}$(c_1,c_2)$ states that if $c_1$ holds then $c_2$ holds:
$$c_1\implies c_2$$
\end{notedef}

\begin{itemize}
	\item \textbf{API} : \mylst{implies(Constraint c1, Constraint c2)}
	\item \textbf{return type} : \texttt{Constraint}
	\item \textbf{options} : \emph{n/a}
	\item \textbf{favorite domain} : \emph{n/a}
\end{itemize}

\textbf{Example}:
\lstinputlisting{java/cimplies.j2t}

%\part{increasingnvalue}
\label{increasingnvalue}
\hypertarget{increasingnvalue}{}

\section{increasingnvalue (constraint)}\label{increasingnvalue:increasingnvalueconstraint}\hypertarget{increasingnvalue:increasingnvalueconstraint}{}
\begin{notedef}
  \texttt{increasing\_nvalue}$(nval, variables)$ states that $variables$ are increasing. In addition, $nval$ is the number of distinct values taken by $variables$.
\end{notedef}

\begin{itemize}
	\item \textbf{API} : \mylst{increasing\_nvalue(IntegerVariable nval, IntegerVariable[] variables)}
	\item \textbf{return type} : \texttt{Constraint}
	\item \textbf{options} :
	\begin{itemize}
		\item \emph{no option} filter on lower bound and on lower bound
		\item \hyperlink{cinvatleast:cinvatleastoptions}{\tt CPOptions.C\_INCREASING\_NVALUE\_ATLEAST} filter on lower bound only
		\item \hyperlink{cinvatmost:cinvatmostoptions}{\tt CPOptions.C\_INCREASING\_NVALUE\_ATMOST} filter on upper bound only
		\item \hyperlink{cinvboth:cinvbothoptions}{\tt CPOptions.C\_INCREASING\_NVALUE\_BOTH} \textit{--default value--} filter on lower bound and on upper bound
	\end{itemize}
	\item \textbf{favorite domain} : \emph{to complete}
	\item \textbf{references} :\\
      global constraint catalog: \href{http://www.emn.fr/x-info/sdemasse/gccat/Cincreasing_nvalue.html}{increasing\_nvalue}
\end{itemize}

\textbf{Example}:
\lstinputlisting{java/cincreasingnvalue.j2t} 

%!TEX root = ../../content-doc.tex
%\part{inversechanneling}
\label{increasingsum}
\hypertarget{increasingsum}{}

\section{increasingSum (constraint)}\label{increasingsum:increasingsumconstraint}\hypertarget{increasingsum:increasingsumconstraint}{}
\begin{notedef}
  \texttt{increasingSum}$(\collec{x_1}{x_n},s)$ states that \collec{x_1}{x_n} is sorted in increasing order and that $s$ is equal to the sum of $x$.
$$\forall i=[1,n-1] \quad x_i \le x_{i+1}\quad\wedge \quad \sum_{i =[1,n]}x_i = s $$
%$$x_i = j \wedge j \le |y| \quad\iff\quad y_j = i \wedge i \le |x|,\qquad i \in[1,n],  j\in[1,m]$$
\end{notedef}
\begin{itemize}
	\item \textbf{API} : \mylst{increasingSum(IntegerVariable[] x, IntegerVariable s)}
	\item \textbf{return type} : \texttt{Constraint}
	\item \textbf{options} : \emph{no options}
	\item \textbf{favorite domain} : bounded for x and s
	\item \textbf{references} : --
      %\\ global constraint catalog: \href{http://www.emn.fr/z-info/sdemasse/gccat/Cinverse_within_range.html}{inverse\_within\_range}
\end{itemize}

\textbf{Example}:
\lstinputlisting{java/cincreasingsum.j2t}

%\part{infeaspairac}
\label{infeaspairac}
\hypertarget{infeaspairac}{}

\section{infeasPairAC (constraint)}\label{infeaspairac:infeaspairacconstraint}\hypertarget{infeaspairac:infeaspairacconstraint}{}
\begin{notedef}
  \texttt{infeasPairAC}$(x,y,infeasTuples)$ states an extensional binary constraint on $(x,y)$ defined by the table $infeasTuples$ of forbidden pairs of values, and then enforces arc consistency. Two APIs are available to define the forbidden pairs:
\begin{itemize}
	\item if $infeasTuples$ is encoded as a list of pairs \texttt{List<int[2]>}, then:
      $$\forall \text{ tuple } i\ |\quad (x,y)\neq infeasTuples[i]$$
	\item if $infeasTuples$ is encoded as a boolean matrix \texttt{boolean[][]}, let $\underline{x}$ and  $\underline{y}$ be the initial minimum values of $x$ and $y$, then:
      $$\forall (u,v)\ |\quad (x,y)=(u+\underline{x},v+\underline{y})\ \lor\ \neg infeasTuples[u][v]$$
\end{itemize}
\end{notedef}

The two APIs are duplicated to allow definition of options.
\begin{itemize}
	\item \textbf{API} :
	\begin{itemize}
		\item \mylst{infeasPairAC(IntegerVariable x, IntegerVariable y, List<int[]> infeasTuples)}
		\item \mylst{infeasPairAC(String options, IntegerVariable x, IntegerVariable y, List<int[]> infeasTuples)}
		\item \mylst{infeasPairAC(IntegerVariable x, IntegerVariable y, boolean[][] infeasTuples)}
		\item \mylst{infeasPairAC(String options, IntegerVariable x, IntegerVariable y, boolean[][] infeasTuples)}
	\end{itemize}
	\item \textbf{return type} : \texttt{Constraint}
	\item \textbf{options} :
	\begin{itemize}
		\item \emph{no option}: use AC3 (default arc consistency)
		\item \texttt{cp:ac3}: to get AC3 algorithm (searching from scratch for supports on all values)
		\item \texttt{cp:ac2001}: to get AC2001 algorithm (maintaining the current support of each value)
		\item \texttt{cp:ac32}: to get AC3rm algorithm (maintaining the current support of each value in a non backtrackable way)
		\item \texttt{cp:ac322}: to get AC3 with the used of \texttt{BitSet} to know if a support still exists
	\end{itemize}
	\item \textbf{favorite domain} : \emph{to complete}
\end{itemize}

\textbf{Example}:
\begin{lstlisting}
	Model m = new CPModel();
	Solver s = new CPSolver();
	
	boolean[][] matrice2 = new boolean[][]{
	                {false, true, true, false},
	                {true, false, false, false},
	                {false, false, true, false},
	                {false, true, false, false}};
	
	IntegerVariable v1 = makeIntVar("v1", 1, 4);
	IntegerVariable v2 = makeIntVar("v2", 1, 4);
	m.addConstraint(feasPairAC("cp:ac32",v1, v2, matrice2));
	s.read(m);
	s.solveAll();
\end{lstlisting} 

%\part{infeastupleac}
\label{infeastupleac}
\hypertarget{infeastupleac}{}

\section{infeasTupleAC (constraint)}\label{infeastupleac:infeastupleacconstraint}\hypertarget{infeastupleac:infeastupleacconstraint}{}
\begin{notedef}
  \texttt{infeasTupleAC}$(x,feasTuples)$ states an extensional constraint on $(x_1,\ldots,x_n)$ defined by the table $infeasTuples$ of compatible tuples of values, and then enforces arc consistency:
      $$\forall \text{ tuple } i\ |\quad (x_1,\ldots,x_n)\neq infeasTuples[i]$$
\end{notedef}

The API is duplicated to define options.
\begin{itemize}
	\item \textbf{API} :
	\begin{itemize}
		\item \mylst{infeasTupleAC(List<int[]> infeasTuples, IntegerVariable... x)}
		\item \mylst{infeasTupleAC(String options, List<int[]> infeasTuples, IntegerVariable... x)}
	\end{itemize}
	\item \textbf{return type}: \texttt{Constraint}
	\item \textbf{options} :
	\begin{itemize}
		\item \emph{no option}: use AC32 (default arc consistency)
		\item \texttt{cp:ac32}: to get AC3rm algorithm (maintaining the current support of each value in a non backtrackable way)
		\item \texttt{cp:ac2001}: to get AC2001 algorithm (maintaining the current support of each value)
		\item \texttt{cp:ac2008}: to get AC2008 algorithm (maintained by STR)
	\end{itemize}
	\item \textbf{favorite domain} : \emph{to complete}
\end{itemize}

\textbf{Example}:
\lstinputlisting{java/cinfeastupleac.j2t}

%\part{infeastuplefc}
\label{infeastuplefc}
\hypertarget{infeastuplefc}{}

\section{infeasTupleFC (constraint)}\label{infeastuplefc:infeastuplefcconstraint}\hypertarget{infeastuplefc:infeastuplefcconstraint}{}
\begin{notedef}
  \texttt{infeasTupleFC}$(\collec{x_1}{x_n},feasTuples)$ states an extensional constraint on \collec{x_1}{x_n} defined by the table $infeasTuples$ of compatible tuples of values, and then performs Forward Checking:
      $$\forall \text{ tuple } i\ |\quad \collec{x_1}{x_n}\neq infeasTuples[i]$$
\end{notedef}

\begin{itemize}
	\item \textbf{API} : \mylst{infeasTupleFC(List<int[]> infeasTuples, IntegerVariable... x)}
	\item \textbf{return type}: \texttt{Constraint}
	\item \textbf{options} : \emph{n/a}
	\item \textbf{favorite domain}: \emph{to complete}
\end{itemize}

\textbf{Example}:
\lstinputlisting{java/cinfeastuplefc.j2t}

%\part{intdiv}
\label{intdiv}
\hypertarget{intdiv}{}

\section{intDiv (constraint)}\label{intdiv:intdivconstraint}\hypertarget{intdiv:intdivconstraint}{}
\begin{notedef}
  \texttt{intDiv}$(x,y,z)$ states that $z$ is equal to the integer quotient of $x$ by $y$:
$$z = \lfloor x / y \rfloor$$
\end{notedef}

\begin{itemize}
	\item \textbf{API}: \mylst{intDiv(IntegerVariable x, IntegerVariable y, IntegerVariable z)}
	\item \textbf{return type} : \texttt{Constraint}
	\item \textbf{option} : \emph{n/a}
	\item \textbf{favorite domain}: bound
\end{itemize}

\textbf{Example}:
\lstinputlisting{java/cintdiv.j2t}

%\part{inversechanneling}
\label{inversechanneling}
\hypertarget{inversechanneling}{}

\section{inverseChanneling (constraint)}\label{inversechanneling:inversechannelingconstraint}\hypertarget{inversechanneling:inversechannelingconstraint}{}
\begin{notedef}
  \texttt{inverseChanneling}$(x,y)$ states a channeling between two arrays  $x$ and $y$ of integer variables with the same domain.It enforces that if the $i$-th element of $x$ is equal to $j$ then the $j$-th element of $y$ is equal to $i$ and conversely:
$$x_i = j\quad\iff\quad y_j = i$$
\end{notedef}
\begin{itemize}
	\item \textbf{API} : \mylst{inverseChanneling(IntegerVariable[] x, IntegerVariable[] y)}
	\item \textbf{return type} : \texttt{Constraint}
	\item \textbf{options} : \emph{no options}
	\item \textbf{favorite domain} : enumerated for x
	\item \textbf{references} :\\
      global constraint catalog: \href{http://www.emn.fr/x-info/sdemasse/gccat/Cinverse.html}{inverse}
\end{itemize}

\textbf{Example}:
\lstinputlisting{java/cinversechanneling.j2t}

%!TEX root = ../../content-doc.tex
%\part{inversechanneling}
\label{inversechannelingwithinrange}
\hypertarget{inversechannelingwithinrange}{}

\section{inverseChannelingWithinRange (constraint)}\label{inversechannelingwithinrange:inversechannelingconstraintwithinrange}\hypertarget{inversechannelingwithinrange:inversechannelingconstraintwithinrange}{}
\begin{notedef}
  \texttt{inverseChannelingWithinRange}$(\collec{x_1}{x_n},\collec{y_1}{y_m})$ states that
%a channeling between two arrays  $x$ and $y$ of integer variables with the same domain.It enforces 
if $x_i$ is assigned to $j$ and if $j$ is less than or equal to the number of items of the collection Y then $y_j$ is assigned to $i$.
Conversely, if $y_j$ is assigned to $i$ and if $i$ is less than or equal to the number of items of the collection X then $x_i$ is assigned to $j$.

$$x_i = j \wedge j < |y| \quad\iff\quad y_j = i \wedge i < |x|,\qquad i \in[1,n],  j\in[1,m]$$
\end{notedef}
\begin{itemize}
	\item \textbf{API} : \mylst{inverseChanneling(IntegerVariable[] x, IntegerVariable[] y)}
	\item \textbf{return type} : \texttt{Constraint}
	\item \textbf{options} : \emph{no options}
	\item \textbf{favorite domain} : enumerated for x and y
	\item \textbf{references} :\\
      global constraint catalog: \href{http://www.emn.fr/z-info/sdemasse/gccat/Cinverse_within_range.html}{inverse\_within\_range}
\end{itemize}

\textbf{Example}:
\lstinputlisting{java/cinversechannelingwithinrange.j2t}

%\part{inversechanneling}
\label{inverseset}
\hypertarget{inverseset}{}

\section{inverseSet (constraint)}\label{inverseset:inversesetconstraint}\hypertarget{inverseset:inversesetconstraint}{}
\begin{notedef}
  \texttt{inverseSet}$(\collec{x_1}{x_n},\collec{s_1}{s_m})$ states 
that $x_i$ has value $j$ if and only if $s_j$ contains value $i$:
$$x_i = j\quad\iff\quad i\in s_j,\qquad\forall i=1..n,j=1..m$$
\end{notedef}
\begin{itemize}
	\item \textbf{API} : \mylst{inverseSet(IntegerVariable[] x, SetVariable[] s)}
	\item \textbf{return type} : \texttt{Constraint}
	\item \textbf{options} : \emph{no options}
	\item \textbf{favorite domain} : enumerated for x
	\item \textbf{references} :\\
      global constraint catalog: \href{http://www.emn.fr/x-info/sdemasse/gccat/Cinverse_set.html}{inverse\_set}
\end{itemize}

\textbf{Example}:
\lstinputlisting{java/cinverseset.j2t}

%\part{isincluded}
\label{isincluded}
\hypertarget{isincluded}{}

\section{isIncluded (constraint)}\label{isincluded:isincludedconstraint}\hypertarget{isincluded:isincludedconstraint}{}
\begin{notedef}
  \texttt{isIncluded}$(x,y)$ states that the second set $y$ contains the first set $x$:
 $$x\subseteq y$$
\end{notedef}

\begin{itemize}
	\item \textbf{API} : \mylst{isIncluded(SetVariable x, SetVariable y)}
	\item \textbf{return type} : \texttt{Constraint}
	\item \textbf{options} :\emph{n/a}
	\item \textbf{favorite domain} : \emph{to complete}
\end{itemize}

\textbf{Example}:
\begin{lstlisting}
	Model m = new CPModel();
	Solver s = new CPSolver();
	SetVariable v1 = makeSetVar("v1", 3, 4);
	SetVariable v2 = makeSetVar("v2", 3, 8);
	m.addConstraint(isIncluded(v1, v2));
	s.read(m);
	s.solveAll();
\end{lstlisting} 

%\part{isnotincluded}
\label{isnotincluded}
\hypertarget{isnotincluded}{}

\section{isNotIncluded (constraint)}\label{isnotincluded:isnotincludedconstraint}\hypertarget{isnotincluded:isnotincludedconstraint}{}
\begin{notedef}
  \texttt{isNotIncluded}$(s_1,s_2)$ states that set $s_1$ is not included in set $s_2$:
 $$s_2\not\subseteq s_2$$
\end{notedef}

\begin{itemize}
	\item \textbf{API} : \mylst{isNotIncluded(SetVariable s1, SetVariable s2)}
	\item \textbf{return type} : \texttt{Constraint}
	\item \textbf{options} :\emph{n/a}
	\item \textbf{favorite domain} : \emph{to complete}
\end{itemize}

\textbf{Example}:
\lstinputlisting{java/cisnotincluded.j2t} 

\section{knapsackProblem (constraint)}\label{knapsackproblem:knapsackproblemconstraint}\hypertarget{knapsackproblem:knapsackproblemconstraint}{}
\begin{notedef}
  \texttt{knapsackProblem}$(x, W, V, w, v)$ states that :
$$\sum_{i=1}^{n}x_{i}w_{i}=W \wedge \sum_{i=1}^{n}x_{i}v_{i}=V$$
$x_i$ is the number of copies of the item $i$, each kind of item $i$ has a value $v_i$ and a weight $w_i$.
\end{notedef}
It constructs a knapsack problem constraint with an underlying costregular constraint, where it simulates M. Triks dynamic programming approach.

\begin{itemize}
	\item \textbf{API} :\item \mylst{knapsackProblem(IntegerVariable[] vars, IntegerVariable cVar1, IntegerVariable cVar2, int[] costs1, int[] costs2)}
	\item \textbf{return type} : \texttt{Constraint}
	\item \textbf{options} : \emph{n/a}
	\item \textbf{favorite domain} : \emph{n/a}
	\item \textbf{references} : 
	\begin{itemize}
		\item \cite{DemasseyC06}: \emph{A \texttt{Cost-Regular} based hybrid column generation approach}
		\item \cite{TrickAOR03}: \emph{A Dynamic Programming Approach for Consistency and Propagation for Knapsack Constraints}
		\item \hyperlink{costregular:costregularconstraint}{costRegular}
	\end{itemize}
\end{itemize}

\textbf{Example}:
\lstinputlisting{java/cknapsack.j2t}


%\part{leq}
\label{leq}
\hypertarget{leq}{}

\section{leq (constraint)}\label{leq:leqconstraint}\hypertarget{leq:leqconstraint}{}
\begin{notedef}
  \texttt{leq}$(x,y)$ states that $x$ is less than or equal to $y$:
$$x \le y$$
\end{notedef}

\begin{itemize}
	\item \textbf{API} :
	\begin{itemize}
		\item \mylst{leq(IntegerExpressionVariable x, IntegerExpressionVariable y)}
		\item \mylst{leq(IntegerExpressionVariable x, int y)}
		\item \mylst{leq(int x, IntegerExpressionVariable y)}
		\item \mylst{leq(RealExpressionVariable x, RealExpressionVariable y)}
		\item \mylst{leq(RealExpressionVariable x, double y)}
		\item \mylst{leq(double x, RealExpressionVariable y)}
	\end{itemize}
	\item \textbf{return type} : \texttt{Constraint}
	\item \textbf{options} : \emph{n/a}
	\item \textbf{favorite domain} : \emph{to complete}.
	\item \textbf{references} :\\
      global constraint catalog: \href{http://www.emn.fr/x-info/sdemasse/gccat/Cleq.html}{leq}
\end{itemize}

\textbf{Example:}
\lstinputlisting{java/cleq.j2t}

%\part{leqcard}
\label{leqcard}
\hypertarget{leqcard}{}

\section{leqCard (constraint)}\label{leqcard:leqcardconstraint}\hypertarget{leqcard:leqcardconstraint}{}
\begin{notedef}
  \texttt{leqCard}$(s,x)$ states that the cardinality of set $s$ is less than or equal to $x$:
$$|s| \le x$$
\end{notedef}

\begin{itemize}
	\item \textbf{API} :
	\begin{itemize}
		\item \mylst{leqCard(SetVariable s, IntegerVariable x)}
		\item \mylst{leqCard(SetVariable s, int x)}
	\end{itemize}
	\item \textbf{return type} : \texttt{Constraint}
	\item \textbf{options} : \emph{n/a}
	\item \textbf{favorite domain} : \emph{to complete}
\end{itemize}

\textbf{Example}:
\lstinputlisting{java/cleqcard.j2t}

%\part{lex}
\label{lex}
\hypertarget{lex}{}

\section{lex (constraint)}\label{lex:lexconstraint}\hypertarget{lex:lexconstraint}{}
\begin{notedef}
  \texttt{lex}$(x,y)$ enforces a strict lexicographic ordering  $x <_{lex} y$ between two arrays of same size $n$:
$$\exists\ j\in\{1,\ldots,n\}\ |\qquad x_j<y_j\quad \land\quad x_i=y_i\ (\forall\  i<j)$$
\end{notedef}

\begin{itemize}
	\item \textbf{API} : \mylst{lex(IntegerVariable[] x, IntegerVariable[] y)}
	\item \textbf{return type} : \texttt{Constraint}
	\item \textbf{options} :\emph{n/a}
	\item \textbf{favorite domain} : \emph{to complete}
	\item \textbf{references} :
      \begin{itemize}
      \item \cite{FrischCP02}: \emph{Global Constraints for Lexicographic Orderings}
      \item global constraint catalog: \href{http://www.emn.fr/x-info/sdemasse/gccat/Clex_less.html}{lex\_less}
      \end{itemize}
\end{itemize}

\textbf{Example}:
\begin{lstlisting}
	Model m = new CPModel();
	Solver s = new CPSolver();
	
	int n1 = 8;
	int k = 2;
	IntegerVariable[] vs1 = new IntegerVariable[n1 / 2];
	IntegerVariable[] vs2 = new IntegerVariable[n1 / 2];
	
	for (int i = 0; i < n1 / 2; i++) {
	   vs1[i] = makeIntVar("" + i, 0, k);
	   vs2[i] = makeIntVar("" + i, 0, k);
	}
	m.addConstraint(lex(vs1, vs2));
	
	s.read(m);
	s.solve();
\end{lstlisting} 

%\part{lexchain}
\label{lexchain}
\hypertarget{lexchain}{}

\section{lexChain (constraint)}\label{lexchain:lexchainconstraint}\hypertarget{lexchain:lexchainconstraint}{}
\begin{notedef}
  \texttt{lexChain}$(\collec{x^1_1}{x^1_n},\ldots,\collec{x^p_1}{x^p_n})$ states a strict lexicographic ordering on a chain of $p$ integer vectors:
$$x^1 <_{lex} x^2 <_{lex}\cdots <_{lex} x^p$$
%where $X^1$ contains up to $n$ variables. 
\end{notedef}

\begin{itemize}
	\item \textbf{API} : \mylst{lexChain(IntegerVariable[]... x)}
	\item \textbf{return type} : \texttt{Constraint}
	\item \textbf{options} : \emph{n/a}
	\item \textbf{favorite domain} : \emph{to complete}
	\item \textbf{references} :
      \begin{itemize}
      \item \cite{BeldiceanuSICS02} \emph{Arc-Consistency for a chain of Lexicographic Ordering Constraints} 
      \item global constraint catalog: \href{http://www.emn.fr/x-info/sdemasse/gccat/Clex_chain_less.html}{lex\_chain\_less}
      \end{itemize}
\end{itemize}

\textbf{Example}:
\lstinputlisting{java/clexchain.j2t}

%\part{lexchaineq}
\label{lexchaineq}
\hypertarget{lexchaineq}{}

\section{lexChainEq (constraint)}\label{lexchaineq:lexchaineqconstraint}\hypertarget{lexchaineq:lexchaineqconstraint}{}
\begin{notedef}
\texttt{lexChainEq}$(x^1 ,x^2 ,x^3,\ldots)$ enforces a lexicographic ordering on a chain of integer vectors:
$$x^1 \le_{lex} x^2 \le_{lex} x^3 \le_{lex}\cdots$$
%where $X^1$ contains up to $n$ variables. 
\end{notedef}

\begin{itemize}
	\item \textbf{API} : \mylst{lexChainEq(IntegerVariable[]... arrayOfVectors)}
	\item \textbf{return type} : \texttt{Constraint}
	\item \textbf{options} : \emph{n/a}
	\item \textbf{favorite domain} : \emph{to complete}
	\item \textbf{references} :
      \begin{itemize}
      \item \cite{BeldiceanuSICS02} \emph{Arc-Consistency for a chain of Lexicographic Ordering Constraints} 
      \item global constraint catalog: \href{http://www.emn.fr/x-info/sdemasse/gccat/Clex_chain_lesseq.html}{lex\_chain\_lesseq}
      \end{itemize}
\end{itemize}

\textbf{Example}:
\lstinputlisting{java/clexchaineq.j2t}
%\part{lexeq}
\label{lexeq}
\hypertarget{lexeq}{}

\section{lexEq (constraint)}\label{lexeq:lexeqconstraint}\hypertarget{lexeq:lexeqconstraint}{}
\begin{notedef}
  \texttt{lexEq}$(\collec{x_1}{x_n},\collec{y_1}{y_n})$ states a lexicographic ordering  $x \le_{lex} y$:
$$\exists\ j=1..n\ |\qquad x_j\le y_j\quad \land\quad x_i=y_i\ (\forall\  i<j)$$
\end{notedef}

\begin{itemize}
	\item \textbf{API} : \mylst{lexEq(IntegerVariable[] x, IntegerVariable[] y)}
	\item \textbf{return type} : \texttt{Constraint}
	\item \textbf{options} :\emph{n/a}
	\item \textbf{favorite domain} : \emph{to complete}
	\item \textbf{references} :
      \begin{itemize}
      \item \cite{FrischCP02}: \emph{Global Constraints for Lexicographic Orderings}
      \item global constraint catalog: \href{http://www.emn.fr/x-info/sdemasse/gccat/Clex_lesseq.html}{lex\_lesseq}
      \end{itemize}
\end{itemize}

\textbf{Example}:
\textbf{Example}:
\lstinputlisting{java/clexeq.j2t}

%\part{leximin}
\label{leximin}
\hypertarget{leximin}{}

\section{leximin (constraint)}\label{leximin:leximinconstraint}\hypertarget{leximin:leximinconstraint}{}

\emph{TODO: verify the specifications of the implemented version.}

\begin{notedef}
Let $x = (x_1,\ldots, x_n)$ and $y = (y_1,\ldots, y_n)$ be two vectors of $n$ integers, and let $x'$ and $y'$ be respectively permutations of vectors $x$ and $y$ sorted by increasing order of the components.
Constraint \texttt{leximin(x, y)} holds if and only if $x'<_{lex} y'$:
$$\exists\ j\in\{1,\ldots,n\}\ |\qquad x'_j<y'_j\quad \land\quad x'_i=y'_i\ (\forall\  i<j)$$
  \end{notedef}

\begin{itemize}
	\item \textbf{API} :
	\begin{itemize}
		\item \mylst{leximin(IntegerVariable[] x, IntegerVariable[] y)}
		\item \mylst{leximin(int[] x, IntegerVariable[] y)}
	\end{itemize}
	\item \textbf{return type} : \texttt{Constraint}
	\item \textbf{options} :\emph{n/a}
	\item \textbf{favorite domain} : \emph{to complete}
	\item \textbf{references} :
      \begin{itemize}
      \item \cite{FrischIJCAI03}: \emph{Multiset ordering constraints} 
      \item global constraint catalog: \href{http://www.emn.fr/x-info/sdemasse/gccat/Clex_lesseq_allperm.html}{lex\_lesseq\_allperm} (variant)
      \end{itemize}
\end{itemize}

\textbf{Example}:
\lstinputlisting{java/cleximin.j2t}

%\part{lt}
\label{lt}
\hypertarget{lt}{}

\section{lt (constraint)}\label{lt:ltconstraint}\hypertarget{lt:ltconstraint}{}
\begin{notedef}
  \texttt{lt}$(x,y)$ states that $x$ is strictly smaller than $y$:
$$x<y$$
\end{notedef}

\begin{itemize}
	\item \textbf{API} :
	\begin{itemize}
		\item \mylst{lt(IntegerExpressionVariable x, IntegerExpressionVariable y)}
		\item \mylst{lt(IntegerExpressionVariable x, int y)}
		\item \mylst{lt(int x, IntegerExpressionVariable y)}
	\end{itemize}
	\item \textbf{return type} : \texttt{Constraint}
	\item \textbf{options} : \emph{n/a}
	\item \textbf{favorite domain} : \emph{to complete}.
	\item \textbf{references} :\\
      global constraint catalog: \href{http://www.emn.fr/x-info/sdemasse/gccat/Clt.html}{lt}
\end{itemize}

\textbf{Example}:
\lstinputlisting{java/clt.j2t}

%\part{max}
\label{max}
\hypertarget{max}{}

\section{max (constraint)}\label{max:maxconstraint}\hypertarget{max:maxconstraint}{}

\subsection{max of a list}\label{max:maxofalist}\hypertarget{max:maxofalist}{}

\begin{notedef}
\texttt{max}$(x,z)$ states that $z$ is equal to the greater element of vector $x$:
$$z = \max(x_1, x_2, ..., x_n)$$  
\end{notedef}

\begin{itemize}
	\item \textbf{API}:
	\begin{itemize}
		\item \mylst{max(IntegerVariable[] x, IntegerVariable z)}
		\item \mylst{max(IntegerVariable x1, IntegerVariable x2, IntegerVariable z)}
		\item \mylst{max(int x1, IntegerVariable x2, IntegerVariable z)}
		\item \mylst{max(IntegerVariable x1, int x2, IntegerVariable z)}
	\end{itemize}
	\item \textbf{return type}: \texttt{Constraint}
	\item \textbf{options} : \emph{n/a}
	\item \textbf{favorite domain} : \emph{to complete}
	\item \textbf{references} :\\
      global constraint catalog: \href{http://www.emn.fr/x-info/sdemasse/gccat/Cmaximum.html}{maximum}
\end{itemize}

\textbf{Example}:
\begin{lstlisting}
	Model m = new CPModel();
	Solver s= new CPSolver();
	IntegerVariable x = makeIntVar("x", 1, 5);
	IntegerVariable y = makeIntVar("y", 1, 5);
	IntegerVariable z = makeIntVar("z", 1, 5);
	m.addVariable("cp:bound", x, y, z);
	m.addConstraint(max(y, z, x));
	s.read(m);
	s.solve();
\end{lstlisting}

\subsection{max of a set}\label{max:maxofaset}\hypertarget{max:maxofaset}{}

\begin{notedef}
\texttt{max}$(s,x,z)$ states that $z$ is equal to the greater element of vector $x$ whose index is in set $s$:
$$z = \max_{i\in s}( x_i )$$
  \end{notedef}

\begin{itemize}
	\item \textbf{API}:
	\begin{itemize}
		\item \mylst{max(SetVariable s, IntegerVariable[] x, IntegerVariable z)}
	\end{itemize}
	\item \textbf{return type}: \texttt{Constraint}
	\item \textbf{options} : \emph{n/a}
	\item \textbf{favorite domain} : \emph{to complete}
\end{itemize}

\begin{lstlisting}
	Model m = new CPModel();
	Solver s= new CPSolver();
	IntegerVariable[] x = constantArray(new int[]{5,7,9,10,12,3,2});
	IntegerVariable max = makeIntVar("max", 1, 100);
	SetVariable set = makeSetVar("set", 0, x.length-1);
	m.addConstraints(max(set, x, max), leqCard(set, constant(5)));
	s.read(m);
	s.solve();
\end{lstlisting}

%\part{member}
\label{member}
\hypertarget{member}{}

\section{member (constraint)}\label{member:memberconstraint}\hypertarget{member:memberconstraint}{}

\begin{notedef}
  \texttt{member}$(x,s)$ states that integer $x$ belongs to set $s$:
$$x\in s$$
\end{notedef}

\begin{itemize}
	\item \textbf{API} :
	\begin{itemize}
		\item \mylst{member(int x, SetVariable s)}
		\item \mylst{member(SetVariable s, int x)}
		\item \mylst{member(SetVariable s, IntegerVariable x)}
		\item \mylst{member(IntegerVariable x, SetVariable s)}
		\item \mylst{member(member(SetVariable sv, IntegerVariable... vars)}
		\item \mylst{member(IntegerVariable x, int[] s)}
		\item \mylst{member(IntegerVariable x, int lower, int upper)}
	\end{itemize}
	\item \textbf{return type} : \texttt{Constraint}
	\item \textbf{options} :\emph{n/a}
	\item \textbf{favorite domain} : \emph{to complete}
	\item \textbf{references} :\\
      global constraint catalog: \href{http://www.emn.fr/x-info/sdemasse/gccat/Cin_set.html}{in\_set}
\end{itemize}

\textbf{Examples}:
1. using a set variable
\lstinputlisting{java/cmember.j2t}

2. using an array of integers
\lstinputlisting{java/camong1.j2t}

%\part{min}
\label{min}
\hypertarget{min}{}

\section{min (constraint)}\label{min:minconstraint}\hypertarget{min:minconstraint}{}

\subsection{min of a list}\label{min:minofalist}\hypertarget{min:minofalist}{}

\begin{notedef}
  \texttt{mix}$(x,z)$ states that $z$ is equal to the smaller element
  of vector $x$:
$$z = \min(x_1, x_2, ..., x_n).$$
\end{notedef}
\begin{itemize}
	\item \textbf{API}:
	\begin{itemize}
		\item \mylst{min(IntegerVariable[] x, IntegerVariable z)}
		\item \mylst{min(IntegerVariable x1, IntegerVariable x2, IntegerVariable z)}
		\item \mylst{min(int x1, IntegerVariable x2, IntegerVariable z)}
		\item \mylst{min(IntegerVariable x1, int x2, IntegerVariable z)}
	\end{itemize}
	\item \textbf{return type}: \texttt{Constraint}
	\item \textbf{options} : \emph{n/a}
	\item \textbf{favorite domain} : \emph{to complete}
	\item \textbf{references} :\\
      global constraint catalog: \href{http://www.emn.fr/x-info/sdemasse/gccat/Cminimum.html}{minimum}
\end{itemize}

\textbf{Example}:
\begin{lstlisting}
	Model m = new CPModel();
	Solver s= new CPSolver();
	IntegerVariable x = makeIntVar("x", 1, 5);
	IntegerVariable y = makeIntVar("y", 1, 5);
	IntegerVariable z = makeIntVar("z", 1, 5);
	m.addVariable("cp:bound", x, y, z);
	m.addConstraint(min(y, z, x));
	s.read(m);
	s.solve();
\end{lstlisting}

\subsection{min of a set}\label{min:minofaset}\hypertarget{min:minofaset}{}

\begin{notedef}
  \texttt{min}$(s,x,z)$ states that $z$ is equal to the smaller
  element of vector $x$ whose index is in set $s$:
$$z = \min_{i\in s}( x_i ).$$
\end{notedef}
\begin{itemize}
	\item \textbf{API}:
	\begin{itemize}
		\item \mylst{min(SetVariable s,IntegerVariable[] x, IntegerVariable z)}
	\end{itemize}
	\item \textbf{return type}: \texttt{Constraint}
	\item \textbf{options} : \emph{n/a}
	\item \textbf{favorite domain} : \emph{to complete}
\end{itemize}

\begin{lstlisting}
	Model m = new CPModel();
	Solver s= new CPSolver();
	IntegerVariable[] x = constantArray(new int[]{5,7,9,10,12,3,2});
	IntegerVariable min = makeIntVar("min", 1, 100);
	SetVariable set = makeSetVar("set", 0, x.length-1);
	m.addConstraints(min(set, x, max), leqCard(set, constant(5)));
	s.read(m);
	s.solve();
\end{lstlisting}

%\part{mod}
\label{mod}
\hypertarget{mod}{}

\section{mod (constraint)}\label{mod:modconstraint}\hypertarget{mod:modconstraint}{}
\begin{notedef}
  \texttt{mod}$(x_1,x_2,x_3)$ states that $x_1$ is congruent to $x_2$
  modulo $x_3$:
$$x_1 \equiv x_2 \mod x_3$$
\end{notedef}
\begin{itemize}
	\item \textbf{API} : \mylst{mod(IntegerVariable x1, IntegerVariable x2, int x3)}
	\item \textbf{return type} : \texttt{Constraint}
	\item \textbf{options} : \emph{n/a}
	\item \textbf{favorite domain} : \emph{n/a}
\end{itemize}

\textbf{Example}:
\begin{lstlisting}
	Model m = new CPModel();
	Solver s = new CPSolver();
	
	IntegerVariable x = makeIntVar("x", 0, 10);
	IntegerVariable w = makeIntVar("w", 0, 10);
	
	m.addConstraint(mod(w,x, 1));
	
	s.read(m);
	s.solve();
\end{lstlisting}

%\part{multicostregular}
\label{multicostregular}
\hypertarget{multicostregular}{}

\section{multiCostRegular (constraint)}\label{multicostregular:multicostregularconstraint}\hypertarget{multicostregular:multicostregularconstraint}{}
\begin{notedef}
  \texttt{multiCostRegular}$(x,z,\mathcal{L}(\Pi),c)$ states that sequence $x$ is a word belonging to the regular language $\mathcal{L}(\Pi)$,
% recognized by a deterministic finite automaton (DFA) or a multicostregular expression $\Pi$:
$$(x_1,\ldots,x_n)\in\mathcal{L}(\Pi)$$
and that the bounded vector $z$ is equal to the costs of $x$ according to the assigment cost matrix $c$:
$$\sum_{i=1}^{n} c[r][i][x_i]=z[r],\quad \forall r\in\{0,\ldots,R\}$$
\end{notedef}
\texttt{multiCostRegular} is a conjunction of a \hyperlink{regular}{\texttt{regular}} constraint with $R+1$ cost functions.
It may be used in the context of personnel scheduling problems, handling complex work regulations by the mean of regular expressions, together with cardinality or financial constraints by the mean of cost functions.
The filtering algorithm associated with \texttt{multiCostRegular} is based on lagrangian relaxation and computations of shortest/longest pathes in a layered digraph~\cite{MenanaCPAIOR09}. It typically performs more filtering than the conjunction of \texttt{costRegular} and \texttt{globalCardinality} or than multiple \texttt{costRegular}.

The accepting language is specified by a deterministic finite automaton (DFA):
Automaton $\Pi$ is defined on a given \emph{alphabet} $\Sigma\subseteq\Z$ by a set $Q=\{0,\ldots,m\}$ of \emph{states}, a subset $A\subseteq Q$ of \emph{final} or \emph{accepting states} and a table $\Delta\subseteq Q\!\times\!\Sigma\!\times Q$ of \emph{transitions} between states. $\Pi$ is encoded as an object of class \texttt{Automaton}, whose API contains:
\begin{lstlisting}
  Automaton();
  int addState();
  void setStartingState(int state); 
  void setAcceptingState(int state); 
  void addTransition(int state1, int state2, int label);
  int getNbStates();
\end{lstlisting}
The cost functions are encoded as one matrix \texttt{int cost[nTime][nAct][auto.getNbStates()][nRes]} such that
\texttt{cost[i][j][s][r]} is the cost of assigning variable $x_i$ to activity $j$ at state $s$ on dimension $r+1$.

\begin{itemize}
	\item \textbf{API} : \mylst{multiCostRegular(IntegerVariable[] x, IntegerVariable[] z, Automaton P, int[][][][] c)}
	\item \textbf{return type} : \texttt{Constraint}
	\item \textbf{options} :
      \begin{itemize}
      \item \texttt{MultiCostRegular.DATA\_STRUCT} is  \texttt{MultiCostRegular.BITSET} or \texttt{MultiCostRegular.LIST}: a parameter stating which backtrable data structure to use for storing the outgoing arcs of the layered digraph. The observed behaviour is until $1000$ arcs the bipartite list is much more efficient, afterwards the memory efficiency of the bitset representation allow faster operations. 
      \item \texttt{MultiCostRegular.U0}, \texttt{MultiCostRegular.R0}, \texttt{MultiCostRegular.MAXNONIMPROVEITER}, and \texttt{MultiCostRegular.MAXBOUNDITER} are value parameters of the subgradient algorithm used for solving the lagrangean relaxation.
      \item \texttt{MultiCostRegular.D\_PREC} is a double parameter stating the precision of float computation. It is set by default to $10^{-5}$.
      \end{itemize}
	\item \textbf{favorite domain} : \emph{to complete}
	\item \textbf{references} :\\
       \cite{MenanaCPAIOR09}: \emph{Sequencing and Counting with the {\tt multicost-regular} Constraint}
\end{itemize}
%\begin{notedef}
%  For further informations, see the multicost-regular description.
%\end{notedef}

\textbf{Example}:
\lstinputlisting{java/cmulticosteregular_import.j2t}
\lstinputlisting{java/cmulticostregular.j2t}


%\part{and}
\label{nand}
\hypertarget{nand}{}

\section{nand (constraint)}\label{nand:nandconstraint}\hypertarget{nand:nandconstraint}{}
\begin{notedef}
  \texttt{nand}$(\collec{C_1}{C_n})$ states that at least one constraint in arguments is not satisfied:
$$ C_1 \uparrow C_2 \uparrow\ldots\uparrow C_n$$

  \texttt{nand}$(\collec{b_1}{b_n})$ states that at least one boolean in arguments is false:
$$ (b_1=1)\ \uparrow\ (b_2=1)\ \uparrow\ \ldots\ \uparrow\ (b_n=1)$$
\end{notedef}

\begin{itemize}
\item \textbf{API} : 
\begin{itemize}
\item \mylst{nand(Constraint... c)}
\item \mylst{nand(IntegerVariable... b)}
\end{itemize}
\item \textbf{return type} : \texttt{Constraint}
\item \textbf{options} : \emph{n/a}
\item \textbf{favorite domain} : \emph{n/a}
\item \textbf{references} :\\
  global constraint catalog: \href{http://www.emn.fr/z-info/sdemasse/gccat/Cnand.html}{\tt nand}
\end{itemize}

\textbf{Examples:}
\begin{itemize}
	\item example1:
\end{itemize}
\lstinputlisting{java/cnand1.j2t}
\begin{itemize}
	\item example2
\end{itemize}

\lstinputlisting{java/cnand2.j2t}


%\part{neq}
\label{neq}
\hypertarget{neq}{}

\section{neq (constraint)}\label{neq:neqconstraint}\hypertarget{neq:neqconstraint}{}

\begin{notedef}
  \texttt{neq} states that the two arguments are different:
$$x \neq y.$$
\end{notedef}
\begin{itemize}
	\item \textbf{API} :
	\begin{itemize}
		\item \mylst{neq(IntegerExpressionVariable x, IntegerExpressionVariable y)}
		\item \mylst{neq(IntegerExpressionVariable x, int y)}
		\item \mylst{neq(int x, IntegerExpressionVariable y)}
	\end{itemize}
	\item \textbf{return type} : \texttt{Constraint}
	\item \textbf{options} : \emph{n/a}
	\item \textbf{favorite domain} : \emph{to complete}.
	\item \textbf{references} :\\
      global constraint catalog: \href{http://www.emn.fr/x-info/sdemasse/gccat/Cneq.html}{neq}
\end{itemize}

\textbf{Examples:}
\begin{itemize}
	\item example1:
\end{itemize}

\lstinputlisting{java/cneq1.j2t}

\begin{itemize}
	\item example2
\end{itemize}

\lstinputlisting{java/cneq2.j2t}
%\part{neqcard}
\label{neqcard}
\hypertarget{neqcard}{}

\section{neqCard (constraint)}\label{neqcard:neqcardconstraint}\hypertarget{neqcard:neqcardconstraint}{}
\begin{notedef}
  \texttt{neqCard}$(s,z)$ states that the cardinality of set $s$ is not equal to $z$:
$$|s| \neq z$$
\end{notedef}

\begin{itemize}
	\item \textbf{API} :
	\begin{itemize}
		\item \mylst{neqCard(SetVariable s, IntegerVariable z)}
		\item \mylst{neqCard(SetVariable s, int z)}
	\end{itemize}
	\item \textbf{return type} : \texttt{Constraint}
	\item \textbf{options} : \emph{n/a}
	\item \textbf{favorite domain} : \emph{to complete}
\end{itemize}

\textbf{Example}:
\lstinputlisting{java/cneqcard.j2t}

%\part{not}
\label{not}
\hypertarget{not}{}

\section{not (constraint)}\label{not:notconstraint}\hypertarget{not:notconstraint}{}
\begin{notedef}
  \texttt{not}$(c)$ holds if and only if constraint $c$ does not hold:
$$\neg c$$
\end{notedef}
\begin{itemize}
	\item \textbf{API} : \mylst{not(Constraint c)}
	\item \textbf{return type} : \texttt{Constraint}
	\item \textbf{options} : \emph{n/a}
	\item \textbf{favorite domain} : \emph{n/a}
\end{itemize}

\textbf{Example} : 
\begin{lstlisting}
	Model m = new CPModel();
	Solver s = new CPSolver();
	IntegerVariable x = makeIntVar("x", 1, 10);
	// x < 3
	m.addConstraint(not(geq(x, 3)));
	
	s.read(m);
	s.solve();
\end{lstlisting}

%\part{notmember}
\label{notmember}
\hypertarget{notmember}{}

\section{notMember (constraint)}\label{notmember:notmemberconstraint}\hypertarget{notmember:notmemberconstraint}{}
\begin{notedef}
\texttt{notMember}$(x,s)$ states that integer $x$ is not contained in set $s$:
$$x\not\in s$$  
\end{notedef}

\begin{itemize}
	\item \textbf{API} :
	\begin{itemize}
		\item \mylst{notMember(int x, SetVariable s)}
		\item \mylst{notMember(SetVariable s, int x)}
		\item \mylst{notMember(SetVariable s, IntegerVariable x)}
		\item \mylst{notMember(IntegerVariable x, SetVariable s)}
	\end{itemize}
	\item \textbf{return type} : \texttt{Constraint}
	\item \textbf{options} :\emph{n/a}
	\item \textbf{favorite domain} : \emph{to complete}
\end{itemize}

\textbf{Example}:
\lstinputlisting{java/cnotmember.j2t}

%\part{nor}
\label{nor}
\hypertarget{nor}{}

\section{nor (constraint)}\label{nor:norconstraint}\hypertarget{nor:norconstraint}{}
\begin{notedef}
  \texttt{nor}$(c_1,\ldots,c_n)$ states that at least one constraint in arguments is satisfied:
$$ c_1 \downarrow c_2 \downarrow\ldots\downarrow c_n$$

  \texttt{nor}$(b_1,\ldots,b_n)$ states that at least one boolean variable in argument is true:
$$ (b_1=1) \downarrow (b_2=1) \downarrow\ldots\downarrow (b_n=1)$$
\end{notedef}

\begin{itemize}
\item \textbf{API} : 
\begin{itemize}
\item \mylst{nor(Constraint... c)}
\item \mylst{nor(IntegerVariable... b)}
\end{itemize}
	\item \textbf{return type} : \texttt{Constraint}
	\item \textbf{options} : \emph{n/a}
	\item \textbf{favorite domain} : \emph{n/a}
	\item \textbf{references} :\\
  global constraint catalog: \href{http://www.emn.fr/z-info/sdemasse/gccat/Cnor.html}{\tt nor}
\end{itemize}

\textbf{Examples:}
\begin{itemize}
	\item example1:
\end{itemize}
\lstinputlisting{java/cnor1.j2t}
\begin{itemize}
	\item example2
\end{itemize}
\lstinputlisting{java/cnor2.j2t}

%\part{nth}
\label{nth}
\hypertarget{nth}{}

\section{nth (constraint)}\label{nth:nthconstraint}\hypertarget{nth:nthconstraint}{}
\texttt{nth} is the well known \emph{element} constraint.
Several APIs are available: 
\begin{notedef}
\begin{itemize}
\item \texttt{nth}$(i,x,y)$ ensures that $x[i]=y$
\item \texttt{nth}$(i,x,y,o)$ ensures that $x[i+o]=y$ ($o$ is an \emph{offset} for shifting values)
\item \texttt{nth}$(i,j,x,y)$ ensures that $x[i][j]=y$
\end{itemize}
\end{notedef}

\begin{itemize}
	\item \textbf{API} :
	\begin{itemize}
		\item \mylst{nth(IntegerVariable i, int[] x, IntegerVariable y)}
		\item \mylst{nth(String option, IntegerVariable i, int[] x, IntegerVariable y)}
		\item \mylst{nth(IntegerVariable i, IntegerVariable[] x, IntegerVariable y)}
		\item \mylst{nth(IntegerVariable i, int[] x, IntegerVariable y, int offset)}
		\item \mylst{nth(String option, IntegerVariable i, int[] x, IntegerVariable y, int offset)}		
		\item \mylst{nth(IntegerVariable i, IntegerVariable[] x, IntegerVariable y, int offset)}
		\item \mylst{nth(String option, IntegerVariable i, IntegerVariable[] x, IntegerVariable y, int offset)}
		\item \mylst{nth(IntegerVariable i, IntegerVariable j, int[][] x, IntegerVariable y)}
	\end{itemize}
	\item \textbf{return type} : \texttt{Constraint}
	\item \textbf{options} :
	\begin{itemize}
		\item \emph{no option} 
		\item \hyperlink{cnthg:cnthgoptions}{\tt Options.C\_NTH\_G} for global consistency
	\end{itemize}
	\item \textbf{favorite domain} : \emph{to complete}
	\item \textbf{references} :\\
      global constraint catalog: \href{http://www.emn.fr/x-info/sdemasse/gccat/Celement.html}{element}
\end{itemize}

\textbf{Example}:
\lstinputlisting{java/cnth.j2t} 

%\part{occurrence}
\label{occurrence}
\hypertarget{occurrence}{}

\section{occurrence (constraint)}\label{occurrence:occurrenceconstraint}\hypertarget{occurrence:occurrenceconstraint}{}
\begin{notedef}
  \texttt{occurrence}$(v,z,x)$ states that $z$ is equal to the number of elements in $x$ with value $v$:
$$z=|\{i\ |\ x_i=v\}|$$   
\end{notedef}
  This is a specialization of the \texttt{globalCardinality} constraint.

\begin{itemize}
	\item \textbf{API}: \mylst{occurrence(int v, IntegerVariable z, IntegerVariable... x)}
	\item \textbf{return type} : \texttt{Constraint}
	\item \textbf{options} :\emph{n/a}
	\item \textbf{favorite domain} : \emph{to complete}
	\item \textbf{references} :\\
      global constraint catalog: \href{http://www.emn.fr/x-info/sdemasse/gccat/Ccount.html}{count}
\end{itemize}

\textbf{Example}:
\begin{lstlisting}
	Model m = new CPModel();
	Solver s = new CPSolver();
	
	IntegerVariable x1 = makeIntVar("X1", 0, 10);
	IntegerVariable x2 = makeIntVar("X2", 0, 10);
	IntegerVariable x3 = makeIntVar("X3", 0, 10);
	IntegerVariable x4 = makeIntVar("X4", 0, 10);
	IntegerVariable x5 = makeIntVar("X5", 0, 10);
	IntegerVariable x6 = makeIntVar("X6", 0, 10);
	IntegerVariable x7 = makeIntVar("X7", 0, 10);
	IntegerVariable y1 = makeIntVar("Y1", 0, 10);
	
	m.addConstraint(occurrence(3, y1, new IntegerVariable[]{x1, x2, x3, x4, x5, x6, x7}));
	
	s.read(m);
	s.solve();
\end{lstlisting} 

%\part{occurrencemax}
\label{occurrencemax}
\hypertarget{occurrencemax}{}

\section{occurrenceMax (constraint)}\label{occurrencemax:occurrencemaxconstraint}\hypertarget{occurrencemax:occurrencemaxconstraint}{}
\begin{notedef}
  \texttt{occurrenceMax}$(v,z,x)$ states that $z$ is at most equal to the number of elements in $x$ with value $v$:
$$z\le|\{i\ |\ x_i=v\}|$$   
\end{notedef}
  This is a specialization of the \texttt{globalCardinality} constraint.

\begin{itemize}
	\item \textbf{API}: \mylst{occurrenceMax(int v, IntegerVariable z, IntegerVariable... x)}
	\item \textbf{return type} : \texttt{Constraint}
	\item \textbf{options} :\emph{n/a}
	\item \textbf{favorite domain} : \emph{to complete}
	\item \textbf{references} :\\
      global constraint catalog: \href{http://www.emn.fr/x-info/sdemasse/gccat/Ccount.html}{count}
\end{itemize}

\textbf{Example}:
\begin{lstlisting}
	Model m = new CPModel();
	Solver s = new CPSolver();
	 
	IntegerVariable x1 = makeIntVar("X1", 0, 10);
	IntegerVariable x2 = makeIntVar("X2", 0, 10);
	IntegerVariable x3 = makeIntVar("X3", 0, 10);
	IntegerVariable x4 = makeIntVar("X4", 0, 10);
	IntegerVariable x5 = makeIntVar("X5", 0, 10);
	IntegerVariable x6 = makeIntVar("X6", 0, 10);
	IntegerVariable x7 = makeIntVar("X7", 0, 10);
	IntegerVariable y1 = makeIntVar("Y1", 0, 10);
	 
	m.addConstraint(occurrenceMax(3, y1, new IntegerVariable[]{x1, x2, x3, x4, x5, x6, x7}));
	 
	s.read(m);
	s.solve();
\end{lstlisting} 

%\part{occurrencemin}
\label{occurrencemin}
\hypertarget{occurrencemin}{}

\section{occurrenceMin (constraint)}\label{occurrencemin:occurrenceminconstraint}\hypertarget{occurrencemin:occurrenceminconstraint}{}
\begin{notedef}
  \texttt{occurrenceMin}$(v,z,x)$ states that $z$ is at least equal to the number of elements in $x$ with value $v$:
$$z\ge|\{i\ |\ x_i=v\}|$$   
\end{notedef}
  This is a specialization of the \texttt{globalCardinality} constraint.

\begin{itemize}
	\item \textbf{API}: \mylst{occurrenceMin(int v, IntegerVariable z, IntegerVariable... x)}
	\item \textbf{return type} : \texttt{Constraint}
	\item \textbf{options} :\emph{n/a}
	\item \textbf{favorite domain} : \emph{to complete}
	\item \textbf{references} :\\
      global constraint catalog: \href{http://www.emn.fr/x-info/sdemasse/gccat/Ccount.html}{count}
\end{itemize}

\textbf{Example}:
\begin{lstlisting}
	Model m = new CPModel();
	Solver s = new CPSolver();
	 
	IntegerVariable x1 = makeIntVar("X1", 0, 10);
	IntegerVariable x2 = makeIntVar("X2", 0, 10);
	IntegerVariable x3 = makeIntVar("X3", 0, 10);
	IntegerVariable x4 = makeIntVar("X4", 0, 10);
	IntegerVariable x5 = makeIntVar("X5", 0, 10);
	IntegerVariable x6 = makeIntVar("X6", 0, 10);
	IntegerVariable x7 = makeIntVar("X7", 0, 10);
	IntegerVariable y1 = makeIntVar("Y1", 0, 10);
	 
	m.addConstraint(occurrenceMin(3, y1, new IntegerVariable[]{x1, x2, x3, x4, x5, x6, x7}));
	 
	s.read(m);
	s.solve();
\end{lstlisting} 

%\part{oppositesign}
\label{oppositesign}
\hypertarget{oppositesign}{}

\section{oppositeSign (constraint)}\label{oppositesign:oppositesignconstraint}\hypertarget{oppositesign:oppositesignconstraint}{}


\begin{notedef}
  \texttt{oppositeSign}$(x,y)$ states that the two arguments have opposite signs:
$$xy\le 0$$
\end{notedef}

\begin{notedef}
  0 is considered as both sign, if one argument is equal to 0, the constraint is not satisfied.
\end{notedef}


\begin{itemize}
	\item \textbf{API} : \mylst{oppositeSign(IntegerExpressionVariable x, IntegerExpressionVariable y)}
	\item \textbf{return type} : \texttt{Constraint}
	\item \textbf{options} :\emph{n/a}
	\item \textbf{favorite domain} : \emph{to complete}
\end{itemize}

\textbf{Example}:
\lstinputlisting{java/coppositesign.j2t} 

%\part{or}
\label{or}
\hypertarget{or}{}

\section{or (constraint)}\label{or:orconstraint}\hypertarget{or:orconstraint}{}
\begin{notedef}
  \texttt{or}$(c_1,\ldots,c_n)$ states that at least one constraint in arguments is satisfied:
$$ c_1 \lor c_2 \lor\ldots\lor c_n$$

  \texttt{or}$(b_1,\ldots,b_n)$ states that at least one boolean variable in argument is true:
$$ (b_1=1) \lor (b_2=1) \lor\ldots\lor (b_n=1)$$
\end{notedef}

\begin{itemize}
\item \textbf{API} : 
\begin{itemize}
\item \mylst{or(Constraint... c)}
\item \mylst{or(IntegerVariable... b)}
\end{itemize}
	\item \textbf{return type} : \texttt{Constraint}
	\item \textbf{options} : \emph{n/a}
	\item \textbf{favorite domain} : \emph{n/a}
	\item \textbf{references} :\\
  global constraint catalog: \href{http://www.emn.fr/z-info/sdemasse/gccat/Cor.html}{\tt or}
\end{itemize}

\textbf{Examples:}
\begin{itemize}
	\item example1:
\end{itemize}
\lstinputlisting{java/cor1.j2t}
\begin{itemize}
	\item example2
\end{itemize}
\lstinputlisting{java/cor2.j2t}

%\part{pack}
\label{pack}
\hypertarget{pack}{}

\section{pack (constraint)}\label{pack:packconstraint}\hypertarget{pack:packconstraint}{}

\begin{notedef}
  \texttt{pack(items, load, bin, size)} states that a collection of items is packed into different bins, such that the total size of the items in each bin does not exceed the bin capacity:
$$ \mathtt{load}[b] = \sum_{i\in\mathtt{items}[b]} \mathtt{size}[i],\quad\forall \text{ bin } b $$
%and
$$ i\in\mathtt{items}[b]\ \iff\ \mathtt{bin}[i]=b,\quad\forall \text{ bin } b,\ \forall \text{ item } i $$
\end{notedef}
%\texttt{pack}$(sizes, n, )$ states a collection of items (each of them having a specific size) is packed into different bins of given capacity such that the total weight of the items in each bin does not exceed the bin capacity.
\texttt{pack} is a \href{http://www.emn.fr/x-info/sdemasse/gccat/Cbin_packing.html}{bin packing constraint} based on \cite{ShawCP04}. 

\begin{itemize}
	\item \textbf{API} :
	\begin{itemize}
		\item \mylst{pack(SetVariable[] items, IntegerVariable[] load, IntegerVariable[] bin, IntegerConstantVariable[] size, String... options)}
		\item \mylst{pack(PackModeler modeler,String... options)}: PackModeler is a high-level modeling object.
		\item \mylst{pack(int[] sizes, int nbBins, int capacity, String... options)}: build instance with PackModeler.
	\end{itemize}
	\item \textbf{Variables}:
	\begin{itemize}
		\item \texttt{SetVariable[] items: items}$[b]$ is the set of items packed into bin $b$.
		\item \texttt{IntegerVariable[] load: load}$[b]$ is the total size of the items packed into bin $b$.
		\item \texttt{IntegerVariable[] bin: bin}$[i]$ is the bin where item $i$ is packed into.
		\item \texttt{IntegerConstantVariable[] size: size}$[i]$ is the size of item $i$.
	\end{itemize}
	\item \textbf{return type} : \texttt{Constraint}
	\item \textbf{options} : 	
      \begin{itemize}
      \item \hyperlink{cpackar:cpackaroptions}{SettingType.ADDITIONAL\_RULES.getOptionName()}: additional filtering rules \emph{recommended}
      \item \hyperlink{cpackdlb:cpackdlboptions}{SettingType.DYNAMIC\_LB.getOptionName()}: feasibility tests based on dynamic lower bounds for 1D-bin packing
      \item \hyperlink{cpackfill:cpackfilloptions}{SettingType.FILL\_BIN.getOptionName()}: dominance rule: fill a bin when an item fit into pertfectly equal-sized items and bins must be equivalent
      \item \hyperlink{cpacklbe:cpacklbeoptions}{SettingType.LAST\_BINS\_EMPTY.getOptionName()}: empty bins are the last ones 
      \end{itemize}
	\item \textbf{favorite domain} : \emph{to complete}
	\item \textbf{references} :
      \begin{itemize}
      \item \cite{ShawCP04}: \emph{A constraint for bin packing}
      \item global constraint catalog: \href{http://www.emn.fr/x-info/sdemasse/gccat/Cbin_packing.html}{bin\_packing} (variant)
      \end{itemize}
\end{itemize}

\textbf{Example}:

Take a look at \emph{samples.pack} to see advanced use of the constraint.
\lstinputlisting{java/cpack_import.j2t}
\lstinputlisting{java/cpack.j2t}

%\part{precedencereified}
\label{precedencereified}
\hypertarget{precedencereified}{}

\section{precedenceReified (constraint)}\label{precedencereified:precedencereifiedconstraint}\hypertarget{precedencereified:precedencereifiedconstraint}{}
\begin{notedef}
  \texttt{precedenceReified}$(x_1,d,x_2,b)$ states that $x_1$ plus duration $d$ is less than or equal to $x_2$ requires boolean $b$ to be true:
  $$b\quad\iff\quad x_1 + d \le x_2$$
\end{notedef}

\begin{itemize}
	\item \textbf{API} : \mylst{precedenceReified(IntegerVariable x1, int d, IntegerVariable x2, IntegerVariable b)}
	\item \textbf{return type} : \texttt{Constraint}
	\item \textbf{options} :\emph{n/a}
	\item \textbf{favorite domain} : \emph{to complete}
\end{itemize}

\textbf{Example}:
\mylst{//TODO: complete} 

\section{precedenceimplied (constraint)}\label{precedenceimplied:precedenceimpliedconstraint}\hypertarget{precedenceimplied:precedenceimpliedconstraint}{}
\emph{To complete}

\section{precedence (constraint)}\label{precedence:precedenceconstraint}\hypertarget{precedence:precedenceconstraint}{}
\emph{To complete}

\section{precedencedisjoint (constraint)}\label{precedencedisjoint:precedencedisjointconstraint}\hypertarget{precedencedisjoint:precedencedisjointconstraint}{}
\emph{To complete}

%\part{regular}
\label{regular}
\hypertarget{regular}{}

\section{regular (constraint)}\label{regular:regularconstraint}\hypertarget{regular:regularconstraint}{}
\begin{notedef}
  \texttt{regular}$(x,\mathcal{L}(\Pi))$ states that sequence $x$ is a word belonging to the regular language $\mathcal{L}(\Pi)$:
% recognized by a deterministic finite automaton (DFA) or a regular expression $\Pi$:
$$(x_1,\ldots,x_n)\in\mathcal{L}(\Pi)$$
\end{notedef}

The accepting language can be specified either by a deterministic finite automaton (DFA), a list of feasible or infeasible tuples, or a regular expression:
\begin{description}
\item[DFA:] Automaton $\Pi$ is defined on a given \emph{alphabet} $\Sigma\subseteq\Z$ by a set $Q=\{0,\ldots,m\}$ of \emph{states}, a subset $A\subseteq Q$ of \emph{final} or \emph{accepting states} and a table $\Delta\subseteq Q\!\times\!\Sigma\!\times Q$ of \emph{transitions} between states. $\Delta$ is encoded as \texttt{List<Transition>} where a Transition object $\delta=\texttt{new Transition}(q_i,\sigma,q_j)$ is made of three integers expressing the ingoing state $q_i$, the label $\sigma$, and the outgoing state $q_j$.
Automaton $\Pi$ is a DFA if $\Delta$ is finite and if it has only one initial state (here, state $0$ is considered as the unique initial state) and no two transitions sharing the same ingoing state and the same label.
\item[FiniteAutomaton:] is another API for building a DFA (manually, or from a regular expression, or from a \mylst{dk.brics.Automaton}) and operating on them (intersection, union, complement) in a more flexible way. Using this API leads to another implementation of the constraint: \mylst{FastRegular}. See \hyperlink{costregular}{\texttt{costRegular}} for a short API of \texttt{FiniteAutomaton}. 
\item[feasible tuples:] \emph{regular} can be used as an extensional constraint. Given the list of \emph{feasible} tuples for sequence $x$, this API builds a DFA from the list, and then enforces GAC on the constraint. Using \texttt{regular} can be more efficient than a standard GAC algorithm on tables of tuples if the tuples are structured so that the resulting DFA is compact. The DFA is built from the list of tuples by computing incrementally the minimal DFA after each addition of tuple. 
\item[infeasible tuples:] An another API allows to specify the list of \emph{infeasible} tuples and then builds the corresponding feasible DFA. This operation requires to know the entire alphabet, hence this API has two mandatory table fields \emph{min} and \emph{max} defining the minimum and maximum values of each variable $x_i$.
\item[regular expression:] Finally, the \texttt{regular} constraint can be based on a \href{http://en.wikipedia.org/wiki/regularexpression}{regular expression}, such as \mylst{String regexp = "(1\|2)3\{4\}5*";}. This expression recognizes any sequences starting by one 1 or one 2, then four consecutive 3 followed by any (possibly empty) sequences of 5.
\end{description}

\todo{Warning ! DFA and FiniteAutomaton are both based on the dk.brics library. The construction of these objects is non-deterministic and the order the filtering occur (not the result) may vary at each execution. This may results in different first solutions when branching dynamically using  weighted degrees-base heuristics for example.}

\begin{itemize}
	\item \textbf{API} :
	\begin{itemize}
		\item \mylst{regular(IntegerVariable[] x, FiniteAutomaton pi)}
		\item \mylst{regular(IntegerVariable[] x, DFA pi)}
		\item \mylst{regular(IntegerVariable[] x, List<int[]> feasTuples)}
		\item \mylst{regular(IntegerVariable[] x, List<int[]> infeasTuples, int[] min, int[] max)}
		\item \mylst{regular(IntegerVariable[] x, String regexp)}
	\end{itemize}
	\item \textbf{return type} : \texttt{Constraint}
	\item \textbf{options} :\emph{n/a}
	\item \textbf{favorite domain} : \emph{to complete}
	\item \textbf{references} :\\
       \cite{PesantCP04}: \emph{A regular language membership constraint}
\end{itemize}

\textbf{Examples}:
\begin{itemize}
	\item example with \texttt{FiniteAutomaton}: see \hyperlink{costregular:costregularconstraint}{\texttt{costRegular}}.
	\item example 1 with DFA:
\end{itemize}
\lstinputlisting{java/cregular1_import.j2t}
\lstinputlisting{java/cregular1.j2t}

\begin{itemize}
	\item example 2 with feasible tuples:
\end{itemize}
\lstinputlisting{java/cregular2.j2t}

\begin{itemize}
	\item example 3 with regular expression:
\end{itemize}
\lstinputlisting{java/cregular3.j2t}

\section{reifiedAnd (constraint)}\label{reifiedand:reifiedandconstraint}\hypertarget{reifiedand:reifiedandconstraint}{}
\emph{To complete}

\section{reifiedConstraint (constraint)}\label{reifiedconstraint:reifiedconstraintconstraint}\hypertarget{reifiedconstraint:reifiedconstraintconstraint}{}
\begin{notedef}
  \begin{itemize}
  \item \texttt{reifiedConstraint}$(b,C)$ states that boolean $b$ is true if and only if constraint $C$ holds:
  $$(b=1)\ \iff\ C$$
  \item \texttt{reifiedConstraint}$(b,C_1,C_2)$ states that boolean $b$ is true if and only if $C_1$ holds, and $b$ is false if and only if $C_2$ holds ($C_2$ must be the negation of constraint of $C_1$):
$$(b\land C_1) \lor (\neg b \land C_2)$$
  \end{itemize}
\end{notedef}

\begin{itemize}
	\item \textbf{API} :
	\begin{itemize}
		\item \mylst{reifiedConstraint(IntegerVariable b, Constraint c)}
		\item \mylst{reifiedConstraint(IntegerVariable b, Constraint c1, Constraint c2)}
	\end{itemize}
	\item \textbf{return type} : \texttt{Constraint}
	\item \textbf{options} : \emph{n/a}
	\item \textbf{favorite domain} : \emph{n/a}
\end{itemize}

The constraint $C$ to reify has to provide its negation $\neg C$ (the negation is needed for propagation). 
Most basic constraints of Choco provides their negation by default, and can then be reified using the first API.
The second API attends to reify user-defined constraints as it allows the user to directly specify the negation constraint. 

The {\tt  reifiedConstraint} filter algorithm:
\begin{enumerate}
	\item if $b$ is instantiated to 1 (resp. to 0), then $C$ (resp. $\neg C$) is propagated
	\item otherwise
	\begin{enumerate}
		\item if $C$ is entailed, $b$ is set to 1
		\item else if $C$ is failed, $b$ is set to 0.
	\end{enumerate}
\end{enumerate}


\textbf{Example}:
\lstinputlisting{java/creifiedintconstraint.j2t}


%%\part{reifiedintconstraint}
\label{reifiedintconstraint}
\hypertarget{reifiedintconstraint}{}

\section{reifiedIntConstraint (constraint)}\label{reifiedintconstraint:reifiedintconstraintconstraint}\hypertarget{reifiedintconstraint:reifiedintconstraintconstraint}{}
\begin{notedef}
  \begin{itemize}
  \item \texttt{reifiedIntConstraint}$(b,c)$ states that boolean $b$ is true if and only if constraint $c$ holds:
  $$b\ \iff\ c$$
  \item \texttt{reifiedIntConstraint}$(b,c_1,c_2)$ states that boolean $b$ is true if and only if $c_1$ holds, and $b$ is false if and only if $c_2$ holds ($c_2$ must be the opposite constraint of $c_1$):
$$(b\land c_1) \lor (\neg b \land c_2)$$
  \end{itemize}
\end{notedef}

\begin{itemize}
	\item \textbf{API} :
	\begin{itemize}
		\item \mylst{reifiedIntConstraint(IntegerVariable b, Constraint c)}
		\item \mylst{reifiedIntConstraint(IntegerVariable b, Constraint c1, Constraint c2)}
	\end{itemize}
	\item \textbf{return type} : \texttt{Constraint}
	\item \textbf{options} : \emph{n/a}
	\item \textbf{favorite domain} : \emph{n/a}
\end{itemize}

Parameter \emph{b} is a boolean variable (enumerated domain with two values $\{0,1\}$) and \emph{c} is a constraint over Integer variables.

The constraint $c$ to reify has to provide its opposite (the opposite is needed for propagation). Most basic constraints of Choco provides their opposite by default, and can then be reified using the first API.
The second API attends to reify user-defined constraints as it allows the user to directly specify the opposite constraint.

\textbf{Example}:

\lstinputlisting{java/creifiedintconstraint.j2t}

\section{reifiedLeftImp (constraint)}\label{reifiedleftimp:reifiedleftimpconstraint}\hypertarget{reifiedleftimp:reifiedleftimpconstraint}{}
\emph{To complete}

\section{reifiedNot (constraint)}\label{reifiednot:reifiednotconstraint}\hypertarget{reifiednot:reifiednotconstraint}{}
\emph{To complete}

\section{reifiedOr (constraint)}\label{reifiedor:reifiedorconstraint}\hypertarget{reifiedor:reifiedorconstraint}{}
\emph{To complete}

\section{reifiedRightImp (constraint)}\label{reifiedrightimp:reifiedrightimpconstraint}\hypertarget{reifiedrightimp:reifiedrightimpconstraint}{}
\emph{To complete}

\section{reifiedXnor (constraint)}\label{reifiedxnor:reifiedxnorconstraint}\hypertarget{reifiedxnor:reifiedxnorconstraint}{}
\emph{To complete}

\section{reifiedXor (constraint)}\label{reifiedxor:reifiedxorconstraint}\hypertarget{reifiedxor:reifiedxorconstraint}{}
\emph{To complete}

%\part{relationpairac}
\label{relationpairac}
\hypertarget{relationpairac}{}

\section{relationPairAC (constraint)}\label{relationpairac:relationpairacconstraint}\hypertarget{relationpairac:relationpairacconstraint}{}
\begin{notedef}
  \texttt{relationPairAC}$(x,y,rel)$ states an extensional binary constraint on $(x,y)$ defined by the binary relation $rel$:
$$(x,y)\in rel$$
\end{notedef}
Many constraints of the same kind often appear in a model. Relations can therefore often be shared among many constraints to spare memory.

The API is duplicated to allow definition of options.

\begin{itemize}
	\item \textbf{API} :
	\begin{itemize}
		\item \mylst{relationPairAC(IntegerVariable x, IntegerVariable y, BinRelation rel)}
		\item \mylst{relationPairAC(String options, IntegerVariable x, IntegerVariable y, BinRelation rel)}
	\end{itemize}
	\item \textbf{return type} : \texttt{Constraint}
	\item \textbf{options} :
	\begin{itemize}
		\item \emph{no option} : use AC3 (default arc consistency)
		\item \texttt{CPOptions.C_EXT_AC3}: to get AC3 algorithm (searching from scratch for supports on all values)
		\item \texttt{CPOptions.C_EXT_AC2001}: to get AC2001 algorithm (maintaining the current support of each value)
		\item \texttt{CPOptions.C_EXT_AC32}: to get AC3rm algorithm (maintaining the current support of each value in a non backtrackable way)
		\item \texttt{CPOptions.C_EXT_AC322}: to get AC3 with the used of \texttt{BitSet} to know if a support still exists
	\end{itemize}
	\item \textbf{favorite domain} : \emph{to complete}
\end{itemize}

\textbf{Example}:
\lstinputlisting{java/crelationpairac_import.j2t}
\lstinputlisting{java/ccoupletest.j2t}
\lstinputlisting{java/crelationpairac.j2t}

%\part{relationtupleac}
\label{relationtupleac}
\hypertarget{relationtupleac}{}

\section{relationTupleAC (constraint)}\label{relationtupleac:relationtupleacconstraint}\hypertarget{relationtupleac:relationtupleacconstraint}{}
\begin{notedef}
  \texttt{relationTupleAC}$(x,rel)$ states an extensional constraint on $(x_1,\ldots,x_n)$ defined by the $n$-ary relation $rel$, and then enforces arc consistency:
$$(x_1,\ldots,x_n)\in rel$$
\end{notedef}
Many constraints of the same kind often appear in a model. Relations can therefore often be shared among many constraints to spare memory.
The API is duplicated to define options.

\begin{itemize}
	\item \textbf{API}:
	\begin{itemize}
		\item \mylst{relationTupleAC(IntegerVariable[] x, LargeRelation rel)}
		\item \mylst{relationTupleAC(String options, IntegerVariable[] x, LargeRelation rel)}
	\end{itemize}
	\item \textbf{return type}: \texttt{Constraint}
	\item \textbf{options} :
	\begin{itemize}
		\item \emph{no option}: use AC32 (default arc consistency)
		\item \texttt{cp:ac32}: to get AC3rm algorithm (maintaining the current support of each value in a non backtrackable way)
		\item \texttt{cp:ac2001}: to get AC2001 algorithm (maintaining the current support of each value)
		\item \texttt{cp:ac2008}: to get AC2008 algorithm (maintained by STR)
	\end{itemize}
	\item \textbf{favorite domain} : \emph{to complete}
\end{itemize}

\textbf{Example} :
\mylst{// TODO : add example}

%\part{relationtuplefc}
\label{relationtuplefc}
\hypertarget{relationtuplefc}{}

\section{relationTupleFC (constraint)}\label{relationtuplefc:relationtuplefcconstraint}\hypertarget{relationtuplefc:relationtuplefcconstraint}{}
\begin{notedef}
  \texttt{relationTupleFC}$(x,rel)$ states an extensional constraint on $(x_1,\ldots,x_n)$ defined by the $n$-ary relation $rel$, and then enforces forward checking:
$$(x_1,\ldots,x_n)\in rel$$
\end{notedef}
Many constraints of the same kind often appear in a model. Relations can therefore often be shared among many constraints to spare memory.

\begin{itemize}
	\item \textbf{API}: \mylst{relationTupleFC(IntegerVariable[] x, LargeRelation rel)}
	\item \textbf{return type}: \texttt{Constraint}
	\item \textbf{options} : \emph{n/a}
	\item \textbf{favorite domain} : \emph{to complete}
\end{itemize}

\textbf{Example} :
\mylst{// TODO : add example}

%\part{samesign}
\label{samesign}
\hypertarget{samesign}{}

\section{sameSign (constraint)}\label{samesign:samesignconstraint}\hypertarget{samesign:samesignconstraint}{}
\todo{verify case 0}

\begin{notedef}
  \texttt{sameSign}$(x,y)$ states that the two arguments have the same sign:
$$xy\ge 0$$
\end{notedef}

\begin{itemize}
	\item \textbf{API} : \mylst{sameSign(IntegerExpressionVariable x, IntegerExpressionVariable y)}
	\item \textbf{return type} : \texttt{Constraint}
	\item \textbf{options} :\emph{n/a}
	\item \textbf{favorite domain} : \emph{to complete}
\end{itemize}

\textbf{Example}:
\lstinputlisting{java/csamesign.j2t}

%\part{setdisjoint}
\label{setdisjoint}
\hypertarget{setdisjoint}{}

\section{setDisjoint (constraint)}\label{setdisjoint:setdisjointconstraint}\hypertarget{setdisjoint:setdisjointconstraint}{}
\begin{notedef}
  \texttt{setDisjoint}$(s_1,s_2)$ states that the two set arguments are disjoint:
$$s_1\cap s_2=\emptyset$$
\end{notedef}

\begin{itemize}
	\item \textbf{API} : \mylst{setDisjoint(SetVariable s1, SetVariable s2)}
	\item \textbf{return type} : \texttt{Constraint}
	\item \textbf{options} :\emph{n/a}
	\item \textbf{favorite domain} : \emph{to complete}
\end{itemize}

\textbf{Example}:
\begin{lstlisting}
	Model m = new CPModel();
	Solver s = new CPSolver();
	SetVar x = makeSetVar("X", 1, 3);
	SetVar y = makeSetVar("Y", 1, 3);
	Constraint c1 = setDisjoint(x, y);
	m.addConstraint(c1);
	s.read(m);
	s.solveAll();
\end{lstlisting} 

%\part{setinter}
\label{setinter}
\hypertarget{setinter}{}

\section{setInter (constraint)}\label{setinter:setinterconstraint}\hypertarget{setinter:setinterconstraint}{}
\begin{notedef}
  \texttt{setInter}$(s_1,s_2,s_3)$ states that the third set $s_3$ is exactly the intersection of the two first sets:
$$s_1\cap s_2=s_3$$
\end{notedef}

\begin{itemize}
	\item \textbf{API} : \mylst{setInter(SetVariable s1, SetVariable s2, SetVariable s3)}
	\item \textbf{return type} : \texttt{Constraint}
	\item \textbf{options} :\emph{n/a}
	\item \textbf{favorite domain} : \emph{to complete}
\end{itemize}

\textbf{Example}:
\lstinputlisting{java/csetinter.j2t}
%\part{setunion}
\label{setunion}
\hypertarget{setunion}{}

\section{setUnion (constraint)}\label{setunion:setunionconstraint}\hypertarget{setunion:setunionconstraint}{}
\begin{notedef}
  \texttt{setUnion}$(sv,s_{union})$ states that the $s_{union}$ set is exactly the union of the sets $sv$:
$$sv_1\cup sv_2 \cup \ldots sv_i \cup sv_{i+1} \ldots \cup sv_n=s_{union}$$
\end{notedef}

\begin{itemize}
	\item \textbf{API} : 
	\begin{itemize}
		\item \mylst{setUnion(SetVariable s1, SetVariable s2, SetVariable union)}
		\item \mylst{setUnion(SetVariable[] sv, SetVariable union)}
	\end{itemize}
	\item \textbf{return type} : \texttt{Constraint}
	\item \textbf{options} :\emph{n/a}
	\item \textbf{favorite domain} : \emph{to complete}
\end{itemize}

\textbf{Example}:
\lstinputlisting{java/csetunion.j2t}
%\part{sorting}
\label{sorting}
\hypertarget{sorting}{}

\section{sorting (constraint)}\label{sorting:sortingconstraint}\hypertarget{sorting:sortingconstraint}{}

\section{startsAfter (constraint)}\label{startsafter:startsafterconstraint}\hypertarget{startsafter:startsafterconstraint}{}
\begin{notedef}
  \texttt{startsAfter}$(T,c)$ states that the task variable $t$ starts after time $c$:
  $$T.start \ge c$$
\end{notedef}

\begin{itemize}
	\item \textbf{API} :\mylst{startsAfter(final TaskVariable t, final int c) }
	\item \textbf{return type} : \texttt{Constraint}
	\item \textbf{options} : \emph{n/a}
	\item \textbf{favorite domain} : \emph{n/a}.
\end{itemize}

\textbf{Examples:}
%\lstinputlisting{java/ceq1.j2t}
\emph{to complete}

\section{startsAfterBegin (constraint)}\label{startsafterbegin:startsafterbeginconstraint}\hypertarget{startsafterbegin:startsafterbeginconstraint}{}
\begin{notedef}
  \texttt{startsAfterBegin}$(t_1,t_2,c)$ states that the task variable $t_1$ starts after the beginning of $t_2$, minus $c$:
  $$s_{t_{1}} \ge s_{t_{2}} - c$$
\end{notedef}

\begin{itemize}
	\item \textbf{API} :\mylst{startsAfterBegin(TaskVariable t1, TaskVariable t2, int c)}
	\item \textbf{return type} : \texttt{Constraint}
	\item \textbf{options} : \emph{n/a}
	\item \textbf{favorite domain} : \emph{n/a}.
\end{itemize}

\textbf{Examples:}
%\lstinputlisting{java/ceq1.j2t}
\emph{to complete}

\section{startsAfterEnd (constraint)}\label{startsafterend:startsafterendconstraint}\hypertarget{startsafterend:startsafterendconstraint}{}
\begin{notedef}
\texttt{startsAfterEnd}$(T_1,T_2,c)$  states that task $T_1$ starts after the end time of $T_2$ minus $c$:
  $$T_{1}.start \ge T_{2}.end - c$$
\end{notedef}

\begin{itemize}
	\item \textbf{API} :\mylst{startsAfterEnd(TaskVariable t1, TaskVariable t2, int c)}
	\item \textbf{return type} : \texttt{Constraint}
	\item \textbf{options} : \emph{n/a}
	\item \textbf{favorite domain} : \emph{n/a}.
\end{itemize}

\textbf{Examples:}
%\lstinputlisting{java/ceq1.j2t}
\emph{to complete}

\section{startsBefore (constraint)}\label{startsbefore:startsbeforeconstraint}\hypertarget{startsbefore:startsbeforeconstraint}{}
\begin{notedef}
  \texttt{startsBefore}$(T,c)$ states that the task variable $T$ starts before time $c$:
  $$T.start \le c$$
\end{notedef}

\begin{itemize}
	\item \textbf{API} :\mylst{startsBefore(final TaskVariable t, final int c) }
	\item \textbf{return type} : \texttt{Constraint}
	\item \textbf{options} : \emph{n/a}
	\item \textbf{favorite domain} : \emph{n/a}.
\end{itemize}

\textbf{Examples:}
%\lstinputlisting{java/ceq1.j2t}
\emph{to complete}

\section{startsBeforeBegin (constraint)}\label{startsbeforebegin:startsbeforebeginconstraint}\hypertarget{startsbeforebegin:startsbeforebeginconstraint}{}
\begin{notedef}
\texttt{startsBeforeBegin}$(T_1,T_2,c)$  states that task $T_1$ starts before the start time of $T_2$ minus $c$:
  $$T_{1}.start \le T_{2}.start - c$$
\end{notedef}

\begin{itemize}
	\item \textbf{API} :\mylst{startsBeforeBegin(TaskVariable t1, TaskVariable t2, int c)}
	\item \textbf{return type} : \texttt{Constraint}
	\item \textbf{options} : \emph{n/a}
	\item \textbf{favorite domain} : \emph{n/a}.
\end{itemize}

\textbf{Examples:}
%\lstinputlisting{java/ceq1.j2t}
\emph{to complete}

\section{startsBeforeEnd (constraint)}\label{startsbeforeend:startsbeforeendconstraint}\hypertarget{startsbeforeend:startsbeforeendconstraint}{}
\begin{notedef}
  \texttt{startsBeforeEnd}$(t_1,t_2,c)$ states that the task variable $t_1$ starts before the end of $t_2$, minus $c$:
  $$s_{t_{1}} \le e_{t_{2}} - c$$
\end{notedef}

\begin{itemize}
	\item \textbf{API} :\mylst{startsBeforeEnd(TaskVariable t1, TaskVariable t2, int c)}
	\item \textbf{return type} : \texttt{Constraint}
	\item \textbf{options} : \emph{n/a}
	\item \textbf{favorite domain} : \emph{n/a}.
\end{itemize}

\textbf{Examples:}
%\lstinputlisting{java/ceq1.j2t}
\emph{to complete}

\section{startsBetween (constraint)}\label{startsbetween:startsbetweenconstraint}\hypertarget{startsbetween:startsbetweenconstraint}{}
\begin{notedef}
  \texttt{startsBetween}$(T, c_1, c_2)$ states that task $T$ starts between times $c_1$ and $c_2$:
  $$c_1 \le T.start \le c_2$$
\end{notedef}

\begin{itemize}
	\item \textbf{API} :\mylst{startsBetween(TaskVariable t, int min, int max)}
	\item \textbf{return type} : \texttt{Constraint}
	\item \textbf{options} : \emph{n/a}
	\item \textbf{favorite domain} : \emph{n/a}.
\end{itemize}

\textbf{Examples:}
%\lstinputlisting{java/ceq1.j2t}
\emph{to complete}

%\part{stretchpath}
\label{stretchpath}
\hypertarget{stretchpath}{}

\section{stretchPath (constraint)}\label{stretchpath:stretchpathconstraint}\hypertarget{stretchpath:stretchpathconstraint}{}


\begin{notedef}
  A \emph{stretch} in a sequence $x$ is a maximum subsequence of (consecutive) identical values.  
  \texttt{stretchPath}$(param,x)$ enforces the minimal and maximal length of the stretches in sequence $x$ of any values given in $param$:
Consider the sequence $x$ as a concatenation of stretches $x^1.x^2\ldots x^k$ with $v^i$ and $l^i$ being respectively the value and the length of stretch $x^i$,
$$\forall i\in\{1,\ldots,k\},\ \forall j,\quad param[j][0]=v^i\quad\implies\quad param[j][1]\le l^i\le param[j][2]$$
\end{notedef}


Useful for Rostering Problems. \texttt{stretchPath} is implemented by a \hyperlink{regular}{\texttt{regular}} constraint that performs GAC. The bounds on the stretch lengths are defined by $param$ a list of triples of integers: $[value, min, max]$ specifying the minimal and maximal lengths of any stretch of the corresponding value. 

This API requires a Java library on automaton available on \href{http://www.brics.dk/automaton/}{http://www.brics.dk/automaton/}. (It is contained in the Choco jar file.)

\begin{itemize}
	\item \textbf{API} : \mylst{stretchPath(List<int[]> param, IntegerVariable... x)}
	\item \textbf{return type} : Constraint
	\item \textbf{options} :\emph{n/a}
	\item \textbf{favorite domain} : \emph{to complete}
	\item \textbf{references} :
      \begin{itemize}
      \item \cite{PesantCP04}: \emph{A regular language membership constraint}
      \item global constraint catalog: \href{http://www.emn.fr/x-info/sdemasse/gccat/Cstretch_path.html}{stretch\_path}
      \end{itemize}
\end{itemize}


\textbf{Example}:
\begin{lstlisting}
  Model m = new CPModel();
  int n = 7;
  IntegerVariable[] vars = makeIntVarArray("v", n, 0, 2);
	
  //define the stretches
  ArrayList<int[]> lgt = new ArrayList<int[]>();
  lgt.add(new int[]{0, 2, 2}); // stretches of value 0 are of length 2
  lgt.add(new int[]{1, 2, 3}); // stretches of value 1 are of length 2 or 3
  lgt.add(new int[]{2, 2, 2}); // stretches of value 2 are of length 2
	
  m.addConstraint(stretchPath(lgt, vars));
	
  Solver s = new CPSolver();
  s.read(m);
  s.solve();
\end{lstlisting}

%\part{times}
\label{times}
\hypertarget{times}{}

\section{times (constraint)}\label{times:timesconstraint}\hypertarget{times:timesconstraint}{}
\begin{notedef}
  \texttt{times}$(x_1, x_2, x_3)$ states that the third argument is equal to the product of the two arguments:
$$x_3=x_1\times x_2.$$
\end{notedef}

\begin{itemize}
	\item \textbf{API}:
	\begin{itemize}
		\item \mylst{times(IntegerVariable x1, IntegerVariable x2, IntegerVariable x3)}
		\item \mylst{times(int x1, IntegerVariable x2, IntegerVariable x3)}
		\item \mylst{times(IntegerVariable x1, int x2, IntegerVariable x3)}
	\end{itemize}
	\item \textbf{return type} : \texttt{Constraint}
	\item \textbf{option} : \emph{n/a}
	\item \textbf{favorite domain}: bound
\end{itemize}

\textbf{Example}:
\lstinputlisting{java/ctimes.j2t}
%\part{tree}
\label{tree}
\hypertarget{tree}{}

\section{tree (constraint)}\label{tree:treeconstraint}\hypertarget{tree:treeconstraint}{}

Let $G=(V,A)$ be a digraph on $V=\{1,\ldots,n\}$. $G$ can be modeled by a sequence of domain variables $x=(x_1,\dots,x_n)\in V^n$ -- the \emph{successors} variables -- whose respective domains are given by $D_i=\{j\in V\ |\ (i,j)\in A\}$. Conversely, when instantiated, $x$ defines a subgraph $G_x=(V,A_x)$ of $G$ with $A_x=\{(i,x_i)\ |\ i\in V\}\subseteq A$. Such a subgraph has one particularity: any connected component of $G_x$ contains either no loop -- and then it contains a cycle -- or exactly one loop $x_i=i$ and then it is a \emph{tree} of root $i$ (literally, it is an anti-arborescence as there exists a path from each node to $i$ and $i$ has a loop).

\begin{notedef}
  \texttt{tree}$(x,restrictions)$ is a vertex-disjoint graph partitioning constraint. It states that $G_x$ is a forest (its connected components are trees) that satisfies some conditions specified by $restrictions$.
\texttt{tree} deals with several kinds of graph restrictions on:
\begin{itemize}
	\item the number of trees
	\item the number of proper trees (a tree is proper if it contains more than 2 nodes)
    \item the weigth of the partition: the sum of the weights of the edges
	\item incomparability: some nodes in pairs have to belong to distinct trees
	\item precedence: some nodes in pairs have to belong to the same tree in a given order
	\item conditional precedence: some nodes in pairs have to respect a given order if they belong to the same tree
	\item the in-degree of the nodes
	\item the time windows on nodes (given travelling times on arcs)
\end{itemize}
\end{notedef}

Many applications require to partition a graph such that each component contains exactly one \emph{resource} node and several \emph{task} nodes. A typical example is a routing problem where vehicle routes are paths (a path is a special case of tree) starting from a depot and delivering goods to several clients. Another example is a local network where each computer has to be connected to one shared printer. Last, one can cite the problem of reconstructing plylogeny trees.
The constraint \texttt{tree} can handle these kinds of problems with many additional constraints on the structure of the partition.

\begin{itemize}
	\item \textbf{API} : \mylst{tree(TreeParametersObject param)}
	\item \textbf{return type} : \texttt{Constraint}
	\item \textbf{options} :\emph{n/a}
	\item \textbf{favorite domain} : \emph{to complete}
	\item \textbf{references} :
      \begin{itemize}
      \item \cite{beldiceanuCONSTRAINTS08}: \emph{Combining tree partitioning, precedence, and incomparability constraints}
      \item global constraint catalog: \href{http://www.emn.fr/x-info/sdemasse/gccat/Cproper_forest.html}{proper\_forest} (variant)
      \end{itemize}

\end{itemize}

The tree constraint API requires a particular Model object, named \textbf{\tt TreeParametersObject}.
It can be created with the following parameters:

\begin{tabular}{p{3cm}p{3cm}p{7cm}}
parameter &type &description\\
\hline
$n$ &\texttt{int} &number of nodes in the initial graph $G$\\
$nTree$ &\texttt{IntegerVariable} &number of trees in the resulting forest $G_x$\\
$nProper$ &\texttt{IntegerVariable} &number of proper trees in $G_x$\\
$objective$ &\texttt{IntegerVariable} &(bounded) total \todo{cost} of $G_x$\\
%$objective$ &\texttt{IntegerVariable} &(bounded) total weight of $G_x$\\
$graphs$ &\texttt{List<BitSet[]>} &
\begin{minipage}[t]{7cm}
graphs encoded as successor lists,\\
  \texttt{graphs[0]} the initial graph $G$,\\
  \texttt{graphs[1]} a precedence graph,\\
  \texttt{graphs[2]} a conditional precedence graph,\\
  \texttt{graphs[3]} an incomparability graph
\end{minipage}\\
$matrix$ &\texttt{List<int[][]>} &\texttt{matrix[0]} the indegree of each node, and \texttt{matrix[1]} the starting time from each node\\
$travel$ &\texttt{int[][]} &the travel time of each arc
\end{tabular}

\textbf{Example}:
\lstinputlisting{java/ctree_import.j2t}
\lstinputlisting{java/ctree.j2t}

%\part{true}
\label{true}
\hypertarget{true}{}

\section{TRUE (constraint)}\label{true:trueconstraint}\hypertarget{true:trueconstraint}{}
\(TRUE\) always returns \emph{true}.

%\part{xnor}
\label{xnor}
\hypertarget{xnor}{}

\section{xnor (constraint)}\label{xnor:xnorconstraint}\hypertarget{xnor:xnorconstraint}{}
\begin{notedef}
    \texttt{xnor}$(b_1,b_2)$ states that two booleans are either both true, or both false:
$$ (b_1=1)\quad\iff\quad(b_2=1)$$
\end{notedef}

\begin{itemize}
    \item \textbf{API} : \mylst{xnor(IntegerVariable b1, IntegerVariable b2)}
	\item \textbf{return type} : \texttt{Constraint}
	\item \textbf{options} : \emph{n/a}
	\item \textbf{favorite domain} : \emph{n/a}
\end{itemize}

\textbf{Examples:}
\lstinputlisting{java/cxnor.j2t}

%\part{xor}
\label{xor}
\hypertarget{xor}{}

\section{xor (constraint)}\label{xor:xorconstraint}\hypertarget{xor:xorconstraint}{}
\begin{notedef}
    \texttt{xor}$(b_1,b_2)$ states that the 0-1 variables in arguments take distinct value:
$$ (b_1=1 \land b_2=0) \lor (b_1=0 \land b_2=1)$$
\end{notedef}

\begin{itemize}
    \item \textbf{API} : \mylst{xor(IntegerVariable... b)}
	\item \textbf{return type} : \texttt{Constraint}
	\item \textbf{options} : \emph{n/a}
	\item \textbf{favorite domain} : \emph{n/a}
\end{itemize}

\textbf{Examples:}
\lstinputlisting{java/cxor.j2t}


\chapter{Options (Model)}\label{ch:options}\hypertarget{ch:options}{}
This section lists and details the options that can be declared on variables or constraints within a Choco Model.
\input{chapters/Options.tex}

\chapter{Branching strategies (Solver)}\label{ch:branchstrat}\hypertarget{ch:branchstrat}{}
This section lists and details the \hyperlink{solver:searchstrategy}{branching strategies} currently available in Choco.

\section{AssignInterval (Branching strategy)}\label{assigninterval:assignintervalbranchstrat}\hypertarget{assigninterval:assignintervalbranchstrat}{}
\begin{notedef}
  \texttt{AssignInterval} is a \textbf{binary branching} assigning interval to subinterval. It splits the interval of the selected real variable in two parts. The first decision is to reduced the variable's domain to the first half, the second and last decision is to reduced the variable's domain to its second half. 
\end{notedef}

\begin{itemize}
	\item \textbf{Constructor} :\mylst{AssignInterval(VarSelector<RealVar> varSelector, ValIterator<RealVar> valIterator)}
	\item \textbf{type of variable} : real
	\item \textbf{references} : \emph{n/a}
\end{itemize}

\textbf{Example}:
%\lstinputlisting{java/cabs.j2t}


\section{AssignOrForbidIntVarVal (Branching strategy)}\label{assignorforbidintvarval:assignorforbidintvarvalbranchstrat}\hypertarget{assignorforbidintvarval:assignorforbidintvarvalbranchstrat}{}
\begin{notedef}
  \texttt{AssignOrForbidIntVarVal} is a \textbf{binary branching} assigning value to variable in the left branch and forbidding the assignment in the right branch.  
\end{notedef}

\begin{itemize}
	\item \textbf{Constructor} :\mylst{AssignOrForbidIntVarVal(VarSelector<IntDomainVar> varsel, ValSelector<IntDomainVar> valsel)}
	\item \textbf{type of variable} : integer
	\item \textbf{references} : \emph{n/a}
\end{itemize}

\textbf{Example}:
%\lstinputlisting{java/cabs.j2t}


\section{AssignOrForbidIntVarValPair (Branching strategy)}\label{assignorforbidintvarvalpair:assignorforbidintvarvalpairbranchstrat}\hypertarget{assignorforbidintvarvalpair:assignorforbidintvarvalpairbranchstrat}{}
\begin{notedef}
\texttt{AssignOrForbidIntVarValPair} is a  \textbf{binary branching} assigning a value to an integer variable. In the first branch, the selected value is assigned to the selected variable; in the second branch, the value is removed from the variable domain.
It requires a \texttt{VarValPairSelector} selecting both variable and value at the same time.
$$B_1: x=v,\quad B_2: x\neq v$$
\end{notedef}


\begin{itemize}
	\item \textbf{Constructor} :\mylst{AssignOrForbidIntVarValPair(VarValPairSelector varValSel)}
	\item \textbf{type of variable} : integer
	\item \textbf{references} : \emph{n/a}
	\item \textbf{see also} : \hyperlink{assignorforbidintvarval:assignorforbidintvarvalbranchstrat}{\texttt{AssignOrForbidIntVarVal}} based on distinct variable/value selectors
\end{itemize}

\textbf{Example}:
%\lstinputlisting{java/cabs.j2t}


\section{AssignSetVar (Branching strategy)}\label{assignsetvar:assignsetvarbranchstrat}\hypertarget{assignsetvar:assignsetvarbranchstrat}{}
\begin{notedef}
  \texttt{AssignSetVar} is a \textbf{n-ary branching} assigning a value to a variable.
\end{notedef}

\begin{itemize}
	\item \textbf{Constructor} :\mylst{VarSelector<SetVar> varselect, ValSelector<SetVar> valselect)}
	\item \textbf{type of variable} : set
	\item \textbf{references} : \emph{n/a}
\end{itemize}

\textbf{Example}:
%\lstinputlisting{java/cabs.j2t}


\section{AssignVar (Branching strategy)}\label{assignvar:assignvarbranchstrat}\hypertarget{assignvar:assignvarbranchstrat}{}
\begin{notedef}
  \texttt{AssignVar} is a \textbf{n-ary branching} assigning a value to a variable. \todo{to complete}
\end{notedef}

\begin{itemize}
	\item \textbf{Constructor} :\mylst{VarSelector<SetVar> varselect, ValSelector<SetVar> valselect)}
	\item \textbf{type of variable} : integer
	\item \textbf{references} : \emph{n/a}
\end{itemize}

\textbf{Example}:
%\lstinputlisting{java/cabs.j2t}


\section{DomOverWDegBranchingNew (Branching strategy)}\label{domoverwdeg:domoverwdegbranchstrat}\hypertarget{domoverwdeg:domoverwdegbranchstrat}{}
\begin{notedef}
  \texttt{DomOverWDegBranchingNew} is a \textbf{n-ary branching} assigning distinct values to an integer variable. It maintains (incrementally or dynamically) on each constraint, the count of the failures caused by the constraint from the beginning of the search.
To each variable are then associated, at any time, three values: $dom$ the current domain size, $deg$ the current number of uninstantiated constraints involving the variable, and $w$ the sum of the counters associated to these $deg$ constraints.
The strategy selects the variable with the smallest ratio $r_i=dom/w*deg$.  Ties are randomly broken when \mylst{seed != null}.
The variable is then successively assigned, in each branch, to a next selected value.
% er a counter on each constraint. Whenever a dead-end occurs, the counter associated to the contraint that just failed is updated during search whenever a dead-end occurs. Then, the weight, $w$, of a variable $X_i$ is the sum of the weights of the constraints involving $X_i$ and at least one uninstantiated variable. And combining $w$ with the domain size, $dom$, and the degree, $deg$, of $X_i$, this strategy selects the variable with the smallest ratio $r_i=dom/w*deg$.
$$B_1: x=v_1,\quad B_2: x= v_2,\quad\ldots,\quad B_m: x= v_m,\qquad\text{ with } x=\arg\min\bigl(\frac{dom_x}{w_x*deg_x}\bigr)$$
\end{notedef}

\begin{itemize}
	\item \textbf{Constructor} :\mylst{DomOverWDegBranchingNew(Solver s, IntDomainVar[] vars, ValIterator valIt, Number seed)}
	\item \textbf{type of variable} : integer
	\item \textbf{references} : \cite{Boussemart04}: \emph{Boosting systematic search by weighting constraints}.
\end{itemize}

\textbf{Example}:
%\lstinputlisting{java/cabs.j2t}


\section{DomOverWDegBinBranchingNew (Branching strategy)}\label{domoverwdegbin:domoverwdegbinbranchstrat}\hypertarget{domoverwdegbin:domoverwdegbinbranchstrat}{}
\begin{notedef}
  \texttt{DomOverWDegBranchingNew} is a \textbf{binary branching} assigning distinct values to an integer variable. It maintains (incrementally or dynamically) on each constraint, the count of the failures caused by the constraint from the beginning of the search.
To each variable are then associated, at any time, three values: $dom$ the current domain size, $deg$ the current number of uninstantiated constraints involving the variable, and $w$ the sum of the counters associated to these $deg$ constraints.
The strategy selects the variable with the smallest ratio $r_i=dom/w*deg$.  Ties are randomly broken when \mylst{seed != null}.
The variable is then assigned, in the first branch, to the selected value; in the second branch, the value is removed from the variable domain.
$$B_1: x=v,\quad B_2: x\neq v,\qquad\text{ with } x=\arg\min\bigl(\frac{dom_x}{w_x*deg_x}\bigr)$$\end{notedef}

\begin{itemize}
	\item \textbf{Constructor}: \mylst{DomOverWDegBinBranchingNew(Solver s, IntDomainVar[] vars, ValSelector valSel, Number seed)}
	\item \textbf{type of variable} : integer
	\item \textbf{references} : \cite{Boussemart04}: \emph{Boosting systematic search by weighting constraints}.
\end{itemize}

\textbf{Example}:
%\lstinputlisting{java/cabs.j2t}


\section{ImpactBasedBranching (Branching strategy)}\label{impact:impactbranchstrat}\hypertarget{impact:impactbranchstrat}{}
\begin{notedef}
  \texttt{ImpactBasedBranching} is a \textbf{n-ary branching} assigning all distinct values to an integer variable. 
  The impact of a branching decision measures the reduction of the search space induced when the decision was posted and propagated since the beginning of the search.
  Here, branching decisions are variable-value assignments. 
The strategy selects an integer variable maximizing the total impact minus the domain size. The variable is then successively assigned, in each branch, to its possible values selected by decreasing order of their impacts. 
  \emph{Restarting search can dramatically improve performance.}
$$B_1: x=v_1,\quad B_2: x= v_2,\quad\ldots,\quad B_m: x= v_m,$$
with $x=\arg\max\sum_{v\in D(x)}(impact(x,v)-1)$  and $impact(x,v_i) > impact(x,v_j)$, $\forall j<i$.
\end{notedef}

\begin{itemize}
	\item \textbf{Constructor} :
	\begin{itemize}
	\item \mylst{ImpactBasedBranching(Solver s, IntDomainVar[] vars)}
	\item \mylst{ImpactBasedBranching(Solver s, IntDomainVar[] vars, AbstractImpactStrategy ibs)}
	\end{itemize}	
	\item \textbf{type of variable} : integer
	\item \textbf{references} : \cite{Refalo04}: \emph{Impact-Based Search Strategies for Constraint Programming}. 
\end{itemize}

\textbf{Example}:
%\lstinputlisting{java/cabs.j2t}


\section{PackDynRemovals (Branching strategy)}\label{packdynremovals:packdynremovalsbranchstrat}\hypertarget{packdynremovals:packdynremovalsbranchstrat}{}
\begin{notedef}
  \texttt{PackDynRemovals} is an extension of \hyperlink{assignvar:assignvarbranchstrat}{AssignVar} specialized for packing problem. At every backtrack, we state that the item can not be packed in equivalent bins. Two bins are equivalent if they have the same remaining space.
\end{notedef}

\begin{itemize}
	\item \textbf{Constructor} :\mylst{PackDynRemovals(VarSelector varSel, ValSelector valHeuri, PackSConstraint pack)}
	\item \textbf{type of variable} : integer
	\item \textbf{references} : \emph{n/a}
\end{itemize}

\textbf{Example}:
%\lstinputlisting{java/cabs.j2t}


\section{SetTimes (Branching strategy)}\label{settimes:settimesbranchstrat}\hypertarget{settimes:settimesbranchstrat}{}
\begin{notedef}
  \texttt{SetTimes} is a \textbf{n-ary branching} for scheduling problems with makespan minimization objective, fixing distinct task variables to their earliest starting time.
At each node, a set of available tasks is considered then, in each branch, one task variable is selected and its starting time is set to its smallest value. 
The task variables are selected in order according to a \mylst{TaskVarSelector}, or in the decreasing order defined by a \mylst{Comparator}, ties being randomly broken when boolean \mylst{rand} is set. 
%The selected variable is either chosen by a \mylst{TaskVarSelector} or as the maximal element according to a \mylst{Comparator}. Ties in the variable selection are randomly broken when boolean \mylst{rand} is set. 
\emph{The search is not complete: do not use within a} \texttt{solveAll}. 
$$B_1: T_1.start=\underline{T_1.start},\quad B_2: T_2.start=\underline{T_2.start},\quad \ldots,\quad B_n: T_n.start=\underline{T_n.start}$$
  \end{notedef}

\begin{itemize}
	\item \textbf{Constructor} :
      \begin{itemize}
      \item \mylst{SetTimes(Solver solver, List<TaskVar> tasks, TaskVarSelector varSel)}
      \item \mylst{SetTimes(Solver solver, List<TaskVar> tasks, Comparator<ITask> comp, boolean rand)}
      \end{itemize}
	\item \textbf{type of variable} : task
	\item \textbf{references} : \emph{n/a}
\end{itemize}

\textbf{Example}:
%\lstinputlisting{java/cabs.j2t}


\section{TaskOverWDegBinBranching (Branching strategy)}\label{taskdomoverwdeg:taskdomoverwdegbranchstrat}\hypertarget{taskdomoverwdeg:taskdomoverwdegbranchstrat}{}
\begin{notedef}
  \texttt{TaskOverWDegBinBranching} is a \textbf{n-ary branching} assigning a value to a variable. \todo{to complete}.
\end{notedef}

\begin{itemize}
	\item \textbf{Constructor} :\mylst{TaskOverWDegBinBranching(Solver s, IntDomainVar[] vars, ValIterator valHeuri, Number seed)}
	\item \textbf{type of variable} : integer
	\item \textbf{references} : \cite{Boussemart04}: \emph{Boosting systematic search by weighting constraints}.
\end{itemize}

\textbf{Example}:
%\lstinputlisting{java/cabs.j2t}



\chapter{Variable selectors (Solver)}\label{ch:varsel}\hypertarget{ch:varsel}{}
This section lists and details the \hyperlink{solver:variableselector}{variable selectors} currently available in Choco.

\section{CompositeIntVarSelector (Variable selector)}\label{compositeintvarselector:compositeintvarselectorvarselector}\hypertarget{compositeintvarselector:compositeintvarselectorvarselector}{}
\begin{notedef}
  \texttt{CompositeIntVarSelector} composes two heuristics for selecting a variable: a first heuristic is applied for selecting a constraint, then from that constraint a second heuristic is applied for selecting the variable
  \end{notedef}

\begin{itemize}
	\item \textbf{Constructor} :\mylst{CompositeIntVarSelector(ConstraintSelector cs, HeuristicIntVarSelector cvs)}
	\begin{itemize}
\item \texttt{ConstraintSelector} is an \texttt{interface} not implemented. You should define your own implementation. 

\item \texttt{HeuristicIntVarSelector} is an \texttt{abstract class}  implemented by: \hyperlink{mindomain:mindomainvarselector}{MinDomain}, \hyperlink{maxdomain:maxdomainvarselector}{MaxDomain}, \hyperlink{maxregret:maxregretvarselector}{MaxRegret}, \hyperlink{maxvaldomain:maxvaldomainvarselector}{MaxValueDomain}, \hyperlink{minvaldomain:minvaldomainvarselector}{MinValueDomain}, \hyperlink{mostconstrained:mostconstrainedvarselector}{MostConstrained}. //TODO to complete

\end{itemize}

	\item \textbf{type of variable} : integer
	\item \textbf{references} : \emph{n/a}
\end{itemize}

\textbf{Example}:
%\lstinputlisting{java/cabs.j2t}


\section{CyclicRealVarSelector (Variable selector)}\label{cyclicrealvarselector:cyclicrealvarselectorvarselector}\hypertarget{cyclicrealvarselector:cyclicrealvarselectorvarselector}{}
\begin{notedef}
  \texttt{CyclicRealVarSelector}$(xs)$: since a dichotomy algorithm is used, cyclic assigning is needed to instantiate a real interval variable. A variable is selected several times to split its domain until it reaches the desired precision.
\end{notedef}

\begin{itemize}
	\item \textbf{Constructor} : 
	\begin{itemize}
	\item \mylst{CyclicRealVarSelector(Solver solver)}
	\item \mylst{CyclicRealVarSelector(Solver solver, RealVar[] vs)}
	\end{itemize}	
	\item \textbf{type of variable} : real
	\item \textbf{references} : \emph{n/a}
\end{itemize}

\textbf{Example}:
%\lstinputlisting{java/cabs.j2t}


\section{LexIntVarSelector (Variable selector)}\label{lexintvarselector:lexintvarselectorvarselector}\hypertarget{lexintvarselector:lexintvarselectorvarselector}{}
\begin{notedef}
  \texttt{LexIntVarSelector}$(h_1,h_2)$ selects the integer variable according to the first heuristic $h_1$, ties being broken by the second heuristic $h_2$:
%applies two heuristics lexicographically for selecting a variable: a first heuristic is applied finding the best variables, ties are broken with the second heuristic 
$$\begin{cases}
  h_1(x) &\quad\text{ if } |h_1(x)|=1,\\
  h_2(h_1(x)) &\quad\text{ otherwise}.
\end{cases}$$
\end{notedef}

\begin{itemize}
	\item \textbf{Constructor}: \mylst{LexIntVarSelector(TiedIntVarSelector h1, HeuristicIntVarSelector h2)}
	\begin{itemize}
\item \texttt{TiedIntVarSelector} is an \texttt{interface} implemented by : \hyperlink{mindomain:mindomainvarselector}{MinDomain}, \hyperlink{maxdomain:maxdomainvarselector}{MaxDomain}, \hyperlink{maxregret:maxregretvarselector}{MaxRegret}, \hyperlink{maxvaldomain:maxvaldomainvarselector}{MaxValueDomain}, \hyperlink{minvaldomain:minvaldomainvarselector}{MinValueDomain}, \hyperlink{mostconstrained:mostconstrainedvarselector}{MostConstrained}, \hyperlink{randomvarint:randomvarintvarselector}{RandomIntVarSelector}.
\item \texttt{HeuristicIntVarSelector} is an \texttt{abstract class}  implemented by: \hyperlink{mindomain:mindomainvarselector}{MinDomain}, \hyperlink{maxdomain:maxdomainvarselector}{MaxDomain}, \hyperlink{maxregret:maxregretvarselector}{MaxRegret}, \hyperlink{maxvaldomain:maxvaldomainvarselector}{MaxValueDomain}, \hyperlink{minvaldomain:minvaldomainvarselector}{MinValueDomain}, \hyperlink{mostconstrained:mostconstrainedvarselector}{MostConstrained}, \hyperlink{randomvarint:randomvarintvarselector}{RandomIntVarSelector}.
\end{itemize}
	\item \textbf{type of variable} : integer
	\item \textbf{references} : \emph{n/a}
\end{itemize}

\textbf{Example}:
%\lstinputlisting{java/cabs.j2t}


\section{MaxDomain (Variable selector)}\label{maxdomain:maxdomainvarselector}\hypertarget{maxdomain:maxdomainvarselector}{}
\begin{notedef}
  \texttt{MaxDomain}$(x)$ selects the integer variable with the largest domain (\textbf{dynamic}):
$$\max |D(x)|$$
\end{notedef}

\begin{itemize}
	\item \textbf{Constructor} : 
	\begin{itemize}
	\item \mylst{MaxDomain(Solver s)}
	\item \mylst{MaxDomain(Solver s, IntDomainVar[] vars)}
	\end{itemize}	
	\item \textbf{type of variable} : integer
	\item \textbf{references} : \emph{n/a}
\end{itemize}

\textbf{Example}:
%\lstinputlisting{java/cabs.j2t}


\section{MaxDomSet (Variable selector)}\label{maxdomset:maxdomsetvarselector}\hypertarget{maxdomset:maxdomsetvarselector}{}
\begin{notedef}
  \texttt{MaxDomSet}$(xs)$ chooses the variable with the largest domain :
\end{notedef}

\begin{itemize}
	\item \textbf{Constructor} : 
	\begin{itemize}
	\item \mylst{MaxDomSet(Solver solver)}
	\item \mylst{MaxDomSet(Solver solver, SetVar[] vs)}
	\end{itemize}	
	\item \textbf{type of variable} : set
	\item \textbf{references} : \emph{n/a}
\end{itemize}

\textbf{Example}:
%\lstinputlisting{java/cabs.j2t}


\section{MaxRegret (Variable selector)}\label{maxregret:maxregretvarselector}\hypertarget{maxregret:maxregretvarselector}{}
\begin{notedef}
  \texttt{MaxRegret}$(x)$ selects the integer variable with the largest difference between the two smallest values in its domain (\textbf{dynamic}):
$$\max(\underline{x}_2 - \underline{x}),\qquad\text{with } \underline{x}_2=\min(D(x)\setminus\underline{x})$$
\end{notedef}

\begin{itemize}
	\item \textbf{Constructor} : 
	\begin{itemize}
	\item \mylst{MaxRegret(Solver s)}
	\item \mylst{MaxRegret(Solver s, IntDomainVar[] vars)}
	\end{itemize}	
	\item \textbf{type of variable} : integer
	\item \textbf{references} : \emph{n/a}
\end{itemize}

\textbf{Example}:
%\lstinputlisting{java/cabs.j2t}


\section{MaxRegretSet (Variable selector)}\label{maxregretset:maxregretsetvarselector}\hypertarget{maxregretset:maxregretsetvarselector}{}
\begin{notedef}
  \texttt{MaxRegretSet}$(xs)$ chooses the variable with the largest difference between the two smallest values in its domain :
\end{notedef}

\begin{itemize}
	\item \textbf{Constructor} : 
	\begin{itemize}
	\item \mylst{MaxRegretSet(Solver solver)}
	\item \mylst{MaxRegretSet(Solver solver, SetVar[] vs)}
	\end{itemize}	
	\item \textbf{type of variable} : set
	\item \textbf{references} : \emph{n/a}
\end{itemize}

\textbf{Example}:
%\lstinputlisting{java/cabs.j2t}


\section{MaxValueDomain (Variable selector)}\label{maxvaldomain:maxvaldomainvarselector}\hypertarget{maxvaldomain:maxvaldomainvarselector}{}
\begin{notedef}
  \texttt{MaxValueDomain}$(x)$ selects the integer variable with the largest value in its domain (\textbf{dynamic}):
$$\max(\bar{x})$$
\end{notedef}

\begin{itemize}
	\item \textbf{Constructor} : 
	\begin{itemize}
	\item \mylst{MaxValueDomain(Solver s)}
	\item \mylst{MaxValueDomain(Solver s, IntDomainVar[] vars)}
	\end{itemize}	
	\item \textbf{type of variable} : integer
	\item \textbf{references} : \emph{n/a}
\end{itemize}

\textbf{Example}:
%\lstinputlisting{java/cabs.j2t}


\section{MaxValueDomSet (Variable selector)}\label{maxvaldomset:maxvaldomsetvarselector}\hypertarget{maxvaldomset:maxvaldomsetvarselector}{}
\begin{notedef}
  \texttt{MaxValueDomSet}$(x)$ selects the set variable with the largest value in its envelope (\textbf{dynamic}):
$$\max(\bar{x}),\qquad\text{with } \bar{x}=\max(Env(x))$$
\end{notedef}

\begin{itemize}
	\item \textbf{Constructor} : 
	\begin{itemize}
	\item \mylst{MaxValueDomSet(Solver s)}
	\item \mylst{MaxValueDomSet(Solver s, SetVar[] vars)}
	\end{itemize}	
	\item \textbf{type of variable} : set
	\item \textbf{references} : \emph{n/a}
\end{itemize}

\textbf{Example}:
%\lstinputlisting{java/cabs.j2t}


\section{MinDomain (Variable selector)}\label{mindomain:mindomainvarselector}\hypertarget{mindomain:mindomainvarselector}{}
\begin{notedef}
  \texttt{MinDomain}$(x)$ selects the integer variable with the smallest domain (\textbf{dynamic}):
$$\min |D(x)|$$
\end{notedef}

\begin{itemize}
	\item \textbf{Constructor} : 
	\begin{itemize}
	\item \mylst{MinDomain(Solver s)}
	\item \mylst{MinDomain(Solver s, IntDomainVar[] vars)}
	\end{itemize}	
	\item \textbf{type of variable} : integer
	\item \textbf{references} : \emph{n/a}
\end{itemize}

\textbf{Example}:
%\lstinputlisting{java/cabs.j2t}


\section{MinDomSet (Variable selector)}\label{mindomset:mindomsetvarselector}\hypertarget{mindomset:mindomsetvarselector}{}
\begin{notedef}
  \texttt{MinDomSet}$(x)$ selects the set variable with the smallest open domain (\textbf{dynamic}):
$$\min |Env(x)\setminus Ker(x)|$$
\end{notedef}

\begin{itemize}
	\item \textbf{Constructor} : 
	\begin{itemize}
	\item \mylst{MinDomSet(Solver s)}
	\item \mylst{MinDomSet(Solver s, SetVar[] vars)}
	\end{itemize}	
	\item \textbf{type of variable} : set
	\item \textbf{references} : \emph{n/a}
\end{itemize}

\textbf{Example}:
%\lstinputlisting{java/cabs.j2t}


\section{MinValueDomain (Variable selector)}\label{minvaldomain:minvaldomainvarselector}\hypertarget{minvaldomain:minvaldomainvarselector}{}
\begin{notedef}
  \texttt{MinValueDomain}$(x)$ selects the integer variable with the smallest value in its domain (\textbf{dynamic}):
$$\min(\underline{x})$$
\end{notedef}

\begin{itemize}
	\item \textbf{Constructor} : 
	\begin{itemize}
	\item \mylst{MinValueDomain(Solver s)}
	\item \mylst{MinValueDomain(Solver s, IntDomainVar[] vars)}
	\end{itemize}	
	\item \textbf{type of variable} : integer
	\item \textbf{references} : \emph{n/a}
\end{itemize}

\textbf{Example}:
%\lstinputlisting{java/cabs.j2t}


\section{MinValueDomSet (Variable selector)}\label{minvaldomset:minvaldomsetvarselector}\hypertarget{minvaldomset:minvaldomsetvarselector}{}
\begin{notedef}
  \texttt{MinValueDomSet}$(x)$ selects the set variable with the smallest value in its envelope (\textbf{dynamic}):
$$\min(\underline{x}),\qquad\text{with } \underline{x}=\min(Env(x))$$
\end{notedef}

\begin{itemize}
	\item \textbf{Constructor} : 
	\begin{itemize}
	\item \mylst{MinValueDomSet(Solver s)}
	\item \mylst{MinValueDomSet(Solver s, SetVar[] vars)}
	\end{itemize}	
	\item \textbf{type of variable} : set
	\item \textbf{references} : \emph{n/a}
\end{itemize}

\textbf{Example}:
%\lstinputlisting{java/cabs.j2t}


\section{MostConstrained (Variable selector)}\label{mostconstrained:mostconstrainedvarselector}\hypertarget{mostconstrained:mostconstrainedvarselector}{}
\begin{notedef}
  \texttt{MostConstrained}$(x)$ selects the integer variable involved in the largest number of constraints initially present in the solver (\textbf{static}):
$$\max(initDeg(x))$$
\end{notedef}

\begin{itemize}
	\item \textbf{Constructor} : 
	\begin{itemize}
	\item \mylst{MostConstrained(Solver s)}
	\item \mylst{MostConstrained(Solver s, IntDomainVar[] vars)}
	\end{itemize}	
	\item \textbf{type of variable} : integer
	\item \textbf{references} : \emph{n/a}
\end{itemize}

\textbf{Example}:
%\lstinputlisting{java/cabs.j2t}


\section{MostConstrainedSet (Variable selector)}\label{mostconstrainedset:mostconstrainedsetvarselector}\hypertarget{mostconstrainedset:mostconstrainedsetvarselector}{}
\begin{notedef}
  \texttt{MostConstrainedSet}$(xs)$ chooses the variable with the largest number of constraints:
$$\{x | \}$$
\end{notedef}

\begin{itemize}
	\item \textbf{Constructor} : 
	\begin{itemize}
	\item \mylst{MostConstrainedSet(Solver solver)}
	\item \mylst{MostConstrainedSet(Solver solver, SetVar[] vs)}
	\end{itemize}	
	\item \textbf{type of variable} : set
	\item \textbf{references} : \emph{n/a}
\end{itemize}

\textbf{Example}:
%\lstinputlisting{java/cabs.j2t}


\section{RandomIntVarSelector (Variable selector)}\label{randomvarint:randomvarintvarselector}\hypertarget{randomvarint:randomvarintvarselector}{}
\begin{notedef}
  \texttt{RandomIntVarSelector}$(xs)$ chooses a random variable in $$xs$$:
$$\{x | \}$$
\end{notedef}

\begin{itemize}
	\item \textbf{Constructor} : 
	\begin{itemize}
	\item \mylst{RandomIntVarSelector(Solver solver)}
	\item \mylst{RandomIntVarSelector(Solver solver, long seed)}
	\item \mylst{RandomIntVarSelector(Solver solver, IntDomainVar[] vs, long seed)}
	\end{itemize}	
	\item \textbf{type of variable} : integer
	\item \textbf{references} : \emph{n/a}
\end{itemize}

\textbf{Example}:
%\lstinputlisting{java/cabs.j2t}


\section{RandomSetVarSelector (Variable selector)}\label{randomvarset:randomvarsetvarselector}\hypertarget{randomvarset:randomvarsetvarselector}{}
\begin{notedef}
  \texttt{RandomSetVarSelector}$(x)$ selects a set variable randomly (\textbf{dynamic}). The random \texttt{seed} can be fixed.
\end{notedef}

\begin{itemize}
	\item \textbf{Constructor} : 
	\begin{itemize}
	\item \mylst{RandomSetVarSelector(Solver s)}
	\item \mylst{RandomSetVarSelector(Solver s, long seed)}
	\item \mylst{RandomSetVarSelector(Solver s, SetVar[] vars, long seed)}
	\end{itemize}	
	\item \textbf{type of variable} : set
	\item \textbf{references} : \emph{n/a}
\end{itemize}

\textbf{Example}:
%\lstinputlisting{java/cabs.j2t}


\section{StaticSetVarOrder (Variable selector)}\label{staticsetvarorder:staticsetvarordervarselector}\hypertarget{staticsetvarorder:staticsetvarordervarselector}{}
\begin{notedef}
  \texttt{StaticSetVarOrder}$(xs)$ chooses the variable in the order they appears in $$xs$$:
$$\{x | \}$$
\end{notedef}

\begin{itemize}
	\item \textbf{Constructor} : 
	\begin{itemize}
	\item \mylst{StaticSetVarOrder(Solver solver)}
	\item \mylst{StaticSetVarOrder(Solver solver, SetVar[] vs)}
	\end{itemize}	
	\item \textbf{type of variable} : set
	\item \textbf{references} : \emph{n/a}
\end{itemize}

\textbf{Example}:
%\lstinputlisting{java/cabs.j2t}


\section{StaticVarOrder (Variable selector)}\label{staticvarorder:staticvarordervarselector}\hypertarget{staticvarorder:staticvarordervarselector}{}
\begin{notedef}
  \texttt{StaticVarOrder}$(x)$ selects the integer variables in the order they appear in array $x$ (\textbf{static}).
$$x_1$$
\end{notedef}

\begin{itemize}
	\item \textbf{Constructor} : 
	\begin{itemize}
	\item \mylst{StaticVarOrder(Solver s)}
	\item \mylst{StaticVarOrder(Solver s, IntDomainVar[] vars)}
	\end{itemize}	
	\item \textbf{type of variable} : integer
	\item \textbf{references} : \emph{n/a}
\end{itemize}

\textbf{Example}:
%\lstinputlisting{java/cabs.j2t}



\chapter{Value iterators (Solver)}\label{ch:valite}\hypertarget{ch:valite}{}
This section lists and details the \hyperlink{solver:valueiterator}{value iterators} currently available in Choco.

\section{DecreasingDomain (Value iterator)}\label{decreasingdomain:decreasingdomainvaliterator}\hypertarget{decreasingdomain:decreasingdomainvaliterator}{}
\begin{notedef}
  \texttt{DecreasingDomain} selects the integer variable largest value:
$$\max(D(x))$$
\end{notedef}

\begin{itemize}
	\item \textbf{Constructor} : \mylst{DecreasingDomain()}
	\item \textbf{type of variable} : integer
%	\item \textbf{references} : \emph{n/a}
\end{itemize}

%\textbf{Example}:
%\lstinputlisting{java/cabs.j2t}


\section{IncreasingDomain (Value iterator)}\label{increasingdomain:increasingdomainvaliterator}\hypertarget{increasingdomain:increasingdomainvaliterator}{}
\begin{notedef}
  \texttt{IncreasingDomain} selects the integer variable smallest value:
$$\min(D(x))$$
\end{notedef}

\begin{itemize}
	\item \textbf{Constructor} : \mylst{IncreasingDomain()}
	\item \textbf{type of variable} : integer
%	\item \textbf{references} : \emph{n/a}
\end{itemize}

%\textbf{Example}:
%\lstinputlisting{java/cabs.j2t}


\section{RealIncreasingDomain (Value iterator)}\label{realincreasingdomain:realincreasingdomainvaliterator}\hypertarget{realincreasingdomain:realincreasingdomainvaliterator}{}
\begin{notedef}
  \texttt{RealIncreasingDomain} selects the real variable smallest value:
$$\min(D(x))$$
\end{notedef}

\begin{itemize}
	\item \textbf{Constructor} : \mylst{RealIncreasingDomain()}
	\item \textbf{type of variable} : real
%	\item \textbf{references} : \emph{n/a}
\end{itemize}

%\textbf{Example}:
%\lstinputlisting{java/cabs.j2t}



\chapter{Value selector (Solver)}\label{ch:valsel}\hypertarget{ch:valsel}{}
This section lists and details the \hyperlink{solver:valueselector}{value selectors} currently available in Choco.

\section{BestFit (Value selector)}\label{bestfit:bestfitvalselector}\hypertarget{bestfit:bestfitvalselector}{}
\begin{notedef}
  \texttt{BestFit}, associated with a \hyperlink{pack:packconstraint}{pack}(items, load, bin, size) constraint, selects the bin $v$ with the minimum residual space $rcap_v$ for the bin assignment integer variable $\mathtt{bin}[i]$ of item $i$:
$$v\in D(\mathtt{bin}[i]):\quad\min\{rcap_v\}\qquad\text{with } rcap_v=\overline{\mathtt{load}[v]} -\!\sum_{j\in Ker(\mathtt{items}[v])}\!\mathtt{size}[j]$$
\end{notedef}

\begin{itemize}
	\item \textbf{Constructor} : \mylst{BestFit(PackSConstraint cstr)}
	\item \textbf{type of variable} : integer (\emph{bin-to-item assignment variable} $\mathtt{bin}[i]$)
%	\item \textbf{references} : \emph{n/a}
\end{itemize}

%\textbf{Example}:
%\lstinputlisting{java/cabs.j2t}


\section{CostRegularValSelector (Value selector)}\label{costregularvalselector:costregularvalselectorvalselector}\hypertarget{costregularvalselector:costregularvalselectorvalselector}{}
\begin{notedef}
  \texttt{CostRegularValSelector}, \todo{to complete}
\end{notedef}

\begin{itemize}
	\item \textbf{Constructor} : \mylst{CostRegularValSelector(CostRegular cr,boolean max)}
	\item \textbf{type of variable} : integer
	\item \textbf{references} : \emph{n/a}
\end{itemize}

\textbf{Example}:
%\lstinputlisting{java/cabs.j2t}


\section{FCostRegularValSelector (Value selector)}\label{fcostregularvalselector:fcostregularvalselectorvalselector}\hypertarget{fcostregularvalselector:fcostregularvalselectorvalselector}{}
\begin{notedef}
  \texttt{FCostRegularValSelector}, associated with a \hyperlink{costregular:costregularconstraint}{\texttt{CostRegular}}$(z, \collec{x_1}{x_n},\mathcal{L}(\Pi), \coll{c_{i,j}})$ constraint, selects a value $v_i$ in the domain of the sequence integer variable $x_i$ that can be extended to a solution for the constraint $\collec{v_1}{v_n}\in\mathcal{L}(\Pi)$ of maximum cost (if \texttt{max}, resp. minimum cost if $\neg$\texttt{max}):
$$ v_i\in D(x_i):\quad \exists v_j\in D(x_j)\ \forall j\neq i, \text{ s.t. } \sum_{j=1}^n c_{j,v_j} = \bar{z}\ (\text{resp. } =\underline{z})$$   
\end{notedef}
Remember that the underlying data structure of a \texttt{CostRegular} constraint is a layered acyclic digraph, where each arc in the $i$-th layer corresponds to a possible assignment of variable $x_i$ to a value $v\in D(x_i)$ and has length $c_{iv}$, and where each path corresponds to a solution of the constraint, i.e. a word $\collec{v_1}{v_n}$ of language $\mathcal{L}(\Pi)$, whose cost $z$ is equal to the path length $\sum_{j=1}^n c_{j,v_j}$.

This heuristic works with the \mylst{ConstraintType.FASTCOSTREGULAR} implementation of the constraint (preferred): see \hyperlink{costregularvalselector:costregularvalselectorvalselector}{CostRegularValSelector} for use with the \mylst{ConstraintType.COSTREGULAR} implementation.

\begin{itemize}
	\item \textbf{Constructor} : \mylst{FCostRegularValSelector(FastCostRegular cr, boolean max)}
	\item \textbf{type of variable} : integer
%	\item \textbf{references} : \emph{n/a}
\end{itemize}

%\textbf{Example}:
%\lstinputlisting{java/cabs.j2t}


\section{MaxVal (Value selector)}\label{maxval:maxvalvalselector}\hypertarget{maxval:maxvalvalselector}{}
\begin{notedef}
  \texttt{MaxVal} chooses the variable's domain largest value.
\end{notedef}

\begin{itemize}
	\item \textbf{Constructor} : \mylst{MaxVal()}
	\item \textbf{type of variable} : integer
	\item \textbf{references} : \emph{n/a}
\end{itemize}

\textbf{Example}:
%\lstinputlisting{java/cabs.j2t}


%\section{MCRValSelector (Value selector)}\label{mcrvalselector:mcrvalselectorvalselector}\hypertarget{mcrvalselector:mcrvalselectorvalselector}{}
\begin{notedef}
  \texttt{MCRValSelector}, \todo{to complete}
\end{notedef}

\begin{itemize}
	\item \textbf{Constructor} : \mylst{MCRValSelector(FastMultiCostRegular[] cons, boolean max)}
	\item \textbf{type of variable} : integer
	\item \textbf{references} : \emph{n/a}
\end{itemize}

\textbf{Example}:
%\lstinputlisting{java/cabs.j2t}


\section{MidVal (Value selector)}\label{midval:midvalvalselector}\hypertarget{midval:midvalvalselector}{}
\begin{notedef}
  \texttt{MidVal} chooses the value closest (equals or right after) to the variable's domain mean of its current bounds.
\end{notedef}

\begin{itemize}
	\item \textbf{Constructor} : \mylst{MidVal()}
	\item \textbf{type of variable} : integer
	\item \textbf{references} : \emph{n/a}
\end{itemize}

\textbf{Example}:
%\lstinputlisting{java/cabs.j2t}


\section{MinEnv (Value selector)}\label{minenv:minenvvalselector}\hypertarget{minenv:minenvvalselector}{}
\begin{notedef}
  \texttt{MinEnv} chooses the variable's domain smallest value.
\end{notedef}

\begin{itemize}
	\item \textbf{Constructor} : \mylst{MinEnv()}
	\item \textbf{type of variable} : set
	\item \textbf{references} : \emph{n/a}
\end{itemize}

\textbf{Example}:
%\lstinputlisting{java/cabs.j2t}


\section{MinVal (Value selector)}\label{minval:minvalvalselector}\hypertarget{minval:minvalvalselector}{}
\begin{notedef}
  \texttt{MinVal} selects the smallest value in the domain of the integer variable:
$$\min(D(x))$$
\end{notedef}

\begin{itemize}
	\item \textbf{Constructor} : \mylst{MinVal()}
	\item \textbf{type of variable} : integer
%	\item \textbf{references} : \emph{n/a}
\end{itemize}

%\textbf{Example}:
%\lstinputlisting{java/cabs.j2t}


\section{RandomIntValSelector (Value selector)}\label{randomintvalselector:randomintvalselectorvalselector}\hypertarget{randomintvalselector:randomintvalselectorvalselector}{}
\begin{notedef}
  \texttt{RandomIntValSelector} selects a random value in the domain of the integer variable. The random \texttt{seed} can be specified:
$$rand(D(x))$$
In case of bounded domain, it selects one of the bounds of the domain. 

\end{notedef}

\begin{itemize}
	\item \textbf{Constructor} : 
		\begin{itemize}
		\item \mylst{RandomIntValSelector()}
		\item \mylst{RandomIntValSelector(long seed)}
		\end{itemize}
	\item \textbf{type of variable} : set
%	\item \textbf{references} : \emph{n/a}
\end{itemize}

%\textbf{Example}:
%\lstinputlisting{java/cabs.j2t}


\section{RandomSetValSelector (Value selector)}\label{randomsetvalselector:randomsetvalselectorvalselector}\hypertarget{randomsetvalselector:randomsetvalselectorvalselector}{}
\begin{notedef}
  \texttt{RandomSetValSelector} selects a random value in the open domain of the set variable. The random \texttt{seed} can be specified:
$$rand(Env(x)\setminus Ker(x))$$
\end{notedef}

\begin{itemize}
	\item \textbf{Constructor} : 
		\begin{itemize}
		\item \mylst{RandomSetValSelector()}
		\item \mylst{RandomSetValSelector(long seed)}
		\end{itemize}
	\item \textbf{type of variable} : set
%	\item \textbf{references} : \emph{n/a}
\end{itemize}

%\textbf{Example}:
%\lstinputlisting{java/cabs.j2t}



\chapter{Visualizer (Visualization)}\label{ch:visualizer}\hypertarget{ch:visualizer}{}
\section{Variable-oriented Visualizers}\label{vov:visualizers}\hypertarget{vov:visualizers}{}
This section lists and details the \hyperlink{chocoandcpviz:visualizers}{variable-oriented visualizers} currently available in Choco.
\subsection{BinaryMatrix (visualizer)}\label{binarymatrix:visu}\hypertarget{binarymatrix:visu}{}
\begin{notedef}
  \texttt{BinaryMatrix}$(bool)$ is a specialized visualizer for a matrix of 0/1 variables $bool$.
\end{notedef}

\begin{itemize}
	\item \textbf{API} : 
	\begin{itemize}
	\item \mylst{BinaryMatrix(IntDomainVar[][] bool, String display, int width, int height)}
	\item \mylst{BinaryMatrix(IntDomainVar[][] bool, String display, int x, int y, int width, int height, String group, int min, int max)}
	\end{itemize}
\end{itemize}

%\textbf{Example}:
%\lstinputlisting{java/cabs.j2t}


\subsection{BinaryVector (visualizer)}\label{binaryvector:visu}\hypertarget{binaryvector:visu}{}
\begin{notedef}
  \texttt{BinaryVector}$(bool)$  is a specialized visualizer for a vector of 0/1 variables $bool$.
\end{notedef}

\begin{itemize}
	\item \textbf{API} : 
	\begin{itemize}
	\item \mylst{BinaryVector(IntDomainVar[] bool, String display, int width, int height)}
	\item \mylst{BinaryVector(IntDomainVar[] bool, String display, int x, int y, int width, int height, String group, int min, int max)}
	\end{itemize}
\end{itemize}

%\textbf{Example}:
%\lstinputlisting{java/cabs.j2t}


\subsection{DomainMatrix (visualizer)}\label{domainmatrix:visu}\hypertarget{domainmatrix:visu}{}
\begin{notedef}
  \texttt{DomainMatrix}$(xs)$ is a specialized visualizer for a matrix of variables $xs$.
\end{notedef}

\begin{itemize}
	\item \textbf{API} : 
	\begin{itemize}
	\item \mylst{DomainMatrix(Var[][] vars, String display, int width, int height)}
	\item \mylst{DomainMatrix(Var[][] vars, String display, int x, int y, int width, int height, String group, int min, int max)}
	\end{itemize}
\end{itemize}

%\textbf{Example}:
%\lstinputlisting{java/cabs.j2t}


\subsection{Vector (visualizer)}\label{vector:visu}\hypertarget{vector:visu}{}
\begin{notedef}
  \texttt{Vector}$(vars)$ is a specialized visualizer for variables $vars$.
\end{notedef}

\begin{itemize}
	\item \textbf{API} : 
	\begin{itemize}
	\item \mylst{Vector(Var[] vars, String display, int width, int height)}
	\item \mylst{Vector(Var[] vars, String display, int x, int y, int width, int height, String group, int min, int max)}
	\end{itemize}
\end{itemize}

%\textbf{Example}:
%\lstinputlisting{java/cabs.j2t}


\subsection{VectorSize (visualizer)}\label{vectorsize:visu}\hypertarget{vectorsize:visu}{}
\begin{notedef}
  \texttt{VectorSize}$(vars)$ is specialized visualizer for the sum of cardinalities of variables $vars$.
\end{notedef}

\begin{itemize}
	\item \textbf{API} : 
	\begin{itemize}
	\item \mylst{VectorSize(Var[] vars, String display, int width, int height)}
	\item \mylst{VectorSize(Var[] vars, String display, int x, int y, int width, int height, String group, int min, int max)}
	\end{itemize}
\end{itemize}

%\textbf{Example}:
%\lstinputlisting{java/cabs.j2t}


\subsection{VectorWaterfall (visualizer)}\label{vectorwaterfall:visu}\hypertarget{vectorwaterfall:visu}{}
\begin{notedef}
  \texttt{VectorWaterfall}$(vars$ is a specialized visualizer for a waterfall representation of a vector of variables $vars$.
\end{notedef}

\begin{itemize}
	\item \textbf{API} : 
	\begin{itemize}
	\item \mylst{VectorWaterfall(Var[] vars, String display, int width, int height)}
	\item \mylst{VectorWaterfall(Var[] vars, String display, int x, int y, int width, int height, String group, int min, int max)}
	\end{itemize}
\end{itemize}

%\textbf{Example}:
%\lstinputlisting{java/cabs.j2t}



\section{Constraint-oriented Visualizers}\label{cov:visualizers}\hypertarget{cov:visualizers}{}
This section lists and details the \hyperlink{chocoandcpviz:visualizers}{constraint-oriented visualizers} currently available in Choco.
\subsection{AllDifferent (visualizer)}\label{alldiff:visu}\hypertarget{alldiff:visu}{}
\begin{notedef}
  \texttt{AllDifferent}$(vars)$ a specialized visualizer for the \mylst{AllDifferent} constraint, where $vars$ is an array of variables involved in one AllDifferent constraint.
\end{notedef}

\begin{itemize}
	\item \textbf{API} : 
	\begin{itemize}
	\item \mylst{AllDifferent(IntDomainVar[] vars, String display, int width, int height)}
	\item \mylst{AllDifferent(IntDomainVar[] vars, String display, int x, int y, int width, int height, String group, int min, int max)}
	\end{itemize}
\end{itemize}

%\textbf{Example}:
%\lstinputlisting{java/cabs.j2t}


\subsection{AllDifferentMatrix (visualizer)}\label{alldiffmatrix:visu}\hypertarget{alldiffmatrix:visu}{}
\begin{notedef}
  \texttt{AllDifferentMatrix}$(vars)$ a specialized visualizer for a list of \mylst{AllDifferent} constraints, where $vars$ is a matrix of variables involved in $n$ AllDifferent constraints.
\end{notedef}

\begin{itemize}
	\item \textbf{API} : 
	\begin{itemize}
	\item \mylst{AllDifferentMatrix(IntDomainVar[][] vars, String display, int width, int height)}
	\item \mylst{AllDifferentMatrix(IntDomainVar[][] vars, String type, String display, int x, int y, int width, int height, String group, int min, int max)}
	\end{itemize}
\end{itemize}

%\textbf{Example}:
%\lstinputlisting{java/cabs.j2t}


\subsection{BooleanChanneling (visualizer)}\label{boolchan:visu}\hypertarget{boolchan:visu}{}
\begin{notedef}
  \texttt{BooleanChanneling}$(x, bs, o)$ is a specialized visualizer for the boolean channeling constraint, where $x$ is a variable, $bs$ an array of boolean variables and $o$ the first value (in Choco: \mylst{domainChanneling}).
\end{notedef}

\begin{itemize}
	\item \textbf{API} : 
	\begin{itemize}
	\item \mylst{BoolChanneling(IntDomainVar x, IntDomainVar[] bs, int o, String display, int width, int height)}
	\item \mylst{BoolChanneling(IntDomainVar x, IntDomainVar[] bs, String display, int width, int height)}, where $o$ is computed from the lower bound $x$'s domain.
	\item \mylst{BoolChanneling(IntDomainVar var, IntDomainVar[] bool, int o, String display, int x, int y, int width, int height, String group, int min, int max)}
	\item \mylst{BoolChanneling(IntDomainVar var, IntDomainVar[] bool, String display, int x, int y, int width, int height, String group, int min, int max)}, where $o$ is computed from the lower bound $x$'s domain.
	\end{itemize}
\end{itemize}

%\textbf{Example}:
%\lstinputlisting{java/cabs.j2t}


\subsection{Cumulative (visualizer)}\label{cumulative:visu}\hypertarget{cumulative:visu}{}
\begin{notedef}
  \texttt{Cumulative}$(ts, l, e)$ is a specialized visualizer for the \mylst{Cumulative} constraint, where $ts$ is an array of task variables, $l$ the limit variable and $e$ the end variable.
\end{notedef}

\begin{itemize}
	\item \textbf{API} : 
	\begin{itemize}
	\item \mylst{Cumulative(TaskVar[] tasks, IntDomainVar limit, IntDomainVar end, String display, int width, int height)}, where \mylst{display} can take its value in \mylst{\{"compact", "expanded", "gantt"\}}
	\item \mylst{Cumulative(TaskVar[] tasks, IntDomainVar limit, IntDomainVar end, String display, int x, int y, int width, int height, String group, int min, int max)}, where \mylst{display} can take its value in \mylst{\{"compact", "expanded", "gantt"\}}
	\end{itemize}
\end{itemize}

%\textbf{Example}:
%\lstinputlisting{java/cabs.j2t}


\subsection{Element (visualizer)}\label{element:visu}\hypertarget{element:visu}{}
\begin{notedef}
  \texttt{Element}$(index, values, value)$ is a specialized visualizer for the \mylst{Element} constraint, where $index$ is a variable, $values$ an array of integer values and $value$ a variable.
\end{notedef}

\begin{itemize}
	\item \textbf{API} : 
	\begin{itemize}
	\item \mylst{Element(IntDomainVar index, int[] values, IntDomainVar value, String display, int width, int height)}
	\item \mylst{Element(IntDomainVar index, int[] values, IntDomainVar value, String display, int x, int y, int width, int height, String group, int min, int max)}
	\end{itemize}
\end{itemize}

%\textbf{Example}:
%\lstinputlisting{java/cabs.j2t}


\subsection{Gcc (visualizer)}\label{gcc:visu}\hypertarget{gcc:visu}{}
\begin{notedef}
  \texttt{Gcc}$(vars, values, low, high)$ is a specialized visualizer for the \mylst{Gcc} constraint, where $vars$ is an array of variables, $values$ an array of integer values, $low$ minimum number of occurrences of values and $high$ maximum number of occurrences of values.
\end{notedef}

\begin{itemize}
	\item \textbf{API} : 
	\begin{itemize}
	\item \mylst{Gcc(IntDomainVar[] vars, int[] values, int[] low, int[] high, String display, int width, int height)}
	\item \mylst{Gcc(IntDomainVar[] vars, int[] values, int[] low, int[] high, String display, int x, int y, int width, int height, String group, int min, int max) }
	\end{itemize}
\end{itemize}

%\textbf{Example}:
%\lstinputlisting{java/cabs.j2t}


\section{Inverse (visualizer)}\label{inverse:visu}\hypertarget{inverse:visu}{}
\begin{notedef}
  \texttt{Inverse}$(X, Y)$ is a specialized visualizer for the \mylst{Inverse} constraint, where $X$ is an array of variables and $Y$ is an other array of variables.
\end{notedef}

\begin{itemize}
	\item \textbf{API} : 
	\begin{itemize}
	\item \mylst{Inverse(IntDomainVar[] X, IntDomainVar[] Y, String display, int width, int height)}
	\item \mylst{Inverse(IntDomainVar[] X, IntDomainVar[] Y, String display, int x, int y, int width, int height, String group, int min, int max)}
	\end{itemize}
\end{itemize}

%\textbf{Example}:
%\lstinputlisting{java/cabs.j2t}


\subsection{LexLe (visualizer)}\label{lexle:visu}\hypertarget{lexle:visu}{}
\begin{notedef}
  \texttt{LexLe}$(X, Y)$ is a specialized visualizer for the \mylst{LexLe} constraint, where $X$ is an array of variables and $Y$ is an other array of variables.
\end{notedef}

\begin{itemize}
	\item \textbf{API} : 
	\begin{itemize}
	\item \mylst{LexLe(IntDomainVar[] X, IntDomainVar[] Y, String display, int width, int height)}
	\item \mylst{LexLe(IntDomainVar[] X, IntDomainVar[] Y, String display, int x, int y, int width, int height, String group, int min, int max)}
	\end{itemize}
\end{itemize}

%\textbf{Example}:
%\lstinputlisting{java/cabs.j2t}


\subsection{LexLt (visualizer)}\label{lexlt:visu}\hypertarget{lexlt:visu}{}
\begin{notedef}
  \texttt{LexLt}$(X, Y)$ is a specialized visualizer for the \mylst{LexLt} constraint, where $X$ is an array of variables and $Y$ is an other array of variables.
\end{notedef}

\begin{itemize}
	\item \textbf{API} : 
	\begin{itemize}
	\item \mylst{LexLt(IntDomainVar[] X, IntDomainVar[] Y, String display, int width, int height)}
	\item \mylst{LexLt(IntDomainVar[] X, IntDomainVar[] Y, String display, int x, int y, int width, int height, String group, int min, int max)}
	\end{itemize}
\end{itemize}

%\textbf{Example}:
%\lstinputlisting{java/cabs.j2t}



\part{Extras}\label{ch:extra}\hypertarget{ch:extra}{}
%\input{chapters/Choco_and_visu.tex}
%\part{sudoku and cp}
\label{sudokuandcp}
\hypertarget{sudokuandcp}{}

\chapter{Sudoku and Constraint Programming}\label{sudokuandcp:sudokuandconstraintprogramming}\hypertarget{sudokuandcp:sudokuandconstraintprogramming}{}

\section{Sudoku ?!?}\label{sudokuandcp:sudoku!}\hypertarget{sudokuandcp:sudoku!}{}

\insertGraphique{5cm}{media/sudokuillustration.jpg}{A sudoku grid}

Everybody knows those grids that appeared last year in the subway, in wating lounges, on colleague's desks, etc. In Japanese \emph{su} means digit and \emph{doku}, unique. But this game has been discovered by an American ! The first grids appeared in the USA in 1979 (they were hand crafted). \href{http://en.wikipedia.org/wiki/sudoku}{Wikipedia} tells us that they were designed by Howard Garns a retired architect. He died in 1989 well before the success story of sudoku initiated by Wayne Gould, a retired judge from Hong-Kong. The rules are really simple: a 81 cells square grid is divided in 9 smaller blocks of 9 cells (3 x 3). Some of the 81 are filled with one digit. The aim of the puzzle is to fill in the other cells, using digits except 0, such as each digit appears once and only once in each row, each column and each smaller block. The solution is unique.

\subsection{Solving sudokus}\label{sudokuandcp:solvingsudokus}\hypertarget{sudokuandcp:solvingsudokus}{}

Many computer techniques exist to quickly solve a sudoku puzzle. Mainly, they are based on backtracking algorithms. The idea is the following: give a free cell a value and continue as long as choices remain consistent. As soon as an inconsistency is detected, the computer program backtracks to its earliest past choice et tries another value. If no more value is available, the program keeps backtracking until it can go forward again. This systematic technique make it sure to solve a sudoku grid. However, no human player plays this way: this needs too much memory ! 

\begin{note}
see \href{http://en.wikipedia.org/wiki/sudoku}{Wikipedia} for a panel of solving techniques.
\end{note}

\section{Sudoku and Artificial Intelligence}\label{sudokuandcp:sudokuandartificialintelligence}\hypertarget{sudokuandcp:sudokuandartificialintelligence}{}

Many techniques and rules have been designed and discovered to solve sudoku grids. Some are really simple, some need to use some useful tools: pencil and eraser. 

\subsection{Simple rules: single candidate and single position}\label{sudokuandcp:simplerules:singlecandidateandsingleposition}\hypertarget{sudokuandcp:simplerules:singlecandidateandsingleposition}{}

\insertGraphique{5cm}{media/sudokuillustrationscsp.jpg}{Simple rules: single candidates and single position} 

Let consider the grid on Figure~\ref{fig:media/sudokuillustrationscsp.jpg} and the cell with the red dot. In the same line, we find: 3, 4, 6, 7, and 9. In the same column: 2, 3, 5, and 8. In the same block: 2, 7, 8, and 9. There remain only one possibility: \textbf{1}. This is the \textbf{single candidate} rule. This cell should be filled in with \textbf{1}. 

Now let consider a given digit: let's say 4. In the block with a blue dot, there is no 4. Where can it be ? The 4's in the surrounding blocks heavily constrain the problem. There is a \textbf{single position} possible: the blue dot. This another simple rule to apply.

Alternatively using these two rules allows a player to fill in many cells and even solve the simplest grids. 
But, limits are easily reached. More subtle approaches are needed: but an important tool is now needed ... an eraser ! 

\subsection{Human reasoning principles}\label{sudokuandcp:humanreasoningprinciples}\hypertarget{sudokuandcp:humanreasoningprinciples}{}

\insertGraphique{7cm}{media/sudokuillustrationmarks.jpg}{Introducing marks}

Many techniques do exist but a vaste amount of them rely on simple principles. The first one is: do not try to find the value of a cell but instead focus on values that \textbf{will never be assigned} to it. The space of possibility is then reduced. This is where the eraser comes handy. Many players marks the remaining possibilities as in the grid on the left. 

Using this information, rather subtle reasoning is possible. For example, consider the seventh column on the grid on the left. Two cells contain as possible values the two values 5 and 7. This means that those two values cannot appear elsewhere in that very same column. Therefore, the other unassigned cell on the column can only contain a 6. We have \emph{deduced} something.

\noindent\begin{minipage}[b]{0.8\linewidth}
This was an easy to spot inference. This is not always the case. Consider the part of the grid on the right. Let us consider the third column. For cells 4 and 5, only two values are available: 4 and 8. Those values cannot be assigned to any other cell in that column. Therefore, in cell 6 we have a 3, and thus and 7 in cell 2 and finally a 1 in cell 3. This can be a very powerful rule.

Such a reasoning (sometimes called \emph{Naked Pairs}) is easily generalized to any number of cells (always in the same region: row, column or block) presenting this same configuration. This local reasoning can be applied to any region of the grid. It is important to notice that the inferred information can (and should) be used from a region to another.   
\end{minipage}%
\begin{minipage}[m]{0.2\linewidth}
~~\Graph{media/sudokuillustrationpart.jpg}{width=3cm}
\end{minipage}

\noindent The following principles of \emph{human} reasoning can be listed: 
\begin{itemize}
	\item reasoning on \emph{possible} values for a cell (by erasing impossible ones)
	\item systematically applying an evolved local reasoning (such as the \emph{Naked Pairs} rule)
	\item transmitting inferred information from a region to another related through a given a set of cells
\end{itemize}

\subsection{Towards Constraint Programming}\label{sudokuandcp:towardsconstraintprogramming}\hypertarget{sudokuandcp:towardsconstraintprogramming}{}

Those three principles are at the core of \textbf{constraint programming} a recent technique coming from both \emph{artificial intelligence} and \emph{operations research}.

\begin{itemize}
	\item The first principle is called \textbf{domain reduction} or \emph{filtering}
	\item The second considers its region as a \textbf{constraint} (a relation to be verified by the solution of the problem): here we consider an \emph{all different} constraint (all the values must be different in a given region). Constraints are considered \textbf{locally} for reasoning
	\item The third principle is called \textbf{propagation}: constraints (regions) communicate with one another through the available values in variables (cells)
\end{itemize}

Constraint programming is able to solve this problem as a human would do. Moreover, a large majority of the rules and techniques described on the Internet amount to a well-known problem: the alldifferent problem. A \textbf{constraint solver} (as \textbf{Choco}) is therefore able to reason on this problem allowing the solving of sudoku grid as a human would do although it has not be specifically designed to.

Ideally, iterating local reasoning will lead to a solution. However, for exceptionnaly hard grids, an enumerating phase (all constraint solvers provide tools for that) relying on backtracking may be necessary.

\section{See also}\label{sudokuandcp:seealso}\hypertarget{sudokuandcp:seealso}{}

\begin{itemize}
	\item \href{http://njussien.e-constraints.net/sudoku/eng-jouer.html}{SudokuHelper} a sudoku solver and helper applet developed with \emph{Choco}.
	\item \href{http://www.palmsudoku.com}{PalmSudoku} a rather complete list of rules and tips for solving sudokus
\end{itemize}



\backmatter

\printglossaries
\addcontentsline{toc}{part}{Glossary}

\addcontentsline{toc}{part}{Bibliography}
\bibliographystyle{apalike}
\bibliography{choco-doc}

\label{licence}
\hypertarget{licence}{}


%---------------------------------------------------------------------
\chapter*{\rlap{GNU Free Documentation License}}
\phantomsection  % so hyperref creates bookmarks
\addcontentsline{toc}{chapter}{GNU Free Documentation License}
%\label{label_fdl}

 \begin{center}

       Version 1.3, 3 November 2008


 Copyright \copyright{} 2000, 2001, 2002, 2007, 2008  Free Software Foundation, Inc.
 
 \bigskip
 
     \url{http://fsf.org/}
  
 \bigskip
 
 Everyone is permitted to copy and distribute verbatim copies
 of this license document, but changing it is not allowed.
\end{center}


\begin{center}
{\bf\large Preamble}
\end{center}

The purpose of this License is to make a manual, textbook, or other
functional and useful document ``free'' in the sense of freedom: to
assure everyone the effective freedom to copy and redistribute it,
with or without modifying it, either commercially or noncommercially.
Secondarily, this License preserves for the author and publisher a way
to get credit for their work, while not being considered responsible
for modifications made by others.

This License is a kind of ``copyleft'', which means that derivative
works of the document must themselves be free in the same sense.  It
complements the GNU General Public License, which is a copyleft
license designed for free software.

We have designed this License in order to use it for manuals for free
software, because free software needs free documentation: a free
program should come with manuals providing the same freedoms that the
software does.  But this License is not limited to software manuals;
it can be used for any textual work, regardless of subject matter or
whether it is published as a printed book.  We recommend this License
principally for works whose purpose is instruction or reference.


\begin{center}
{\Large\bf 1. APPLICABILITY AND DEFINITIONS\par}
\phantomsection
\addcontentsline{toc}{section}{1. APPLICABILITY AND DEFINITIONS}
\end{center}

This License applies to any manual or other work, in any medium, that
contains a notice placed by the copyright holder saying it can be
distributed under the terms of this License.  Such a notice grants a
world-wide, royalty-free license, unlimited in duration, to use that
work under the conditions stated herein.  The ``\textbf{Document}'', below,
refers to any such manual or work.  Any member of the public is a
licensee, and is addressed as ``\textbf{you}''.  You accept the license if you
copy, modify or distribute the work in a way requiring permission
under copyright law.

A ``\textbf{Modified Version}'' of the Document means any work containing the
Document or a portion of it, either copied verbatim, or with
modifications and/or translated into another language.

A ``\textbf{Secondary Section}'' is a named appendix or a front-matter section of
the Document that deals exclusively with the relationship of the
publishers or authors of the Document to the Document's overall subject
(or to related matters) and contains nothing that could fall directly
within that overall subject.  (Thus, if the Document is in part a
textbook of mathematics, a Secondary Section may not explain any
mathematics.)  The relationship could be a matter of historical
connection with the subject or with related matters, or of legal,
commercial, philosophical, ethical or political position regarding
them.

The ``\textbf{Invariant Sections}'' are certain Secondary Sections whose titles
are designated, as being those of Invariant Sections, in the notice
that says that the Document is released under this License.  If a
section does not fit the above definition of Secondary then it is not
allowed to be designated as Invariant.  The Document may contain zero
Invariant Sections.  If the Document does not identify any Invariant
Sections then there are none.

The ``\textbf{Cover Texts}'' are certain short passages of text that are listed,
as Front-Cover Texts or Back-Cover Texts, in the notice that says that
the Document is released under this License.  A Front-Cover Text may
be at most 5 words, and a Back-Cover Text may be at most 25 words.

A ``\textbf{Transparent}'' copy of the Document means a machine-readable copy,
represented in a format whose specification is available to the
general public, that is suitable for revising the document
straightforwardly with generic text editors or (for images composed of
pixels) generic paint programs or (for drawings) some widely available
drawing editor, and that is suitable for input to text formatters or
for automatic translation to a variety of formats suitable for input
to text formatters.  A copy made in an otherwise Transparent file
format whose markup, or absence of markup, has been arranged to thwart
or discourage subsequent modification by readers is not Transparent.
An image format is not Transparent if used for any substantial amount
of text.  A copy that is not ``Transparent'' is called ``\textbf{Opaque}''.

Examples of suitable formats for Transparent copies include plain
ASCII without markup, Texinfo input format, LaTeX input format, SGML
or XML using a publicly available DTD, and standard-conforming simple
HTML, PostScript or PDF designed for human modification.  Examples of
transparent image formats include PNG, XCF and JPG.  Opaque formats
include proprietary formats that can be read and edited only by
proprietary word processors, SGML or XML for which the DTD and/or
processing tools are not generally available, and the
machine-generated HTML, PostScript or PDF produced by some word
processors for output purposes only.

The ``\textbf{Title Page}'' means, for a printed book, the title page itself,
plus such following pages as are needed to hold, legibly, the material
this License requires to appear in the title page.  For works in
formats which do not have any title page as such, ``Title Page'' means
the text near the most prominent appearance of the work's title,
preceding the beginning of the body of the text.

The ``\textbf{publisher}'' means any person or entity that distributes
copies of the Document to the public.

A section ``\textbf{Entitled XYZ}'' means a named subunit of the Document whose
title either is precisely XYZ or contains XYZ in parentheses following
text that translates XYZ in another language.  (Here XYZ stands for a
specific section name mentioned below, such as ``\textbf{Acknowledgements}'',
``\textbf{Dedications}'', ``\textbf{Endorsements}'', or ``\textbf{History}''.)  
To ``\textbf{Preserve the Title}''
of such a section when you modify the Document means that it remains a
section ``Entitled XYZ'' according to this definition.

The Document may include Warranty Disclaimers next to the notice which
states that this License applies to the Document.  These Warranty
Disclaimers are considered to be included by reference in this
License, but only as regards disclaiming warranties: any other
implication that these Warranty Disclaimers may have is void and has
no effect on the meaning of this License.


\begin{center}
{\Large\bf 2. VERBATIM COPYING\par}
\phantomsection
\addcontentsline{toc}{section}{2. VERBATIM COPYING}
\end{center}

You may copy and distribute the Document in any medium, either
commercially or noncommercially, provided that this License, the
copyright notices, and the license notice saying this License applies
to the Document are reproduced in all copies, and that you add no other
conditions whatsoever to those of this License.  You may not use
technical measures to obstruct or control the reading or further
copying of the copies you make or distribute.  However, you may accept
compensation in exchange for copies.  If you distribute a large enough
number of copies you must also follow the conditions in section~3.

You may also lend copies, under the same conditions stated above, and
you may publicly display copies.


\begin{center}
{\Large\bf 3. COPYING IN QUANTITY\par}
\phantomsection
\addcontentsline{toc}{section}{3. COPYING IN QUANTITY}
\end{center}


If you publish printed copies (or copies in media that commonly have
printed covers) of the Document, numbering more than 100, and the
Document's license notice requires Cover Texts, you must enclose the
copies in covers that carry, clearly and legibly, all these Cover
Texts: Front-Cover Texts on the front cover, and Back-Cover Texts on
the back cover.  Both covers must also clearly and legibly identify
you as the publisher of these copies.  The front cover must present
the full title with all words of the title equally prominent and
visible.  You may add other material on the covers in addition.
Copying with changes limited to the covers, as long as they preserve
the title of the Document and satisfy these conditions, can be treated
as verbatim copying in other respects.

If the required texts for either cover are too voluminous to fit
legibly, you should put the first ones listed (as many as fit
reasonably) on the actual cover, and continue the rest onto adjacent
pages.

If you publish or distribute Opaque copies of the Document numbering
more than 100, you must either include a machine-readable Transparent
copy along with each Opaque copy, or state in or with each Opaque copy
a computer-network location from which the general network-using
public has access to download using public-standard network protocols
a complete Transparent copy of the Document, free of added material.
If you use the latter option, you must take reasonably prudent steps,
when you begin distribution of Opaque copies in quantity, to ensure
that this Transparent copy will remain thus accessible at the stated
location until at least one year after the last time you distribute an
Opaque copy (directly or through your agents or retailers) of that
edition to the public.

It is requested, but not required, that you contact the authors of the
Document well before redistributing any large number of copies, to give
them a chance to provide you with an updated version of the Document.


\begin{center}
{\Large\bf 4. MODIFICATIONS\par}
\phantomsection
\addcontentsline{toc}{section}{4. MODIFICATIONS}
\end{center}

You may copy and distribute a Modified Version of the Document under
the conditions of sections 2 and 3 above, provided that you release
the Modified Version under precisely this License, with the Modified
Version filling the role of the Document, thus licensing distribution
and modification of the Modified Version to whoever possesses a copy
of it.  In addition, you must do these things in the Modified Version:

\begin{itemize}
\item[A.] 
   Use in the Title Page (and on the covers, if any) a title distinct
   from that of the Document, and from those of previous versions
   (which should, if there were any, be listed in the History section
   of the Document).  You may use the same title as a previous version
   if the original publisher of that version gives permission.
   
\item[B.]
   List on the Title Page, as authors, one or more persons or entities
   responsible for authorship of the modifications in the Modified
   Version, together with at least five of the principal authors of the
   Document (all of its principal authors, if it has fewer than five),
   unless they release you from this requirement.
   
\item[C.]
   State on the Title page the name of the publisher of the
   Modified Version, as the publisher.
   
\item[D.]
   Preserve all the copyright notices of the Document.
   
\item[E.]
   Add an appropriate copyright notice for your modifications
   adjacent to the other copyright notices.
   
\item[F.]
   Include, immediately after the copyright notices, a license notice
   giving the public permission to use the Modified Version under the
   terms of this License, in the form shown in the Addendum below.
   
\item[G.]
   Preserve in that license notice the full lists of Invariant Sections
   and required Cover Texts given in the Document's license notice.
   
\item[H.]
   Include an unaltered copy of this License.
   
\item[I.]
   Preserve the section Entitled ``History'', Preserve its Title, and add
   to it an item stating at least the title, year, new authors, and
   publisher of the Modified Version as given on the Title Page.  If
   there is no section Entitled ``History'' in the Document, create one
   stating the title, year, authors, and publisher of the Document as
   given on its Title Page, then add an item describing the Modified
   Version as stated in the previous sentence.
   
\item[J.]
   Preserve the network location, if any, given in the Document for
   public access to a Transparent copy of the Document, and likewise
   the network locations given in the Document for previous versions
   it was based on.  These may be placed in the ``History'' section.
   You may omit a network location for a work that was published at
   least four years before the Document itself, or if the original
   publisher of the version it refers to gives permission.
   
\item[K.]
   For any section Entitled ``Acknowledgements'' or ``Dedications'',
   Preserve the Title of the section, and preserve in the section all
   the substance and tone of each of the contributor acknowledgements
   and/or dedications given therein.
   
\item[L.]
   Preserve all the Invariant Sections of the Document,
   unaltered in their text and in their titles.  Section numbers
   or the equivalent are not considered part of the section titles.
   
\item[M.]
   Delete any section Entitled ``Endorsements''.  Such a section
   may not be included in the Modified Version.
   
\item[N.]
   Do not retitle any existing section to be Entitled ``Endorsements''
   or to conflict in title with any Invariant Section.
   
\item[O.]
   Preserve any Warranty Disclaimers.
\end{itemize}

If the Modified Version includes new front-matter sections or
appendices that qualify as Secondary Sections and contain no material
copied from the Document, you may at your option designate some or all
of these sections as invariant.  To do this, add their titles to the
list of Invariant Sections in the Modified Version's license notice.
These titles must be distinct from any other section titles.

You may add a section Entitled ``Endorsements'', provided it contains
nothing but endorsements of your Modified Version by various
parties---for example, statements of peer review or that the text has
been approved by an organization as the authoritative definition of a
standard.

You may add a passage of up to five words as a Front-Cover Text, and a
passage of up to 25 words as a Back-Cover Text, to the end of the list
of Cover Texts in the Modified Version.  Only one passage of
Front-Cover Text and one of Back-Cover Text may be added by (or
through arrangements made by) any one entity.  If the Document already
includes a cover text for the same cover, previously added by you or
by arrangement made by the same entity you are acting on behalf of,
you may not add another; but you may replace the old one, on explicit
permission from the previous publisher that added the old one.

The author(s) and publisher(s) of the Document do not by this License
give permission to use their names for publicity for or to assert or
imply endorsement of any Modified Version.


\begin{center}
{\Large\bf 5. COMBINING DOCUMENTS\par}
\phantomsection
\addcontentsline{toc}{section}{5. COMBINING DOCUMENTS}
\end{center}


You may combine the Document with other documents released under this
License, under the terms defined in section~4 above for modified
versions, provided that you include in the combination all of the
Invariant Sections of all of the original documents, unmodified, and
list them all as Invariant Sections of your combined work in its
license notice, and that you preserve all their Warranty Disclaimers.

The combined work need only contain one copy of this License, and
multiple identical Invariant Sections may be replaced with a single
copy.  If there are multiple Invariant Sections with the same name but
different contents, make the title of each such section unique by
adding at the end of it, in parentheses, the name of the original
author or publisher of that section if known, or else a unique number.
Make the same adjustment to the section titles in the list of
Invariant Sections in the license notice of the combined work.

In the combination, you must combine any sections Entitled ``History''
in the various original documents, forming one section Entitled
``History''; likewise combine any sections Entitled ``Acknowledgements'',
and any sections Entitled ``Dedications''.  You must delete all sections
Entitled ``Endorsements''.

\begin{center}
{\Large\bf 6. COLLECTIONS OF DOCUMENTS\par}
\phantomsection
\addcontentsline{toc}{section}{6. COLLECTIONS OF DOCUMENTS}
\end{center}

You may make a collection consisting of the Document and other documents
released under this License, and replace the individual copies of this
License in the various documents with a single copy that is included in
the collection, provided that you follow the rules of this License for
verbatim copying of each of the documents in all other respects.

You may extract a single document from such a collection, and distribute
it individually under this License, provided you insert a copy of this
License into the extracted document, and follow this License in all
other respects regarding verbatim copying of that document.


\begin{center}
{\Large\bf 7. AGGREGATION WITH INDEPENDENT WORKS\par}
\phantomsection
\addcontentsline{toc}{section}{7. AGGREGATION WITH INDEPENDENT WORKS}
\end{center}


A compilation of the Document or its derivatives with other separate
and independent documents or works, in or on a volume of a storage or
distribution medium, is called an ``aggregate'' if the copyright
resulting from the compilation is not used to limit the legal rights
of the compilation's users beyond what the individual works permit.
When the Document is included in an aggregate, this License does not
apply to the other works in the aggregate which are not themselves
derivative works of the Document.

If the Cover Text requirement of section~3 is applicable to these
copies of the Document, then if the Document is less than one half of
the entire aggregate, the Document's Cover Texts may be placed on
covers that bracket the Document within the aggregate, or the
electronic equivalent of covers if the Document is in electronic form.
Otherwise they must appear on printed covers that bracket the whole
aggregate.


\begin{center}
{\Large\bf 8. TRANSLATION\par}
\phantomsection
\addcontentsline{toc}{section}{8. TRANSLATION}
\end{center}


Translation is considered a kind of modification, so you may
distribute translations of the Document under the terms of section~4.
Replacing Invariant Sections with translations requires special
permission from their copyright holders, but you may include
translations of some or all Invariant Sections in addition to the
original versions of these Invariant Sections.  You may include a
translation of this License, and all the license notices in the
Document, and any Warranty Disclaimers, provided that you also include
the original English version of this License and the original versions
of those notices and disclaimers.  In case of a disagreement between
the translation and the original version of this License or a notice
or disclaimer, the original version will prevail.

If a section in the Document is Entitled ``Acknowledgements'',
``Dedications'', or ``History'', the requirement (section~4) to Preserve
its Title (section~1) will typically require changing the actual
title.


\begin{center}
{\Large\bf 9. TERMINATION\par}
\phantomsection
\addcontentsline{toc}{section}{9. TERMINATION}
\end{center}


You may not copy, modify, sublicense, or distribute the Document
except as expressly provided under this License.  Any attempt
otherwise to copy, modify, sublicense, or distribute it is void, and
will automatically terminate your rights under this License.

However, if you cease all violation of this License, then your license
from a particular copyright holder is reinstated (a) provisionally,
unless and until the copyright holder explicitly and finally
terminates your license, and (b) permanently, if the copyright holder
fails to notify you of the violation by some reasonable means prior to
60 days after the cessation.

Moreover, your license from a particular copyright holder is
reinstated permanently if the copyright holder notifies you of the
violation by some reasonable means, this is the first time you have
received notice of violation of this License (for any work) from that
copyright holder, and you cure the violation prior to 30 days after
your receipt of the notice.

Termination of your rights under this section does not terminate the
licenses of parties who have received copies or rights from you under
this License.  If your rights have been terminated and not permanently
reinstated, receipt of a copy of some or all of the same material does
not give you any rights to use it.


\begin{center}
{\Large\bf 10. FUTURE REVISIONS OF THIS LICENSE\par}
\phantomsection
\addcontentsline{toc}{section}{10. FUTURE REVISIONS OF THIS LICENSE}
\end{center}


The Free Software Foundation may publish new, revised versions
of the GNU Free Documentation License from time to time.  Such new
versions will be similar in spirit to the present version, but may
differ in detail to address new problems or concerns.  See
http://www.gnu.org/copyleft/.

Each version of the License is given a distinguishing version number.
If the Document specifies that a particular numbered version of this
License ``or any later version'' applies to it, you have the option of
following the terms and conditions either of that specified version or
of any later version that has been published (not as a draft) by the
Free Software Foundation.  If the Document does not specify a version
number of this License, you may choose any version ever published (not
as a draft) by the Free Software Foundation.  If the Document
specifies that a proxy can decide which future versions of this
License can be used, that proxy's public statement of acceptance of a
version permanently authorizes you to choose that version for the
Document.


\begin{center}
{\Large\bf 11. RELICENSING\par}
\phantomsection
\addcontentsline{toc}{section}{11. RELICENSING}
\end{center}


``Massive Multiauthor Collaboration Site'' (or ``MMC Site'') means any
World Wide Web server that publishes copyrightable works and also
provides prominent facilities for anybody to edit those works.  A
public wiki that anybody can edit is an example of such a server.  A
``Massive Multiauthor Collaboration'' (or ``MMC'') contained in the
site means any set of copyrightable works thus published on the MMC
site.

``CC-BY-SA'' means the Creative Commons Attribution-Share Alike 3.0
license published by Creative Commons Corporation, a not-for-profit
corporation with a principal place of business in San Francisco,
California, as well as future copyleft versions of that license
published by that same organization.

``Incorporate'' means to publish or republish a Document, in whole or
in part, as part of another Document.

An MMC is ``eligible for relicensing'' if it is licensed under this
License, and if all works that were first published under this License
somewhere other than this MMC, and subsequently incorporated in whole
or in part into the MMC, (1) had no cover texts or invariant sections,
and (2) were thus incorporated prior to November 1, 2008.

The operator of an MMC Site may republish an MMC contained in the site
under CC-BY-SA on the same site at any time before August 1, 2009,
provided the MMC is eligible for relicensing.


\begin{center}
{\Large\bf ADDENDUM: How to use this License for your documents\par}
\phantomsection
\addcontentsline{toc}{section}{ADDENDUM: How to use this License for your documents}
\end{center}

To use this License in a document you have written, include a copy of
the License in the document and put the following copyright and
license notices just after the title page:

\bigskip
\begin{quote}
    Copyright \copyright{}  YEAR  YOUR NAME.
    Permission is granted to copy, distribute and/or modify this document
    under the terms of the GNU Free Documentation License, Version 1.3
    or any later version published by the Free Software Foundation;
    with no Invariant Sections, no Front-Cover Texts, and no Back-Cover Texts.
    A copy of the license is included in the section entitled ``GNU
    Free Documentation License''.
\end{quote}
\bigskip
    
If you have Invariant Sections, Front-Cover Texts and Back-Cover Texts,
replace the ``with \dots\ Texts.'' line with this:

\bigskip
\begin{quote}
    with the Invariant Sections being LIST THEIR TITLES, with the
    Front-Cover Texts being LIST, and with the Back-Cover Texts being LIST.
\end{quote}
\bigskip
    
If you have Invariant Sections without Cover Texts, or some other
combination of the three, merge those two alternatives to suit the
situation.

If your document contains nontrivial examples of program code, we
recommend releasing these examples in parallel under your choice of
free software license, such as the GNU General Public License,
to permit their use in free software.

%---------------------------------------------------------------------


\end{document}

%%% Local Variables: 
%%% mode: latex
%%% TeX-master: t
%%% End: 
